% !TEX program = pdflatex
% !TEX encoding = UTF-8 Unicode

% Trabajo de Fin de Grado: Una perspectiva algorímtica a la incompletitud de las Matemáticas
% Miguel Ángel Fernández Gutiérrez

\documentclass{scrbook}

\KOMAoptions{%
  fontsize=10pt,        % Tamaño de fuente
  paper=a4,             % Tamaño del papel
  headings=normal,      % Tamaño de letra para los títulos: small, normal, big
  parskip=half,         % Espacio entre párrafos: full (una línea) o half (media línea)
  headsepline=false,    % Una linea separa la cabecera del texto
  cleardoublepage=empty,% No imprime cabecera ni pie en páginas en blanco 
  chapterprefix=false,  % No antepone el texto "capítulo" antes del número
  appendixprefix=false, % No antepone el texto "Apéndice" antes de la letra
  listof=totoc,         % Añade a la tabla de contenidos la lista de tablas y figuras
  index=totoc,          % Añade a la talba de contenidos una entrada para el índice
  bibliography=totoc,	  % Añade a la tabla de contenidos una entrada para bibliografía
  %BCOR=5mm,            % Reserva margen interior para la encuadernación. 
                        % El valor dependerá el tipo de encuadernado y del grosor del libro.
  DIV=10,               % Cálcula el diseño de página según ciertos 
                        % parámetros. Al aumentar el número aumentamos el ancho de texto y disminuimos el ancho del margen. Una opción de 14 producirá márgenes estrechos y texto ancho.
  twoside=false
}

% INFORMACIÓN PARA LA VERSIÓN IMPRESA
% Si el documento ha de ser impreso en papel de tamaño a4 pero el tamaño del documento (elegido en \KOMAoptions con la ocpión paper) no es a4 descomentar la línea que carga el paquete `crop` más abajo. El paquete crop se encargará de centrar el documento en un a4 e imprimir unas guías de corte. El procedimiento completo para imprenta sería el siguiente:
% 0. Determinar, según el tipo de encuadernación del documento, el ancho reservado para el proceso de encuadernación (preguntar en la imprenta), es decir, la anchura del área del papel que se pierde durante el proceso de encuadernación. Fijar la varibale BCOR de \KOMAoptions a dicho valor.
% 1. Descomentar la siguiente línea e imprimir una única página con las guías de corte
% 2. Cambiar la opción `cross` por `cam` (o `off`) en el paquete crop y volver a compilar. Imprimir el documento (las guías de corte impresas no inferfieren con el texto).
% 3. Usar la página con las guías impresas en el punto 1 para cortar todas las páginas.

% \usepackage[a4, odd, center, pdflatex, cross]{crop} % Permite imprimir el documento en un a4 (si el tamaño es más pequeño) mostrando unas guías de corte. Útil para imprenta.

% VERSIÓN ELECTRÓNICA PARA TABLETA
% Las opciones siguientes seleccionan un tamaño de impresión similar a una tableta de 9 pulgadas con márgenes estrechos. Útil para producir una versión en pdf para ser leída en una tableta en lugar de impresa.
% Para que la portada quede centrada correctamente hay que editar el
% archivo `portada.tex` y eliminar el entorno `addmargin`

% \KOMAoptions{fontsize=10pt, paper=19.7104cm:14.7828cm, twoside=false, BCOR=0cm, DIV=14}

% ---------------------------------------------------------------------
%	PAQUETES 
% ---------------------------------------------------------------------

% CODIFICACIÓN E IDIOMA
% ---------------------------------------------------------------------
\usepackage[utf8]{inputenc} 			    % Codificación de caracteres

% Selección del idioma: cargamos por defecto inglés y español (aunque este último es el idioma por defecto para el documento). Cuando queramos cambiar de idioma escribiremos:
% \selectlanguage{english} o \selectlanguage{spanish}

\usepackage[english, spanish, es-nodecimaldot, es-noindentfirst, es-tabla]{babel}
\selectlanguage{spanish}

% Opciones cargadas para el paquete babel:
  % es-nodecimaldot: No cambia el punto decimal por una coma en modo matemático.
  % es-noindentfirst: No sangra los párrafos tras los títulos.
  % es-tabla: cambia el título del entorno `table` de "Cuadro" a "Tabla"

% Otras opciones del paquete spanish-babel:
  \unaccentedoperators % Desactiva los acentos en los operadores matemáticso (p.e. \lim, \max, ...). Eliminar esta opción si queremos que vayan acentuados

% MATEMÁTICAS
% ---------------------------------------------------------------------
\usepackage{amsmath, amsthm, amssymb, etoolbox, thmtools} % Paquetes matemáticas
\usepackage{mathtools}                % Añade mejoras a amsmath
\mathtoolsset{showonlyrefs=true}      % sólo se numeran las ecuaciones que se usan
\usepackage[mathscr]{eucal} 					% Proporciona el comando \mathscr para
                                      % fuentes de tipo manuscrito en modo matemático sin sobreescribir el comando \mathcal

% TIPOGRAFÍA 
% ---------------------------------------------------------------------
% El paquete microtype mejora la tipografía del documento.
\usepackage[activate={true,nocompatibility},final,tracking=true,kerning=true,spacing=true,factor=1100,stretch=10,shrink=10]{microtype}

% Las tipografías elegidas para el documento, siguiendo la guía de estilo de la UGR,
% son las siguientes
% Normal font: 			URW Palladio typeface. 
% Sans-serif font: 	Gill Sans
% Monospace font: 	Inconsolata
\usepackage[LGR, T1]{fontenc}
\usepackage[sc, osf]{mathpazo} \linespread{1.05}         
\usepackage[scaled=.95,type1]{cabin} % sans serif in style of Gill Sans
% Si el paquete cabin da error usar el siguiente comando en su lugar
% \renewcommand{\sfdefault}{iwona} 
\usepackage{inconsolata}
\usepackage{enumitem}
\setitemize{itemsep=4pt,topsep=4pt,parsep=4pt,partopsep=0pt}
\usepackage{dirtree}
\usepackage{caption}

% Selecciona el tipo de fuente para los títulos (capítulo, sección, subsección) del documento.
\setkomafont{disposition}{\sffamily\bfseries}

% Cambia el ancho de la cita. Al inicio de un capítulo podemos usar el comando \dictum[autor]{cita} para añadir una cita famosa de un autor.
\renewcommand{\dictumwidth}{0.45\textwidth} 

\recalctypearea % Necesario tras definir la tipografía a usar.

\usepackage{setspace}
\usepackage{changepage}   % for the adjustwidth environment

% TABLAS, GRÁFICOS Y LISTADOS DE CÓDIGO
% ---------------------------------------------------------------------
\usepackage{booktabs}
% \renewcommand{\arraystretch}{1.5} % Aumenta el espacio vertical entre las filas de un entorno tabular

\usepackage{tikz}
\usepackage{tikz-cd}
\usepackage{float}
\usepackage{rotating}
\usepackage{svg}
\usepackage{xcolor, graphicx}
\usepackage{pdflscape}
\usepackage{stfloats,framed}
\usepackage{booktabs}
\usepackage{mathabx}

\usepackage{tabularx}

\newcolumntype{L}[1]{>{\hsize=#1\hsize\RaggedRight} X}


% Carpeta donde buscar los archivos de imagen por defecto
\graphicspath{{img/}}

% IMAGEN DE LA PORTADA
% Existen varias opciones para la imagen de fondo de la portada del TFG. Todas ellas tienen en logotipo de la universidad de Granada en la cabecera. Las opciones son las siguientes:
% 1. portada-ugr y portada-ugr-color: diseño con marca de agua basada en el logo de la UGR (en escala de grises y color).
% 2. portada-ugr-sencilla y portada-ugr-sencilla-color: portada únicamente con el logotipo de la UGR en la cabecera.
\usepackage{eso-pic}
\newcommand\BackgroundPic{%
	\put(0,-1){%
		\parbox[b][\paperheight]{\paperwidth}{%
			\vfill
			\centering
      % Indicar la imagen de fondo en el siguiente comando
			\includegraphics[width=\paperwidth,height=\paperheight,%
			keepaspectratio]{portada-ugr-color-no-margin-bottom}%
			\vfill
}}}

\usepackage{listings} % Para la inclusión de trozos de código
\lstset{
     literate=%
         {á}{{\'a}}1
         {í}{{\'i}}1
         {é}{{\'e}}1
         {ý}{{\'y}}1
         {ú}{{\'u}}1
         {ó}{{\'o}}1
         {ě}{{\v{e}}}1
         {š}{{\v{s}}}1
         {č}{{\v{c}}}1
         {ř}{{\v{r}}}1
         {ž}{{\v{z}}}1
         {ď}{{\v{d}}}1
         {ť}{{\v{t}}}1
         {ň}{{\v{n}}}1                
         {ů}{{\r{u}}}1
         {Á}{{\'A}}1
         {Í}{{\'I}}1
         {É}{{\'E}}1
         {Ý}{{\'Y}}1
         {Ú}{{\'U}}1
         {Ó}{{\'O}}1
         {Ě}{{\v{E}}}1
         {Š}{{\v{S}}}1
         {Č}{{\v{C}}}1
         {Ř}{{\v{R}}}1
         {Ž}{{\v{Z}}}1
         {Ď}{{\v{D}}}1
         {Ť}{{\v{T}}}1
         {Ň}{{\v{N}}}1                
         {Ů}{{\r{U}}}1    
}
\usepackage{lstautogobble}  % Fix relative indenting
\usepackage{zi4}            % Nice font

\definecolor{bluekeywords}{rgb}{0.13, 0.13, 1}
\definecolor{greencomments}{rgb}{0, 0.5, 0}
\definecolor{redstrings}{rgb}{0.9, 0, 0}
\definecolor{graynumbers}{rgb}{0.5, 0.5, 0.5}

\lstset{
    autogobble,
    columns=fullflexible,
    showspaces=false,
    showtabs=false,
    breaklines=true,
    showstringspaces=false,
    breakatwhitespace=true,
    escapeinside={(*@}{@*)},
    commentstyle=\color{greencomments},
    keywordstyle=\color{bluekeywords},
    stringstyle=\color{redstrings},
    numberstyle=\color{graynumbers},
    basicstyle=\ttfamily\footnotesize,
    frame=l,
    framesep=30pt,
    xleftmargin=30pt,
    tabsize=4,
    captionpos=b,
    basicstyle=\ttfamily,
    numbers=left,
    numberstyle=\footnotesize\color{graynumbers},
    numbersep=10pt,
    escapechar=|,
    numberblanklines=false,
    columns=flexible
}

\renewcommand{\lstlistingname}{Programa}% Listing -> Algorithm
\renewcommand{\lstlistlistingname}{List of \lstlistingname s}% List of Listings -> List of Algorithms

\usepackage{newfloat}
%\DeclareCaptionType[fileext=ext, listname={Lista de problemas}, name=Problema, placement=tbph!]{problema}
%\captionsetup[scheme]{labelsep=period, font={small, bf}, skip=8pt}
%\newfloat{problema}{H}{flt}[chapter]
%\floatname{problema}{Problema}
\DeclareFloatingEnvironment[listname={Índice de problemas},placement=H]{problema}

%\newfloat{tabla}{H}{flt}[chapter]
%\floatname{tabla}{Tabla}
\DeclareFloatingEnvironment[listname={Índice de tablas},placement=H]{tabla}


%\DeclareFloatingEnvironment[
%    fileext=los,
%    listname={Lista de problemas},
%    name=Problema,
%    placement=tbhp!,
%    within=chapter,
%]{problema}

% CABECERAS
% ---------------------------------------------------------------------
% Si queremos modificar las cabeceras del documento podemos usar el paquete
% `scrlayer-scrpage` de KOMA-Script. Consultar la documentación al respecto.
% \usepackage[automark]{scrlayer-scrpage}

% VARIOS
% ---------------------------------------------------------------------

%\usepackage{showkeys}	% Muestra las etiquetas del documento. Útil para revisar las referencias cruzadas.

\usepackage[style=apa,backend=biber,sorting=nty,citestyle=numeric]{biblatex}
\usepackage{hyperref}
\hypersetup{%
  % hidelinks,            % Enlaces sin color ni borde. El borde no se imprime
  linkbordercolor=0.8 0 0,
  citebordercolor=0 0.8 0,
  citebordercolor=0 0.8 0,
  colorlinks = true,            % Color en texto de los enlaces. Comentar esta línea o cambiar `true` por `false` para imprimir el documento.
  linkcolor = [rgb]{0.5, 0, 0}, % Color de los enlaces internos
  urlcolor = [rgb]{0, 0, 0.5},  % Color de los hipervínculos
  citecolor = [rgb]{0, 0.5, 0}, % Color de las referencias bibliográficas
	pdftitle={\miTitulo},%
	pdfauthor={\textcopyright\ \miNombre, \miFacultad, \miUniversidad},%
  pdfsubject={Trabajo de Fin de Grado},%
	pdfkeywords={},%
	pdfcreator={pdfLaTeX},%
}
\addbibresource{library.bib}

\DeclareFieldFormat{labelnumberwidth}{\mkbibbrackets{#1}}

\defbibenvironment{bibliography}
  {\list
     {\printtext[labelnumberwidth]{%
      \printfield{labelprefix}%
      \printfield{labelnumber}}}
     {\setlength{\labelwidth}{\labelnumberwidth}%
      \setlength{\leftmargin}{\labelwidth}%
      \setlength{\labelsep}{\biblabelsep}%
      \addtolength{\leftmargin}{\labelsep}%
      \setlength{\itemsep}{8pt}%
      \setlength{\parsep}{\bibparsep}}%
      \renewcommand*{\makelabel}[1]{\hss##1}}
  {\endlist}
  {\item}


% ÍNDICE 
% Para generar el índice hay que compilar el documento con MakeIndex. Generalmente los editores se encargan de ello automáticamente.
% ----------------------------------------------------------------------
% \index{} para añadir un elemento
% \index{main!sub} para añadir un elementos "sub" bajo la categoría "main".
% \index{termino|textbf} para dar formato al número de página (negrita).
% \index{termino|see{termino relacionado}} para crear una referencia cruzada

% Ejemplo: \index{espacio homogéneo}, \index{superficie!mínima}, \index{esfera|see{espacio homogéneo}}

% Activar los siguientes comandos para generar el índice terminológico. Ver también comandos al final de este documento para incluir dicho índice en el pdf final.
% \usepackage{makeidx}
% \makeindex

% Para revisar las entradas al índice conforme las incluimos en el documento es útil el siguiente paquete. Conviene observar que mientras esté cargado no se generará el índice.
%\usepackage{showidx} % Muestra en el margen del documento las entradas añadidas al índice. Útil para revisar el documento. Si está activo el índice no se genera

% ---------------------------------------------------------------------
% COMANDOS Y ENTORNOS
% ---------------------------------------------------------------------

% DEFINICIÓN DE COMANDOS Y ENTORNOS

% CONJUNTOS DE NÚMEROS

\newcommand{\N}{\mathbb{N}}     % Naturales
\newcommand{\R}{\mathbb{R}}     % Reales
\newcommand{\Z}{\mathbb{Z}}     % Enteros
\newcommand{\Q}{\mathbb{Q}}     % Racionales
\newcommand{\C}{\mathbb{C}}     % Complejos

% TEOREMAS Y ENTORNOS

\newtheorem*{teorema*}{Teorema}
\newtheorem{teorema}{Teorema}[chapter]
\newtheorem{proposicion}{Proposición}[chapter]
\newtheorem{lema}{Lema}[chapter]
\newtheorem{corolario}{Corolario}[chapter]

\theoremstyle{definition}
\newtheorem{definicion}{Definición}[chapter]
\newtheorem{ejemplo}{Ejemplo}[chapter]
\newtheorem{sistemaformal}{Sistema formal}[chapter]
\newtheorem{sistemalogico}[sistemaformal]{Sistema lógico}

\theoremstyle{remark}
\newtheorem{observacion}{Observación}[chapter]
\newtheorem{nota}{Nota}[definicion]

\newcommand\Vtextvisiblespace[1][.3em]{%
\mbox{\kern.06em\vrule height.3ex}%
\vbox{\hrule width#1}%
\hbox{\vrule height.3ex}}

\newcommand{\cabezal}[1]{\overset{\blacktriangledown}{#1}}
\newcommand{\palabra}[1]{\texttt{\textquotesingle{}{#1}\textquotesingle{}}}

\makeatletter
\newenvironment{proofw}{\par
\pushQED{\qed}%
\normalfont \topsep6\p@\@plus6\p@\relax
\trivlist
\item[]\ignorespaces
}{%
\popQED\endtrivlist\@endpefalse
}
\makeatother

\makeatletter
\def\ll@definicion{%
\protect\numberline{\csname the\thmt@envname\endcsname}%
\ifx\@empty\thmt@shortoptarg
  \thmt@thmname
\else
  \thmt@shortoptarg
\fi}
\def\l@thmt@definicion{} 
\makeatother

\makeatletter
\def\ll@sistemaformal{%
\protect\numberline{\csname the\thmt@envname\endcsname}%
\ifx\@empty\thmt@shortoptarg
  \thmt@thmname
\else
  \thmt@shortoptarg
\fi}
\def\l@thmt@sistemaformal{} 
\makeatother

\makeatletter
\def\ll@sistemalogico{%
\protect\numberline{\csname the\thmt@envname\endcsname}%
\ifx\@empty\thmt@shortoptarg
  \thmt@thmname
\else
  \thmt@shortoptarg
\fi}
\def\l@thmt@sistemalogico{} 
\makeatother

% --------------------------------------------------------------------
% INFORMACIÓN DEL TFG Y EL AUTOR
% --------------------------------------------------------------------
\usepackage{xspace} % Para problemas de espaciado al definir comandos

\newcommand{\miBorrador}{BORRADOR 5}
\newcommand{\miTitulo}{Una perspectiva algorítmica a la incompletitud de las Matemáticas\xspace}
\newcommand{\miNombre}{Miguel Ángel Fernández Gutiérrez\xspace}
\newcommand{\miGrado}{Doble Grado en \\ Ingeniería Informática y Matemáticas}
\newcommand{\miFacultad}{Escuela Técnica Superior de \\ Ingenierías Informática y de Telecomunicación \\[4pt] Facultad de Ciencias}
\newcommand{\miFacultadNoSpace}{Escuela Técnica Superior de Ingenierías Informática y de Telecomunicación \\[4pt] Facultad de Ciencias}
\newcommand{\miUniversidad}{Universidad de Granada}
% Añadir tantos tutores como sea necesario separando cada uno de ellos
% mediante el comando `\medskip` y una línea en blanco
\newcommand{\miTutor}{
  Serafín Moral Callejón \\ \emph{Departamento de Ciencias de la Computación e Inteligencia Artificial}
}
\newcommand{\miCurso}{2022/2023\xspace}

% HYPERREFERENCES
% --------------------------------------------------------------------
\usepackage{xurl}
\usepackage[noabbrev,capitalize,nameinlink]{cleveref}

\crefname{section}{sección}{secciones}
\Crefname{section}{sec.}{secs.}
\crefname{chapter}{capítulo}{capítulos}
\Crefname{chapter}{cap.}{caps.}
\crefname{part}{parte}{partes}
\Crefname{part}{parte}{partes}

\crefname{programa}{programa}{programas}
\Crefname{programa}{Programa}{Programas}
\crefname{problema}{problema}{problemas}
\Crefname{problema}{Problema}{Problemas}
\crefname{ejemplo}{ejemplo}{ejemplos}
\Crefname{ejemplo}{Ejemplo}{Ejemplos}
\crefname{proposicion}{proposición}{proposiciones}
\Crefname{proposicion}{prop.}{props.}
\crefname{lema}{lema}{lemas}
\Crefname{lema}{Lema}{Lemas}
\crefname{corolario}{corolario}{corolarios}
\Crefname{corolario}{Corolario}{Corolarios}
\crefname{teorema}{teorema}{teoremas}
\Crefname{teorema}{Teorema}{Teoremas}
\crefname{programa}{programa}{programas}
\Crefname{programa}{Programa}{Programas}
\crefname{definicion}{definición}{definiciones}
\Crefname{definicion}{Definición}{Definiciones}
\crefname{figure}{figura}{figuras}
\Crefname{figure}{Figura}{Figuras}
\crefname{tabla}{tabla}{tablas}
\Crefname{tabla}{Tabla}{Tablas}
\crefname{appendix}{apéndice}{apéndices}
\Crefname{appendix}{Apéndice}{Apéndices}
\crefname{listing}{programa}{programas}
\Crefname{listing}{Programa}{Programas}
\crefname{line}{línea}{líneas}
\Crefname{line}{Línea}{Líneas}
\crefname{page}{página}{páginas}
\Crefname{page}{pág.}{págs.}
\crefname{sistemalogico}{sistema lógico}{sistemas lógicos}
\Crefname{sistemalogico}{Sistema Lógico}{Sistemas Lógicos}
\crefname{sistemaformal}{sistema formal}{sistemas formales}
\Crefname{sistemaformal}{Sistema Formal}{Sistemas Formales}
\newcommand{\crefpairconjunction}{ y }
\newcommand{\crefrangeconjunction}{ a }


%index
\usepackage{makeidx}
\makeindex

\usepackage[firstpageonly=true, angle=58]{draftwatermark}
\SetWatermarkText{\sffamily \miBorrador}
\SetWatermarkScale{1}


\hypersetup{
 pdfauthor={Fernández Gutiérrez, Miguel Ángel},
 pdftitle={[\miBorrador] — \miTitulo},
 pdfsubject={Trabajo de Fin de Grado — \miGrado — Curso \miCurso},
 pdfkeywords={incompletitud, informática, matemáticas, TFG},
 pdfproducer=mianfg,
 pdfcreator=pdflatex,
 breaklinks=true
}

\begin{document}

% --------------------------------------------------------------------
% FRONTMATTER
% -------------------------------------------------------------------
\frontmatter % Desactiva la numeración de capítulos y usa numeración romana para las páginas
%\flushbottom
% \pagestyle{plain} % No imprime cabeceras

% !TeX root = ../libro.tex
% !TeX encoding = utf8

%*******************************************************
% Titlepage
%*******************************************************
\begin{titlepage}
  \AddToShipoutPicture*{\BackgroundPic}
  \phantomsection 
  \pdfbookmark[1]{Portada}{title}

  % Para que el título esté centrado en la página.
  % Los valores numéricos deberán elegirse de acuerdo con el diseño de
  % página (sobre todo si se cambia la opción BCOR o DIV).
  % antes 2.575cm
  \begin{addmargin}[2cm]{0cm}
  \begin{flushleft}
    \Large  
    \hfill\vfil

    \large{\textsf{\miFacultad}}
    \vfill

    {\large\textsc\miGrado} \vfill


    {\large\textsc{trabajo de fin de grado}}

    \begin{flushleft}
      \Huge
      \setstretch{0.8}
      Una perspectiva algorítmica\\a la incompletitud de\\las Matemáticas\\[18pt]
      \huge
      An algorithmic perspective into the incompleteness of Mathematics
    \end{flushleft}

    \vfill\vfill\vfill\vfill

    \textsf{\normalsize{Presentado por:}}\\
    {\normalsize\textrm{\miNombre}} 
    \bigskip

    \textsf{\normalsize{Tutor:}}\\
    {\normalsize\rmfamily\miTutor}

    \bigskip
    \textsf{\normalsize{Curso académico \miCurso}}
    \vspace*{-2cm}
  \end{flushleft}  
  \end{addmargin}       

\end{titlepage}   
\cleardoublepage
\endinput

% !TeX root = ../libro.tex
% !TeX encoding = utf8

%*******************************************************
% Little Dirty Titlepage
%*******************************************************

\thispagestyle{empty}

\begin{center}
  \large  

  \vspace*{\stretch{1}}

  \begingroup
  \huge{\miTitulo} \\[18pt]
  \Large{An algorithmic perspective into the\\ incompleteness of Mathematics}
  \bigskip
  \endgroup

  \textrm{\miNombre}

  \vspace{\stretch{5}}

\end{center}  

\newpage
\thispagestyle{empty}

\hfill

\vfill

\textbf{Alumno}\\
\miNombre
\medskip

\emph{Una perspectiva algorítmica a la incompletitud de las Matemáticas\\(An algorithmic perspective into the incompleteness of Matematics)}\\
Trabajo de Fin de Grado. Curso académico \miCurso
\bigskip

\begin{minipage}[t]{0.45\textwidth} %0.25
  \flushleft
  \textbf{Responsable de tutorización}\\
  \medskip
  \miTutor
\end{minipage}
\begin{minipage}[t]{0.10\textwidth} %0.25
  \hspace{1cm}
\end{minipage}
%\begin{minipage}[t]{0.45\textwidth}
%  \flushleft
%  \miTutor
%\end{minipage}
\begin{minipage}[t]{0.45\textwidth} %0.30
  \flushright
  \textbf{\miGrado}
  \medskip

  \miFacultadNoSpace
  \medskip

  \miUniversidad
\end{minipage}

\newpage
\endinput

% !TeX root = ../libro.tex
% !TeX encoding = utf8
%
%*******************************************************
% Declaración de originalidad
%*******************************************************

\thispagestyle{empty}

\hfill\vfill

%\textsc{Declaración de originalidad}\\\bigskip

%D. \miNombre \\\medskip

Declaro explícitamente que el trabajo presentado como Trabajo de Fin de Grado (TFG), correspondiente al curso académico \miCurso, es original, entendida esta, en el sentido de que no he utilizado para la elaboración del trabajo fuentes sin citarlas debidamente.
\medskip

En Granada, a \today 
\begin{flushleft} 
Fdo: \miNombre 

\end{flushleft}

\vfill

\cleardoublepage
\endinput

%% !TeX root = ../libro.tex
% !TeX encoding = utf8

%*******************************************************
% Dedication
%*******************************************************
\thispagestyle{empty}
\phantomsection 
%\pdfbookmark[1]{Dedicatoria}{Dedicatoria}

\hfill
\vfill

\begin{flushright}
\itshape
A Sonia, por acompañarme\\
en este año tan diferente
\end{flushright}

\vfill
\vfill

\cleardoublepage
\endinput
                % Opcional
% !TeX root = ../libro.tex
% !TeX encoding = utf8

%*******************************************************
% Table of Contents
%*******************************************************
\phantomsection
\pdfbookmark[0]{\contentsname}{toc}

\setcounter{tocdepth}{2} % <-- 2 includes up to subsections in the ToC
\setcounter{secnumdepth}{3} % <-- 3 numbers up to subsubsections

% \manualmark
% \markboth{\textsc{\contentsname}}{\textsc{\contentsname}}
\tableofcontents 

%*******************************************************
% List of Figures and of the Tables
%*******************************************************

    % *******************************************************
    %  List of Figures
    % *******************************************************    
    \phantomsection 
    % \listoffigures

    %*******************************************************
    % List of Tables
    %*******************************************************
    \phantomsection 
    % \listoftables
    
    %*******************************************************
    % List of Listings
    % The package \usepackage{listings} is needed
    %*******************************************************      
	  % \phantomsection 
    % \renewcommand{\lstlistlistingname}{Listados de código}
    % \lstlistoflistings 

\cleardoublepage


%% !TeX root = ../libro.tex
% !TeX encoding = utf8

%*******************************************************
% Agradecimientos
%*******************************************************

\chapter{Agradecimientos}

Agradecimientos del libro (opcional, ver archivo \texttt{preliminares/agradecimiento.tex}).

\cleardoublepage
\endinput
            % Opcional

% \pagestyle{scrheadings} % A partir de ahora sí imprime cabeceras

% !TeX root = ../libro.tex
% !TeX encoding = utf8
%
%*******************************************************
% Summary
%*******************************************************

\selectlanguage{english}
\chapter{Summary}

%An english summary of the project (around 800 and 1500 words are recommended).

%File: \texttt{preliminares/summary.tex}

We begin this Bachelor's Thesis with \hyperref[ch:motivacion]{chapter 1}, which highlights the motivation behind this work, as well as the main goals to be achieved. Immediately after, in \hyperref[ch:historia]{chapter 2}, we give a brief historical contextualization of the topics covered in this writing.

The theoretical development starts in \hyperref[ch:programas-python-maquinas-turing]{chapter 3}, whose \hyperref[sec:programas-python]{section 3.1} introduces a specific subset of Python programs, namely \emph{SISO programs}, and we observe this class is as expressive as the class of all Python programs. Next, in \hyperref[sec:maquinas-turing]{section 3.2}, we delve into \emph{Turing machines}, the main mathematical tool of our analysis. In \hyperref[sec:equivalencia]{section 3.3} we introduce new and equivalent computing models, until we arrive at the main result of this chapter: the fact that Turing machines can emulate Python programs, and vice versa; in other words, Python programs and Turing machines are equivalent. This will allow us to use Python programs conveniently, instead of having to specify precise descriptions of Turing machines. We conclude the chapter introducing the Church-Turing thesis in \hyperref[sec:church-turing]{section 3.4}, and realizing that the equivalence result proved can be an argument in its favor.

We define mathematically in \hyperref[sec:problemas-decidibles]{section 4.1} the notion of \emph{computational problem}, and we discuss Python programs (and, equivalently, Turing machines), can solve them. We call \emph{decidable} any problem that can be solved by a program that never halts. In \hyperref[sec:universalidad]{section 4.2} we observe a Python program can emulate itself, which leads us in \hyperref[sec:problemas-no-decidibles]{section 4.3} to the existence of undecidable problems. In \hyperref[sec:problemas-semidecidibles]{section 4.4} we relax the halting condition to obtain \emph{semidecidable} problems. Afterwards, we introduce in \hyperref[sec:reducciones]{section 4.5} the \emph{reduction} technique, which allows us to obtain new, undecidable problems, including the most important one for this thesis in \hyperref[sec:problema-parada]{section 4.6}: the \emph{halting problem}.

In order to prove the main result of this work, it is necessary to introduce a formalization of the concepts of \emph{proof} and \emph{truth}. This is precisely what we do when introducing \emph{formal systems} \hyperref[sec:sistemas-formales]{section 4.1} and \emph{logical systems} in \hyperref[sec:sistemas-logicos]{section 4.2}. Formal systems are \emph{syntactic}: they formalize what a \emph{theorem} is, a provable \emph{formula} from a set of \emph{axioms} and \emph{inference rules}. On the other hand, logical systems incorporate to formal systems a semantic component: \emph{truth}. In \hyperref[sec:solidez-completitud-decidibilidad]{sections 5.3} and \hyperref[sec:consistencia]{5.5} we present essential properties of the aforementioned systems, which relate semantic and syntactic concepts. We highlight in \hyperref[sec:aritmetica-peano]{section 5.4} the introduction of \emph{Peano arithmetic}: a logical system of especial significance.

As colophon, we proceed to prove the First Incompleteness Theorem in \hyperref[ch:teorema-incompletitud]{chapter 6}: we do so first from Peano arithmetic in \hyperref[sec:primera-aproximacion]{section 6.1}, and then we generalize it by means of two versions. The first one, of a semantic nature, is proven in \hyperref[sec:version-semantica]{section 6.2}, and allows us (under some computational hypotheses) to find true but unprovable statements in a logical system. Next, the syntactic version of \hyperref[sec:version-sintactica]{section 6.3}, closer to Gödel's, allows us to find unprovable statements in formal systems. We finish the study of this theorem in \hyperref[sec:consecuencias]{section 6.4}, with some important mathematical (\hyperref[subsec:segundo-problema-hilbert]{section 6.4.1}) and philosophical (\hyperref[subsec:mente-humana-tesis-church-turing]{section 6.4.2}) consequences.

Finally, we wrap up with some conclusions in \hyperref[ch:conclusiones]{chapter 7}, and some final words for further work in \hyperref[ch:trabajo-futuro]{chapter 8}.

%\vfill
% ====================
\begin{sidewaysfigure}
\centering
%\vspace{8pt}
%\hspace*{-1cm}


\tikzset{every picture/.style={line width=0.75pt}} %set default line width to 0.75pt        

\begin{tikzpicture}[x=0.75pt,y=0.75pt,yscale=-1,xscale=1]
%uncomment if require: \path (0,694); %set diagram left start at 0, and has height of 694

%Straight Lines [id:da2056540562354816] 
\draw  [dash pattern={on 4.5pt off 4.5pt}]  (672,64.47) -- (810,64.47) ;
%Straight Lines [id:da23044041717189923] 
\draw  [dash pattern={on 4.5pt off 4.5pt}]  (412.67,64.47) -- (610,64.47) ;
\draw [shift={(410.67,64.47)}, rotate = 0] [color={rgb, 255:red, 0; green, 0; blue, 0 }  ][line width=0.75]    (10.93,-3.29) .. controls (6.95,-1.4) and (3.31,-0.3) .. (0,0) .. controls (3.31,0.3) and (6.95,1.4) .. (10.93,3.29)   ;
%Straight Lines [id:da8398435684999113] 
\draw  [dash pattern={on 4.5pt off 4.5pt}]  (810,64.47) -- (810,530) ;
%Rounded Rect [id:dp027018508739026448] 
\draw  [draw opacity=0][fill={rgb, 255:red, 230; green, 230; blue, 230 }  ,fill opacity=0.5 ] (20,439.33) .. controls (20,445.22) and (24.78,450) .. (30.67,450) -- (409.33,450) .. controls (415.22,450) and (420,445.22) .. (420,439.33) -- (420,140.67) .. controls (420,134.78) and (415.22,130) .. (409.33,130) -- (30.67,130) .. controls (24.78,130) and (20,134.78) .. (20,140.67) -- cycle ;
%Rounded Rect [id:dp9415679498487528] 
\draw  [dash pattern={on 4.5pt off 4.5pt}] (29.83,292.77) .. controls (29.83,285.88) and (35.41,280.3) .. (42.3,280.3) -- (398.54,280.3) .. controls (405.42,280.3) and (411,285.88) .. (411,292.77) -- (411,337.17) .. controls (411,344.05) and (405.42,349.64) .. (398.54,349.64) -- (42.3,349.64) .. controls (35.41,349.64) and (29.83,344.05) .. (29.83,337.17) -- cycle ;
%Rounded Rect [id:dp05100022295548956] 
\draw  [draw opacity=0][fill={rgb, 255:red, 230; green, 230; blue, 230 }  ,fill opacity=0.5 ] (130.5,575) .. controls (130.5,583.28) and (137.22,590) .. (145.5,590) -- (615,590) .. controls (623.28,590) and (630,583.28) .. (630,575) -- (630,395) .. controls (630,386.72) and (623.28,380) .. (615,380) -- (145.5,380) .. controls (137.22,380) and (130.5,386.72) .. (130.5,395) -- cycle ;
%Rounded Rect [id:dp6496809750851564] 
\draw  [draw opacity=0][fill={rgb, 255:red, 230; green, 230; blue, 230 }  ,fill opacity=0.5 ] (450.61,440.79) .. controls (450.61,445.82) and (454.69,449.9) .. (459.72,449.9) -- (780.89,449.9) .. controls (785.92,449.9) and (790,445.82) .. (790,440.79) -- (790,139.11) .. controls (790,134.08) and (785.92,130) .. (780.89,130) -- (459.72,130) .. controls (454.69,130) and (450.61,134.08) .. (450.61,139.11) -- cycle ;
%Rounded Rect [id:dp549636265291094] 
\draw  [fill={rgb, 255:red, 230; green, 230; blue, 230 }  ,fill opacity=0.5 ][dash pattern={on 4.5pt off 4.5pt}] (29.83,42.77) .. controls (29.83,35.88) and (35.41,30.3) .. (42.3,30.3) -- (398.54,30.3) .. controls (405.42,30.3) and (411,35.88) .. (411,42.77) -- (411,87.17) .. controls (411,94.05) and (405.42,99.64) .. (398.54,99.64) -- (42.3,99.64) .. controls (35.41,99.64) and (29.83,94.05) .. (29.83,87.17) -- cycle ;
%Rounded Rect [id:dp1378581119146769] 
\draw  [fill={rgb, 255:red, 255; green, 255; blue, 255 }  ,fill opacity=1 ] (40.4,49.82) .. controls (40.4,44.29) and (44.89,39.8) .. (50.42,39.8) -- (169.58,39.8) .. controls (175.11,39.8) and (179.6,44.29) .. (179.6,49.82) -- (179.6,79.88) .. controls (179.6,85.41) and (175.11,89.9) .. (169.58,89.9) -- (50.42,89.9) .. controls (44.89,89.9) and (40.4,85.41) .. (40.4,79.88) -- cycle ;
%Straight Lines [id:da8729153203483555] 
\draw    (188.42,63.77) -- (250.99,63.77)(188.42,66.77) -- (250.99,66.77) ;
\draw [shift={(257.99,65.27)}, rotate = 180] [color={rgb, 255:red, 0; green, 0; blue, 0 }  ][line width=0.75]    (10.93,-4.9) .. controls (6.95,-2.3) and (3.31,-0.67) .. (0,0) .. controls (3.31,0.67) and (6.95,2.3) .. (10.93,4.9)   ;
\draw [shift={(181.42,65.27)}, rotate = 0] [color={rgb, 255:red, 0; green, 0; blue, 0 }  ][line width=0.75]    (10.93,-4.9) .. controls (6.95,-2.3) and (3.31,-0.67) .. (0,0) .. controls (3.31,0.67) and (6.95,2.3) .. (10.93,4.9)   ;
%Straight Lines [id:da845243590120559] 
\draw    (220.17,99.97) -- (220.17,137.64) ;
\draw [shift={(220.17,139.64)}, rotate = 270] [color={rgb, 255:red, 0; green, 0; blue, 0 }  ][line width=0.75]    (10.93,-3.29) .. controls (6.95,-1.4) and (3.31,-0.3) .. (0,0) .. controls (3.31,0.3) and (6.95,1.4) .. (10.93,3.29)   ;
%Straight Lines [id:da7497001546502899] 
\draw    (380,420.33) -- (380,478) ;
\draw [shift={(380,480)}, rotate = 270] [color={rgb, 255:red, 0; green, 0; blue, 0 }  ][line width=0.75]    (10.93,-3.29) .. controls (6.95,-1.4) and (3.31,-0.3) .. (0,0) .. controls (3.31,0.3) and (6.95,1.4) .. (10.93,3.29)   ;
%Straight Lines [id:da5689893955334542] 
\draw    (300,420) -- (460,420) ;
%Rounded Rect [id:dp18582062333838767] 
\draw  [fill={rgb, 255:red, 255; green, 255; blue, 255 }  ,fill opacity=1 ] (260.4,49.82) .. controls (260.4,44.29) and (264.89,39.8) .. (270.42,39.8) -- (389.58,39.8) .. controls (395.11,39.8) and (399.6,44.29) .. (399.6,49.82) -- (399.6,79.88) .. controls (399.6,85.41) and (395.11,89.9) .. (389.58,89.9) -- (270.42,89.9) .. controls (264.89,89.9) and (260.4,85.41) .. (260.4,79.88) -- cycle ;
%Rounded Rect [id:dp9523169763067778] 
\draw  [fill={rgb, 255:red, 255; green, 255; blue, 255 }  ,fill opacity=1 ] (470,299.92) .. controls (470,294.39) and (474.49,289.9) .. (480.02,289.9) -- (599.18,289.9) .. controls (604.71,289.9) and (609.2,294.39) .. (609.2,299.92) -- (609.2,329.98) .. controls (609.2,335.51) and (604.71,340) .. (599.18,340) -- (480.02,340) .. controls (474.49,340) and (470,335.51) .. (470,329.98) -- cycle ;
%Rounded Rect [id:dp48544030815291705] 
\draw  [fill={rgb, 255:red, 255; green, 255; blue, 255 }  ,fill opacity=1 ] (220,490) .. controls (220,484.48) and (224.48,480) .. (230,480) -- (530,480) .. controls (535.52,480) and (540,484.48) .. (540,490) -- (540,570) .. controls (540,575.52) and (535.52,580) .. (530,580) -- (230,580) .. controls (224.48,580) and (220,575.52) .. (220,570) -- cycle ;
%Rounded Rect [id:dp08222627508455749] 
\draw  [fill={rgb, 255:red, 255; green, 255; blue, 255 }  ,fill opacity=1 ] (540.4,49.82) .. controls (540.4,44.29) and (544.89,39.8) .. (550.42,39.8) -- (690.65,39.8) .. controls (696.18,39.8) and (700.67,44.29) .. (700.67,49.82) -- (700.67,79.88) .. controls (700.67,85.41) and (696.18,89.9) .. (690.65,89.9) -- (550.42,89.9) .. controls (544.89,89.9) and (540.4,85.41) .. (540.4,79.88) -- cycle ;
%Straight Lines [id:da5868346364154593] 
\draw  [dash pattern={on 4.5pt off 4.5pt}]  (810,530) -- (542,530) ;
\draw [shift={(540,530)}, rotate = 360] [color={rgb, 255:red, 0; green, 0; blue, 0 }  ][line width=0.75]    (10.93,-3.29) .. controls (6.95,-1.4) and (3.31,-0.3) .. (0,0) .. controls (3.31,0.3) and (6.95,1.4) .. (10.93,3.29)   ;
%Rounded Rect [id:dp7392784471609453] 
\draw  [fill={rgb, 255:red, 255; green, 255; blue, 255 }  ,fill opacity=1 ] (30.33,150.47) .. controls (30.33,144.58) and (35.11,139.8) .. (41,139.8) -- (400.34,139.8) .. controls (406.23,139.8) and (411,144.58) .. (411,150.47) -- (411,229.34) .. controls (411,235.23) and (406.23,240.01) .. (400.34,240.01) -- (41,240.01) .. controls (35.11,240.01) and (30.33,235.23) .. (30.33,229.34) -- cycle ;
%Rounded Rect [id:dp90901637181607] 
\draw  [fill={rgb, 255:red, 255; green, 255; blue, 255 }  ,fill opacity=1 ] (41,189.82) .. controls (41,184.29) and (45.49,179.8) .. (51.02,179.8) -- (145.02,179.8) .. controls (150.55,179.8) and (155.04,184.29) .. (155.04,189.82) -- (155.04,219.88) .. controls (155.04,225.41) and (150.55,229.9) .. (145.02,229.9) -- (51.02,229.9) .. controls (45.49,229.9) and (41,225.41) .. (41,219.88) -- cycle ;
%Rounded Rect [id:dp06115018179625542] 
\draw  [fill={rgb, 255:red, 255; green, 255; blue, 255 }  ,fill opacity=1 ] (163.65,189.82) .. controls (163.65,184.29) and (168.13,179.8) .. (173.67,179.8) -- (267.67,179.8) .. controls (273.2,179.8) and (277.69,184.29) .. (277.69,189.82) -- (277.69,219.88) .. controls (277.69,225.41) and (273.2,229.9) .. (267.67,229.9) -- (173.67,229.9) .. controls (168.13,229.9) and (163.65,225.41) .. (163.65,219.88) -- cycle ;
%Rounded Rect [id:dp6195275548153432] 
\draw  [fill={rgb, 255:red, 255; green, 255; blue, 255 }  ,fill opacity=1 ] (286.3,189.82) .. controls (286.3,184.29) and (290.78,179.8) .. (296.32,179.8) -- (390.31,179.8) .. controls (395.85,179.8) and (400.33,184.29) .. (400.33,189.82) -- (400.33,219.88) .. controls (400.33,225.41) and (395.85,229.9) .. (390.31,229.9) -- (296.32,229.9) .. controls (290.78,229.9) and (286.3,225.41) .. (286.3,219.88) -- cycle ;
%Rounded Rect [id:dp41283415100159493] 
\draw  [fill={rgb, 255:red, 255; green, 255; blue, 255 }  ,fill opacity=1 ] (40.8,299.92) .. controls (40.8,294.39) and (45.29,289.9) .. (50.82,289.9) -- (169.98,289.9) .. controls (175.51,289.9) and (180,294.39) .. (180,299.92) -- (180,329.98) .. controls (180,335.51) and (175.51,340) .. (169.98,340) -- (50.82,340) .. controls (45.29,340) and (40.8,335.51) .. (40.8,329.98) -- cycle ;
%Rounded Rect [id:dp40154854661776396] 
\draw  [fill={rgb, 255:red, 255; green, 255; blue, 255 }  ,fill opacity=1 ] (260.4,299.82) .. controls (260.4,294.29) and (264.89,289.8) .. (270.42,289.8) -- (389.58,289.8) .. controls (395.11,289.8) and (399.6,294.29) .. (399.6,299.82) -- (399.6,329.88) .. controls (399.6,335.41) and (395.11,339.9) .. (389.58,339.9) -- (270.42,339.9) .. controls (264.89,339.9) and (260.4,335.41) .. (260.4,329.88) -- cycle ;
%Straight Lines [id:da6602630444249074] 
\draw    (180.17,315.8) -- (260,315.8) ;
%Straight Lines [id:da30905866347109745] 
\draw    (219.86,315.97) -- (219.86,386.81) ;
\draw [shift={(219.86,388.81)}, rotate = 270] [color={rgb, 255:red, 0; green, 0; blue, 0 }  ][line width=0.75]    (10.93,-3.29) .. controls (6.95,-1.4) and (3.31,-0.3) .. (0,0) .. controls (3.31,0.3) and (6.95,1.4) .. (10.93,3.29)   ;
%Rounded Rect [id:dp44627062039563214] 
\draw  [fill={rgb, 255:red, 255; green, 255; blue, 255 }  ,fill opacity=1 ] (139.67,399.82) .. controls (139.67,394.29) and (144.15,389.8) .. (149.69,389.8) -- (290.31,389.8) .. controls (295.85,389.8) and (300.33,394.29) .. (300.33,399.82) -- (300.33,429.88) .. controls (300.33,435.41) and (295.85,439.9) .. (290.31,439.9) -- (149.69,439.9) .. controls (144.15,439.9) and (139.67,435.41) .. (139.67,429.88) -- cycle ;
%Straight Lines [id:da8375870493480149] 
\draw    (220,240) -- (220,277.67) ;
\draw [shift={(220,279.67)}, rotate = 270] [color={rgb, 255:red, 0; green, 0; blue, 0 }  ][line width=0.75]    (10.93,-3.29) .. controls (6.95,-1.4) and (3.31,-0.3) .. (0,0) .. controls (3.31,0.3) and (6.95,1.4) .. (10.93,3.29)   ;
%Rounded Rect [id:dp37285624114093485] 
\draw  [fill={rgb, 255:red, 255; green, 255; blue, 255 }  ,fill opacity=1 ] (460,400.02) .. controls (460,394.49) and (464.49,390) .. (470.02,390) -- (610.65,390) .. controls (616.18,390) and (620.67,394.49) .. (620.67,400.02) -- (620.67,430.08) .. controls (620.67,435.61) and (616.18,440.1) .. (610.65,440.1) -- (470.02,440.1) .. controls (464.49,440.1) and (460,435.61) .. (460,430.08) -- cycle ;
%Rounded Rect [id:dp6883616335272991] 
\draw  [fill={rgb, 255:red, 255; green, 255; blue, 255 }  ,fill opacity=1 ] (629.2,299.92) .. controls (629.2,294.39) and (633.69,289.9) .. (639.22,289.9) -- (758.38,289.9) .. controls (763.91,289.9) and (768.4,294.39) .. (768.4,299.92) -- (768.4,329.98) .. controls (768.4,335.51) and (763.91,340) .. (758.38,340) -- (639.22,340) .. controls (633.69,340) and (629.2,335.51) .. (629.2,329.98) -- cycle ;
%Straight Lines [id:da19870419404011752] 
\draw    (540,340) -- (540,388) ;
\draw [shift={(540,390)}, rotate = 270] [color={rgb, 255:red, 0; green, 0; blue, 0 }  ][line width=0.75]    (10.93,-3.29) .. controls (6.95,-1.4) and (3.31,-0.3) .. (0,0) .. controls (3.31,0.3) and (6.95,1.4) .. (10.93,3.29)   ;
%Rounded Rect [id:dp9692784599420237] 
\draw  [fill={rgb, 255:red, 255; green, 255; blue, 255 }  ,fill opacity=1 ] (230,529.92) .. controls (230,524.39) and (234.49,519.9) .. (240.02,519.9) -- (359.18,519.9) .. controls (364.71,519.9) and (369.2,524.39) .. (369.2,529.92) -- (369.2,559.98) .. controls (369.2,565.51) and (364.71,570) .. (359.18,570) -- (240.02,570) .. controls (234.49,570) and (230,565.51) .. (230,559.98) -- cycle ;
%Rounded Rect [id:dp749684165723072] 
\draw  [fill={rgb, 255:red, 255; green, 255; blue, 255 }  ,fill opacity=1 ] (390,530.02) .. controls (390,524.49) and (394.49,520) .. (400.02,520) -- (519.18,520) .. controls (524.71,520) and (529.2,524.49) .. (529.2,530.02) -- (529.2,560.08) .. controls (529.2,565.61) and (524.71,570.1) .. (519.18,570.1) -- (400.02,570.1) .. controls (394.49,570.1) and (390,565.61) .. (390,560.08) -- cycle ;
%Rounded Rect [id:dp28001422897762485] 
\draw  [fill={rgb, 255:red, 255; green, 255; blue, 255 }  ,fill opacity=1 ][dash pattern={on 4.5pt off 4.5pt}] (460,148.92) .. controls (460,143.99) and (463.99,140) .. (468.92,140) -- (771.08,140) .. controls (776.01,140) and (780,143.99) .. (780,148.92) -- (780,241.55) .. controls (780,246.48) and (776.01,250.47) .. (771.08,250.47) -- (468.92,250.47) .. controls (463.99,250.47) and (460,246.48) .. (460,241.55) -- cycle ;
%Shape: Brace [id:dp9411387904179627] 
\draw   (774.67,286.8) .. controls (774.67,282.13) and (772.34,279.8) .. (767.67,279.8) -- (629.96,279.8) .. controls (623.29,279.8) and (619.96,277.47) .. (619.96,272.8) .. controls (619.96,277.47) and (616.63,279.8) .. (609.96,279.8)(612.96,279.8) -- (472.25,279.8) .. controls (467.58,279.8) and (465.25,282.13) .. (465.25,286.8) ;
%Straight Lines [id:da9902399507274253] 
\draw    (620,273) -- (620,253.14) ;
\draw [shift={(620,251.14)}, rotate = 90] [color={rgb, 255:red, 0; green, 0; blue, 0 }  ][line width=0.75]    (10.93,-3.29) .. controls (6.95,-1.4) and (3.31,-0.3) .. (0,0) .. controls (3.31,0.3) and (6.95,1.4) .. (10.93,3.29)   ;

% Text Node
\draw (110,64.85) node  [color={rgb, 255:red, 0; green, 0; blue, 0 }  ,opacity=1 ] [align=left] {\begin{minipage}[lt]{92.93pt}\setlength\topsep{0pt}
\begin{center}
Python program\\
\small{(\hyperref[sec:programas-python]{Section 3.1})}
\end{center}

\end{minipage}};
% Text Node
\draw (219.89,51.97) node  [font=\small,color={rgb, 255:red, 0; green, 0; blue, 0 }  ,opacity=1 ] [align=left] {\begin{minipage}[lt]{54.86pt}\setlength\topsep{0pt}
\begin{center}
\hyperref[sec:equivalencia]{Sect. 3.3}
\end{center}

\end{minipage}};
% Text Node
\draw (69.5,15.5) node  [font=\small,color={rgb, 255:red, 0; green, 0; blue, 0 }  ,opacity=1 ] [align=left] {\begin{minipage}[lt]{41.55pt}\setlength\topsep{0pt}
\hyperref[ch:programas-python-maquinas-turing]{Chapter 3}
\end{minipage}};
% Text Node
\draw (330,64.85) node  [color={rgb, 255:red, 0; green, 0; blue, 0 }  ,opacity=1 ] [align=left] {\begin{minipage}[lt]{92.93pt}\setlength\topsep{0pt}
\begin{center}
Turing machine\\
\small{(\hyperref[sec:maquinas-turing]{Section 3.2})}
\end{center}

\end{minipage}};
% Text Node
\draw (539.6,314.95) node  [color={rgb, 255:red, 0; green, 0; blue, 0 }  ,opacity=1 ] [align=left] {\begin{minipage}[lt]{92.93pt}\setlength\topsep{0pt}
\begin{center}
logical system\\
\small{(\hyperref[sec:sistemas-logicos]{Section 5.2})}
\end{center}

\end{minipage}};
% Text Node
\draw (380.5,500.11) node  [color={rgb, 255:red, 0; green, 0; blue, 0 }  ,opacity=1 ] [align=left] {\begin{minipage}[lt]{159.55pt}\setlength\topsep{0pt}
\begin{center}
\textbf{First Incompleteness Theorem}
\end{center}

\end{minipage}};
% Text Node
\draw (620.53,64.85) node  [color={rgb, 255:red, 0; green, 0; blue, 0 }  ,opacity=1 ] [align=left] {\begin{minipage}[lt]{106.99pt}\setlength\topsep{0pt}
\begin{center}
Church-Turing thesis\\
\small{(\hyperref[sec:church-turing]{Section 3.4})}
\end{center}

\end{minipage}};
% Text Node
\draw (480.28,52.8) node  [font=\small,color={rgb, 255:red, 0; green, 0; blue, 0 }  ,opacity=1 ] [align=left] {\begin{minipage}[lt]{81.9pt}\setlength\topsep{0pt}
\begin{center}
\hyperref[sec:church-turing]{Section 3.4}
\end{center}

\end{minipage}};
% Text Node
\draw (589.5,518.3) node  [font=\small,color={rgb, 255:red, 0; green, 0; blue, 0 }  ,opacity=1 ] [align=left] {\begin{minipage}[lt]{54.86pt}\setlength\topsep{0pt}
\begin{center}
\hyperref[subsec:mente-humana-tesis-church-turing]{Sect. 6.4.2}
\end{center}

\end{minipage}};
% Text Node
\draw (220.2,159.62) node  [color={rgb, 255:red, 0; green, 0; blue, 0 }  ,opacity=1 ] [align=left] {\begin{minipage}[lt]{230.11pt}\setlength\topsep{0pt}
\begin{center}
computational problem
\end{center}

\end{minipage}};
% Text Node
\draw (98.02,204.85) node  [color={rgb, 255:red, 0; green, 0; blue, 0 }  ,opacity=1 ] [align=left] {\begin{minipage}[lt]{76.13pt}\setlength\topsep{0pt}
\begin{center}
decidable\\
\small{(\hyperref[sec:problemas-decidibles]{Section 4.1})}
\end{center}

\end{minipage}};
% Text Node
\draw (220.67,204.85) node  [color={rgb, 255:red, 0; green, 0; blue, 0 }  ,opacity=1 ] [align=left] {\begin{minipage}[lt]{76.13pt}\setlength\topsep{0pt}
\begin{center}
undecidable\\
\small{(\hyperref[sec:problemas-no-decidibles]{Section 4.3})}
\end{center}

\end{minipage}};
% Text Node
\draw (343.31,204.85) node  [color={rgb, 255:red, 0; green, 0; blue, 0 }  ,opacity=1 ] [align=left] {\begin{minipage}[lt]{76.13pt}\setlength\topsep{0pt}
\begin{center}
semidecidable\\
\small{(\hyperref[sec:problemas-semidecidibles]{Section 4.4})}
\end{center}

\end{minipage}};
% Text Node
\draw (110.4,314.95) node  [color={rgb, 255:red, 0; green, 0; blue, 0 }  ,opacity=1 ] [align=left] {\begin{minipage}[lt]{92.93pt}\setlength\topsep{0pt}
\begin{center}
universality\\
\small{(\hyperref[sec:universalidad]{Section 4.2})}
\end{center}

\end{minipage}};
% Text Node
\draw (330,314.85) node  [color={rgb, 255:red, 0; green, 0; blue, 0 }  ,opacity=1 ] [align=left] {\begin{minipage}[lt]{92.93pt}\setlength\topsep{0pt}
\begin{center}
reduction\\
\small{(\hyperref[sec:reducciones]{Section 4.5})}
\end{center}

\end{minipage}};
% Text Node
\draw (220,414.85) node  [color={rgb, 255:red, 0; green, 0; blue, 0 }  ,opacity=1 ] [align=left] {\begin{minipage}[lt]{107.26pt}\setlength\topsep{0pt}
\begin{center}
halting problem\\
\small{(\hyperref[sec:problema-parada]{Section 4.6})}
\end{center}

\end{minipage}};
% Text Node
\draw (540.33,415.05) node  [color={rgb, 255:red, 0; green, 0; blue, 0 }  ,opacity=1 ] [align=left] {\begin{minipage}[lt]{107.26pt}\setlength\topsep{0pt}
\begin{center}
Peano arithmetic\\
\small{(\hyperref[sec:aritmetica-peano]{Section 5.4})}
\end{center}

\end{minipage}};
% Text Node
\draw (698.8,314.95) node  [color={rgb, 255:red, 0; green, 0; blue, 0 }  ,opacity=1 ] [align=left] {\begin{minipage}[lt]{92.93pt}\setlength\topsep{0pt}
\begin{center}
formal system\\
\small{(\hyperref[sec:sistemas-formales]{Section 5.1})}
\end{center}

\end{minipage}};
% Text Node
\draw (299.6,544.95) node  [color={rgb, 255:red, 0; green, 0; blue, 0 }  ,opacity=1 ] [align=left] {\begin{minipage}[lt]{92.93pt}\setlength\topsep{0pt}
\begin{center}
semantic version\\
\small{(\hyperref[sec:version-semantica]{Section 6.2})}
\end{center}

\end{minipage}};
% Text Node
\draw (459.6,545.05) node  [color={rgb, 255:red, 0; green, 0; blue, 0 }  ,opacity=1 ] [align=left] {\begin{minipage}[lt]{92.93pt}\setlength\topsep{0pt}
\begin{center}
syntactic version\\
\small{(\hyperref[sec:version-sintactica]{Section 6.3})}
\end{center}

\end{minipage}};
% Text Node
\draw (60.5,464.5) node  [font=\small,color={rgb, 255:red, 0; green, 0; blue, 0 }  ,opacity=1 ] [align=left] {\begin{minipage}[lt]{41.55pt}\setlength\topsep{0pt}
\hyperref[ch:problema-parada]{Chapter 4}
\end{minipage}};
% Text Node
\draw (749.5,465.5) node  [font=\small,color={rgb, 255:red, 0; green, 0; blue, 0 }  ,opacity=1 ] [align=left] {\begin{minipage}[lt]{41.55pt}\setlength\topsep{0pt}
\begin{flushright}
\hyperref[ch:sistemas-logicos]{Chapter 5}
\end{flushright}

\end{minipage}};
% Text Node
\draw (170.5,604.5) node  [font=\small,color={rgb, 255:red, 0; green, 0; blue, 0 }  ,opacity=1 ] [align=left] {\begin{minipage}[lt]{41.55pt}\setlength\topsep{0pt}
\hyperref[ch:teorema-incompletitud]{Chapter 6}
\end{minipage}};
% Text Node
\draw (658.76,195.55) node  [color={rgb, 255:red, 0; green, 0; blue, 0 }  ,opacity=1 ] [align=left] {\begin{minipage}[lt]{217.92pt}\setlength\topsep{0pt}
-- sound {\small{(\hyperref[sec:solidez-completitud-decidibilidad]{Section 5.3})}}\\
-- decidable {\small{(\hyperref[sec:solidez-completitud-decidibilidad]{Section 5.3})}}\\
-- (semantically) complete {\small{(\hyperref[sec:solidez-completitud-decidibilidad]{Section 5.3})}}\\
-- syntactically complete {\small{(\hyperref[sec:consistencia]{Section 5.5})}}\\
-- sem./synt. consistent {\small{(\hyperref[sec:consistencia]{Section 5.5})}}
\end{minipage}};
% Text Node
\draw (414.84,645.25) node  [color={rgb, 255:red, 0; green, 0; blue, 0 }  ,opacity=1 ] [align=left] {\begin{minipage}[lt]{534.7pt}\setlength\topsep{0pt}
\begin{center}
Summary map of this Bachelor's Thesis
\end{center}
\end{minipage}};
% Text Node
\draw (380.15,407.8) node  [font=\small,color={rgb, 255:red, 0; green, 0; blue, 0 }  ,opacity=1 ] [align=left] {\begin{minipage}[lt]{108.93pt}\setlength\topsep{0pt}
\begin{center}
\hyperref[sec:primera-aproximacion]{Section 6.1}
\end{center}

\end{minipage}};

\end{tikzpicture}
\end{sidewaysfigure}
% ====================
%\vfill

%\begin{landscape}
% ====================
%\begin{figure}[H]
%\captionsetup{margin=200cm}

%\centering
%\vspace{-1.5cm}
%\hspace{-1cm}
%

\tikzset{every picture/.style={line width=0.75pt}} %set default line width to 0.75pt        

\begin{tikzpicture}[x=0.75pt,y=0.75pt,yscale=-1,xscale=1]
%uncomment if require: \path (0,694); %set diagram left start at 0, and has height of 694

%Straight Lines [id:da2056540562354816] 
\draw  [dash pattern={on 4.5pt off 4.5pt}]  (672,64.47) -- (810,64.47) ;
%Straight Lines [id:da23044041717189923] 
\draw  [dash pattern={on 4.5pt off 4.5pt}]  (412.67,64.47) -- (610,64.47) ;
\draw [shift={(410.67,64.47)}, rotate = 0] [color={rgb, 255:red, 0; green, 0; blue, 0 }  ][line width=0.75]    (10.93,-3.29) .. controls (6.95,-1.4) and (3.31,-0.3) .. (0,0) .. controls (3.31,0.3) and (6.95,1.4) .. (10.93,3.29)   ;
%Straight Lines [id:da8398435684999113] 
\draw  [dash pattern={on 4.5pt off 4.5pt}]  (810,64.47) -- (810,530) ;
%Rounded Rect [id:dp027018508739026448] 
\draw  [draw opacity=0][fill={rgb, 255:red, 230; green, 230; blue, 230 }  ,fill opacity=0.5 ] (20,439.33) .. controls (20,445.22) and (24.78,450) .. (30.67,450) -- (409.33,450) .. controls (415.22,450) and (420,445.22) .. (420,439.33) -- (420,140.67) .. controls (420,134.78) and (415.22,130) .. (409.33,130) -- (30.67,130) .. controls (24.78,130) and (20,134.78) .. (20,140.67) -- cycle ;
%Rounded Rect [id:dp9415679498487528] 
\draw  [dash pattern={on 4.5pt off 4.5pt}] (29.83,292.77) .. controls (29.83,285.88) and (35.41,280.3) .. (42.3,280.3) -- (398.54,280.3) .. controls (405.42,280.3) and (411,285.88) .. (411,292.77) -- (411,337.17) .. controls (411,344.05) and (405.42,349.64) .. (398.54,349.64) -- (42.3,349.64) .. controls (35.41,349.64) and (29.83,344.05) .. (29.83,337.17) -- cycle ;
%Rounded Rect [id:dp05100022295548956] 
\draw  [draw opacity=0][fill={rgb, 255:red, 230; green, 230; blue, 230 }  ,fill opacity=0.5 ] (130.5,575) .. controls (130.5,583.28) and (137.22,590) .. (145.5,590) -- (615,590) .. controls (623.28,590) and (630,583.28) .. (630,575) -- (630,395) .. controls (630,386.72) and (623.28,380) .. (615,380) -- (145.5,380) .. controls (137.22,380) and (130.5,386.72) .. (130.5,395) -- cycle ;
%Rounded Rect [id:dp6496809750851564] 
\draw  [draw opacity=0][fill={rgb, 255:red, 230; green, 230; blue, 230 }  ,fill opacity=0.5 ] (450.61,440.79) .. controls (450.61,445.82) and (454.69,449.9) .. (459.72,449.9) -- (780.89,449.9) .. controls (785.92,449.9) and (790,445.82) .. (790,440.79) -- (790,139.11) .. controls (790,134.08) and (785.92,130) .. (780.89,130) -- (459.72,130) .. controls (454.69,130) and (450.61,134.08) .. (450.61,139.11) -- cycle ;
%Rounded Rect [id:dp549636265291094] 
\draw  [fill={rgb, 255:red, 230; green, 230; blue, 230 }  ,fill opacity=0.5 ][dash pattern={on 4.5pt off 4.5pt}] (29.83,42.77) .. controls (29.83,35.88) and (35.41,30.3) .. (42.3,30.3) -- (398.54,30.3) .. controls (405.42,30.3) and (411,35.88) .. (411,42.77) -- (411,87.17) .. controls (411,94.05) and (405.42,99.64) .. (398.54,99.64) -- (42.3,99.64) .. controls (35.41,99.64) and (29.83,94.05) .. (29.83,87.17) -- cycle ;
%Rounded Rect [id:dp1378581119146769] 
\draw  [fill={rgb, 255:red, 255; green, 255; blue, 255 }  ,fill opacity=1 ] (40.4,49.82) .. controls (40.4,44.29) and (44.89,39.8) .. (50.42,39.8) -- (169.58,39.8) .. controls (175.11,39.8) and (179.6,44.29) .. (179.6,49.82) -- (179.6,79.88) .. controls (179.6,85.41) and (175.11,89.9) .. (169.58,89.9) -- (50.42,89.9) .. controls (44.89,89.9) and (40.4,85.41) .. (40.4,79.88) -- cycle ;
%Straight Lines [id:da8729153203483555] 
\draw    (188.42,63.77) -- (250.99,63.77)(188.42,66.77) -- (250.99,66.77) ;
\draw [shift={(257.99,65.27)}, rotate = 180] [color={rgb, 255:red, 0; green, 0; blue, 0 }  ][line width=0.75]    (10.93,-4.9) .. controls (6.95,-2.3) and (3.31,-0.67) .. (0,0) .. controls (3.31,0.67) and (6.95,2.3) .. (10.93,4.9)   ;
\draw [shift={(181.42,65.27)}, rotate = 0] [color={rgb, 255:red, 0; green, 0; blue, 0 }  ][line width=0.75]    (10.93,-4.9) .. controls (6.95,-2.3) and (3.31,-0.67) .. (0,0) .. controls (3.31,0.67) and (6.95,2.3) .. (10.93,4.9)   ;
%Straight Lines [id:da845243590120559] 
\draw    (220.17,99.97) -- (220.17,137.64) ;
\draw [shift={(220.17,139.64)}, rotate = 270] [color={rgb, 255:red, 0; green, 0; blue, 0 }  ][line width=0.75]    (10.93,-3.29) .. controls (6.95,-1.4) and (3.31,-0.3) .. (0,0) .. controls (3.31,0.3) and (6.95,1.4) .. (10.93,3.29)   ;
%Straight Lines [id:da7497001546502899] 
\draw    (380,420.33) -- (380,478) ;
\draw [shift={(380,480)}, rotate = 270] [color={rgb, 255:red, 0; green, 0; blue, 0 }  ][line width=0.75]    (10.93,-3.29) .. controls (6.95,-1.4) and (3.31,-0.3) .. (0,0) .. controls (3.31,0.3) and (6.95,1.4) .. (10.93,3.29)   ;
%Straight Lines [id:da5689893955334542] 
\draw    (300,420) -- (460,420) ;
%Rounded Rect [id:dp18582062333838767] 
\draw  [fill={rgb, 255:red, 255; green, 255; blue, 255 }  ,fill opacity=1 ] (260.4,49.82) .. controls (260.4,44.29) and (264.89,39.8) .. (270.42,39.8) -- (389.58,39.8) .. controls (395.11,39.8) and (399.6,44.29) .. (399.6,49.82) -- (399.6,79.88) .. controls (399.6,85.41) and (395.11,89.9) .. (389.58,89.9) -- (270.42,89.9) .. controls (264.89,89.9) and (260.4,85.41) .. (260.4,79.88) -- cycle ;
%Rounded Rect [id:dp9523169763067778] 
\draw  [fill={rgb, 255:red, 255; green, 255; blue, 255 }  ,fill opacity=1 ] (470,299.92) .. controls (470,294.39) and (474.49,289.9) .. (480.02,289.9) -- (599.18,289.9) .. controls (604.71,289.9) and (609.2,294.39) .. (609.2,299.92) -- (609.2,329.98) .. controls (609.2,335.51) and (604.71,340) .. (599.18,340) -- (480.02,340) .. controls (474.49,340) and (470,335.51) .. (470,329.98) -- cycle ;
%Rounded Rect [id:dp48544030815291705] 
\draw  [fill={rgb, 255:red, 255; green, 255; blue, 255 }  ,fill opacity=1 ] (220,490) .. controls (220,484.48) and (224.48,480) .. (230,480) -- (530,480) .. controls (535.52,480) and (540,484.48) .. (540,490) -- (540,570) .. controls (540,575.52) and (535.52,580) .. (530,580) -- (230,580) .. controls (224.48,580) and (220,575.52) .. (220,570) -- cycle ;
%Rounded Rect [id:dp08222627508455749] 
\draw  [fill={rgb, 255:red, 255; green, 255; blue, 255 }  ,fill opacity=1 ] (540.4,49.82) .. controls (540.4,44.29) and (544.89,39.8) .. (550.42,39.8) -- (690.65,39.8) .. controls (696.18,39.8) and (700.67,44.29) .. (700.67,49.82) -- (700.67,79.88) .. controls (700.67,85.41) and (696.18,89.9) .. (690.65,89.9) -- (550.42,89.9) .. controls (544.89,89.9) and (540.4,85.41) .. (540.4,79.88) -- cycle ;
%Straight Lines [id:da5868346364154593] 
\draw  [dash pattern={on 4.5pt off 4.5pt}]  (810,530) -- (542,530) ;
\draw [shift={(540,530)}, rotate = 360] [color={rgb, 255:red, 0; green, 0; blue, 0 }  ][line width=0.75]    (10.93,-3.29) .. controls (6.95,-1.4) and (3.31,-0.3) .. (0,0) .. controls (3.31,0.3) and (6.95,1.4) .. (10.93,3.29)   ;
%Rounded Rect [id:dp7392784471609453] 
\draw  [fill={rgb, 255:red, 255; green, 255; blue, 255 }  ,fill opacity=1 ] (30.33,150.47) .. controls (30.33,144.58) and (35.11,139.8) .. (41,139.8) -- (400.34,139.8) .. controls (406.23,139.8) and (411,144.58) .. (411,150.47) -- (411,229.34) .. controls (411,235.23) and (406.23,240.01) .. (400.34,240.01) -- (41,240.01) .. controls (35.11,240.01) and (30.33,235.23) .. (30.33,229.34) -- cycle ;
%Rounded Rect [id:dp90901637181607] 
\draw  [fill={rgb, 255:red, 255; green, 255; blue, 255 }  ,fill opacity=1 ] (41,189.82) .. controls (41,184.29) and (45.49,179.8) .. (51.02,179.8) -- (145.02,179.8) .. controls (150.55,179.8) and (155.04,184.29) .. (155.04,189.82) -- (155.04,219.88) .. controls (155.04,225.41) and (150.55,229.9) .. (145.02,229.9) -- (51.02,229.9) .. controls (45.49,229.9) and (41,225.41) .. (41,219.88) -- cycle ;
%Rounded Rect [id:dp06115018179625542] 
\draw  [fill={rgb, 255:red, 255; green, 255; blue, 255 }  ,fill opacity=1 ] (163.65,189.82) .. controls (163.65,184.29) and (168.13,179.8) .. (173.67,179.8) -- (267.67,179.8) .. controls (273.2,179.8) and (277.69,184.29) .. (277.69,189.82) -- (277.69,219.88) .. controls (277.69,225.41) and (273.2,229.9) .. (267.67,229.9) -- (173.67,229.9) .. controls (168.13,229.9) and (163.65,225.41) .. (163.65,219.88) -- cycle ;
%Rounded Rect [id:dp6195275548153432] 
\draw  [fill={rgb, 255:red, 255; green, 255; blue, 255 }  ,fill opacity=1 ] (286.3,189.82) .. controls (286.3,184.29) and (290.78,179.8) .. (296.32,179.8) -- (390.31,179.8) .. controls (395.85,179.8) and (400.33,184.29) .. (400.33,189.82) -- (400.33,219.88) .. controls (400.33,225.41) and (395.85,229.9) .. (390.31,229.9) -- (296.32,229.9) .. controls (290.78,229.9) and (286.3,225.41) .. (286.3,219.88) -- cycle ;
%Rounded Rect [id:dp41283415100159493] 
\draw  [fill={rgb, 255:red, 255; green, 255; blue, 255 }  ,fill opacity=1 ] (40.8,299.92) .. controls (40.8,294.39) and (45.29,289.9) .. (50.82,289.9) -- (169.98,289.9) .. controls (175.51,289.9) and (180,294.39) .. (180,299.92) -- (180,329.98) .. controls (180,335.51) and (175.51,340) .. (169.98,340) -- (50.82,340) .. controls (45.29,340) and (40.8,335.51) .. (40.8,329.98) -- cycle ;
%Rounded Rect [id:dp40154854661776396] 
\draw  [fill={rgb, 255:red, 255; green, 255; blue, 255 }  ,fill opacity=1 ] (260.4,299.82) .. controls (260.4,294.29) and (264.89,289.8) .. (270.42,289.8) -- (389.58,289.8) .. controls (395.11,289.8) and (399.6,294.29) .. (399.6,299.82) -- (399.6,329.88) .. controls (399.6,335.41) and (395.11,339.9) .. (389.58,339.9) -- (270.42,339.9) .. controls (264.89,339.9) and (260.4,335.41) .. (260.4,329.88) -- cycle ;
%Straight Lines [id:da6602630444249074] 
\draw    (180.17,315.8) -- (260,315.8) ;
%Straight Lines [id:da30905866347109745] 
\draw    (219.86,315.97) -- (219.86,386.81) ;
\draw [shift={(219.86,388.81)}, rotate = 270] [color={rgb, 255:red, 0; green, 0; blue, 0 }  ][line width=0.75]    (10.93,-3.29) .. controls (6.95,-1.4) and (3.31,-0.3) .. (0,0) .. controls (3.31,0.3) and (6.95,1.4) .. (10.93,3.29)   ;
%Rounded Rect [id:dp44627062039563214] 
\draw  [fill={rgb, 255:red, 255; green, 255; blue, 255 }  ,fill opacity=1 ] (139.67,399.82) .. controls (139.67,394.29) and (144.15,389.8) .. (149.69,389.8) -- (290.31,389.8) .. controls (295.85,389.8) and (300.33,394.29) .. (300.33,399.82) -- (300.33,429.88) .. controls (300.33,435.41) and (295.85,439.9) .. (290.31,439.9) -- (149.69,439.9) .. controls (144.15,439.9) and (139.67,435.41) .. (139.67,429.88) -- cycle ;
%Straight Lines [id:da8375870493480149] 
\draw    (220,240) -- (220,277.67) ;
\draw [shift={(220,279.67)}, rotate = 270] [color={rgb, 255:red, 0; green, 0; blue, 0 }  ][line width=0.75]    (10.93,-3.29) .. controls (6.95,-1.4) and (3.31,-0.3) .. (0,0) .. controls (3.31,0.3) and (6.95,1.4) .. (10.93,3.29)   ;
%Rounded Rect [id:dp37285624114093485] 
\draw  [fill={rgb, 255:red, 255; green, 255; blue, 255 }  ,fill opacity=1 ] (460,400.02) .. controls (460,394.49) and (464.49,390) .. (470.02,390) -- (610.65,390) .. controls (616.18,390) and (620.67,394.49) .. (620.67,400.02) -- (620.67,430.08) .. controls (620.67,435.61) and (616.18,440.1) .. (610.65,440.1) -- (470.02,440.1) .. controls (464.49,440.1) and (460,435.61) .. (460,430.08) -- cycle ;
%Rounded Rect [id:dp6883616335272991] 
\draw  [fill={rgb, 255:red, 255; green, 255; blue, 255 }  ,fill opacity=1 ] (629.2,299.92) .. controls (629.2,294.39) and (633.69,289.9) .. (639.22,289.9) -- (758.38,289.9) .. controls (763.91,289.9) and (768.4,294.39) .. (768.4,299.92) -- (768.4,329.98) .. controls (768.4,335.51) and (763.91,340) .. (758.38,340) -- (639.22,340) .. controls (633.69,340) and (629.2,335.51) .. (629.2,329.98) -- cycle ;
%Straight Lines [id:da19870419404011752] 
\draw    (540,340) -- (540,388) ;
\draw [shift={(540,390)}, rotate = 270] [color={rgb, 255:red, 0; green, 0; blue, 0 }  ][line width=0.75]    (10.93,-3.29) .. controls (6.95,-1.4) and (3.31,-0.3) .. (0,0) .. controls (3.31,0.3) and (6.95,1.4) .. (10.93,3.29)   ;
%Rounded Rect [id:dp9692784599420237] 
\draw  [fill={rgb, 255:red, 255; green, 255; blue, 255 }  ,fill opacity=1 ] (230,529.92) .. controls (230,524.39) and (234.49,519.9) .. (240.02,519.9) -- (359.18,519.9) .. controls (364.71,519.9) and (369.2,524.39) .. (369.2,529.92) -- (369.2,559.98) .. controls (369.2,565.51) and (364.71,570) .. (359.18,570) -- (240.02,570) .. controls (234.49,570) and (230,565.51) .. (230,559.98) -- cycle ;
%Rounded Rect [id:dp749684165723072] 
\draw  [fill={rgb, 255:red, 255; green, 255; blue, 255 }  ,fill opacity=1 ] (390,530.02) .. controls (390,524.49) and (394.49,520) .. (400.02,520) -- (519.18,520) .. controls (524.71,520) and (529.2,524.49) .. (529.2,530.02) -- (529.2,560.08) .. controls (529.2,565.61) and (524.71,570.1) .. (519.18,570.1) -- (400.02,570.1) .. controls (394.49,570.1) and (390,565.61) .. (390,560.08) -- cycle ;
%Rounded Rect [id:dp28001422897762485] 
\draw  [fill={rgb, 255:red, 255; green, 255; blue, 255 }  ,fill opacity=1 ][dash pattern={on 4.5pt off 4.5pt}] (460,148.92) .. controls (460,143.99) and (463.99,140) .. (468.92,140) -- (771.08,140) .. controls (776.01,140) and (780,143.99) .. (780,148.92) -- (780,241.55) .. controls (780,246.48) and (776.01,250.47) .. (771.08,250.47) -- (468.92,250.47) .. controls (463.99,250.47) and (460,246.48) .. (460,241.55) -- cycle ;
%Shape: Brace [id:dp9411387904179627] 
\draw   (774.67,286.8) .. controls (774.67,282.13) and (772.34,279.8) .. (767.67,279.8) -- (629.96,279.8) .. controls (623.29,279.8) and (619.96,277.47) .. (619.96,272.8) .. controls (619.96,277.47) and (616.63,279.8) .. (609.96,279.8)(612.96,279.8) -- (472.25,279.8) .. controls (467.58,279.8) and (465.25,282.13) .. (465.25,286.8) ;
%Straight Lines [id:da9902399507274253] 
\draw    (620,273) -- (620,253.14) ;
\draw [shift={(620,251.14)}, rotate = 90] [color={rgb, 255:red, 0; green, 0; blue, 0 }  ][line width=0.75]    (10.93,-3.29) .. controls (6.95,-1.4) and (3.31,-0.3) .. (0,0) .. controls (3.31,0.3) and (6.95,1.4) .. (10.93,3.29)   ;

% Text Node
\draw (110,64.85) node  [color={rgb, 255:red, 0; green, 0; blue, 0 }  ,opacity=1 ] [align=left] {\begin{minipage}[lt]{92.93pt}\setlength\topsep{0pt}
\begin{center}
programa en Python\\
\small{(\cref{sec:programas-python})}
\end{center}

\end{minipage}};
% Text Node
\draw (219.89,51.97) node  [font=\small,color={rgb, 255:red, 0; green, 0; blue, 0 }  ,opacity=1 ] [align=left] {\begin{minipage}[lt]{54.86pt}\setlength\topsep{0pt}
\begin{center}
\Cref{sec:equivalencia}
\end{center}

\end{minipage}};
% Text Node
\draw (69.5,15.5) node  [font=\small,color={rgb, 255:red, 0; green, 0; blue, 0 }  ,opacity=1 ] [align=left] {\begin{minipage}[lt]{41.55pt}\setlength\topsep{0pt}
\cref{ch:programas-python-maquinas-turing}
\end{minipage}};
% Text Node
\draw (330,64.85) node  [color={rgb, 255:red, 0; green, 0; blue, 0 }  ,opacity=1 ] [align=left] {\begin{minipage}[lt]{92.93pt}\setlength\topsep{0pt}
\begin{center}
máquina de Turing\\
\small{(\cref{sec:maquinas-turing})}
\end{center}

\end{minipage}};
% Text Node
\draw (539.6,314.95) node  [color={rgb, 255:red, 0; green, 0; blue, 0 }  ,opacity=1 ] [align=left] {\begin{minipage}[lt]{92.93pt}\setlength\topsep{0pt}
\begin{center}
sistema lógico\\
\small{(\cref{sec:sistemas-logicos})}
\end{center}

\end{minipage}};
% Text Node
\draw (380.5,500.11) node  [color={rgb, 255:red, 0; green, 0; blue, 0 }  ,opacity=1 ] [align=left] {\begin{minipage}[lt]{159.55pt}\setlength\topsep{0pt}
\begin{center}
\textbf{Primer Teorema de Incompletitud}
\end{center}

\end{minipage}};
% Text Node
\draw (620.53,64.85) node  [color={rgb, 255:red, 0; green, 0; blue, 0 }  ,opacity=1 ] [align=left] {\begin{minipage}[lt]{106.99pt}\setlength\topsep{0pt}
\begin{center}
tesis de Church-Turing\\
\small{(\cref{sec:church-turing})}
\end{center}

\end{minipage}};
% Text Node
\draw (480.28,52.8) node  [font=\small,color={rgb, 255:red, 0; green, 0; blue, 0 }  ,opacity=1 ] [align=left] {\begin{minipage}[lt]{81.9pt}\setlength\topsep{0pt}
\begin{center}
\cref{sec:church-turing}
\end{center}

\end{minipage}};
% Text Node
\draw (589.5,518.3) node  [font=\small,color={rgb, 255:red, 0; green, 0; blue, 0 }  ,opacity=1 ] [align=left] {\begin{minipage}[lt]{54.86pt}\setlength\topsep{0pt}
\begin{center}
\Cref{subsec:mente-humana-tesis-church-turing}
\end{center}

\end{minipage}};
% Text Node
\draw (220.2,159.62) node  [color={rgb, 255:red, 0; green, 0; blue, 0 }  ,opacity=1 ] [align=left] {\begin{minipage}[lt]{230.11pt}\setlength\topsep{0pt}
\begin{center}
problema computacional
\end{center}

\end{minipage}};
% Text Node
\draw (98.02,204.85) node  [color={rgb, 255:red, 0; green, 0; blue, 0 }  ,opacity=1 ] [align=left] {\begin{minipage}[lt]{76.13pt}\setlength\topsep{0pt}
\begin{center}
decidible\\
\small{(\cref{sec:problemas-decidibles})}
\end{center}

\end{minipage}};
% Text Node
\draw (220.67,204.85) node  [color={rgb, 255:red, 0; green, 0; blue, 0 }  ,opacity=1 ] [align=left] {\begin{minipage}[lt]{76.13pt}\setlength\topsep{0pt}
\begin{center}
no decidible\\
\small{(\cref{sec:problemas-no-decidibles})}
\end{center}

\end{minipage}};
% Text Node
\draw (343.31,204.85) node  [color={rgb, 255:red, 0; green, 0; blue, 0 }  ,opacity=1 ] [align=left] {\begin{minipage}[lt]{76.13pt}\setlength\topsep{0pt}
\begin{center}
semidecidible\\
\small{(\cref{sec:problemas-semidecidibles})}
\end{center}

\end{minipage}};
% Text Node
\draw (110.4,314.95) node  [color={rgb, 255:red, 0; green, 0; blue, 0 }  ,opacity=1 ] [align=left] {\begin{minipage}[lt]{92.93pt}\setlength\topsep{0pt}
\begin{center}
universalidad\\
\small{(\cref{sec:universalidad})}
\end{center}

\end{minipage}};
% Text Node
\draw (330,314.85) node  [color={rgb, 255:red, 0; green, 0; blue, 0 }  ,opacity=1 ] [align=left] {\begin{minipage}[lt]{92.93pt}\setlength\topsep{0pt}
\begin{center}
reducción\\
\small{(\cref{sec:reducciones})}
\end{center}

\end{minipage}};
% Text Node
\draw (220,414.85) node  [color={rgb, 255:red, 0; green, 0; blue, 0 }  ,opacity=1 ] [align=left] {\begin{minipage}[lt]{107.26pt}\setlength\topsep{0pt}
\begin{center}
problema de la parada\\
\small{(\cref{sec:problema-parada})}
\end{center}

\end{minipage}};
% Text Node
\draw (540.33,415.05) node  [color={rgb, 255:red, 0; green, 0; blue, 0 }  ,opacity=1 ] [align=left] {\begin{minipage}[lt]{107.26pt}\setlength\topsep{0pt}
\begin{center}
aritmética de Peano\\
\small{(\cref{sec:aritmetica-peano})}
\end{center}

\end{minipage}};
% Text Node
\draw (698.8,314.95) node  [color={rgb, 255:red, 0; green, 0; blue, 0 }  ,opacity=1 ] [align=left] {\begin{minipage}[lt]{92.93pt}\setlength\topsep{0pt}
\begin{center}
sistema formal\\
\small{(\cref{sec:sistemas-formales})}
\end{center}

\end{minipage}};
% Text Node
\draw (299.6,544.95) node  [color={rgb, 255:red, 0; green, 0; blue, 0 }  ,opacity=1 ] [align=left] {\begin{minipage}[lt]{92.93pt}\setlength\topsep{0pt}
\begin{center}
versión semántica\\
\small{(\cref{sec:version-semantica})}
\end{center}

\end{minipage}};
% Text Node
\draw (459.6,545.05) node  [color={rgb, 255:red, 0; green, 0; blue, 0 }  ,opacity=1 ] [align=left] {\begin{minipage}[lt]{92.93pt}\setlength\topsep{0pt}
\begin{center}
versión sintáctica\\
\small{(\cref{sec:version-sintactica})}
\end{center}

\end{minipage}};
% Text Node
\draw (60.5,464.5) node  [font=\small,color={rgb, 255:red, 0; green, 0; blue, 0 }  ,opacity=1 ] [align=left] {\begin{minipage}[lt]{41.55pt}\setlength\topsep{0pt}
\cref{ch:problema-parada}
\end{minipage}};
% Text Node
\draw (749.5,465.5) node  [font=\small,color={rgb, 255:red, 0; green, 0; blue, 0 }  ,opacity=1 ] [align=left] {\begin{minipage}[lt]{41.55pt}\setlength\topsep{0pt}
\begin{flushright}
\cref{ch:sistemas-logicos}
\end{flushright}

\end{minipage}};
% Text Node
\draw (170.5,604.5) node  [font=\small,color={rgb, 255:red, 0; green, 0; blue, 0 }  ,opacity=1 ] [align=left] {\begin{minipage}[lt]{41.55pt}\setlength\topsep{0pt}
\cref{ch:teorema-incompletitud}
\end{minipage}};
% Text Node
\draw (643.76,195.55) node  [color={rgb, 255:red, 0; green, 0; blue, 0 }  ,opacity=1 ] [align=left] {\begin{minipage}[lt]{217.92pt}\setlength\topsep{0pt}
-- sólido {\small{(\cref{sec:solidez-completitud-decidibilidad})}}\\
-- decidible {\small{(\cref{sec:solidez-completitud-decidibilidad})}}\\
-- (semánticamente) completo {\small{(\cref{sec:solidez-completitud-decidibilidad})}}\\
-- sintácticamente completo {\small{(\cref{sec:consistencia})}}\\
-- sem./sint. consistente {\small{(\cref{sec:consistencia})}}
\end{minipage}};
% Text Node
\draw (414.84,645.25) node  [color={rgb, 255:red, 0; green, 0; blue, 0 }  ,opacity=1 ] [align=left] {\begin{minipage}[lt]{534.7pt}\setlength\topsep{0pt}
\begin{center}
Esquema resumen de este Trabajo de Fin de Grado
\end{center}
\end{minipage}};
% Text Node
\draw (380.15,407.8) node  [font=\small,color={rgb, 255:red, 0; green, 0; blue, 0 }  ,opacity=1 ] [align=left] {\begin{minipage}[lt]{108.93pt}\setlength\topsep{0pt}
\begin{center}
\cref{sec:primera-aproximacion}
\end{center}

\end{minipage}};

\end{tikzpicture}
%\raisebox{-15.5cm}[0pt][0pt]{%
%\resizebox{\linewidth}{!}{

\tikzset{every picture/.style={line width=0.75pt}} %set default line width to 0.75pt        

\begin{tikzpicture}[x=0.75pt,y=0.75pt,yscale=-1,xscale=1]
%uncomment if require: \path (0,694); %set diagram left start at 0, and has height of 694

%Straight Lines [id:da2056540562354816] 
\draw  [dash pattern={on 4.5pt off 4.5pt}]  (672,64.47) -- (810,64.47) ;
%Straight Lines [id:da23044041717189923] 
\draw  [dash pattern={on 4.5pt off 4.5pt}]  (412.67,64.47) -- (610,64.47) ;
\draw [shift={(410.67,64.47)}, rotate = 0] [color={rgb, 255:red, 0; green, 0; blue, 0 }  ][line width=0.75]    (10.93,-3.29) .. controls (6.95,-1.4) and (3.31,-0.3) .. (0,0) .. controls (3.31,0.3) and (6.95,1.4) .. (10.93,3.29)   ;
%Straight Lines [id:da8398435684999113] 
\draw  [dash pattern={on 4.5pt off 4.5pt}]  (810,64.47) -- (810,530) ;
%Rounded Rect [id:dp027018508739026448] 
\draw  [draw opacity=0][fill={rgb, 255:red, 230; green, 230; blue, 230 }  ,fill opacity=0.5 ] (20,439.33) .. controls (20,445.22) and (24.78,450) .. (30.67,450) -- (409.33,450) .. controls (415.22,450) and (420,445.22) .. (420,439.33) -- (420,140.67) .. controls (420,134.78) and (415.22,130) .. (409.33,130) -- (30.67,130) .. controls (24.78,130) and (20,134.78) .. (20,140.67) -- cycle ;
%Rounded Rect [id:dp9415679498487528] 
\draw  [dash pattern={on 4.5pt off 4.5pt}] (29.83,292.77) .. controls (29.83,285.88) and (35.41,280.3) .. (42.3,280.3) -- (398.54,280.3) .. controls (405.42,280.3) and (411,285.88) .. (411,292.77) -- (411,337.17) .. controls (411,344.05) and (405.42,349.64) .. (398.54,349.64) -- (42.3,349.64) .. controls (35.41,349.64) and (29.83,344.05) .. (29.83,337.17) -- cycle ;
%Rounded Rect [id:dp05100022295548956] 
\draw  [draw opacity=0][fill={rgb, 255:red, 230; green, 230; blue, 230 }  ,fill opacity=0.5 ] (130.5,575) .. controls (130.5,583.28) and (137.22,590) .. (145.5,590) -- (615,590) .. controls (623.28,590) and (630,583.28) .. (630,575) -- (630,395) .. controls (630,386.72) and (623.28,380) .. (615,380) -- (145.5,380) .. controls (137.22,380) and (130.5,386.72) .. (130.5,395) -- cycle ;
%Rounded Rect [id:dp6496809750851564] 
\draw  [draw opacity=0][fill={rgb, 255:red, 230; green, 230; blue, 230 }  ,fill opacity=0.5 ] (450.61,440.79) .. controls (450.61,445.82) and (454.69,449.9) .. (459.72,449.9) -- (780.89,449.9) .. controls (785.92,449.9) and (790,445.82) .. (790,440.79) -- (790,139.11) .. controls (790,134.08) and (785.92,130) .. (780.89,130) -- (459.72,130) .. controls (454.69,130) and (450.61,134.08) .. (450.61,139.11) -- cycle ;
%Rounded Rect [id:dp549636265291094] 
\draw  [fill={rgb, 255:red, 230; green, 230; blue, 230 }  ,fill opacity=0.5 ][dash pattern={on 4.5pt off 4.5pt}] (29.83,42.77) .. controls (29.83,35.88) and (35.41,30.3) .. (42.3,30.3) -- (398.54,30.3) .. controls (405.42,30.3) and (411,35.88) .. (411,42.77) -- (411,87.17) .. controls (411,94.05) and (405.42,99.64) .. (398.54,99.64) -- (42.3,99.64) .. controls (35.41,99.64) and (29.83,94.05) .. (29.83,87.17) -- cycle ;
%Rounded Rect [id:dp1378581119146769] 
\draw  [fill={rgb, 255:red, 255; green, 255; blue, 255 }  ,fill opacity=1 ] (40.4,49.82) .. controls (40.4,44.29) and (44.89,39.8) .. (50.42,39.8) -- (169.58,39.8) .. controls (175.11,39.8) and (179.6,44.29) .. (179.6,49.82) -- (179.6,79.88) .. controls (179.6,85.41) and (175.11,89.9) .. (169.58,89.9) -- (50.42,89.9) .. controls (44.89,89.9) and (40.4,85.41) .. (40.4,79.88) -- cycle ;
%Straight Lines [id:da8729153203483555] 
\draw    (188.42,63.77) -- (250.99,63.77)(188.42,66.77) -- (250.99,66.77) ;
\draw [shift={(257.99,65.27)}, rotate = 180] [color={rgb, 255:red, 0; green, 0; blue, 0 }  ][line width=0.75]    (10.93,-4.9) .. controls (6.95,-2.3) and (3.31,-0.67) .. (0,0) .. controls (3.31,0.67) and (6.95,2.3) .. (10.93,4.9)   ;
\draw [shift={(181.42,65.27)}, rotate = 0] [color={rgb, 255:red, 0; green, 0; blue, 0 }  ][line width=0.75]    (10.93,-4.9) .. controls (6.95,-2.3) and (3.31,-0.67) .. (0,0) .. controls (3.31,0.67) and (6.95,2.3) .. (10.93,4.9)   ;
%Straight Lines [id:da845243590120559] 
\draw    (220.17,99.97) -- (220.17,137.64) ;
\draw [shift={(220.17,139.64)}, rotate = 270] [color={rgb, 255:red, 0; green, 0; blue, 0 }  ][line width=0.75]    (10.93,-3.29) .. controls (6.95,-1.4) and (3.31,-0.3) .. (0,0) .. controls (3.31,0.3) and (6.95,1.4) .. (10.93,3.29)   ;
%Straight Lines [id:da7497001546502899] 
\draw    (380,420.33) -- (380,478) ;
\draw [shift={(380,480)}, rotate = 270] [color={rgb, 255:red, 0; green, 0; blue, 0 }  ][line width=0.75]    (10.93,-3.29) .. controls (6.95,-1.4) and (3.31,-0.3) .. (0,0) .. controls (3.31,0.3) and (6.95,1.4) .. (10.93,3.29)   ;
%Straight Lines [id:da5689893955334542] 
\draw    (300,420) -- (460,420) ;
%Rounded Rect [id:dp18582062333838767] 
\draw  [fill={rgb, 255:red, 255; green, 255; blue, 255 }  ,fill opacity=1 ] (260.4,49.82) .. controls (260.4,44.29) and (264.89,39.8) .. (270.42,39.8) -- (389.58,39.8) .. controls (395.11,39.8) and (399.6,44.29) .. (399.6,49.82) -- (399.6,79.88) .. controls (399.6,85.41) and (395.11,89.9) .. (389.58,89.9) -- (270.42,89.9) .. controls (264.89,89.9) and (260.4,85.41) .. (260.4,79.88) -- cycle ;
%Rounded Rect [id:dp9523169763067778] 
\draw  [fill={rgb, 255:red, 255; green, 255; blue, 255 }  ,fill opacity=1 ] (470,299.92) .. controls (470,294.39) and (474.49,289.9) .. (480.02,289.9) -- (599.18,289.9) .. controls (604.71,289.9) and (609.2,294.39) .. (609.2,299.92) -- (609.2,329.98) .. controls (609.2,335.51) and (604.71,340) .. (599.18,340) -- (480.02,340) .. controls (474.49,340) and (470,335.51) .. (470,329.98) -- cycle ;
%Rounded Rect [id:dp48544030815291705] 
\draw  [fill={rgb, 255:red, 255; green, 255; blue, 255 }  ,fill opacity=1 ] (220,490) .. controls (220,484.48) and (224.48,480) .. (230,480) -- (530,480) .. controls (535.52,480) and (540,484.48) .. (540,490) -- (540,570) .. controls (540,575.52) and (535.52,580) .. (530,580) -- (230,580) .. controls (224.48,580) and (220,575.52) .. (220,570) -- cycle ;
%Rounded Rect [id:dp08222627508455749] 
\draw  [fill={rgb, 255:red, 255; green, 255; blue, 255 }  ,fill opacity=1 ] (540.4,49.82) .. controls (540.4,44.29) and (544.89,39.8) .. (550.42,39.8) -- (690.65,39.8) .. controls (696.18,39.8) and (700.67,44.29) .. (700.67,49.82) -- (700.67,79.88) .. controls (700.67,85.41) and (696.18,89.9) .. (690.65,89.9) -- (550.42,89.9) .. controls (544.89,89.9) and (540.4,85.41) .. (540.4,79.88) -- cycle ;
%Straight Lines [id:da5868346364154593] 
\draw  [dash pattern={on 4.5pt off 4.5pt}]  (810,530) -- (542,530) ;
\draw [shift={(540,530)}, rotate = 360] [color={rgb, 255:red, 0; green, 0; blue, 0 }  ][line width=0.75]    (10.93,-3.29) .. controls (6.95,-1.4) and (3.31,-0.3) .. (0,0) .. controls (3.31,0.3) and (6.95,1.4) .. (10.93,3.29)   ;
%Rounded Rect [id:dp7392784471609453] 
\draw  [fill={rgb, 255:red, 255; green, 255; blue, 255 }  ,fill opacity=1 ] (30.33,150.47) .. controls (30.33,144.58) and (35.11,139.8) .. (41,139.8) -- (400.34,139.8) .. controls (406.23,139.8) and (411,144.58) .. (411,150.47) -- (411,229.34) .. controls (411,235.23) and (406.23,240.01) .. (400.34,240.01) -- (41,240.01) .. controls (35.11,240.01) and (30.33,235.23) .. (30.33,229.34) -- cycle ;
%Rounded Rect [id:dp90901637181607] 
\draw  [fill={rgb, 255:red, 255; green, 255; blue, 255 }  ,fill opacity=1 ] (41,189.82) .. controls (41,184.29) and (45.49,179.8) .. (51.02,179.8) -- (145.02,179.8) .. controls (150.55,179.8) and (155.04,184.29) .. (155.04,189.82) -- (155.04,219.88) .. controls (155.04,225.41) and (150.55,229.9) .. (145.02,229.9) -- (51.02,229.9) .. controls (45.49,229.9) and (41,225.41) .. (41,219.88) -- cycle ;
%Rounded Rect [id:dp06115018179625542] 
\draw  [fill={rgb, 255:red, 255; green, 255; blue, 255 }  ,fill opacity=1 ] (163.65,189.82) .. controls (163.65,184.29) and (168.13,179.8) .. (173.67,179.8) -- (267.67,179.8) .. controls (273.2,179.8) and (277.69,184.29) .. (277.69,189.82) -- (277.69,219.88) .. controls (277.69,225.41) and (273.2,229.9) .. (267.67,229.9) -- (173.67,229.9) .. controls (168.13,229.9) and (163.65,225.41) .. (163.65,219.88) -- cycle ;
%Rounded Rect [id:dp6195275548153432] 
\draw  [fill={rgb, 255:red, 255; green, 255; blue, 255 }  ,fill opacity=1 ] (286.3,189.82) .. controls (286.3,184.29) and (290.78,179.8) .. (296.32,179.8) -- (390.31,179.8) .. controls (395.85,179.8) and (400.33,184.29) .. (400.33,189.82) -- (400.33,219.88) .. controls (400.33,225.41) and (395.85,229.9) .. (390.31,229.9) -- (296.32,229.9) .. controls (290.78,229.9) and (286.3,225.41) .. (286.3,219.88) -- cycle ;
%Rounded Rect [id:dp41283415100159493] 
\draw  [fill={rgb, 255:red, 255; green, 255; blue, 255 }  ,fill opacity=1 ] (40.8,299.92) .. controls (40.8,294.39) and (45.29,289.9) .. (50.82,289.9) -- (169.98,289.9) .. controls (175.51,289.9) and (180,294.39) .. (180,299.92) -- (180,329.98) .. controls (180,335.51) and (175.51,340) .. (169.98,340) -- (50.82,340) .. controls (45.29,340) and (40.8,335.51) .. (40.8,329.98) -- cycle ;
%Rounded Rect [id:dp40154854661776396] 
\draw  [fill={rgb, 255:red, 255; green, 255; blue, 255 }  ,fill opacity=1 ] (260.4,299.82) .. controls (260.4,294.29) and (264.89,289.8) .. (270.42,289.8) -- (389.58,289.8) .. controls (395.11,289.8) and (399.6,294.29) .. (399.6,299.82) -- (399.6,329.88) .. controls (399.6,335.41) and (395.11,339.9) .. (389.58,339.9) -- (270.42,339.9) .. controls (264.89,339.9) and (260.4,335.41) .. (260.4,329.88) -- cycle ;
%Straight Lines [id:da6602630444249074] 
\draw    (180.17,315.8) -- (260,315.8) ;
%Straight Lines [id:da30905866347109745] 
\draw    (219.86,315.97) -- (219.86,386.81) ;
\draw [shift={(219.86,388.81)}, rotate = 270] [color={rgb, 255:red, 0; green, 0; blue, 0 }  ][line width=0.75]    (10.93,-3.29) .. controls (6.95,-1.4) and (3.31,-0.3) .. (0,0) .. controls (3.31,0.3) and (6.95,1.4) .. (10.93,3.29)   ;
%Rounded Rect [id:dp44627062039563214] 
\draw  [fill={rgb, 255:red, 255; green, 255; blue, 255 }  ,fill opacity=1 ] (139.67,399.82) .. controls (139.67,394.29) and (144.15,389.8) .. (149.69,389.8) -- (290.31,389.8) .. controls (295.85,389.8) and (300.33,394.29) .. (300.33,399.82) -- (300.33,429.88) .. controls (300.33,435.41) and (295.85,439.9) .. (290.31,439.9) -- (149.69,439.9) .. controls (144.15,439.9) and (139.67,435.41) .. (139.67,429.88) -- cycle ;
%Straight Lines [id:da8375870493480149] 
\draw    (220,240) -- (220,277.67) ;
\draw [shift={(220,279.67)}, rotate = 270] [color={rgb, 255:red, 0; green, 0; blue, 0 }  ][line width=0.75]    (10.93,-3.29) .. controls (6.95,-1.4) and (3.31,-0.3) .. (0,0) .. controls (3.31,0.3) and (6.95,1.4) .. (10.93,3.29)   ;
%Rounded Rect [id:dp37285624114093485] 
\draw  [fill={rgb, 255:red, 255; green, 255; blue, 255 }  ,fill opacity=1 ] (460,400.02) .. controls (460,394.49) and (464.49,390) .. (470.02,390) -- (610.65,390) .. controls (616.18,390) and (620.67,394.49) .. (620.67,400.02) -- (620.67,430.08) .. controls (620.67,435.61) and (616.18,440.1) .. (610.65,440.1) -- (470.02,440.1) .. controls (464.49,440.1) and (460,435.61) .. (460,430.08) -- cycle ;
%Rounded Rect [id:dp6883616335272991] 
\draw  [fill={rgb, 255:red, 255; green, 255; blue, 255 }  ,fill opacity=1 ] (629.2,299.92) .. controls (629.2,294.39) and (633.69,289.9) .. (639.22,289.9) -- (758.38,289.9) .. controls (763.91,289.9) and (768.4,294.39) .. (768.4,299.92) -- (768.4,329.98) .. controls (768.4,335.51) and (763.91,340) .. (758.38,340) -- (639.22,340) .. controls (633.69,340) and (629.2,335.51) .. (629.2,329.98) -- cycle ;
%Straight Lines [id:da19870419404011752] 
\draw    (540,340) -- (540,388) ;
\draw [shift={(540,390)}, rotate = 270] [color={rgb, 255:red, 0; green, 0; blue, 0 }  ][line width=0.75]    (10.93,-3.29) .. controls (6.95,-1.4) and (3.31,-0.3) .. (0,0) .. controls (3.31,0.3) and (6.95,1.4) .. (10.93,3.29)   ;
%Rounded Rect [id:dp9692784599420237] 
\draw  [fill={rgb, 255:red, 255; green, 255; blue, 255 }  ,fill opacity=1 ] (230,529.92) .. controls (230,524.39) and (234.49,519.9) .. (240.02,519.9) -- (359.18,519.9) .. controls (364.71,519.9) and (369.2,524.39) .. (369.2,529.92) -- (369.2,559.98) .. controls (369.2,565.51) and (364.71,570) .. (359.18,570) -- (240.02,570) .. controls (234.49,570) and (230,565.51) .. (230,559.98) -- cycle ;
%Rounded Rect [id:dp749684165723072] 
\draw  [fill={rgb, 255:red, 255; green, 255; blue, 255 }  ,fill opacity=1 ] (390,530.02) .. controls (390,524.49) and (394.49,520) .. (400.02,520) -- (519.18,520) .. controls (524.71,520) and (529.2,524.49) .. (529.2,530.02) -- (529.2,560.08) .. controls (529.2,565.61) and (524.71,570.1) .. (519.18,570.1) -- (400.02,570.1) .. controls (394.49,570.1) and (390,565.61) .. (390,560.08) -- cycle ;
%Rounded Rect [id:dp28001422897762485] 
\draw  [fill={rgb, 255:red, 255; green, 255; blue, 255 }  ,fill opacity=1 ][dash pattern={on 4.5pt off 4.5pt}] (460,148.92) .. controls (460,143.99) and (463.99,140) .. (468.92,140) -- (771.08,140) .. controls (776.01,140) and (780,143.99) .. (780,148.92) -- (780,241.55) .. controls (780,246.48) and (776.01,250.47) .. (771.08,250.47) -- (468.92,250.47) .. controls (463.99,250.47) and (460,246.48) .. (460,241.55) -- cycle ;
%Shape: Brace [id:dp9411387904179627] 
\draw   (774.67,286.8) .. controls (774.67,282.13) and (772.34,279.8) .. (767.67,279.8) -- (629.96,279.8) .. controls (623.29,279.8) and (619.96,277.47) .. (619.96,272.8) .. controls (619.96,277.47) and (616.63,279.8) .. (609.96,279.8)(612.96,279.8) -- (472.25,279.8) .. controls (467.58,279.8) and (465.25,282.13) .. (465.25,286.8) ;
%Straight Lines [id:da9902399507274253] 
\draw    (620,273) -- (620,253.14) ;
\draw [shift={(620,251.14)}, rotate = 90] [color={rgb, 255:red, 0; green, 0; blue, 0 }  ][line width=0.75]    (10.93,-3.29) .. controls (6.95,-1.4) and (3.31,-0.3) .. (0,0) .. controls (3.31,0.3) and (6.95,1.4) .. (10.93,3.29)   ;

% Text Node
\draw (110,64.85) node  [color={rgb, 255:red, 0; green, 0; blue, 0 }  ,opacity=1 ] [align=left] {\begin{minipage}[lt]{92.93pt}\setlength\topsep{0pt}
\begin{center}
Python program\\
\small{(\hyperref[sec:programas-python]{Section 3.1})}
\end{center}

\end{minipage}};
% Text Node
\draw (219.89,51.97) node  [font=\small,color={rgb, 255:red, 0; green, 0; blue, 0 }  ,opacity=1 ] [align=left] {\begin{minipage}[lt]{54.86pt}\setlength\topsep{0pt}
\begin{center}
\hyperref[sec:equivalencia]{Sect. 3.3}
\end{center}

\end{minipage}};
% Text Node
\draw (69.5,15.5) node  [font=\small,color={rgb, 255:red, 0; green, 0; blue, 0 }  ,opacity=1 ] [align=left] {\begin{minipage}[lt]{41.55pt}\setlength\topsep{0pt}
\hyperref[ch:programas-python-maquinas-turing]{Chapter 3}
\end{minipage}};
% Text Node
\draw (330,64.85) node  [color={rgb, 255:red, 0; green, 0; blue, 0 }  ,opacity=1 ] [align=left] {\begin{minipage}[lt]{92.93pt}\setlength\topsep{0pt}
\begin{center}
Turing machine\\
\small{(\hyperref[sec:maquinas-turing]{Section 3.2})}
\end{center}

\end{minipage}};
% Text Node
\draw (539.6,314.95) node  [color={rgb, 255:red, 0; green, 0; blue, 0 }  ,opacity=1 ] [align=left] {\begin{minipage}[lt]{92.93pt}\setlength\topsep{0pt}
\begin{center}
logical system\\
\small{(\hyperref[sec:sistemas-logicos]{Section 5.2})}
\end{center}

\end{minipage}};
% Text Node
\draw (380.5,500.11) node  [color={rgb, 255:red, 0; green, 0; blue, 0 }  ,opacity=1 ] [align=left] {\begin{minipage}[lt]{159.55pt}\setlength\topsep{0pt}
\begin{center}
\textbf{First Incompleteness Theorem}
\end{center}

\end{minipage}};
% Text Node
\draw (620.53,64.85) node  [color={rgb, 255:red, 0; green, 0; blue, 0 }  ,opacity=1 ] [align=left] {\begin{minipage}[lt]{106.99pt}\setlength\topsep{0pt}
\begin{center}
Church-Turing thesis\\
\small{(\hyperref[sec:church-turing]{Section 3.4})}
\end{center}

\end{minipage}};
% Text Node
\draw (480.28,52.8) node  [font=\small,color={rgb, 255:red, 0; green, 0; blue, 0 }  ,opacity=1 ] [align=left] {\begin{minipage}[lt]{81.9pt}\setlength\topsep{0pt}
\begin{center}
\hyperref[sec:church-turing]{Section 3.4}
\end{center}

\end{minipage}};
% Text Node
\draw (589.5,518.3) node  [font=\small,color={rgb, 255:red, 0; green, 0; blue, 0 }  ,opacity=1 ] [align=left] {\begin{minipage}[lt]{54.86pt}\setlength\topsep{0pt}
\begin{center}
\hyperref[subsec:mente-humana-tesis-church-turing]{Sect. 6.4.2}
\end{center}

\end{minipage}};
% Text Node
\draw (220.2,159.62) node  [color={rgb, 255:red, 0; green, 0; blue, 0 }  ,opacity=1 ] [align=left] {\begin{minipage}[lt]{230.11pt}\setlength\topsep{0pt}
\begin{center}
computational problem
\end{center}

\end{minipage}};
% Text Node
\draw (98.02,204.85) node  [color={rgb, 255:red, 0; green, 0; blue, 0 }  ,opacity=1 ] [align=left] {\begin{minipage}[lt]{76.13pt}\setlength\topsep{0pt}
\begin{center}
decidable\\
\small{(\hyperref[sec:problemas-decidibles]{Section 4.1})}
\end{center}

\end{minipage}};
% Text Node
\draw (220.67,204.85) node  [color={rgb, 255:red, 0; green, 0; blue, 0 }  ,opacity=1 ] [align=left] {\begin{minipage}[lt]{76.13pt}\setlength\topsep{0pt}
\begin{center}
undecidable\\
\small{(\hyperref[sec:problemas-no-decidibles]{Section 4.3})}
\end{center}

\end{minipage}};
% Text Node
\draw (343.31,204.85) node  [color={rgb, 255:red, 0; green, 0; blue, 0 }  ,opacity=1 ] [align=left] {\begin{minipage}[lt]{76.13pt}\setlength\topsep{0pt}
\begin{center}
semidecidable\\
\small{(\hyperref[sec:problemas-semidecidibles]{Section 4.4})}
\end{center}

\end{minipage}};
% Text Node
\draw (110.4,314.95) node  [color={rgb, 255:red, 0; green, 0; blue, 0 }  ,opacity=1 ] [align=left] {\begin{minipage}[lt]{92.93pt}\setlength\topsep{0pt}
\begin{center}
universality\\
\small{(\hyperref[sec:universalidad]{Section 4.2})}
\end{center}

\end{minipage}};
% Text Node
\draw (330,314.85) node  [color={rgb, 255:red, 0; green, 0; blue, 0 }  ,opacity=1 ] [align=left] {\begin{minipage}[lt]{92.93pt}\setlength\topsep{0pt}
\begin{center}
reduction\\
\small{(\hyperref[sec:reducciones]{Section 4.5})}
\end{center}

\end{minipage}};
% Text Node
\draw (220,414.85) node  [color={rgb, 255:red, 0; green, 0; blue, 0 }  ,opacity=1 ] [align=left] {\begin{minipage}[lt]{107.26pt}\setlength\topsep{0pt}
\begin{center}
halting problem\\
\small{(\hyperref[sec:problema-parada]{Section 4.6})}
\end{center}

\end{minipage}};
% Text Node
\draw (540.33,415.05) node  [color={rgb, 255:red, 0; green, 0; blue, 0 }  ,opacity=1 ] [align=left] {\begin{minipage}[lt]{107.26pt}\setlength\topsep{0pt}
\begin{center}
Peano arithmetic\\
\small{(\hyperref[sec:aritmetica-peano]{Section 5.4})}
\end{center}

\end{minipage}};
% Text Node
\draw (698.8,314.95) node  [color={rgb, 255:red, 0; green, 0; blue, 0 }  ,opacity=1 ] [align=left] {\begin{minipage}[lt]{92.93pt}\setlength\topsep{0pt}
\begin{center}
formal system\\
\small{(\hyperref[sec:sistemas-formales]{Section 5.1})}
\end{center}

\end{minipage}};
% Text Node
\draw (299.6,544.95) node  [color={rgb, 255:red, 0; green, 0; blue, 0 }  ,opacity=1 ] [align=left] {\begin{minipage}[lt]{92.93pt}\setlength\topsep{0pt}
\begin{center}
semantic version\\
\small{(\hyperref[sec:version-semantica]{Section 6.2})}
\end{center}

\end{minipage}};
% Text Node
\draw (459.6,545.05) node  [color={rgb, 255:red, 0; green, 0; blue, 0 }  ,opacity=1 ] [align=left] {\begin{minipage}[lt]{92.93pt}\setlength\topsep{0pt}
\begin{center}
syntactic version\\
\small{(\hyperref[sec:version-sintactica]{Section 6.3})}
\end{center}

\end{minipage}};
% Text Node
\draw (60.5,464.5) node  [font=\small,color={rgb, 255:red, 0; green, 0; blue, 0 }  ,opacity=1 ] [align=left] {\begin{minipage}[lt]{41.55pt}\setlength\topsep{0pt}
\hyperref[ch:problema-parada]{Chapter 4}
\end{minipage}};
% Text Node
\draw (749.5,465.5) node  [font=\small,color={rgb, 255:red, 0; green, 0; blue, 0 }  ,opacity=1 ] [align=left] {\begin{minipage}[lt]{41.55pt}\setlength\topsep{0pt}
\begin{flushright}
\hyperref[ch:sistemas-logicos]{Chapter 5}
\end{flushright}

\end{minipage}};
% Text Node
\draw (170.5,604.5) node  [font=\small,color={rgb, 255:red, 0; green, 0; blue, 0 }  ,opacity=1 ] [align=left] {\begin{minipage}[lt]{41.55pt}\setlength\topsep{0pt}
\hyperref[ch:teorema-incompletitud]{Chapter 6}
\end{minipage}};
% Text Node
\draw (658.76,195.55) node  [color={rgb, 255:red, 0; green, 0; blue, 0 }  ,opacity=1 ] [align=left] {\begin{minipage}[lt]{217.92pt}\setlength\topsep{0pt}
-- sound {\small{(\hyperref[sec:solidez-completitud-decidibilidad]{Section 5.3})}}\\
-- decidable {\small{(\hyperref[sec:solidez-completitud-decidibilidad]{Section 5.3})}}\\
-- (semantically) complete {\small{(\hyperref[sec:solidez-completitud-decidibilidad]{Section 5.3})}}\\
-- syntactically complete {\small{(\hyperref[sec:consistencia]{Section 5.5})}}\\
-- sem./synt. consistent {\small{(\hyperref[sec:consistencia]{Section 5.5})}}
\end{minipage}};
% Text Node
\draw (414.84,645.25) node  [color={rgb, 255:red, 0; green, 0; blue, 0 }  ,opacity=1 ] [align=left] {\begin{minipage}[lt]{534.7pt}\setlength\topsep{0pt}
\begin{center}
Summary map of this Bachelor's Thesis
\end{center}
\end{minipage}};
% Text Node
\draw (380.15,407.8) node  [font=\small,color={rgb, 255:red, 0; green, 0; blue, 0 }  ,opacity=1 ] [align=left] {\begin{minipage}[lt]{108.93pt}\setlength\topsep{0pt}
\begin{center}
\hyperref[sec:primera-aproximacion]{Section 6.1}
\end{center}

\end{minipage}};

\end{tikzpicture}}
%}
%\caption{Esquema resumen (extendido)}

%\label{fig:resumen-extendido}
%\end{figure}
% ====================
%\end{landscape}

% Al finalizar el resumen en inglés, volvemos a seleccionar el idioma español para el documento
\selectlanguage{spanish} 
\endinput

% !TeX root = ../libro.tex
% !TeX encoding = utf8
%
%*******************************************************
% Introducción
%*******************************************************

% \manualmark
% \markboth{\textsc{Introducción}}{\textsc{Introducción}} 
\chapter{Resumen}

%De acuerdo con la comisión de grado, el TFG debe incluir una introducción en la que se describan claramente los objetivos previstos inicialmente en la propuesta de TFG, indicando si han sido o no alcanzados, los antecedentes importantes para el desarrollo, los resultados obtenidos, en su caso y las principales fuentes consultadas.

%En este capítulo, resumiremos el proceso que hemos seguido hasta llegar a la demostración, desde una perspectiva algorítmica, del Primer Teorema de Incompletitud. En la \cref{fig:resumen}, se puede observar un esquema resumen de este trabajo, el cual es extendido en la \cref{fig:resumen-extendido}. En este resumen final se omiten las \cref{part:introduccion,,part:conclusion}.

Comenzamos este Trabajo de Fin de Grado con el \cref{ch:motivacion}, que destaca la motivación que subyace este trabajo, así como los objetivos que se pretenden cumplir. Inmediatamente después, en el \cref{ch:historia}, hacemos una breve contextualización histórica de esta obra.

El desarrollo teórico comienza con el \cref{ch:programas-python-maquinas-turing}, en cuya \cref{sec:programas-python} introducimos una clase específica de programas en Python, los \emph{programas SISO}, y observamos que esta clase es tan expresiva como la clase de todos los programas en Python. A continuación, en la \cref{sec:maquinas-turing}, profundizamos en las \emph{máquinas de Turing}, herramienta matemática principal de nuestro trabajo. En la \cref{sec:equivalencia}, introducimos nuevos modelos de computación equivalentes, hasta llegar al resultado principal del capítulo: que las máquinas de Turing pueden simular programas en Python, y viceversa; en otras palabras, que los programas en Python y las máquinas de Turing son equivalentes. Esto nos permitirá usar programas en Python convenientemente, sin realizar descripciones de máquinas de Turing. Concluimos introduciendo la tesis de Church-Turing en la \cref{sec:church-turing}, y viendo que la equivalencia probada puede ser un argumento a su favor.

En la \cref{sec:problemas-decidibles} definimos los \emph{problemas computacionales}, y observamos que los programas en Python (y, equivalentemente, las máquinas de Turing), pueden resolverlos. Llamamos \emph{decidibles} a aquellos problemas que pueden resolverse mediante un programa que nunca para. En la \cref{sec:universalidad} observamos que un programa puede simularse a sí mismo, lo cual nos conduce en la \cref{sec:problemas-no-decidibles} a la existencia de problemas no decidibles. En la \cref{sec:problemas-semidecidibles} relajamos la condición de parada para obtener problemas \emph{semidecidibles}. A continuación, introducimos en la \cref{sec:reducciones} la técnica de \emph{reducción}, que nos permite obtener nuevos problemas no decidibles, incluyendo el más importante del trabajo en la \cref{sec:problema-parada}: el \emph{problema de la parada}.

Para demostrar el teorema principal de este trabajo, es necesario formalizar los conceptos de \emph{demostración} y \emph{verdad}. Hacemos esto introduciendo los \emph{sistemas formales} en la \cref{sec:sistemas-formales} y los \emph{sistemas lógicos} en la \cref{sec:sistemas-logicos}. Los sistemas formales son \emph{sintácticos}: formalizan qué es un \emph{teorema}, una \emph{fórmula} demostrable a partir de un conjunto de \emph{axiomas} y \emph{reglas de inferencia}. Por otro lado, los sistemas lógicos incorporan a los sistemas formales un componente semántico: la \emph{verdad}. En las \cref{sec:solidez-completitud-decidibilidad,,sec:consistencia} presentamos propiedades esenciales de los sistemas mencionados, que relacionan los conceptos semánticos y sintácticos. Destacamos en la \cref{sec:aritmetica-peano} la introducción de la \emph{aritmética de Peano}: un sistema lógico de especial importancia.

Como culmen, en el \cref{ch:teorema-incompletitud} demostramos el Primer Teorema de Incompletitud: primero desde la aritmética de Peano en la \cref{sec:primera-aproximacion}, para luego generalizarlo mediante dos versiones. La primera, semántica, es demostrada en la \cref{sec:version-semantica}, y nos permite (bajo unas hipótesis computacionales) encontrar afirmaciones ciertas e indemostrables en sistemas lógicos. A continuación, la versión sintáctica de la \cref{sec:version-sintactica}, más cercana a la de Gödel, nos conduce a la existencia de sentencias indemostrables en sistemas formales. Finalizamos con la \cref{sec:consecuencias}, estudiando importantes consecuencias matemáticas (\cref{subsec:segundo-problema-hilbert}) y filosóficas (\cref{subsec:mente-humana-tesis-church-turing}).

Finalmente, terminamos con algunas conclusiones en el \cref{ch:conclusiones}, y dedicando unas palabras finales para el trabajo posterior en el \cref{ch:trabajo-futuro}.


%\vfill
% ====================
%\begin{figure}[H]
%\centering
%\vspace{8pt}
%\hspace*{-1cm}
%\centerline{% !TeX root = ../libro.tex
% !TeX encoding = utf8
%
%*******************************************************
% Introducción
%*******************************************************

% \manualmark
% \markboth{\textsc{Introducción}}{\textsc{Introducción}} 
\chapter{Resumen}

%De acuerdo con la comisión de grado, el TFG debe incluir una introducción en la que se describan claramente los objetivos previstos inicialmente en la propuesta de TFG, indicando si han sido o no alcanzados, los antecedentes importantes para el desarrollo, los resultados obtenidos, en su caso y las principales fuentes consultadas.

%En este capítulo, resumiremos el proceso que hemos seguido hasta llegar a la demostración, desde una perspectiva algorítmica, del Primer Teorema de Incompletitud. En la \cref{fig:resumen}, se puede observar un esquema resumen de este trabajo, el cual es extendido en la \cref{fig:resumen-extendido}. En este resumen final se omiten las \cref{part:introduccion,,part:conclusion}.

Comenzamos este Trabajo de Fin de Grado con el \cref{ch:motivacion}, que destaca la motivación que subyace este trabajo, así como los objetivos que se pretenden cumplir. Inmediatamente después, en el \cref{ch:historia}, hacemos una breve contextualización histórica de esta obra.

El desarrollo teórico comienza con el \cref{ch:programas-python-maquinas-turing}, en cuya \cref{sec:programas-python} introducimos una clase específica de programas en Python, los \emph{programas SISO}, y observamos que esta clase es tan expresiva como la clase de todos los programas en Python. A continuación, en la \cref{sec:maquinas-turing}, profundizamos en las \emph{máquinas de Turing}, herramienta matemática principal de nuestro trabajo. En la \cref{sec:equivalencia}, introducimos nuevos modelos de computación equivalentes, hasta llegar al resultado principal del capítulo: que las máquinas de Turing pueden simular programas en Python, y viceversa; en otras palabras, que los programas en Python y las máquinas de Turing son equivalentes. Esto nos permitirá usar programas en Python convenientemente, sin realizar descripciones de máquinas de Turing. Concluimos introduciendo la tesis de Church-Turing en la \cref{sec:church-turing}, y viendo que la equivalencia probada puede ser un argumento a su favor.

En la \cref{sec:problemas-decidibles} definimos los \emph{problemas computacionales}, y observamos que los programas en Python (y, equivalentemente, las máquinas de Turing), pueden resolverlos. Llamamos \emph{decidibles} a aquellos problemas que pueden resolverse mediante un programa que nunca para. En la \cref{sec:universalidad} observamos que un programa puede simularse a sí mismo, lo cual nos conduce en la \cref{sec:problemas-no-decidibles} a la existencia de problemas no decidibles. En la \cref{sec:problemas-semidecidibles} relajamos la condición de parada para obtener problemas \emph{semidecidibles}. A continuación, introducimos en la \cref{sec:reducciones} la técnica de \emph{reducción}, que nos permite obtener nuevos problemas no decidibles, incluyendo el más importante del trabajo en la \cref{sec:problema-parada}: el \emph{problema de la parada}.

Para demostrar el teorema principal de este trabajo, es necesario formalizar los conceptos de \emph{demostración} y \emph{verdad}. Hacemos esto introduciendo los \emph{sistemas formales} en la \cref{sec:sistemas-formales} y los \emph{sistemas lógicos} en la \cref{sec:sistemas-logicos}. Los sistemas formales son \emph{sintácticos}: formalizan qué es un \emph{teorema}, una \emph{fórmula} demostrable a partir de un conjunto de \emph{axiomas} y \emph{reglas de inferencia}. Por otro lado, los sistemas lógicos incorporan a los sistemas formales un componente semántico: la \emph{verdad}. En las \cref{sec:solidez-completitud-decidibilidad,,sec:consistencia} presentamos propiedades esenciales de los sistemas mencionados, que relacionan los conceptos semánticos y sintácticos. Destacamos en la \cref{sec:aritmetica-peano} la introducción de la \emph{aritmética de Peano}: un sistema lógico de especial importancia.

Como culmen, en el \cref{ch:teorema-incompletitud} demostramos el Primer Teorema de Incompletitud: primero desde la aritmética de Peano en la \cref{sec:primera-aproximacion}, para luego generalizarlo mediante dos versiones. La primera, semántica, es demostrada en la \cref{sec:version-semantica}, y nos permite (bajo unas hipótesis computacionales) encontrar afirmaciones ciertas e indemostrables en sistemas lógicos. A continuación, la versión sintáctica de la \cref{sec:version-sintactica}, más cercana a la de Gödel, nos conduce a la existencia de sentencias indemostrables en sistemas formales. Finalizamos con la \cref{sec:consecuencias}, estudiando importantes consecuencias matemáticas (\cref{subsec:segundo-problema-hilbert}) y filosóficas (\cref{subsec:mente-humana-tesis-church-turing}).

Finalmente, terminamos con algunas conclusiones en el \cref{ch:conclusiones}, y dedicando unas palabras finales para el trabajo posterior en el \cref{ch:trabajo-futuro}.


%\vfill
% ====================
%\begin{figure}[H]
%\centering
%\vspace{8pt}
%\hspace*{-1cm}
%\centerline{% !TeX root = ../libro.tex
% !TeX encoding = utf8
%
%*******************************************************
% Introducción
%*******************************************************

% \manualmark
% \markboth{\textsc{Introducción}}{\textsc{Introducción}} 
\chapter{Resumen}

%De acuerdo con la comisión de grado, el TFG debe incluir una introducción en la que se describan claramente los objetivos previstos inicialmente en la propuesta de TFG, indicando si han sido o no alcanzados, los antecedentes importantes para el desarrollo, los resultados obtenidos, en su caso y las principales fuentes consultadas.

%En este capítulo, resumiremos el proceso que hemos seguido hasta llegar a la demostración, desde una perspectiva algorítmica, del Primer Teorema de Incompletitud. En la \cref{fig:resumen}, se puede observar un esquema resumen de este trabajo, el cual es extendido en la \cref{fig:resumen-extendido}. En este resumen final se omiten las \cref{part:introduccion,,part:conclusion}.

Comenzamos este Trabajo de Fin de Grado con el \cref{ch:motivacion}, que destaca la motivación que subyace este trabajo, así como los objetivos que se pretenden cumplir. Inmediatamente después, en el \cref{ch:historia}, hacemos una breve contextualización histórica de esta obra.

El desarrollo teórico comienza con el \cref{ch:programas-python-maquinas-turing}, en cuya \cref{sec:programas-python} introducimos una clase específica de programas en Python, los \emph{programas SISO}, y observamos que esta clase es tan expresiva como la clase de todos los programas en Python. A continuación, en la \cref{sec:maquinas-turing}, profundizamos en las \emph{máquinas de Turing}, herramienta matemática principal de nuestro trabajo. En la \cref{sec:equivalencia}, introducimos nuevos modelos de computación equivalentes, hasta llegar al resultado principal del capítulo: que las máquinas de Turing pueden simular programas en Python, y viceversa; en otras palabras, que los programas en Python y las máquinas de Turing son equivalentes. Esto nos permitirá usar programas en Python convenientemente, sin realizar descripciones de máquinas de Turing. Concluimos introduciendo la tesis de Church-Turing en la \cref{sec:church-turing}, y viendo que la equivalencia probada puede ser un argumento a su favor.

En la \cref{sec:problemas-decidibles} definimos los \emph{problemas computacionales}, y observamos que los programas en Python (y, equivalentemente, las máquinas de Turing), pueden resolverlos. Llamamos \emph{decidibles} a aquellos problemas que pueden resolverse mediante un programa que nunca para. En la \cref{sec:universalidad} observamos que un programa puede simularse a sí mismo, lo cual nos conduce en la \cref{sec:problemas-no-decidibles} a la existencia de problemas no decidibles. En la \cref{sec:problemas-semidecidibles} relajamos la condición de parada para obtener problemas \emph{semidecidibles}. A continuación, introducimos en la \cref{sec:reducciones} la técnica de \emph{reducción}, que nos permite obtener nuevos problemas no decidibles, incluyendo el más importante del trabajo en la \cref{sec:problema-parada}: el \emph{problema de la parada}.

Para demostrar el teorema principal de este trabajo, es necesario formalizar los conceptos de \emph{demostración} y \emph{verdad}. Hacemos esto introduciendo los \emph{sistemas formales} en la \cref{sec:sistemas-formales} y los \emph{sistemas lógicos} en la \cref{sec:sistemas-logicos}. Los sistemas formales son \emph{sintácticos}: formalizan qué es un \emph{teorema}, una \emph{fórmula} demostrable a partir de un conjunto de \emph{axiomas} y \emph{reglas de inferencia}. Por otro lado, los sistemas lógicos incorporan a los sistemas formales un componente semántico: la \emph{verdad}. En las \cref{sec:solidez-completitud-decidibilidad,,sec:consistencia} presentamos propiedades esenciales de los sistemas mencionados, que relacionan los conceptos semánticos y sintácticos. Destacamos en la \cref{sec:aritmetica-peano} la introducción de la \emph{aritmética de Peano}: un sistema lógico de especial importancia.

Como culmen, en el \cref{ch:teorema-incompletitud} demostramos el Primer Teorema de Incompletitud: primero desde la aritmética de Peano en la \cref{sec:primera-aproximacion}, para luego generalizarlo mediante dos versiones. La primera, semántica, es demostrada en la \cref{sec:version-semantica}, y nos permite (bajo unas hipótesis computacionales) encontrar afirmaciones ciertas e indemostrables en sistemas lógicos. A continuación, la versión sintáctica de la \cref{sec:version-sintactica}, más cercana a la de Gödel, nos conduce a la existencia de sentencias indemostrables en sistemas formales. Finalizamos con la \cref{sec:consecuencias}, estudiando importantes consecuencias matemáticas (\cref{subsec:segundo-problema-hilbert}) y filosóficas (\cref{subsec:mente-humana-tesis-church-turing}).

Finalmente, terminamos con algunas conclusiones en el \cref{ch:conclusiones}, y dedicando unas palabras finales para el trabajo posterior en el \cref{ch:trabajo-futuro}.


%\vfill
% ====================
%\begin{figure}[H]
%\centering
%\vspace{8pt}
%\hspace*{-1cm}
%\centerline{\input{fig/resumen}}
%\caption{Esquema resumen}
%\label{fig:resumen}
%\end{figure}
% ====================
%\vfill

\begin{landscape}
% ====================
\begin{figure}[H]
%\captionsetup{margin=200cm}

\centering
%\vspace{-1.5cm}
%\hspace{-1cm}
%\input{fig/resumen-extendido}
\raisebox{-15.5cm}[0pt][0pt]{%
\resizebox{\linewidth}{!}{\input{fig/resumen-extendido}}
}
%\caption{Esquema resumen (extendido)}

%\label{fig:resumen-extendido}
\end{figure}
% ====================
\end{landscape}

\endinput
}
%\caption{Esquema resumen}
%\label{fig:resumen}
%\end{figure}
% ====================
%\vfill

\begin{landscape}
% ====================
\begin{figure}[H]
%\captionsetup{margin=200cm}

\centering
%\vspace{-1.5cm}
%\hspace{-1cm}
%

\tikzset{every picture/.style={line width=0.75pt}} %set default line width to 0.75pt        

\begin{tikzpicture}[x=0.75pt,y=0.75pt,yscale=-1,xscale=1]
%uncomment if require: \path (0,694); %set diagram left start at 0, and has height of 694

%Straight Lines [id:da2056540562354816] 
\draw  [dash pattern={on 4.5pt off 4.5pt}]  (672,64.47) -- (810,64.47) ;
%Straight Lines [id:da23044041717189923] 
\draw  [dash pattern={on 4.5pt off 4.5pt}]  (412.67,64.47) -- (610,64.47) ;
\draw [shift={(410.67,64.47)}, rotate = 0] [color={rgb, 255:red, 0; green, 0; blue, 0 }  ][line width=0.75]    (10.93,-3.29) .. controls (6.95,-1.4) and (3.31,-0.3) .. (0,0) .. controls (3.31,0.3) and (6.95,1.4) .. (10.93,3.29)   ;
%Straight Lines [id:da8398435684999113] 
\draw  [dash pattern={on 4.5pt off 4.5pt}]  (810,64.47) -- (810,530) ;
%Rounded Rect [id:dp027018508739026448] 
\draw  [draw opacity=0][fill={rgb, 255:red, 230; green, 230; blue, 230 }  ,fill opacity=0.5 ] (20,439.33) .. controls (20,445.22) and (24.78,450) .. (30.67,450) -- (409.33,450) .. controls (415.22,450) and (420,445.22) .. (420,439.33) -- (420,140.67) .. controls (420,134.78) and (415.22,130) .. (409.33,130) -- (30.67,130) .. controls (24.78,130) and (20,134.78) .. (20,140.67) -- cycle ;
%Rounded Rect [id:dp9415679498487528] 
\draw  [dash pattern={on 4.5pt off 4.5pt}] (29.83,292.77) .. controls (29.83,285.88) and (35.41,280.3) .. (42.3,280.3) -- (398.54,280.3) .. controls (405.42,280.3) and (411,285.88) .. (411,292.77) -- (411,337.17) .. controls (411,344.05) and (405.42,349.64) .. (398.54,349.64) -- (42.3,349.64) .. controls (35.41,349.64) and (29.83,344.05) .. (29.83,337.17) -- cycle ;
%Rounded Rect [id:dp05100022295548956] 
\draw  [draw opacity=0][fill={rgb, 255:red, 230; green, 230; blue, 230 }  ,fill opacity=0.5 ] (130.5,575) .. controls (130.5,583.28) and (137.22,590) .. (145.5,590) -- (615,590) .. controls (623.28,590) and (630,583.28) .. (630,575) -- (630,395) .. controls (630,386.72) and (623.28,380) .. (615,380) -- (145.5,380) .. controls (137.22,380) and (130.5,386.72) .. (130.5,395) -- cycle ;
%Rounded Rect [id:dp6496809750851564] 
\draw  [draw opacity=0][fill={rgb, 255:red, 230; green, 230; blue, 230 }  ,fill opacity=0.5 ] (450.61,440.79) .. controls (450.61,445.82) and (454.69,449.9) .. (459.72,449.9) -- (780.89,449.9) .. controls (785.92,449.9) and (790,445.82) .. (790,440.79) -- (790,139.11) .. controls (790,134.08) and (785.92,130) .. (780.89,130) -- (459.72,130) .. controls (454.69,130) and (450.61,134.08) .. (450.61,139.11) -- cycle ;
%Rounded Rect [id:dp549636265291094] 
\draw  [fill={rgb, 255:red, 230; green, 230; blue, 230 }  ,fill opacity=0.5 ][dash pattern={on 4.5pt off 4.5pt}] (29.83,42.77) .. controls (29.83,35.88) and (35.41,30.3) .. (42.3,30.3) -- (398.54,30.3) .. controls (405.42,30.3) and (411,35.88) .. (411,42.77) -- (411,87.17) .. controls (411,94.05) and (405.42,99.64) .. (398.54,99.64) -- (42.3,99.64) .. controls (35.41,99.64) and (29.83,94.05) .. (29.83,87.17) -- cycle ;
%Rounded Rect [id:dp1378581119146769] 
\draw  [fill={rgb, 255:red, 255; green, 255; blue, 255 }  ,fill opacity=1 ] (40.4,49.82) .. controls (40.4,44.29) and (44.89,39.8) .. (50.42,39.8) -- (169.58,39.8) .. controls (175.11,39.8) and (179.6,44.29) .. (179.6,49.82) -- (179.6,79.88) .. controls (179.6,85.41) and (175.11,89.9) .. (169.58,89.9) -- (50.42,89.9) .. controls (44.89,89.9) and (40.4,85.41) .. (40.4,79.88) -- cycle ;
%Straight Lines [id:da8729153203483555] 
\draw    (188.42,63.77) -- (250.99,63.77)(188.42,66.77) -- (250.99,66.77) ;
\draw [shift={(257.99,65.27)}, rotate = 180] [color={rgb, 255:red, 0; green, 0; blue, 0 }  ][line width=0.75]    (10.93,-4.9) .. controls (6.95,-2.3) and (3.31,-0.67) .. (0,0) .. controls (3.31,0.67) and (6.95,2.3) .. (10.93,4.9)   ;
\draw [shift={(181.42,65.27)}, rotate = 0] [color={rgb, 255:red, 0; green, 0; blue, 0 }  ][line width=0.75]    (10.93,-4.9) .. controls (6.95,-2.3) and (3.31,-0.67) .. (0,0) .. controls (3.31,0.67) and (6.95,2.3) .. (10.93,4.9)   ;
%Straight Lines [id:da845243590120559] 
\draw    (220.17,99.97) -- (220.17,137.64) ;
\draw [shift={(220.17,139.64)}, rotate = 270] [color={rgb, 255:red, 0; green, 0; blue, 0 }  ][line width=0.75]    (10.93,-3.29) .. controls (6.95,-1.4) and (3.31,-0.3) .. (0,0) .. controls (3.31,0.3) and (6.95,1.4) .. (10.93,3.29)   ;
%Straight Lines [id:da7497001546502899] 
\draw    (380,420.33) -- (380,478) ;
\draw [shift={(380,480)}, rotate = 270] [color={rgb, 255:red, 0; green, 0; blue, 0 }  ][line width=0.75]    (10.93,-3.29) .. controls (6.95,-1.4) and (3.31,-0.3) .. (0,0) .. controls (3.31,0.3) and (6.95,1.4) .. (10.93,3.29)   ;
%Straight Lines [id:da5689893955334542] 
\draw    (300,420) -- (460,420) ;
%Rounded Rect [id:dp18582062333838767] 
\draw  [fill={rgb, 255:red, 255; green, 255; blue, 255 }  ,fill opacity=1 ] (260.4,49.82) .. controls (260.4,44.29) and (264.89,39.8) .. (270.42,39.8) -- (389.58,39.8) .. controls (395.11,39.8) and (399.6,44.29) .. (399.6,49.82) -- (399.6,79.88) .. controls (399.6,85.41) and (395.11,89.9) .. (389.58,89.9) -- (270.42,89.9) .. controls (264.89,89.9) and (260.4,85.41) .. (260.4,79.88) -- cycle ;
%Rounded Rect [id:dp9523169763067778] 
\draw  [fill={rgb, 255:red, 255; green, 255; blue, 255 }  ,fill opacity=1 ] (470,299.92) .. controls (470,294.39) and (474.49,289.9) .. (480.02,289.9) -- (599.18,289.9) .. controls (604.71,289.9) and (609.2,294.39) .. (609.2,299.92) -- (609.2,329.98) .. controls (609.2,335.51) and (604.71,340) .. (599.18,340) -- (480.02,340) .. controls (474.49,340) and (470,335.51) .. (470,329.98) -- cycle ;
%Rounded Rect [id:dp48544030815291705] 
\draw  [fill={rgb, 255:red, 255; green, 255; blue, 255 }  ,fill opacity=1 ] (220,490) .. controls (220,484.48) and (224.48,480) .. (230,480) -- (530,480) .. controls (535.52,480) and (540,484.48) .. (540,490) -- (540,570) .. controls (540,575.52) and (535.52,580) .. (530,580) -- (230,580) .. controls (224.48,580) and (220,575.52) .. (220,570) -- cycle ;
%Rounded Rect [id:dp08222627508455749] 
\draw  [fill={rgb, 255:red, 255; green, 255; blue, 255 }  ,fill opacity=1 ] (540.4,49.82) .. controls (540.4,44.29) and (544.89,39.8) .. (550.42,39.8) -- (690.65,39.8) .. controls (696.18,39.8) and (700.67,44.29) .. (700.67,49.82) -- (700.67,79.88) .. controls (700.67,85.41) and (696.18,89.9) .. (690.65,89.9) -- (550.42,89.9) .. controls (544.89,89.9) and (540.4,85.41) .. (540.4,79.88) -- cycle ;
%Straight Lines [id:da5868346364154593] 
\draw  [dash pattern={on 4.5pt off 4.5pt}]  (810,530) -- (542,530) ;
\draw [shift={(540,530)}, rotate = 360] [color={rgb, 255:red, 0; green, 0; blue, 0 }  ][line width=0.75]    (10.93,-3.29) .. controls (6.95,-1.4) and (3.31,-0.3) .. (0,0) .. controls (3.31,0.3) and (6.95,1.4) .. (10.93,3.29)   ;
%Rounded Rect [id:dp7392784471609453] 
\draw  [fill={rgb, 255:red, 255; green, 255; blue, 255 }  ,fill opacity=1 ] (30.33,150.47) .. controls (30.33,144.58) and (35.11,139.8) .. (41,139.8) -- (400.34,139.8) .. controls (406.23,139.8) and (411,144.58) .. (411,150.47) -- (411,229.34) .. controls (411,235.23) and (406.23,240.01) .. (400.34,240.01) -- (41,240.01) .. controls (35.11,240.01) and (30.33,235.23) .. (30.33,229.34) -- cycle ;
%Rounded Rect [id:dp90901637181607] 
\draw  [fill={rgb, 255:red, 255; green, 255; blue, 255 }  ,fill opacity=1 ] (41,189.82) .. controls (41,184.29) and (45.49,179.8) .. (51.02,179.8) -- (145.02,179.8) .. controls (150.55,179.8) and (155.04,184.29) .. (155.04,189.82) -- (155.04,219.88) .. controls (155.04,225.41) and (150.55,229.9) .. (145.02,229.9) -- (51.02,229.9) .. controls (45.49,229.9) and (41,225.41) .. (41,219.88) -- cycle ;
%Rounded Rect [id:dp06115018179625542] 
\draw  [fill={rgb, 255:red, 255; green, 255; blue, 255 }  ,fill opacity=1 ] (163.65,189.82) .. controls (163.65,184.29) and (168.13,179.8) .. (173.67,179.8) -- (267.67,179.8) .. controls (273.2,179.8) and (277.69,184.29) .. (277.69,189.82) -- (277.69,219.88) .. controls (277.69,225.41) and (273.2,229.9) .. (267.67,229.9) -- (173.67,229.9) .. controls (168.13,229.9) and (163.65,225.41) .. (163.65,219.88) -- cycle ;
%Rounded Rect [id:dp6195275548153432] 
\draw  [fill={rgb, 255:red, 255; green, 255; blue, 255 }  ,fill opacity=1 ] (286.3,189.82) .. controls (286.3,184.29) and (290.78,179.8) .. (296.32,179.8) -- (390.31,179.8) .. controls (395.85,179.8) and (400.33,184.29) .. (400.33,189.82) -- (400.33,219.88) .. controls (400.33,225.41) and (395.85,229.9) .. (390.31,229.9) -- (296.32,229.9) .. controls (290.78,229.9) and (286.3,225.41) .. (286.3,219.88) -- cycle ;
%Rounded Rect [id:dp41283415100159493] 
\draw  [fill={rgb, 255:red, 255; green, 255; blue, 255 }  ,fill opacity=1 ] (40.8,299.92) .. controls (40.8,294.39) and (45.29,289.9) .. (50.82,289.9) -- (169.98,289.9) .. controls (175.51,289.9) and (180,294.39) .. (180,299.92) -- (180,329.98) .. controls (180,335.51) and (175.51,340) .. (169.98,340) -- (50.82,340) .. controls (45.29,340) and (40.8,335.51) .. (40.8,329.98) -- cycle ;
%Rounded Rect [id:dp40154854661776396] 
\draw  [fill={rgb, 255:red, 255; green, 255; blue, 255 }  ,fill opacity=1 ] (260.4,299.82) .. controls (260.4,294.29) and (264.89,289.8) .. (270.42,289.8) -- (389.58,289.8) .. controls (395.11,289.8) and (399.6,294.29) .. (399.6,299.82) -- (399.6,329.88) .. controls (399.6,335.41) and (395.11,339.9) .. (389.58,339.9) -- (270.42,339.9) .. controls (264.89,339.9) and (260.4,335.41) .. (260.4,329.88) -- cycle ;
%Straight Lines [id:da6602630444249074] 
\draw    (180.17,315.8) -- (260,315.8) ;
%Straight Lines [id:da30905866347109745] 
\draw    (219.86,315.97) -- (219.86,386.81) ;
\draw [shift={(219.86,388.81)}, rotate = 270] [color={rgb, 255:red, 0; green, 0; blue, 0 }  ][line width=0.75]    (10.93,-3.29) .. controls (6.95,-1.4) and (3.31,-0.3) .. (0,0) .. controls (3.31,0.3) and (6.95,1.4) .. (10.93,3.29)   ;
%Rounded Rect [id:dp44627062039563214] 
\draw  [fill={rgb, 255:red, 255; green, 255; blue, 255 }  ,fill opacity=1 ] (139.67,399.82) .. controls (139.67,394.29) and (144.15,389.8) .. (149.69,389.8) -- (290.31,389.8) .. controls (295.85,389.8) and (300.33,394.29) .. (300.33,399.82) -- (300.33,429.88) .. controls (300.33,435.41) and (295.85,439.9) .. (290.31,439.9) -- (149.69,439.9) .. controls (144.15,439.9) and (139.67,435.41) .. (139.67,429.88) -- cycle ;
%Straight Lines [id:da8375870493480149] 
\draw    (220,240) -- (220,277.67) ;
\draw [shift={(220,279.67)}, rotate = 270] [color={rgb, 255:red, 0; green, 0; blue, 0 }  ][line width=0.75]    (10.93,-3.29) .. controls (6.95,-1.4) and (3.31,-0.3) .. (0,0) .. controls (3.31,0.3) and (6.95,1.4) .. (10.93,3.29)   ;
%Rounded Rect [id:dp37285624114093485] 
\draw  [fill={rgb, 255:red, 255; green, 255; blue, 255 }  ,fill opacity=1 ] (460,400.02) .. controls (460,394.49) and (464.49,390) .. (470.02,390) -- (610.65,390) .. controls (616.18,390) and (620.67,394.49) .. (620.67,400.02) -- (620.67,430.08) .. controls (620.67,435.61) and (616.18,440.1) .. (610.65,440.1) -- (470.02,440.1) .. controls (464.49,440.1) and (460,435.61) .. (460,430.08) -- cycle ;
%Rounded Rect [id:dp6883616335272991] 
\draw  [fill={rgb, 255:red, 255; green, 255; blue, 255 }  ,fill opacity=1 ] (629.2,299.92) .. controls (629.2,294.39) and (633.69,289.9) .. (639.22,289.9) -- (758.38,289.9) .. controls (763.91,289.9) and (768.4,294.39) .. (768.4,299.92) -- (768.4,329.98) .. controls (768.4,335.51) and (763.91,340) .. (758.38,340) -- (639.22,340) .. controls (633.69,340) and (629.2,335.51) .. (629.2,329.98) -- cycle ;
%Straight Lines [id:da19870419404011752] 
\draw    (540,340) -- (540,388) ;
\draw [shift={(540,390)}, rotate = 270] [color={rgb, 255:red, 0; green, 0; blue, 0 }  ][line width=0.75]    (10.93,-3.29) .. controls (6.95,-1.4) and (3.31,-0.3) .. (0,0) .. controls (3.31,0.3) and (6.95,1.4) .. (10.93,3.29)   ;
%Rounded Rect [id:dp9692784599420237] 
\draw  [fill={rgb, 255:red, 255; green, 255; blue, 255 }  ,fill opacity=1 ] (230,529.92) .. controls (230,524.39) and (234.49,519.9) .. (240.02,519.9) -- (359.18,519.9) .. controls (364.71,519.9) and (369.2,524.39) .. (369.2,529.92) -- (369.2,559.98) .. controls (369.2,565.51) and (364.71,570) .. (359.18,570) -- (240.02,570) .. controls (234.49,570) and (230,565.51) .. (230,559.98) -- cycle ;
%Rounded Rect [id:dp749684165723072] 
\draw  [fill={rgb, 255:red, 255; green, 255; blue, 255 }  ,fill opacity=1 ] (390,530.02) .. controls (390,524.49) and (394.49,520) .. (400.02,520) -- (519.18,520) .. controls (524.71,520) and (529.2,524.49) .. (529.2,530.02) -- (529.2,560.08) .. controls (529.2,565.61) and (524.71,570.1) .. (519.18,570.1) -- (400.02,570.1) .. controls (394.49,570.1) and (390,565.61) .. (390,560.08) -- cycle ;
%Rounded Rect [id:dp28001422897762485] 
\draw  [fill={rgb, 255:red, 255; green, 255; blue, 255 }  ,fill opacity=1 ][dash pattern={on 4.5pt off 4.5pt}] (460,148.92) .. controls (460,143.99) and (463.99,140) .. (468.92,140) -- (771.08,140) .. controls (776.01,140) and (780,143.99) .. (780,148.92) -- (780,241.55) .. controls (780,246.48) and (776.01,250.47) .. (771.08,250.47) -- (468.92,250.47) .. controls (463.99,250.47) and (460,246.48) .. (460,241.55) -- cycle ;
%Shape: Brace [id:dp9411387904179627] 
\draw   (774.67,286.8) .. controls (774.67,282.13) and (772.34,279.8) .. (767.67,279.8) -- (629.96,279.8) .. controls (623.29,279.8) and (619.96,277.47) .. (619.96,272.8) .. controls (619.96,277.47) and (616.63,279.8) .. (609.96,279.8)(612.96,279.8) -- (472.25,279.8) .. controls (467.58,279.8) and (465.25,282.13) .. (465.25,286.8) ;
%Straight Lines [id:da9902399507274253] 
\draw    (620,273) -- (620,253.14) ;
\draw [shift={(620,251.14)}, rotate = 90] [color={rgb, 255:red, 0; green, 0; blue, 0 }  ][line width=0.75]    (10.93,-3.29) .. controls (6.95,-1.4) and (3.31,-0.3) .. (0,0) .. controls (3.31,0.3) and (6.95,1.4) .. (10.93,3.29)   ;

% Text Node
\draw (110,64.85) node  [color={rgb, 255:red, 0; green, 0; blue, 0 }  ,opacity=1 ] [align=left] {\begin{minipage}[lt]{92.93pt}\setlength\topsep{0pt}
\begin{center}
programa en Python\\
\small{(\cref{sec:programas-python})}
\end{center}

\end{minipage}};
% Text Node
\draw (219.89,51.97) node  [font=\small,color={rgb, 255:red, 0; green, 0; blue, 0 }  ,opacity=1 ] [align=left] {\begin{minipage}[lt]{54.86pt}\setlength\topsep{0pt}
\begin{center}
\Cref{sec:equivalencia}
\end{center}

\end{minipage}};
% Text Node
\draw (69.5,15.5) node  [font=\small,color={rgb, 255:red, 0; green, 0; blue, 0 }  ,opacity=1 ] [align=left] {\begin{minipage}[lt]{41.55pt}\setlength\topsep{0pt}
\cref{ch:programas-python-maquinas-turing}
\end{minipage}};
% Text Node
\draw (330,64.85) node  [color={rgb, 255:red, 0; green, 0; blue, 0 }  ,opacity=1 ] [align=left] {\begin{minipage}[lt]{92.93pt}\setlength\topsep{0pt}
\begin{center}
máquina de Turing\\
\small{(\cref{sec:maquinas-turing})}
\end{center}

\end{minipage}};
% Text Node
\draw (539.6,314.95) node  [color={rgb, 255:red, 0; green, 0; blue, 0 }  ,opacity=1 ] [align=left] {\begin{minipage}[lt]{92.93pt}\setlength\topsep{0pt}
\begin{center}
sistema lógico\\
\small{(\cref{sec:sistemas-logicos})}
\end{center}

\end{minipage}};
% Text Node
\draw (380.5,500.11) node  [color={rgb, 255:red, 0; green, 0; blue, 0 }  ,opacity=1 ] [align=left] {\begin{minipage}[lt]{159.55pt}\setlength\topsep{0pt}
\begin{center}
\textbf{Primer Teorema de Incompletitud}
\end{center}

\end{minipage}};
% Text Node
\draw (620.53,64.85) node  [color={rgb, 255:red, 0; green, 0; blue, 0 }  ,opacity=1 ] [align=left] {\begin{minipage}[lt]{106.99pt}\setlength\topsep{0pt}
\begin{center}
tesis de Church-Turing\\
\small{(\cref{sec:church-turing})}
\end{center}

\end{minipage}};
% Text Node
\draw (480.28,52.8) node  [font=\small,color={rgb, 255:red, 0; green, 0; blue, 0 }  ,opacity=1 ] [align=left] {\begin{minipage}[lt]{81.9pt}\setlength\topsep{0pt}
\begin{center}
\cref{sec:church-turing}
\end{center}

\end{minipage}};
% Text Node
\draw (589.5,518.3) node  [font=\small,color={rgb, 255:red, 0; green, 0; blue, 0 }  ,opacity=1 ] [align=left] {\begin{minipage}[lt]{54.86pt}\setlength\topsep{0pt}
\begin{center}
\Cref{subsec:mente-humana-tesis-church-turing}
\end{center}

\end{minipage}};
% Text Node
\draw (220.2,159.62) node  [color={rgb, 255:red, 0; green, 0; blue, 0 }  ,opacity=1 ] [align=left] {\begin{minipage}[lt]{230.11pt}\setlength\topsep{0pt}
\begin{center}
problema computacional
\end{center}

\end{minipage}};
% Text Node
\draw (98.02,204.85) node  [color={rgb, 255:red, 0; green, 0; blue, 0 }  ,opacity=1 ] [align=left] {\begin{minipage}[lt]{76.13pt}\setlength\topsep{0pt}
\begin{center}
decidible\\
\small{(\cref{sec:problemas-decidibles})}
\end{center}

\end{minipage}};
% Text Node
\draw (220.67,204.85) node  [color={rgb, 255:red, 0; green, 0; blue, 0 }  ,opacity=1 ] [align=left] {\begin{minipage}[lt]{76.13pt}\setlength\topsep{0pt}
\begin{center}
no decidible\\
\small{(\cref{sec:problemas-no-decidibles})}
\end{center}

\end{minipage}};
% Text Node
\draw (343.31,204.85) node  [color={rgb, 255:red, 0; green, 0; blue, 0 }  ,opacity=1 ] [align=left] {\begin{minipage}[lt]{76.13pt}\setlength\topsep{0pt}
\begin{center}
semidecidible\\
\small{(\cref{sec:problemas-semidecidibles})}
\end{center}

\end{minipage}};
% Text Node
\draw (110.4,314.95) node  [color={rgb, 255:red, 0; green, 0; blue, 0 }  ,opacity=1 ] [align=left] {\begin{minipage}[lt]{92.93pt}\setlength\topsep{0pt}
\begin{center}
universalidad\\
\small{(\cref{sec:universalidad})}
\end{center}

\end{minipage}};
% Text Node
\draw (330,314.85) node  [color={rgb, 255:red, 0; green, 0; blue, 0 }  ,opacity=1 ] [align=left] {\begin{minipage}[lt]{92.93pt}\setlength\topsep{0pt}
\begin{center}
reducción\\
\small{(\cref{sec:reducciones})}
\end{center}

\end{minipage}};
% Text Node
\draw (220,414.85) node  [color={rgb, 255:red, 0; green, 0; blue, 0 }  ,opacity=1 ] [align=left] {\begin{minipage}[lt]{107.26pt}\setlength\topsep{0pt}
\begin{center}
problema de la parada\\
\small{(\cref{sec:problema-parada})}
\end{center}

\end{minipage}};
% Text Node
\draw (540.33,415.05) node  [color={rgb, 255:red, 0; green, 0; blue, 0 }  ,opacity=1 ] [align=left] {\begin{minipage}[lt]{107.26pt}\setlength\topsep{0pt}
\begin{center}
aritmética de Peano\\
\small{(\cref{sec:aritmetica-peano})}
\end{center}

\end{minipage}};
% Text Node
\draw (698.8,314.95) node  [color={rgb, 255:red, 0; green, 0; blue, 0 }  ,opacity=1 ] [align=left] {\begin{minipage}[lt]{92.93pt}\setlength\topsep{0pt}
\begin{center}
sistema formal\\
\small{(\cref{sec:sistemas-formales})}
\end{center}

\end{minipage}};
% Text Node
\draw (299.6,544.95) node  [color={rgb, 255:red, 0; green, 0; blue, 0 }  ,opacity=1 ] [align=left] {\begin{minipage}[lt]{92.93pt}\setlength\topsep{0pt}
\begin{center}
versión semántica\\
\small{(\cref{sec:version-semantica})}
\end{center}

\end{minipage}};
% Text Node
\draw (459.6,545.05) node  [color={rgb, 255:red, 0; green, 0; blue, 0 }  ,opacity=1 ] [align=left] {\begin{minipage}[lt]{92.93pt}\setlength\topsep{0pt}
\begin{center}
versión sintáctica\\
\small{(\cref{sec:version-sintactica})}
\end{center}

\end{minipage}};
% Text Node
\draw (60.5,464.5) node  [font=\small,color={rgb, 255:red, 0; green, 0; blue, 0 }  ,opacity=1 ] [align=left] {\begin{minipage}[lt]{41.55pt}\setlength\topsep{0pt}
\cref{ch:problema-parada}
\end{minipage}};
% Text Node
\draw (749.5,465.5) node  [font=\small,color={rgb, 255:red, 0; green, 0; blue, 0 }  ,opacity=1 ] [align=left] {\begin{minipage}[lt]{41.55pt}\setlength\topsep{0pt}
\begin{flushright}
\cref{ch:sistemas-logicos}
\end{flushright}

\end{minipage}};
% Text Node
\draw (170.5,604.5) node  [font=\small,color={rgb, 255:red, 0; green, 0; blue, 0 }  ,opacity=1 ] [align=left] {\begin{minipage}[lt]{41.55pt}\setlength\topsep{0pt}
\cref{ch:teorema-incompletitud}
\end{minipage}};
% Text Node
\draw (643.76,195.55) node  [color={rgb, 255:red, 0; green, 0; blue, 0 }  ,opacity=1 ] [align=left] {\begin{minipage}[lt]{217.92pt}\setlength\topsep{0pt}
-- sólido {\small{(\cref{sec:solidez-completitud-decidibilidad})}}\\
-- decidible {\small{(\cref{sec:solidez-completitud-decidibilidad})}}\\
-- (semánticamente) completo {\small{(\cref{sec:solidez-completitud-decidibilidad})}}\\
-- sintácticamente completo {\small{(\cref{sec:consistencia})}}\\
-- sem./sint. consistente {\small{(\cref{sec:consistencia})}}
\end{minipage}};
% Text Node
\draw (414.84,645.25) node  [color={rgb, 255:red, 0; green, 0; blue, 0 }  ,opacity=1 ] [align=left] {\begin{minipage}[lt]{534.7pt}\setlength\topsep{0pt}
\begin{center}
Esquema resumen de este Trabajo de Fin de Grado
\end{center}
\end{minipage}};
% Text Node
\draw (380.15,407.8) node  [font=\small,color={rgb, 255:red, 0; green, 0; blue, 0 }  ,opacity=1 ] [align=left] {\begin{minipage}[lt]{108.93pt}\setlength\topsep{0pt}
\begin{center}
\cref{sec:primera-aproximacion}
\end{center}

\end{minipage}};

\end{tikzpicture}
\raisebox{-15.5cm}[0pt][0pt]{%
\resizebox{\linewidth}{!}{

\tikzset{every picture/.style={line width=0.75pt}} %set default line width to 0.75pt        

\begin{tikzpicture}[x=0.75pt,y=0.75pt,yscale=-1,xscale=1]
%uncomment if require: \path (0,694); %set diagram left start at 0, and has height of 694

%Straight Lines [id:da2056540562354816] 
\draw  [dash pattern={on 4.5pt off 4.5pt}]  (672,64.47) -- (810,64.47) ;
%Straight Lines [id:da23044041717189923] 
\draw  [dash pattern={on 4.5pt off 4.5pt}]  (412.67,64.47) -- (610,64.47) ;
\draw [shift={(410.67,64.47)}, rotate = 0] [color={rgb, 255:red, 0; green, 0; blue, 0 }  ][line width=0.75]    (10.93,-3.29) .. controls (6.95,-1.4) and (3.31,-0.3) .. (0,0) .. controls (3.31,0.3) and (6.95,1.4) .. (10.93,3.29)   ;
%Straight Lines [id:da8398435684999113] 
\draw  [dash pattern={on 4.5pt off 4.5pt}]  (810,64.47) -- (810,530) ;
%Rounded Rect [id:dp027018508739026448] 
\draw  [draw opacity=0][fill={rgb, 255:red, 230; green, 230; blue, 230 }  ,fill opacity=0.5 ] (20,439.33) .. controls (20,445.22) and (24.78,450) .. (30.67,450) -- (409.33,450) .. controls (415.22,450) and (420,445.22) .. (420,439.33) -- (420,140.67) .. controls (420,134.78) and (415.22,130) .. (409.33,130) -- (30.67,130) .. controls (24.78,130) and (20,134.78) .. (20,140.67) -- cycle ;
%Rounded Rect [id:dp9415679498487528] 
\draw  [dash pattern={on 4.5pt off 4.5pt}] (29.83,292.77) .. controls (29.83,285.88) and (35.41,280.3) .. (42.3,280.3) -- (398.54,280.3) .. controls (405.42,280.3) and (411,285.88) .. (411,292.77) -- (411,337.17) .. controls (411,344.05) and (405.42,349.64) .. (398.54,349.64) -- (42.3,349.64) .. controls (35.41,349.64) and (29.83,344.05) .. (29.83,337.17) -- cycle ;
%Rounded Rect [id:dp05100022295548956] 
\draw  [draw opacity=0][fill={rgb, 255:red, 230; green, 230; blue, 230 }  ,fill opacity=0.5 ] (130.5,575) .. controls (130.5,583.28) and (137.22,590) .. (145.5,590) -- (615,590) .. controls (623.28,590) and (630,583.28) .. (630,575) -- (630,395) .. controls (630,386.72) and (623.28,380) .. (615,380) -- (145.5,380) .. controls (137.22,380) and (130.5,386.72) .. (130.5,395) -- cycle ;
%Rounded Rect [id:dp6496809750851564] 
\draw  [draw opacity=0][fill={rgb, 255:red, 230; green, 230; blue, 230 }  ,fill opacity=0.5 ] (450.61,440.79) .. controls (450.61,445.82) and (454.69,449.9) .. (459.72,449.9) -- (780.89,449.9) .. controls (785.92,449.9) and (790,445.82) .. (790,440.79) -- (790,139.11) .. controls (790,134.08) and (785.92,130) .. (780.89,130) -- (459.72,130) .. controls (454.69,130) and (450.61,134.08) .. (450.61,139.11) -- cycle ;
%Rounded Rect [id:dp549636265291094] 
\draw  [fill={rgb, 255:red, 230; green, 230; blue, 230 }  ,fill opacity=0.5 ][dash pattern={on 4.5pt off 4.5pt}] (29.83,42.77) .. controls (29.83,35.88) and (35.41,30.3) .. (42.3,30.3) -- (398.54,30.3) .. controls (405.42,30.3) and (411,35.88) .. (411,42.77) -- (411,87.17) .. controls (411,94.05) and (405.42,99.64) .. (398.54,99.64) -- (42.3,99.64) .. controls (35.41,99.64) and (29.83,94.05) .. (29.83,87.17) -- cycle ;
%Rounded Rect [id:dp1378581119146769] 
\draw  [fill={rgb, 255:red, 255; green, 255; blue, 255 }  ,fill opacity=1 ] (40.4,49.82) .. controls (40.4,44.29) and (44.89,39.8) .. (50.42,39.8) -- (169.58,39.8) .. controls (175.11,39.8) and (179.6,44.29) .. (179.6,49.82) -- (179.6,79.88) .. controls (179.6,85.41) and (175.11,89.9) .. (169.58,89.9) -- (50.42,89.9) .. controls (44.89,89.9) and (40.4,85.41) .. (40.4,79.88) -- cycle ;
%Straight Lines [id:da8729153203483555] 
\draw    (188.42,63.77) -- (250.99,63.77)(188.42,66.77) -- (250.99,66.77) ;
\draw [shift={(257.99,65.27)}, rotate = 180] [color={rgb, 255:red, 0; green, 0; blue, 0 }  ][line width=0.75]    (10.93,-4.9) .. controls (6.95,-2.3) and (3.31,-0.67) .. (0,0) .. controls (3.31,0.67) and (6.95,2.3) .. (10.93,4.9)   ;
\draw [shift={(181.42,65.27)}, rotate = 0] [color={rgb, 255:red, 0; green, 0; blue, 0 }  ][line width=0.75]    (10.93,-4.9) .. controls (6.95,-2.3) and (3.31,-0.67) .. (0,0) .. controls (3.31,0.67) and (6.95,2.3) .. (10.93,4.9)   ;
%Straight Lines [id:da845243590120559] 
\draw    (220.17,99.97) -- (220.17,137.64) ;
\draw [shift={(220.17,139.64)}, rotate = 270] [color={rgb, 255:red, 0; green, 0; blue, 0 }  ][line width=0.75]    (10.93,-3.29) .. controls (6.95,-1.4) and (3.31,-0.3) .. (0,0) .. controls (3.31,0.3) and (6.95,1.4) .. (10.93,3.29)   ;
%Straight Lines [id:da7497001546502899] 
\draw    (380,420.33) -- (380,478) ;
\draw [shift={(380,480)}, rotate = 270] [color={rgb, 255:red, 0; green, 0; blue, 0 }  ][line width=0.75]    (10.93,-3.29) .. controls (6.95,-1.4) and (3.31,-0.3) .. (0,0) .. controls (3.31,0.3) and (6.95,1.4) .. (10.93,3.29)   ;
%Straight Lines [id:da5689893955334542] 
\draw    (300,420) -- (460,420) ;
%Rounded Rect [id:dp18582062333838767] 
\draw  [fill={rgb, 255:red, 255; green, 255; blue, 255 }  ,fill opacity=1 ] (260.4,49.82) .. controls (260.4,44.29) and (264.89,39.8) .. (270.42,39.8) -- (389.58,39.8) .. controls (395.11,39.8) and (399.6,44.29) .. (399.6,49.82) -- (399.6,79.88) .. controls (399.6,85.41) and (395.11,89.9) .. (389.58,89.9) -- (270.42,89.9) .. controls (264.89,89.9) and (260.4,85.41) .. (260.4,79.88) -- cycle ;
%Rounded Rect [id:dp9523169763067778] 
\draw  [fill={rgb, 255:red, 255; green, 255; blue, 255 }  ,fill opacity=1 ] (470,299.92) .. controls (470,294.39) and (474.49,289.9) .. (480.02,289.9) -- (599.18,289.9) .. controls (604.71,289.9) and (609.2,294.39) .. (609.2,299.92) -- (609.2,329.98) .. controls (609.2,335.51) and (604.71,340) .. (599.18,340) -- (480.02,340) .. controls (474.49,340) and (470,335.51) .. (470,329.98) -- cycle ;
%Rounded Rect [id:dp48544030815291705] 
\draw  [fill={rgb, 255:red, 255; green, 255; blue, 255 }  ,fill opacity=1 ] (220,490) .. controls (220,484.48) and (224.48,480) .. (230,480) -- (530,480) .. controls (535.52,480) and (540,484.48) .. (540,490) -- (540,570) .. controls (540,575.52) and (535.52,580) .. (530,580) -- (230,580) .. controls (224.48,580) and (220,575.52) .. (220,570) -- cycle ;
%Rounded Rect [id:dp08222627508455749] 
\draw  [fill={rgb, 255:red, 255; green, 255; blue, 255 }  ,fill opacity=1 ] (540.4,49.82) .. controls (540.4,44.29) and (544.89,39.8) .. (550.42,39.8) -- (690.65,39.8) .. controls (696.18,39.8) and (700.67,44.29) .. (700.67,49.82) -- (700.67,79.88) .. controls (700.67,85.41) and (696.18,89.9) .. (690.65,89.9) -- (550.42,89.9) .. controls (544.89,89.9) and (540.4,85.41) .. (540.4,79.88) -- cycle ;
%Straight Lines [id:da5868346364154593] 
\draw  [dash pattern={on 4.5pt off 4.5pt}]  (810,530) -- (542,530) ;
\draw [shift={(540,530)}, rotate = 360] [color={rgb, 255:red, 0; green, 0; blue, 0 }  ][line width=0.75]    (10.93,-3.29) .. controls (6.95,-1.4) and (3.31,-0.3) .. (0,0) .. controls (3.31,0.3) and (6.95,1.4) .. (10.93,3.29)   ;
%Rounded Rect [id:dp7392784471609453] 
\draw  [fill={rgb, 255:red, 255; green, 255; blue, 255 }  ,fill opacity=1 ] (30.33,150.47) .. controls (30.33,144.58) and (35.11,139.8) .. (41,139.8) -- (400.34,139.8) .. controls (406.23,139.8) and (411,144.58) .. (411,150.47) -- (411,229.34) .. controls (411,235.23) and (406.23,240.01) .. (400.34,240.01) -- (41,240.01) .. controls (35.11,240.01) and (30.33,235.23) .. (30.33,229.34) -- cycle ;
%Rounded Rect [id:dp90901637181607] 
\draw  [fill={rgb, 255:red, 255; green, 255; blue, 255 }  ,fill opacity=1 ] (41,189.82) .. controls (41,184.29) and (45.49,179.8) .. (51.02,179.8) -- (145.02,179.8) .. controls (150.55,179.8) and (155.04,184.29) .. (155.04,189.82) -- (155.04,219.88) .. controls (155.04,225.41) and (150.55,229.9) .. (145.02,229.9) -- (51.02,229.9) .. controls (45.49,229.9) and (41,225.41) .. (41,219.88) -- cycle ;
%Rounded Rect [id:dp06115018179625542] 
\draw  [fill={rgb, 255:red, 255; green, 255; blue, 255 }  ,fill opacity=1 ] (163.65,189.82) .. controls (163.65,184.29) and (168.13,179.8) .. (173.67,179.8) -- (267.67,179.8) .. controls (273.2,179.8) and (277.69,184.29) .. (277.69,189.82) -- (277.69,219.88) .. controls (277.69,225.41) and (273.2,229.9) .. (267.67,229.9) -- (173.67,229.9) .. controls (168.13,229.9) and (163.65,225.41) .. (163.65,219.88) -- cycle ;
%Rounded Rect [id:dp6195275548153432] 
\draw  [fill={rgb, 255:red, 255; green, 255; blue, 255 }  ,fill opacity=1 ] (286.3,189.82) .. controls (286.3,184.29) and (290.78,179.8) .. (296.32,179.8) -- (390.31,179.8) .. controls (395.85,179.8) and (400.33,184.29) .. (400.33,189.82) -- (400.33,219.88) .. controls (400.33,225.41) and (395.85,229.9) .. (390.31,229.9) -- (296.32,229.9) .. controls (290.78,229.9) and (286.3,225.41) .. (286.3,219.88) -- cycle ;
%Rounded Rect [id:dp41283415100159493] 
\draw  [fill={rgb, 255:red, 255; green, 255; blue, 255 }  ,fill opacity=1 ] (40.8,299.92) .. controls (40.8,294.39) and (45.29,289.9) .. (50.82,289.9) -- (169.98,289.9) .. controls (175.51,289.9) and (180,294.39) .. (180,299.92) -- (180,329.98) .. controls (180,335.51) and (175.51,340) .. (169.98,340) -- (50.82,340) .. controls (45.29,340) and (40.8,335.51) .. (40.8,329.98) -- cycle ;
%Rounded Rect [id:dp40154854661776396] 
\draw  [fill={rgb, 255:red, 255; green, 255; blue, 255 }  ,fill opacity=1 ] (260.4,299.82) .. controls (260.4,294.29) and (264.89,289.8) .. (270.42,289.8) -- (389.58,289.8) .. controls (395.11,289.8) and (399.6,294.29) .. (399.6,299.82) -- (399.6,329.88) .. controls (399.6,335.41) and (395.11,339.9) .. (389.58,339.9) -- (270.42,339.9) .. controls (264.89,339.9) and (260.4,335.41) .. (260.4,329.88) -- cycle ;
%Straight Lines [id:da6602630444249074] 
\draw    (180.17,315.8) -- (260,315.8) ;
%Straight Lines [id:da30905866347109745] 
\draw    (219.86,315.97) -- (219.86,386.81) ;
\draw [shift={(219.86,388.81)}, rotate = 270] [color={rgb, 255:red, 0; green, 0; blue, 0 }  ][line width=0.75]    (10.93,-3.29) .. controls (6.95,-1.4) and (3.31,-0.3) .. (0,0) .. controls (3.31,0.3) and (6.95,1.4) .. (10.93,3.29)   ;
%Rounded Rect [id:dp44627062039563214] 
\draw  [fill={rgb, 255:red, 255; green, 255; blue, 255 }  ,fill opacity=1 ] (139.67,399.82) .. controls (139.67,394.29) and (144.15,389.8) .. (149.69,389.8) -- (290.31,389.8) .. controls (295.85,389.8) and (300.33,394.29) .. (300.33,399.82) -- (300.33,429.88) .. controls (300.33,435.41) and (295.85,439.9) .. (290.31,439.9) -- (149.69,439.9) .. controls (144.15,439.9) and (139.67,435.41) .. (139.67,429.88) -- cycle ;
%Straight Lines [id:da8375870493480149] 
\draw    (220,240) -- (220,277.67) ;
\draw [shift={(220,279.67)}, rotate = 270] [color={rgb, 255:red, 0; green, 0; blue, 0 }  ][line width=0.75]    (10.93,-3.29) .. controls (6.95,-1.4) and (3.31,-0.3) .. (0,0) .. controls (3.31,0.3) and (6.95,1.4) .. (10.93,3.29)   ;
%Rounded Rect [id:dp37285624114093485] 
\draw  [fill={rgb, 255:red, 255; green, 255; blue, 255 }  ,fill opacity=1 ] (460,400.02) .. controls (460,394.49) and (464.49,390) .. (470.02,390) -- (610.65,390) .. controls (616.18,390) and (620.67,394.49) .. (620.67,400.02) -- (620.67,430.08) .. controls (620.67,435.61) and (616.18,440.1) .. (610.65,440.1) -- (470.02,440.1) .. controls (464.49,440.1) and (460,435.61) .. (460,430.08) -- cycle ;
%Rounded Rect [id:dp6883616335272991] 
\draw  [fill={rgb, 255:red, 255; green, 255; blue, 255 }  ,fill opacity=1 ] (629.2,299.92) .. controls (629.2,294.39) and (633.69,289.9) .. (639.22,289.9) -- (758.38,289.9) .. controls (763.91,289.9) and (768.4,294.39) .. (768.4,299.92) -- (768.4,329.98) .. controls (768.4,335.51) and (763.91,340) .. (758.38,340) -- (639.22,340) .. controls (633.69,340) and (629.2,335.51) .. (629.2,329.98) -- cycle ;
%Straight Lines [id:da19870419404011752] 
\draw    (540,340) -- (540,388) ;
\draw [shift={(540,390)}, rotate = 270] [color={rgb, 255:red, 0; green, 0; blue, 0 }  ][line width=0.75]    (10.93,-3.29) .. controls (6.95,-1.4) and (3.31,-0.3) .. (0,0) .. controls (3.31,0.3) and (6.95,1.4) .. (10.93,3.29)   ;
%Rounded Rect [id:dp9692784599420237] 
\draw  [fill={rgb, 255:red, 255; green, 255; blue, 255 }  ,fill opacity=1 ] (230,529.92) .. controls (230,524.39) and (234.49,519.9) .. (240.02,519.9) -- (359.18,519.9) .. controls (364.71,519.9) and (369.2,524.39) .. (369.2,529.92) -- (369.2,559.98) .. controls (369.2,565.51) and (364.71,570) .. (359.18,570) -- (240.02,570) .. controls (234.49,570) and (230,565.51) .. (230,559.98) -- cycle ;
%Rounded Rect [id:dp749684165723072] 
\draw  [fill={rgb, 255:red, 255; green, 255; blue, 255 }  ,fill opacity=1 ] (390,530.02) .. controls (390,524.49) and (394.49,520) .. (400.02,520) -- (519.18,520) .. controls (524.71,520) and (529.2,524.49) .. (529.2,530.02) -- (529.2,560.08) .. controls (529.2,565.61) and (524.71,570.1) .. (519.18,570.1) -- (400.02,570.1) .. controls (394.49,570.1) and (390,565.61) .. (390,560.08) -- cycle ;
%Rounded Rect [id:dp28001422897762485] 
\draw  [fill={rgb, 255:red, 255; green, 255; blue, 255 }  ,fill opacity=1 ][dash pattern={on 4.5pt off 4.5pt}] (460,148.92) .. controls (460,143.99) and (463.99,140) .. (468.92,140) -- (771.08,140) .. controls (776.01,140) and (780,143.99) .. (780,148.92) -- (780,241.55) .. controls (780,246.48) and (776.01,250.47) .. (771.08,250.47) -- (468.92,250.47) .. controls (463.99,250.47) and (460,246.48) .. (460,241.55) -- cycle ;
%Shape: Brace [id:dp9411387904179627] 
\draw   (774.67,286.8) .. controls (774.67,282.13) and (772.34,279.8) .. (767.67,279.8) -- (629.96,279.8) .. controls (623.29,279.8) and (619.96,277.47) .. (619.96,272.8) .. controls (619.96,277.47) and (616.63,279.8) .. (609.96,279.8)(612.96,279.8) -- (472.25,279.8) .. controls (467.58,279.8) and (465.25,282.13) .. (465.25,286.8) ;
%Straight Lines [id:da9902399507274253] 
\draw    (620,273) -- (620,253.14) ;
\draw [shift={(620,251.14)}, rotate = 90] [color={rgb, 255:red, 0; green, 0; blue, 0 }  ][line width=0.75]    (10.93,-3.29) .. controls (6.95,-1.4) and (3.31,-0.3) .. (0,0) .. controls (3.31,0.3) and (6.95,1.4) .. (10.93,3.29)   ;

% Text Node
\draw (110,64.85) node  [color={rgb, 255:red, 0; green, 0; blue, 0 }  ,opacity=1 ] [align=left] {\begin{minipage}[lt]{92.93pt}\setlength\topsep{0pt}
\begin{center}
programa en Python\\
\small{(\cref{sec:programas-python})}
\end{center}

\end{minipage}};
% Text Node
\draw (219.89,51.97) node  [font=\small,color={rgb, 255:red, 0; green, 0; blue, 0 }  ,opacity=1 ] [align=left] {\begin{minipage}[lt]{54.86pt}\setlength\topsep{0pt}
\begin{center}
\Cref{sec:equivalencia}
\end{center}

\end{minipage}};
% Text Node
\draw (69.5,15.5) node  [font=\small,color={rgb, 255:red, 0; green, 0; blue, 0 }  ,opacity=1 ] [align=left] {\begin{minipage}[lt]{41.55pt}\setlength\topsep{0pt}
\cref{ch:programas-python-maquinas-turing}
\end{minipage}};
% Text Node
\draw (330,64.85) node  [color={rgb, 255:red, 0; green, 0; blue, 0 }  ,opacity=1 ] [align=left] {\begin{minipage}[lt]{92.93pt}\setlength\topsep{0pt}
\begin{center}
máquina de Turing\\
\small{(\cref{sec:maquinas-turing})}
\end{center}

\end{minipage}};
% Text Node
\draw (539.6,314.95) node  [color={rgb, 255:red, 0; green, 0; blue, 0 }  ,opacity=1 ] [align=left] {\begin{minipage}[lt]{92.93pt}\setlength\topsep{0pt}
\begin{center}
sistema lógico\\
\small{(\cref{sec:sistemas-logicos})}
\end{center}

\end{minipage}};
% Text Node
\draw (380.5,500.11) node  [color={rgb, 255:red, 0; green, 0; blue, 0 }  ,opacity=1 ] [align=left] {\begin{minipage}[lt]{159.55pt}\setlength\topsep{0pt}
\begin{center}
\textbf{Primer Teorema de Incompletitud}
\end{center}

\end{minipage}};
% Text Node
\draw (620.53,64.85) node  [color={rgb, 255:red, 0; green, 0; blue, 0 }  ,opacity=1 ] [align=left] {\begin{minipage}[lt]{106.99pt}\setlength\topsep{0pt}
\begin{center}
tesis de Church-Turing\\
\small{(\cref{sec:church-turing})}
\end{center}

\end{minipage}};
% Text Node
\draw (480.28,52.8) node  [font=\small,color={rgb, 255:red, 0; green, 0; blue, 0 }  ,opacity=1 ] [align=left] {\begin{minipage}[lt]{81.9pt}\setlength\topsep{0pt}
\begin{center}
\cref{sec:church-turing}
\end{center}

\end{minipage}};
% Text Node
\draw (589.5,518.3) node  [font=\small,color={rgb, 255:red, 0; green, 0; blue, 0 }  ,opacity=1 ] [align=left] {\begin{minipage}[lt]{54.86pt}\setlength\topsep{0pt}
\begin{center}
\Cref{subsec:mente-humana-tesis-church-turing}
\end{center}

\end{minipage}};
% Text Node
\draw (220.2,159.62) node  [color={rgb, 255:red, 0; green, 0; blue, 0 }  ,opacity=1 ] [align=left] {\begin{minipage}[lt]{230.11pt}\setlength\topsep{0pt}
\begin{center}
problema computacional
\end{center}

\end{minipage}};
% Text Node
\draw (98.02,204.85) node  [color={rgb, 255:red, 0; green, 0; blue, 0 }  ,opacity=1 ] [align=left] {\begin{minipage}[lt]{76.13pt}\setlength\topsep{0pt}
\begin{center}
decidible\\
\small{(\cref{sec:problemas-decidibles})}
\end{center}

\end{minipage}};
% Text Node
\draw (220.67,204.85) node  [color={rgb, 255:red, 0; green, 0; blue, 0 }  ,opacity=1 ] [align=left] {\begin{minipage}[lt]{76.13pt}\setlength\topsep{0pt}
\begin{center}
no decidible\\
\small{(\cref{sec:problemas-no-decidibles})}
\end{center}

\end{minipage}};
% Text Node
\draw (343.31,204.85) node  [color={rgb, 255:red, 0; green, 0; blue, 0 }  ,opacity=1 ] [align=left] {\begin{minipage}[lt]{76.13pt}\setlength\topsep{0pt}
\begin{center}
semidecidible\\
\small{(\cref{sec:problemas-semidecidibles})}
\end{center}

\end{minipage}};
% Text Node
\draw (110.4,314.95) node  [color={rgb, 255:red, 0; green, 0; blue, 0 }  ,opacity=1 ] [align=left] {\begin{minipage}[lt]{92.93pt}\setlength\topsep{0pt}
\begin{center}
universalidad\\
\small{(\cref{sec:universalidad})}
\end{center}

\end{minipage}};
% Text Node
\draw (330,314.85) node  [color={rgb, 255:red, 0; green, 0; blue, 0 }  ,opacity=1 ] [align=left] {\begin{minipage}[lt]{92.93pt}\setlength\topsep{0pt}
\begin{center}
reducción\\
\small{(\cref{sec:reducciones})}
\end{center}

\end{minipage}};
% Text Node
\draw (220,414.85) node  [color={rgb, 255:red, 0; green, 0; blue, 0 }  ,opacity=1 ] [align=left] {\begin{minipage}[lt]{107.26pt}\setlength\topsep{0pt}
\begin{center}
problema de la parada\\
\small{(\cref{sec:problema-parada})}
\end{center}

\end{minipage}};
% Text Node
\draw (540.33,415.05) node  [color={rgb, 255:red, 0; green, 0; blue, 0 }  ,opacity=1 ] [align=left] {\begin{minipage}[lt]{107.26pt}\setlength\topsep{0pt}
\begin{center}
aritmética de Peano\\
\small{(\cref{sec:aritmetica-peano})}
\end{center}

\end{minipage}};
% Text Node
\draw (698.8,314.95) node  [color={rgb, 255:red, 0; green, 0; blue, 0 }  ,opacity=1 ] [align=left] {\begin{minipage}[lt]{92.93pt}\setlength\topsep{0pt}
\begin{center}
sistema formal\\
\small{(\cref{sec:sistemas-formales})}
\end{center}

\end{minipage}};
% Text Node
\draw (299.6,544.95) node  [color={rgb, 255:red, 0; green, 0; blue, 0 }  ,opacity=1 ] [align=left] {\begin{minipage}[lt]{92.93pt}\setlength\topsep{0pt}
\begin{center}
versión semántica\\
\small{(\cref{sec:version-semantica})}
\end{center}

\end{minipage}};
% Text Node
\draw (459.6,545.05) node  [color={rgb, 255:red, 0; green, 0; blue, 0 }  ,opacity=1 ] [align=left] {\begin{minipage}[lt]{92.93pt}\setlength\topsep{0pt}
\begin{center}
versión sintáctica\\
\small{(\cref{sec:version-sintactica})}
\end{center}

\end{minipage}};
% Text Node
\draw (60.5,464.5) node  [font=\small,color={rgb, 255:red, 0; green, 0; blue, 0 }  ,opacity=1 ] [align=left] {\begin{minipage}[lt]{41.55pt}\setlength\topsep{0pt}
\cref{ch:problema-parada}
\end{minipage}};
% Text Node
\draw (749.5,465.5) node  [font=\small,color={rgb, 255:red, 0; green, 0; blue, 0 }  ,opacity=1 ] [align=left] {\begin{minipage}[lt]{41.55pt}\setlength\topsep{0pt}
\begin{flushright}
\cref{ch:sistemas-logicos}
\end{flushright}

\end{minipage}};
% Text Node
\draw (170.5,604.5) node  [font=\small,color={rgb, 255:red, 0; green, 0; blue, 0 }  ,opacity=1 ] [align=left] {\begin{minipage}[lt]{41.55pt}\setlength\topsep{0pt}
\cref{ch:teorema-incompletitud}
\end{minipage}};
% Text Node
\draw (643.76,195.55) node  [color={rgb, 255:red, 0; green, 0; blue, 0 }  ,opacity=1 ] [align=left] {\begin{minipage}[lt]{217.92pt}\setlength\topsep{0pt}
-- sólido {\small{(\cref{sec:solidez-completitud-decidibilidad})}}\\
-- decidible {\small{(\cref{sec:solidez-completitud-decidibilidad})}}\\
-- (semánticamente) completo {\small{(\cref{sec:solidez-completitud-decidibilidad})}}\\
-- sintácticamente completo {\small{(\cref{sec:consistencia})}}\\
-- sem./sint. consistente {\small{(\cref{sec:consistencia})}}
\end{minipage}};
% Text Node
\draw (414.84,645.25) node  [color={rgb, 255:red, 0; green, 0; blue, 0 }  ,opacity=1 ] [align=left] {\begin{minipage}[lt]{534.7pt}\setlength\topsep{0pt}
\begin{center}
Esquema resumen de este Trabajo de Fin de Grado
\end{center}
\end{minipage}};
% Text Node
\draw (380.15,407.8) node  [font=\small,color={rgb, 255:red, 0; green, 0; blue, 0 }  ,opacity=1 ] [align=left] {\begin{minipage}[lt]{108.93pt}\setlength\topsep{0pt}
\begin{center}
\cref{sec:primera-aproximacion}
\end{center}

\end{minipage}};

\end{tikzpicture}}
}
%\caption{Esquema resumen (extendido)}

%\label{fig:resumen-extendido}
\end{figure}
% ====================
\end{landscape}

\endinput
}
%\caption{Esquema resumen}
%\label{fig:resumen}
%\end{figure}
% ====================
%\vfill

\begin{landscape}
% ====================
\begin{figure}[H]
%\captionsetup{margin=200cm}

\centering
%\vspace{-1.5cm}
%\hspace{-1cm}
%

\tikzset{every picture/.style={line width=0.75pt}} %set default line width to 0.75pt        

\begin{tikzpicture}[x=0.75pt,y=0.75pt,yscale=-1,xscale=1]
%uncomment if require: \path (0,694); %set diagram left start at 0, and has height of 694

%Straight Lines [id:da2056540562354816] 
\draw  [dash pattern={on 4.5pt off 4.5pt}]  (672,64.47) -- (810,64.47) ;
%Straight Lines [id:da23044041717189923] 
\draw  [dash pattern={on 4.5pt off 4.5pt}]  (412.67,64.47) -- (610,64.47) ;
\draw [shift={(410.67,64.47)}, rotate = 0] [color={rgb, 255:red, 0; green, 0; blue, 0 }  ][line width=0.75]    (10.93,-3.29) .. controls (6.95,-1.4) and (3.31,-0.3) .. (0,0) .. controls (3.31,0.3) and (6.95,1.4) .. (10.93,3.29)   ;
%Straight Lines [id:da8398435684999113] 
\draw  [dash pattern={on 4.5pt off 4.5pt}]  (810,64.47) -- (810,530) ;
%Rounded Rect [id:dp027018508739026448] 
\draw  [draw opacity=0][fill={rgb, 255:red, 230; green, 230; blue, 230 }  ,fill opacity=0.5 ] (20,439.33) .. controls (20,445.22) and (24.78,450) .. (30.67,450) -- (409.33,450) .. controls (415.22,450) and (420,445.22) .. (420,439.33) -- (420,140.67) .. controls (420,134.78) and (415.22,130) .. (409.33,130) -- (30.67,130) .. controls (24.78,130) and (20,134.78) .. (20,140.67) -- cycle ;
%Rounded Rect [id:dp9415679498487528] 
\draw  [dash pattern={on 4.5pt off 4.5pt}] (29.83,292.77) .. controls (29.83,285.88) and (35.41,280.3) .. (42.3,280.3) -- (398.54,280.3) .. controls (405.42,280.3) and (411,285.88) .. (411,292.77) -- (411,337.17) .. controls (411,344.05) and (405.42,349.64) .. (398.54,349.64) -- (42.3,349.64) .. controls (35.41,349.64) and (29.83,344.05) .. (29.83,337.17) -- cycle ;
%Rounded Rect [id:dp05100022295548956] 
\draw  [draw opacity=0][fill={rgb, 255:red, 230; green, 230; blue, 230 }  ,fill opacity=0.5 ] (130.5,575) .. controls (130.5,583.28) and (137.22,590) .. (145.5,590) -- (615,590) .. controls (623.28,590) and (630,583.28) .. (630,575) -- (630,395) .. controls (630,386.72) and (623.28,380) .. (615,380) -- (145.5,380) .. controls (137.22,380) and (130.5,386.72) .. (130.5,395) -- cycle ;
%Rounded Rect [id:dp6496809750851564] 
\draw  [draw opacity=0][fill={rgb, 255:red, 230; green, 230; blue, 230 }  ,fill opacity=0.5 ] (450.61,440.79) .. controls (450.61,445.82) and (454.69,449.9) .. (459.72,449.9) -- (780.89,449.9) .. controls (785.92,449.9) and (790,445.82) .. (790,440.79) -- (790,139.11) .. controls (790,134.08) and (785.92,130) .. (780.89,130) -- (459.72,130) .. controls (454.69,130) and (450.61,134.08) .. (450.61,139.11) -- cycle ;
%Rounded Rect [id:dp549636265291094] 
\draw  [fill={rgb, 255:red, 230; green, 230; blue, 230 }  ,fill opacity=0.5 ][dash pattern={on 4.5pt off 4.5pt}] (29.83,42.77) .. controls (29.83,35.88) and (35.41,30.3) .. (42.3,30.3) -- (398.54,30.3) .. controls (405.42,30.3) and (411,35.88) .. (411,42.77) -- (411,87.17) .. controls (411,94.05) and (405.42,99.64) .. (398.54,99.64) -- (42.3,99.64) .. controls (35.41,99.64) and (29.83,94.05) .. (29.83,87.17) -- cycle ;
%Rounded Rect [id:dp1378581119146769] 
\draw  [fill={rgb, 255:red, 255; green, 255; blue, 255 }  ,fill opacity=1 ] (40.4,49.82) .. controls (40.4,44.29) and (44.89,39.8) .. (50.42,39.8) -- (169.58,39.8) .. controls (175.11,39.8) and (179.6,44.29) .. (179.6,49.82) -- (179.6,79.88) .. controls (179.6,85.41) and (175.11,89.9) .. (169.58,89.9) -- (50.42,89.9) .. controls (44.89,89.9) and (40.4,85.41) .. (40.4,79.88) -- cycle ;
%Straight Lines [id:da8729153203483555] 
\draw    (188.42,63.77) -- (250.99,63.77)(188.42,66.77) -- (250.99,66.77) ;
\draw [shift={(257.99,65.27)}, rotate = 180] [color={rgb, 255:red, 0; green, 0; blue, 0 }  ][line width=0.75]    (10.93,-4.9) .. controls (6.95,-2.3) and (3.31,-0.67) .. (0,0) .. controls (3.31,0.67) and (6.95,2.3) .. (10.93,4.9)   ;
\draw [shift={(181.42,65.27)}, rotate = 0] [color={rgb, 255:red, 0; green, 0; blue, 0 }  ][line width=0.75]    (10.93,-4.9) .. controls (6.95,-2.3) and (3.31,-0.67) .. (0,0) .. controls (3.31,0.67) and (6.95,2.3) .. (10.93,4.9)   ;
%Straight Lines [id:da845243590120559] 
\draw    (220.17,99.97) -- (220.17,137.64) ;
\draw [shift={(220.17,139.64)}, rotate = 270] [color={rgb, 255:red, 0; green, 0; blue, 0 }  ][line width=0.75]    (10.93,-3.29) .. controls (6.95,-1.4) and (3.31,-0.3) .. (0,0) .. controls (3.31,0.3) and (6.95,1.4) .. (10.93,3.29)   ;
%Straight Lines [id:da7497001546502899] 
\draw    (380,420.33) -- (380,478) ;
\draw [shift={(380,480)}, rotate = 270] [color={rgb, 255:red, 0; green, 0; blue, 0 }  ][line width=0.75]    (10.93,-3.29) .. controls (6.95,-1.4) and (3.31,-0.3) .. (0,0) .. controls (3.31,0.3) and (6.95,1.4) .. (10.93,3.29)   ;
%Straight Lines [id:da5689893955334542] 
\draw    (300,420) -- (460,420) ;
%Rounded Rect [id:dp18582062333838767] 
\draw  [fill={rgb, 255:red, 255; green, 255; blue, 255 }  ,fill opacity=1 ] (260.4,49.82) .. controls (260.4,44.29) and (264.89,39.8) .. (270.42,39.8) -- (389.58,39.8) .. controls (395.11,39.8) and (399.6,44.29) .. (399.6,49.82) -- (399.6,79.88) .. controls (399.6,85.41) and (395.11,89.9) .. (389.58,89.9) -- (270.42,89.9) .. controls (264.89,89.9) and (260.4,85.41) .. (260.4,79.88) -- cycle ;
%Rounded Rect [id:dp9523169763067778] 
\draw  [fill={rgb, 255:red, 255; green, 255; blue, 255 }  ,fill opacity=1 ] (470,299.92) .. controls (470,294.39) and (474.49,289.9) .. (480.02,289.9) -- (599.18,289.9) .. controls (604.71,289.9) and (609.2,294.39) .. (609.2,299.92) -- (609.2,329.98) .. controls (609.2,335.51) and (604.71,340) .. (599.18,340) -- (480.02,340) .. controls (474.49,340) and (470,335.51) .. (470,329.98) -- cycle ;
%Rounded Rect [id:dp48544030815291705] 
\draw  [fill={rgb, 255:red, 255; green, 255; blue, 255 }  ,fill opacity=1 ] (220,490) .. controls (220,484.48) and (224.48,480) .. (230,480) -- (530,480) .. controls (535.52,480) and (540,484.48) .. (540,490) -- (540,570) .. controls (540,575.52) and (535.52,580) .. (530,580) -- (230,580) .. controls (224.48,580) and (220,575.52) .. (220,570) -- cycle ;
%Rounded Rect [id:dp08222627508455749] 
\draw  [fill={rgb, 255:red, 255; green, 255; blue, 255 }  ,fill opacity=1 ] (540.4,49.82) .. controls (540.4,44.29) and (544.89,39.8) .. (550.42,39.8) -- (690.65,39.8) .. controls (696.18,39.8) and (700.67,44.29) .. (700.67,49.82) -- (700.67,79.88) .. controls (700.67,85.41) and (696.18,89.9) .. (690.65,89.9) -- (550.42,89.9) .. controls (544.89,89.9) and (540.4,85.41) .. (540.4,79.88) -- cycle ;
%Straight Lines [id:da5868346364154593] 
\draw  [dash pattern={on 4.5pt off 4.5pt}]  (810,530) -- (542,530) ;
\draw [shift={(540,530)}, rotate = 360] [color={rgb, 255:red, 0; green, 0; blue, 0 }  ][line width=0.75]    (10.93,-3.29) .. controls (6.95,-1.4) and (3.31,-0.3) .. (0,0) .. controls (3.31,0.3) and (6.95,1.4) .. (10.93,3.29)   ;
%Rounded Rect [id:dp7392784471609453] 
\draw  [fill={rgb, 255:red, 255; green, 255; blue, 255 }  ,fill opacity=1 ] (30.33,150.47) .. controls (30.33,144.58) and (35.11,139.8) .. (41,139.8) -- (400.34,139.8) .. controls (406.23,139.8) and (411,144.58) .. (411,150.47) -- (411,229.34) .. controls (411,235.23) and (406.23,240.01) .. (400.34,240.01) -- (41,240.01) .. controls (35.11,240.01) and (30.33,235.23) .. (30.33,229.34) -- cycle ;
%Rounded Rect [id:dp90901637181607] 
\draw  [fill={rgb, 255:red, 255; green, 255; blue, 255 }  ,fill opacity=1 ] (41,189.82) .. controls (41,184.29) and (45.49,179.8) .. (51.02,179.8) -- (145.02,179.8) .. controls (150.55,179.8) and (155.04,184.29) .. (155.04,189.82) -- (155.04,219.88) .. controls (155.04,225.41) and (150.55,229.9) .. (145.02,229.9) -- (51.02,229.9) .. controls (45.49,229.9) and (41,225.41) .. (41,219.88) -- cycle ;
%Rounded Rect [id:dp06115018179625542] 
\draw  [fill={rgb, 255:red, 255; green, 255; blue, 255 }  ,fill opacity=1 ] (163.65,189.82) .. controls (163.65,184.29) and (168.13,179.8) .. (173.67,179.8) -- (267.67,179.8) .. controls (273.2,179.8) and (277.69,184.29) .. (277.69,189.82) -- (277.69,219.88) .. controls (277.69,225.41) and (273.2,229.9) .. (267.67,229.9) -- (173.67,229.9) .. controls (168.13,229.9) and (163.65,225.41) .. (163.65,219.88) -- cycle ;
%Rounded Rect [id:dp6195275548153432] 
\draw  [fill={rgb, 255:red, 255; green, 255; blue, 255 }  ,fill opacity=1 ] (286.3,189.82) .. controls (286.3,184.29) and (290.78,179.8) .. (296.32,179.8) -- (390.31,179.8) .. controls (395.85,179.8) and (400.33,184.29) .. (400.33,189.82) -- (400.33,219.88) .. controls (400.33,225.41) and (395.85,229.9) .. (390.31,229.9) -- (296.32,229.9) .. controls (290.78,229.9) and (286.3,225.41) .. (286.3,219.88) -- cycle ;
%Rounded Rect [id:dp41283415100159493] 
\draw  [fill={rgb, 255:red, 255; green, 255; blue, 255 }  ,fill opacity=1 ] (40.8,299.92) .. controls (40.8,294.39) and (45.29,289.9) .. (50.82,289.9) -- (169.98,289.9) .. controls (175.51,289.9) and (180,294.39) .. (180,299.92) -- (180,329.98) .. controls (180,335.51) and (175.51,340) .. (169.98,340) -- (50.82,340) .. controls (45.29,340) and (40.8,335.51) .. (40.8,329.98) -- cycle ;
%Rounded Rect [id:dp40154854661776396] 
\draw  [fill={rgb, 255:red, 255; green, 255; blue, 255 }  ,fill opacity=1 ] (260.4,299.82) .. controls (260.4,294.29) and (264.89,289.8) .. (270.42,289.8) -- (389.58,289.8) .. controls (395.11,289.8) and (399.6,294.29) .. (399.6,299.82) -- (399.6,329.88) .. controls (399.6,335.41) and (395.11,339.9) .. (389.58,339.9) -- (270.42,339.9) .. controls (264.89,339.9) and (260.4,335.41) .. (260.4,329.88) -- cycle ;
%Straight Lines [id:da6602630444249074] 
\draw    (180.17,315.8) -- (260,315.8) ;
%Straight Lines [id:da30905866347109745] 
\draw    (219.86,315.97) -- (219.86,386.81) ;
\draw [shift={(219.86,388.81)}, rotate = 270] [color={rgb, 255:red, 0; green, 0; blue, 0 }  ][line width=0.75]    (10.93,-3.29) .. controls (6.95,-1.4) and (3.31,-0.3) .. (0,0) .. controls (3.31,0.3) and (6.95,1.4) .. (10.93,3.29)   ;
%Rounded Rect [id:dp44627062039563214] 
\draw  [fill={rgb, 255:red, 255; green, 255; blue, 255 }  ,fill opacity=1 ] (139.67,399.82) .. controls (139.67,394.29) and (144.15,389.8) .. (149.69,389.8) -- (290.31,389.8) .. controls (295.85,389.8) and (300.33,394.29) .. (300.33,399.82) -- (300.33,429.88) .. controls (300.33,435.41) and (295.85,439.9) .. (290.31,439.9) -- (149.69,439.9) .. controls (144.15,439.9) and (139.67,435.41) .. (139.67,429.88) -- cycle ;
%Straight Lines [id:da8375870493480149] 
\draw    (220,240) -- (220,277.67) ;
\draw [shift={(220,279.67)}, rotate = 270] [color={rgb, 255:red, 0; green, 0; blue, 0 }  ][line width=0.75]    (10.93,-3.29) .. controls (6.95,-1.4) and (3.31,-0.3) .. (0,0) .. controls (3.31,0.3) and (6.95,1.4) .. (10.93,3.29)   ;
%Rounded Rect [id:dp37285624114093485] 
\draw  [fill={rgb, 255:red, 255; green, 255; blue, 255 }  ,fill opacity=1 ] (460,400.02) .. controls (460,394.49) and (464.49,390) .. (470.02,390) -- (610.65,390) .. controls (616.18,390) and (620.67,394.49) .. (620.67,400.02) -- (620.67,430.08) .. controls (620.67,435.61) and (616.18,440.1) .. (610.65,440.1) -- (470.02,440.1) .. controls (464.49,440.1) and (460,435.61) .. (460,430.08) -- cycle ;
%Rounded Rect [id:dp6883616335272991] 
\draw  [fill={rgb, 255:red, 255; green, 255; blue, 255 }  ,fill opacity=1 ] (629.2,299.92) .. controls (629.2,294.39) and (633.69,289.9) .. (639.22,289.9) -- (758.38,289.9) .. controls (763.91,289.9) and (768.4,294.39) .. (768.4,299.92) -- (768.4,329.98) .. controls (768.4,335.51) and (763.91,340) .. (758.38,340) -- (639.22,340) .. controls (633.69,340) and (629.2,335.51) .. (629.2,329.98) -- cycle ;
%Straight Lines [id:da19870419404011752] 
\draw    (540,340) -- (540,388) ;
\draw [shift={(540,390)}, rotate = 270] [color={rgb, 255:red, 0; green, 0; blue, 0 }  ][line width=0.75]    (10.93,-3.29) .. controls (6.95,-1.4) and (3.31,-0.3) .. (0,0) .. controls (3.31,0.3) and (6.95,1.4) .. (10.93,3.29)   ;
%Rounded Rect [id:dp9692784599420237] 
\draw  [fill={rgb, 255:red, 255; green, 255; blue, 255 }  ,fill opacity=1 ] (230,529.92) .. controls (230,524.39) and (234.49,519.9) .. (240.02,519.9) -- (359.18,519.9) .. controls (364.71,519.9) and (369.2,524.39) .. (369.2,529.92) -- (369.2,559.98) .. controls (369.2,565.51) and (364.71,570) .. (359.18,570) -- (240.02,570) .. controls (234.49,570) and (230,565.51) .. (230,559.98) -- cycle ;
%Rounded Rect [id:dp749684165723072] 
\draw  [fill={rgb, 255:red, 255; green, 255; blue, 255 }  ,fill opacity=1 ] (390,530.02) .. controls (390,524.49) and (394.49,520) .. (400.02,520) -- (519.18,520) .. controls (524.71,520) and (529.2,524.49) .. (529.2,530.02) -- (529.2,560.08) .. controls (529.2,565.61) and (524.71,570.1) .. (519.18,570.1) -- (400.02,570.1) .. controls (394.49,570.1) and (390,565.61) .. (390,560.08) -- cycle ;
%Rounded Rect [id:dp28001422897762485] 
\draw  [fill={rgb, 255:red, 255; green, 255; blue, 255 }  ,fill opacity=1 ][dash pattern={on 4.5pt off 4.5pt}] (460,148.92) .. controls (460,143.99) and (463.99,140) .. (468.92,140) -- (771.08,140) .. controls (776.01,140) and (780,143.99) .. (780,148.92) -- (780,241.55) .. controls (780,246.48) and (776.01,250.47) .. (771.08,250.47) -- (468.92,250.47) .. controls (463.99,250.47) and (460,246.48) .. (460,241.55) -- cycle ;
%Shape: Brace [id:dp9411387904179627] 
\draw   (774.67,286.8) .. controls (774.67,282.13) and (772.34,279.8) .. (767.67,279.8) -- (629.96,279.8) .. controls (623.29,279.8) and (619.96,277.47) .. (619.96,272.8) .. controls (619.96,277.47) and (616.63,279.8) .. (609.96,279.8)(612.96,279.8) -- (472.25,279.8) .. controls (467.58,279.8) and (465.25,282.13) .. (465.25,286.8) ;
%Straight Lines [id:da9902399507274253] 
\draw    (620,273) -- (620,253.14) ;
\draw [shift={(620,251.14)}, rotate = 90] [color={rgb, 255:red, 0; green, 0; blue, 0 }  ][line width=0.75]    (10.93,-3.29) .. controls (6.95,-1.4) and (3.31,-0.3) .. (0,0) .. controls (3.31,0.3) and (6.95,1.4) .. (10.93,3.29)   ;

% Text Node
\draw (110,64.85) node  [color={rgb, 255:red, 0; green, 0; blue, 0 }  ,opacity=1 ] [align=left] {\begin{minipage}[lt]{92.93pt}\setlength\topsep{0pt}
\begin{center}
programa en Python\\
\small{(\cref{sec:programas-python})}
\end{center}

\end{minipage}};
% Text Node
\draw (219.89,51.97) node  [font=\small,color={rgb, 255:red, 0; green, 0; blue, 0 }  ,opacity=1 ] [align=left] {\begin{minipage}[lt]{54.86pt}\setlength\topsep{0pt}
\begin{center}
\Cref{sec:equivalencia}
\end{center}

\end{minipage}};
% Text Node
\draw (69.5,15.5) node  [font=\small,color={rgb, 255:red, 0; green, 0; blue, 0 }  ,opacity=1 ] [align=left] {\begin{minipage}[lt]{41.55pt}\setlength\topsep{0pt}
\cref{ch:programas-python-maquinas-turing}
\end{minipage}};
% Text Node
\draw (330,64.85) node  [color={rgb, 255:red, 0; green, 0; blue, 0 }  ,opacity=1 ] [align=left] {\begin{minipage}[lt]{92.93pt}\setlength\topsep{0pt}
\begin{center}
máquina de Turing\\
\small{(\cref{sec:maquinas-turing})}
\end{center}

\end{minipage}};
% Text Node
\draw (539.6,314.95) node  [color={rgb, 255:red, 0; green, 0; blue, 0 }  ,opacity=1 ] [align=left] {\begin{minipage}[lt]{92.93pt}\setlength\topsep{0pt}
\begin{center}
sistema lógico\\
\small{(\cref{sec:sistemas-logicos})}
\end{center}

\end{minipage}};
% Text Node
\draw (380.5,500.11) node  [color={rgb, 255:red, 0; green, 0; blue, 0 }  ,opacity=1 ] [align=left] {\begin{minipage}[lt]{159.55pt}\setlength\topsep{0pt}
\begin{center}
\textbf{Primer Teorema de Incompletitud}
\end{center}

\end{minipage}};
% Text Node
\draw (620.53,64.85) node  [color={rgb, 255:red, 0; green, 0; blue, 0 }  ,opacity=1 ] [align=left] {\begin{minipage}[lt]{106.99pt}\setlength\topsep{0pt}
\begin{center}
tesis de Church-Turing\\
\small{(\cref{sec:church-turing})}
\end{center}

\end{minipage}};
% Text Node
\draw (480.28,52.8) node  [font=\small,color={rgb, 255:red, 0; green, 0; blue, 0 }  ,opacity=1 ] [align=left] {\begin{minipage}[lt]{81.9pt}\setlength\topsep{0pt}
\begin{center}
\cref{sec:church-turing}
\end{center}

\end{minipage}};
% Text Node
\draw (589.5,518.3) node  [font=\small,color={rgb, 255:red, 0; green, 0; blue, 0 }  ,opacity=1 ] [align=left] {\begin{minipage}[lt]{54.86pt}\setlength\topsep{0pt}
\begin{center}
\Cref{subsec:mente-humana-tesis-church-turing}
\end{center}

\end{minipage}};
% Text Node
\draw (220.2,159.62) node  [color={rgb, 255:red, 0; green, 0; blue, 0 }  ,opacity=1 ] [align=left] {\begin{minipage}[lt]{230.11pt}\setlength\topsep{0pt}
\begin{center}
problema computacional
\end{center}

\end{minipage}};
% Text Node
\draw (98.02,204.85) node  [color={rgb, 255:red, 0; green, 0; blue, 0 }  ,opacity=1 ] [align=left] {\begin{minipage}[lt]{76.13pt}\setlength\topsep{0pt}
\begin{center}
decidible\\
\small{(\cref{sec:problemas-decidibles})}
\end{center}

\end{minipage}};
% Text Node
\draw (220.67,204.85) node  [color={rgb, 255:red, 0; green, 0; blue, 0 }  ,opacity=1 ] [align=left] {\begin{minipage}[lt]{76.13pt}\setlength\topsep{0pt}
\begin{center}
no decidible\\
\small{(\cref{sec:problemas-no-decidibles})}
\end{center}

\end{minipage}};
% Text Node
\draw (343.31,204.85) node  [color={rgb, 255:red, 0; green, 0; blue, 0 }  ,opacity=1 ] [align=left] {\begin{minipage}[lt]{76.13pt}\setlength\topsep{0pt}
\begin{center}
semidecidible\\
\small{(\cref{sec:problemas-semidecidibles})}
\end{center}

\end{minipage}};
% Text Node
\draw (110.4,314.95) node  [color={rgb, 255:red, 0; green, 0; blue, 0 }  ,opacity=1 ] [align=left] {\begin{minipage}[lt]{92.93pt}\setlength\topsep{0pt}
\begin{center}
universalidad\\
\small{(\cref{sec:universalidad})}
\end{center}

\end{minipage}};
% Text Node
\draw (330,314.85) node  [color={rgb, 255:red, 0; green, 0; blue, 0 }  ,opacity=1 ] [align=left] {\begin{minipage}[lt]{92.93pt}\setlength\topsep{0pt}
\begin{center}
reducción\\
\small{(\cref{sec:reducciones})}
\end{center}

\end{minipage}};
% Text Node
\draw (220,414.85) node  [color={rgb, 255:red, 0; green, 0; blue, 0 }  ,opacity=1 ] [align=left] {\begin{minipage}[lt]{107.26pt}\setlength\topsep{0pt}
\begin{center}
problema de la parada\\
\small{(\cref{sec:problema-parada})}
\end{center}

\end{minipage}};
% Text Node
\draw (540.33,415.05) node  [color={rgb, 255:red, 0; green, 0; blue, 0 }  ,opacity=1 ] [align=left] {\begin{minipage}[lt]{107.26pt}\setlength\topsep{0pt}
\begin{center}
aritmética de Peano\\
\small{(\cref{sec:aritmetica-peano})}
\end{center}

\end{minipage}};
% Text Node
\draw (698.8,314.95) node  [color={rgb, 255:red, 0; green, 0; blue, 0 }  ,opacity=1 ] [align=left] {\begin{minipage}[lt]{92.93pt}\setlength\topsep{0pt}
\begin{center}
sistema formal\\
\small{(\cref{sec:sistemas-formales})}
\end{center}

\end{minipage}};
% Text Node
\draw (299.6,544.95) node  [color={rgb, 255:red, 0; green, 0; blue, 0 }  ,opacity=1 ] [align=left] {\begin{minipage}[lt]{92.93pt}\setlength\topsep{0pt}
\begin{center}
versión semántica\\
\small{(\cref{sec:version-semantica})}
\end{center}

\end{minipage}};
% Text Node
\draw (459.6,545.05) node  [color={rgb, 255:red, 0; green, 0; blue, 0 }  ,opacity=1 ] [align=left] {\begin{minipage}[lt]{92.93pt}\setlength\topsep{0pt}
\begin{center}
versión sintáctica\\
\small{(\cref{sec:version-sintactica})}
\end{center}

\end{minipage}};
% Text Node
\draw (60.5,464.5) node  [font=\small,color={rgb, 255:red, 0; green, 0; blue, 0 }  ,opacity=1 ] [align=left] {\begin{minipage}[lt]{41.55pt}\setlength\topsep{0pt}
\cref{ch:problema-parada}
\end{minipage}};
% Text Node
\draw (749.5,465.5) node  [font=\small,color={rgb, 255:red, 0; green, 0; blue, 0 }  ,opacity=1 ] [align=left] {\begin{minipage}[lt]{41.55pt}\setlength\topsep{0pt}
\begin{flushright}
\cref{ch:sistemas-logicos}
\end{flushright}

\end{minipage}};
% Text Node
\draw (170.5,604.5) node  [font=\small,color={rgb, 255:red, 0; green, 0; blue, 0 }  ,opacity=1 ] [align=left] {\begin{minipage}[lt]{41.55pt}\setlength\topsep{0pt}
\cref{ch:teorema-incompletitud}
\end{minipage}};
% Text Node
\draw (643.76,195.55) node  [color={rgb, 255:red, 0; green, 0; blue, 0 }  ,opacity=1 ] [align=left] {\begin{minipage}[lt]{217.92pt}\setlength\topsep{0pt}
-- sólido {\small{(\cref{sec:solidez-completitud-decidibilidad})}}\\
-- decidible {\small{(\cref{sec:solidez-completitud-decidibilidad})}}\\
-- (semánticamente) completo {\small{(\cref{sec:solidez-completitud-decidibilidad})}}\\
-- sintácticamente completo {\small{(\cref{sec:consistencia})}}\\
-- sem./sint. consistente {\small{(\cref{sec:consistencia})}}
\end{minipage}};
% Text Node
\draw (414.84,645.25) node  [color={rgb, 255:red, 0; green, 0; blue, 0 }  ,opacity=1 ] [align=left] {\begin{minipage}[lt]{534.7pt}\setlength\topsep{0pt}
\begin{center}
Esquema resumen de este Trabajo de Fin de Grado
\end{center}
\end{minipage}};
% Text Node
\draw (380.15,407.8) node  [font=\small,color={rgb, 255:red, 0; green, 0; blue, 0 }  ,opacity=1 ] [align=left] {\begin{minipage}[lt]{108.93pt}\setlength\topsep{0pt}
\begin{center}
\cref{sec:primera-aproximacion}
\end{center}

\end{minipage}};

\end{tikzpicture}
\raisebox{-15.5cm}[0pt][0pt]{%
\resizebox{\linewidth}{!}{

\tikzset{every picture/.style={line width=0.75pt}} %set default line width to 0.75pt        

\begin{tikzpicture}[x=0.75pt,y=0.75pt,yscale=-1,xscale=1]
%uncomment if require: \path (0,694); %set diagram left start at 0, and has height of 694

%Straight Lines [id:da2056540562354816] 
\draw  [dash pattern={on 4.5pt off 4.5pt}]  (672,64.47) -- (810,64.47) ;
%Straight Lines [id:da23044041717189923] 
\draw  [dash pattern={on 4.5pt off 4.5pt}]  (412.67,64.47) -- (610,64.47) ;
\draw [shift={(410.67,64.47)}, rotate = 0] [color={rgb, 255:red, 0; green, 0; blue, 0 }  ][line width=0.75]    (10.93,-3.29) .. controls (6.95,-1.4) and (3.31,-0.3) .. (0,0) .. controls (3.31,0.3) and (6.95,1.4) .. (10.93,3.29)   ;
%Straight Lines [id:da8398435684999113] 
\draw  [dash pattern={on 4.5pt off 4.5pt}]  (810,64.47) -- (810,530) ;
%Rounded Rect [id:dp027018508739026448] 
\draw  [draw opacity=0][fill={rgb, 255:red, 230; green, 230; blue, 230 }  ,fill opacity=0.5 ] (20,439.33) .. controls (20,445.22) and (24.78,450) .. (30.67,450) -- (409.33,450) .. controls (415.22,450) and (420,445.22) .. (420,439.33) -- (420,140.67) .. controls (420,134.78) and (415.22,130) .. (409.33,130) -- (30.67,130) .. controls (24.78,130) and (20,134.78) .. (20,140.67) -- cycle ;
%Rounded Rect [id:dp9415679498487528] 
\draw  [dash pattern={on 4.5pt off 4.5pt}] (29.83,292.77) .. controls (29.83,285.88) and (35.41,280.3) .. (42.3,280.3) -- (398.54,280.3) .. controls (405.42,280.3) and (411,285.88) .. (411,292.77) -- (411,337.17) .. controls (411,344.05) and (405.42,349.64) .. (398.54,349.64) -- (42.3,349.64) .. controls (35.41,349.64) and (29.83,344.05) .. (29.83,337.17) -- cycle ;
%Rounded Rect [id:dp05100022295548956] 
\draw  [draw opacity=0][fill={rgb, 255:red, 230; green, 230; blue, 230 }  ,fill opacity=0.5 ] (130.5,575) .. controls (130.5,583.28) and (137.22,590) .. (145.5,590) -- (615,590) .. controls (623.28,590) and (630,583.28) .. (630,575) -- (630,395) .. controls (630,386.72) and (623.28,380) .. (615,380) -- (145.5,380) .. controls (137.22,380) and (130.5,386.72) .. (130.5,395) -- cycle ;
%Rounded Rect [id:dp6496809750851564] 
\draw  [draw opacity=0][fill={rgb, 255:red, 230; green, 230; blue, 230 }  ,fill opacity=0.5 ] (450.61,440.79) .. controls (450.61,445.82) and (454.69,449.9) .. (459.72,449.9) -- (780.89,449.9) .. controls (785.92,449.9) and (790,445.82) .. (790,440.79) -- (790,139.11) .. controls (790,134.08) and (785.92,130) .. (780.89,130) -- (459.72,130) .. controls (454.69,130) and (450.61,134.08) .. (450.61,139.11) -- cycle ;
%Rounded Rect [id:dp549636265291094] 
\draw  [fill={rgb, 255:red, 230; green, 230; blue, 230 }  ,fill opacity=0.5 ][dash pattern={on 4.5pt off 4.5pt}] (29.83,42.77) .. controls (29.83,35.88) and (35.41,30.3) .. (42.3,30.3) -- (398.54,30.3) .. controls (405.42,30.3) and (411,35.88) .. (411,42.77) -- (411,87.17) .. controls (411,94.05) and (405.42,99.64) .. (398.54,99.64) -- (42.3,99.64) .. controls (35.41,99.64) and (29.83,94.05) .. (29.83,87.17) -- cycle ;
%Rounded Rect [id:dp1378581119146769] 
\draw  [fill={rgb, 255:red, 255; green, 255; blue, 255 }  ,fill opacity=1 ] (40.4,49.82) .. controls (40.4,44.29) and (44.89,39.8) .. (50.42,39.8) -- (169.58,39.8) .. controls (175.11,39.8) and (179.6,44.29) .. (179.6,49.82) -- (179.6,79.88) .. controls (179.6,85.41) and (175.11,89.9) .. (169.58,89.9) -- (50.42,89.9) .. controls (44.89,89.9) and (40.4,85.41) .. (40.4,79.88) -- cycle ;
%Straight Lines [id:da8729153203483555] 
\draw    (188.42,63.77) -- (250.99,63.77)(188.42,66.77) -- (250.99,66.77) ;
\draw [shift={(257.99,65.27)}, rotate = 180] [color={rgb, 255:red, 0; green, 0; blue, 0 }  ][line width=0.75]    (10.93,-4.9) .. controls (6.95,-2.3) and (3.31,-0.67) .. (0,0) .. controls (3.31,0.67) and (6.95,2.3) .. (10.93,4.9)   ;
\draw [shift={(181.42,65.27)}, rotate = 0] [color={rgb, 255:red, 0; green, 0; blue, 0 }  ][line width=0.75]    (10.93,-4.9) .. controls (6.95,-2.3) and (3.31,-0.67) .. (0,0) .. controls (3.31,0.67) and (6.95,2.3) .. (10.93,4.9)   ;
%Straight Lines [id:da845243590120559] 
\draw    (220.17,99.97) -- (220.17,137.64) ;
\draw [shift={(220.17,139.64)}, rotate = 270] [color={rgb, 255:red, 0; green, 0; blue, 0 }  ][line width=0.75]    (10.93,-3.29) .. controls (6.95,-1.4) and (3.31,-0.3) .. (0,0) .. controls (3.31,0.3) and (6.95,1.4) .. (10.93,3.29)   ;
%Straight Lines [id:da7497001546502899] 
\draw    (380,420.33) -- (380,478) ;
\draw [shift={(380,480)}, rotate = 270] [color={rgb, 255:red, 0; green, 0; blue, 0 }  ][line width=0.75]    (10.93,-3.29) .. controls (6.95,-1.4) and (3.31,-0.3) .. (0,0) .. controls (3.31,0.3) and (6.95,1.4) .. (10.93,3.29)   ;
%Straight Lines [id:da5689893955334542] 
\draw    (300,420) -- (460,420) ;
%Rounded Rect [id:dp18582062333838767] 
\draw  [fill={rgb, 255:red, 255; green, 255; blue, 255 }  ,fill opacity=1 ] (260.4,49.82) .. controls (260.4,44.29) and (264.89,39.8) .. (270.42,39.8) -- (389.58,39.8) .. controls (395.11,39.8) and (399.6,44.29) .. (399.6,49.82) -- (399.6,79.88) .. controls (399.6,85.41) and (395.11,89.9) .. (389.58,89.9) -- (270.42,89.9) .. controls (264.89,89.9) and (260.4,85.41) .. (260.4,79.88) -- cycle ;
%Rounded Rect [id:dp9523169763067778] 
\draw  [fill={rgb, 255:red, 255; green, 255; blue, 255 }  ,fill opacity=1 ] (470,299.92) .. controls (470,294.39) and (474.49,289.9) .. (480.02,289.9) -- (599.18,289.9) .. controls (604.71,289.9) and (609.2,294.39) .. (609.2,299.92) -- (609.2,329.98) .. controls (609.2,335.51) and (604.71,340) .. (599.18,340) -- (480.02,340) .. controls (474.49,340) and (470,335.51) .. (470,329.98) -- cycle ;
%Rounded Rect [id:dp48544030815291705] 
\draw  [fill={rgb, 255:red, 255; green, 255; blue, 255 }  ,fill opacity=1 ] (220,490) .. controls (220,484.48) and (224.48,480) .. (230,480) -- (530,480) .. controls (535.52,480) and (540,484.48) .. (540,490) -- (540,570) .. controls (540,575.52) and (535.52,580) .. (530,580) -- (230,580) .. controls (224.48,580) and (220,575.52) .. (220,570) -- cycle ;
%Rounded Rect [id:dp08222627508455749] 
\draw  [fill={rgb, 255:red, 255; green, 255; blue, 255 }  ,fill opacity=1 ] (540.4,49.82) .. controls (540.4,44.29) and (544.89,39.8) .. (550.42,39.8) -- (690.65,39.8) .. controls (696.18,39.8) and (700.67,44.29) .. (700.67,49.82) -- (700.67,79.88) .. controls (700.67,85.41) and (696.18,89.9) .. (690.65,89.9) -- (550.42,89.9) .. controls (544.89,89.9) and (540.4,85.41) .. (540.4,79.88) -- cycle ;
%Straight Lines [id:da5868346364154593] 
\draw  [dash pattern={on 4.5pt off 4.5pt}]  (810,530) -- (542,530) ;
\draw [shift={(540,530)}, rotate = 360] [color={rgb, 255:red, 0; green, 0; blue, 0 }  ][line width=0.75]    (10.93,-3.29) .. controls (6.95,-1.4) and (3.31,-0.3) .. (0,0) .. controls (3.31,0.3) and (6.95,1.4) .. (10.93,3.29)   ;
%Rounded Rect [id:dp7392784471609453] 
\draw  [fill={rgb, 255:red, 255; green, 255; blue, 255 }  ,fill opacity=1 ] (30.33,150.47) .. controls (30.33,144.58) and (35.11,139.8) .. (41,139.8) -- (400.34,139.8) .. controls (406.23,139.8) and (411,144.58) .. (411,150.47) -- (411,229.34) .. controls (411,235.23) and (406.23,240.01) .. (400.34,240.01) -- (41,240.01) .. controls (35.11,240.01) and (30.33,235.23) .. (30.33,229.34) -- cycle ;
%Rounded Rect [id:dp90901637181607] 
\draw  [fill={rgb, 255:red, 255; green, 255; blue, 255 }  ,fill opacity=1 ] (41,189.82) .. controls (41,184.29) and (45.49,179.8) .. (51.02,179.8) -- (145.02,179.8) .. controls (150.55,179.8) and (155.04,184.29) .. (155.04,189.82) -- (155.04,219.88) .. controls (155.04,225.41) and (150.55,229.9) .. (145.02,229.9) -- (51.02,229.9) .. controls (45.49,229.9) and (41,225.41) .. (41,219.88) -- cycle ;
%Rounded Rect [id:dp06115018179625542] 
\draw  [fill={rgb, 255:red, 255; green, 255; blue, 255 }  ,fill opacity=1 ] (163.65,189.82) .. controls (163.65,184.29) and (168.13,179.8) .. (173.67,179.8) -- (267.67,179.8) .. controls (273.2,179.8) and (277.69,184.29) .. (277.69,189.82) -- (277.69,219.88) .. controls (277.69,225.41) and (273.2,229.9) .. (267.67,229.9) -- (173.67,229.9) .. controls (168.13,229.9) and (163.65,225.41) .. (163.65,219.88) -- cycle ;
%Rounded Rect [id:dp6195275548153432] 
\draw  [fill={rgb, 255:red, 255; green, 255; blue, 255 }  ,fill opacity=1 ] (286.3,189.82) .. controls (286.3,184.29) and (290.78,179.8) .. (296.32,179.8) -- (390.31,179.8) .. controls (395.85,179.8) and (400.33,184.29) .. (400.33,189.82) -- (400.33,219.88) .. controls (400.33,225.41) and (395.85,229.9) .. (390.31,229.9) -- (296.32,229.9) .. controls (290.78,229.9) and (286.3,225.41) .. (286.3,219.88) -- cycle ;
%Rounded Rect [id:dp41283415100159493] 
\draw  [fill={rgb, 255:red, 255; green, 255; blue, 255 }  ,fill opacity=1 ] (40.8,299.92) .. controls (40.8,294.39) and (45.29,289.9) .. (50.82,289.9) -- (169.98,289.9) .. controls (175.51,289.9) and (180,294.39) .. (180,299.92) -- (180,329.98) .. controls (180,335.51) and (175.51,340) .. (169.98,340) -- (50.82,340) .. controls (45.29,340) and (40.8,335.51) .. (40.8,329.98) -- cycle ;
%Rounded Rect [id:dp40154854661776396] 
\draw  [fill={rgb, 255:red, 255; green, 255; blue, 255 }  ,fill opacity=1 ] (260.4,299.82) .. controls (260.4,294.29) and (264.89,289.8) .. (270.42,289.8) -- (389.58,289.8) .. controls (395.11,289.8) and (399.6,294.29) .. (399.6,299.82) -- (399.6,329.88) .. controls (399.6,335.41) and (395.11,339.9) .. (389.58,339.9) -- (270.42,339.9) .. controls (264.89,339.9) and (260.4,335.41) .. (260.4,329.88) -- cycle ;
%Straight Lines [id:da6602630444249074] 
\draw    (180.17,315.8) -- (260,315.8) ;
%Straight Lines [id:da30905866347109745] 
\draw    (219.86,315.97) -- (219.86,386.81) ;
\draw [shift={(219.86,388.81)}, rotate = 270] [color={rgb, 255:red, 0; green, 0; blue, 0 }  ][line width=0.75]    (10.93,-3.29) .. controls (6.95,-1.4) and (3.31,-0.3) .. (0,0) .. controls (3.31,0.3) and (6.95,1.4) .. (10.93,3.29)   ;
%Rounded Rect [id:dp44627062039563214] 
\draw  [fill={rgb, 255:red, 255; green, 255; blue, 255 }  ,fill opacity=1 ] (139.67,399.82) .. controls (139.67,394.29) and (144.15,389.8) .. (149.69,389.8) -- (290.31,389.8) .. controls (295.85,389.8) and (300.33,394.29) .. (300.33,399.82) -- (300.33,429.88) .. controls (300.33,435.41) and (295.85,439.9) .. (290.31,439.9) -- (149.69,439.9) .. controls (144.15,439.9) and (139.67,435.41) .. (139.67,429.88) -- cycle ;
%Straight Lines [id:da8375870493480149] 
\draw    (220,240) -- (220,277.67) ;
\draw [shift={(220,279.67)}, rotate = 270] [color={rgb, 255:red, 0; green, 0; blue, 0 }  ][line width=0.75]    (10.93,-3.29) .. controls (6.95,-1.4) and (3.31,-0.3) .. (0,0) .. controls (3.31,0.3) and (6.95,1.4) .. (10.93,3.29)   ;
%Rounded Rect [id:dp37285624114093485] 
\draw  [fill={rgb, 255:red, 255; green, 255; blue, 255 }  ,fill opacity=1 ] (460,400.02) .. controls (460,394.49) and (464.49,390) .. (470.02,390) -- (610.65,390) .. controls (616.18,390) and (620.67,394.49) .. (620.67,400.02) -- (620.67,430.08) .. controls (620.67,435.61) and (616.18,440.1) .. (610.65,440.1) -- (470.02,440.1) .. controls (464.49,440.1) and (460,435.61) .. (460,430.08) -- cycle ;
%Rounded Rect [id:dp6883616335272991] 
\draw  [fill={rgb, 255:red, 255; green, 255; blue, 255 }  ,fill opacity=1 ] (629.2,299.92) .. controls (629.2,294.39) and (633.69,289.9) .. (639.22,289.9) -- (758.38,289.9) .. controls (763.91,289.9) and (768.4,294.39) .. (768.4,299.92) -- (768.4,329.98) .. controls (768.4,335.51) and (763.91,340) .. (758.38,340) -- (639.22,340) .. controls (633.69,340) and (629.2,335.51) .. (629.2,329.98) -- cycle ;
%Straight Lines [id:da19870419404011752] 
\draw    (540,340) -- (540,388) ;
\draw [shift={(540,390)}, rotate = 270] [color={rgb, 255:red, 0; green, 0; blue, 0 }  ][line width=0.75]    (10.93,-3.29) .. controls (6.95,-1.4) and (3.31,-0.3) .. (0,0) .. controls (3.31,0.3) and (6.95,1.4) .. (10.93,3.29)   ;
%Rounded Rect [id:dp9692784599420237] 
\draw  [fill={rgb, 255:red, 255; green, 255; blue, 255 }  ,fill opacity=1 ] (230,529.92) .. controls (230,524.39) and (234.49,519.9) .. (240.02,519.9) -- (359.18,519.9) .. controls (364.71,519.9) and (369.2,524.39) .. (369.2,529.92) -- (369.2,559.98) .. controls (369.2,565.51) and (364.71,570) .. (359.18,570) -- (240.02,570) .. controls (234.49,570) and (230,565.51) .. (230,559.98) -- cycle ;
%Rounded Rect [id:dp749684165723072] 
\draw  [fill={rgb, 255:red, 255; green, 255; blue, 255 }  ,fill opacity=1 ] (390,530.02) .. controls (390,524.49) and (394.49,520) .. (400.02,520) -- (519.18,520) .. controls (524.71,520) and (529.2,524.49) .. (529.2,530.02) -- (529.2,560.08) .. controls (529.2,565.61) and (524.71,570.1) .. (519.18,570.1) -- (400.02,570.1) .. controls (394.49,570.1) and (390,565.61) .. (390,560.08) -- cycle ;
%Rounded Rect [id:dp28001422897762485] 
\draw  [fill={rgb, 255:red, 255; green, 255; blue, 255 }  ,fill opacity=1 ][dash pattern={on 4.5pt off 4.5pt}] (460,148.92) .. controls (460,143.99) and (463.99,140) .. (468.92,140) -- (771.08,140) .. controls (776.01,140) and (780,143.99) .. (780,148.92) -- (780,241.55) .. controls (780,246.48) and (776.01,250.47) .. (771.08,250.47) -- (468.92,250.47) .. controls (463.99,250.47) and (460,246.48) .. (460,241.55) -- cycle ;
%Shape: Brace [id:dp9411387904179627] 
\draw   (774.67,286.8) .. controls (774.67,282.13) and (772.34,279.8) .. (767.67,279.8) -- (629.96,279.8) .. controls (623.29,279.8) and (619.96,277.47) .. (619.96,272.8) .. controls (619.96,277.47) and (616.63,279.8) .. (609.96,279.8)(612.96,279.8) -- (472.25,279.8) .. controls (467.58,279.8) and (465.25,282.13) .. (465.25,286.8) ;
%Straight Lines [id:da9902399507274253] 
\draw    (620,273) -- (620,253.14) ;
\draw [shift={(620,251.14)}, rotate = 90] [color={rgb, 255:red, 0; green, 0; blue, 0 }  ][line width=0.75]    (10.93,-3.29) .. controls (6.95,-1.4) and (3.31,-0.3) .. (0,0) .. controls (3.31,0.3) and (6.95,1.4) .. (10.93,3.29)   ;

% Text Node
\draw (110,64.85) node  [color={rgb, 255:red, 0; green, 0; blue, 0 }  ,opacity=1 ] [align=left] {\begin{minipage}[lt]{92.93pt}\setlength\topsep{0pt}
\begin{center}
programa en Python\\
\small{(\cref{sec:programas-python})}
\end{center}

\end{minipage}};
% Text Node
\draw (219.89,51.97) node  [font=\small,color={rgb, 255:red, 0; green, 0; blue, 0 }  ,opacity=1 ] [align=left] {\begin{minipage}[lt]{54.86pt}\setlength\topsep{0pt}
\begin{center}
\Cref{sec:equivalencia}
\end{center}

\end{minipage}};
% Text Node
\draw (69.5,15.5) node  [font=\small,color={rgb, 255:red, 0; green, 0; blue, 0 }  ,opacity=1 ] [align=left] {\begin{minipage}[lt]{41.55pt}\setlength\topsep{0pt}
\cref{ch:programas-python-maquinas-turing}
\end{minipage}};
% Text Node
\draw (330,64.85) node  [color={rgb, 255:red, 0; green, 0; blue, 0 }  ,opacity=1 ] [align=left] {\begin{minipage}[lt]{92.93pt}\setlength\topsep{0pt}
\begin{center}
máquina de Turing\\
\small{(\cref{sec:maquinas-turing})}
\end{center}

\end{minipage}};
% Text Node
\draw (539.6,314.95) node  [color={rgb, 255:red, 0; green, 0; blue, 0 }  ,opacity=1 ] [align=left] {\begin{minipage}[lt]{92.93pt}\setlength\topsep{0pt}
\begin{center}
sistema lógico\\
\small{(\cref{sec:sistemas-logicos})}
\end{center}

\end{minipage}};
% Text Node
\draw (380.5,500.11) node  [color={rgb, 255:red, 0; green, 0; blue, 0 }  ,opacity=1 ] [align=left] {\begin{minipage}[lt]{159.55pt}\setlength\topsep{0pt}
\begin{center}
\textbf{Primer Teorema de Incompletitud}
\end{center}

\end{minipage}};
% Text Node
\draw (620.53,64.85) node  [color={rgb, 255:red, 0; green, 0; blue, 0 }  ,opacity=1 ] [align=left] {\begin{minipage}[lt]{106.99pt}\setlength\topsep{0pt}
\begin{center}
tesis de Church-Turing\\
\small{(\cref{sec:church-turing})}
\end{center}

\end{minipage}};
% Text Node
\draw (480.28,52.8) node  [font=\small,color={rgb, 255:red, 0; green, 0; blue, 0 }  ,opacity=1 ] [align=left] {\begin{minipage}[lt]{81.9pt}\setlength\topsep{0pt}
\begin{center}
\cref{sec:church-turing}
\end{center}

\end{minipage}};
% Text Node
\draw (589.5,518.3) node  [font=\small,color={rgb, 255:red, 0; green, 0; blue, 0 }  ,opacity=1 ] [align=left] {\begin{minipage}[lt]{54.86pt}\setlength\topsep{0pt}
\begin{center}
\Cref{subsec:mente-humana-tesis-church-turing}
\end{center}

\end{minipage}};
% Text Node
\draw (220.2,159.62) node  [color={rgb, 255:red, 0; green, 0; blue, 0 }  ,opacity=1 ] [align=left] {\begin{minipage}[lt]{230.11pt}\setlength\topsep{0pt}
\begin{center}
problema computacional
\end{center}

\end{minipage}};
% Text Node
\draw (98.02,204.85) node  [color={rgb, 255:red, 0; green, 0; blue, 0 }  ,opacity=1 ] [align=left] {\begin{minipage}[lt]{76.13pt}\setlength\topsep{0pt}
\begin{center}
decidible\\
\small{(\cref{sec:problemas-decidibles})}
\end{center}

\end{minipage}};
% Text Node
\draw (220.67,204.85) node  [color={rgb, 255:red, 0; green, 0; blue, 0 }  ,opacity=1 ] [align=left] {\begin{minipage}[lt]{76.13pt}\setlength\topsep{0pt}
\begin{center}
no decidible\\
\small{(\cref{sec:problemas-no-decidibles})}
\end{center}

\end{minipage}};
% Text Node
\draw (343.31,204.85) node  [color={rgb, 255:red, 0; green, 0; blue, 0 }  ,opacity=1 ] [align=left] {\begin{minipage}[lt]{76.13pt}\setlength\topsep{0pt}
\begin{center}
semidecidible\\
\small{(\cref{sec:problemas-semidecidibles})}
\end{center}

\end{minipage}};
% Text Node
\draw (110.4,314.95) node  [color={rgb, 255:red, 0; green, 0; blue, 0 }  ,opacity=1 ] [align=left] {\begin{minipage}[lt]{92.93pt}\setlength\topsep{0pt}
\begin{center}
universalidad\\
\small{(\cref{sec:universalidad})}
\end{center}

\end{minipage}};
% Text Node
\draw (330,314.85) node  [color={rgb, 255:red, 0; green, 0; blue, 0 }  ,opacity=1 ] [align=left] {\begin{minipage}[lt]{92.93pt}\setlength\topsep{0pt}
\begin{center}
reducción\\
\small{(\cref{sec:reducciones})}
\end{center}

\end{minipage}};
% Text Node
\draw (220,414.85) node  [color={rgb, 255:red, 0; green, 0; blue, 0 }  ,opacity=1 ] [align=left] {\begin{minipage}[lt]{107.26pt}\setlength\topsep{0pt}
\begin{center}
problema de la parada\\
\small{(\cref{sec:problema-parada})}
\end{center}

\end{minipage}};
% Text Node
\draw (540.33,415.05) node  [color={rgb, 255:red, 0; green, 0; blue, 0 }  ,opacity=1 ] [align=left] {\begin{minipage}[lt]{107.26pt}\setlength\topsep{0pt}
\begin{center}
aritmética de Peano\\
\small{(\cref{sec:aritmetica-peano})}
\end{center}

\end{minipage}};
% Text Node
\draw (698.8,314.95) node  [color={rgb, 255:red, 0; green, 0; blue, 0 }  ,opacity=1 ] [align=left] {\begin{minipage}[lt]{92.93pt}\setlength\topsep{0pt}
\begin{center}
sistema formal\\
\small{(\cref{sec:sistemas-formales})}
\end{center}

\end{minipage}};
% Text Node
\draw (299.6,544.95) node  [color={rgb, 255:red, 0; green, 0; blue, 0 }  ,opacity=1 ] [align=left] {\begin{minipage}[lt]{92.93pt}\setlength\topsep{0pt}
\begin{center}
versión semántica\\
\small{(\cref{sec:version-semantica})}
\end{center}

\end{minipage}};
% Text Node
\draw (459.6,545.05) node  [color={rgb, 255:red, 0; green, 0; blue, 0 }  ,opacity=1 ] [align=left] {\begin{minipage}[lt]{92.93pt}\setlength\topsep{0pt}
\begin{center}
versión sintáctica\\
\small{(\cref{sec:version-sintactica})}
\end{center}

\end{minipage}};
% Text Node
\draw (60.5,464.5) node  [font=\small,color={rgb, 255:red, 0; green, 0; blue, 0 }  ,opacity=1 ] [align=left] {\begin{minipage}[lt]{41.55pt}\setlength\topsep{0pt}
\cref{ch:problema-parada}
\end{minipage}};
% Text Node
\draw (749.5,465.5) node  [font=\small,color={rgb, 255:red, 0; green, 0; blue, 0 }  ,opacity=1 ] [align=left] {\begin{minipage}[lt]{41.55pt}\setlength\topsep{0pt}
\begin{flushright}
\cref{ch:sistemas-logicos}
\end{flushright}

\end{minipage}};
% Text Node
\draw (170.5,604.5) node  [font=\small,color={rgb, 255:red, 0; green, 0; blue, 0 }  ,opacity=1 ] [align=left] {\begin{minipage}[lt]{41.55pt}\setlength\topsep{0pt}
\cref{ch:teorema-incompletitud}
\end{minipage}};
% Text Node
\draw (643.76,195.55) node  [color={rgb, 255:red, 0; green, 0; blue, 0 }  ,opacity=1 ] [align=left] {\begin{minipage}[lt]{217.92pt}\setlength\topsep{0pt}
-- sólido {\small{(\cref{sec:solidez-completitud-decidibilidad})}}\\
-- decidible {\small{(\cref{sec:solidez-completitud-decidibilidad})}}\\
-- (semánticamente) completo {\small{(\cref{sec:solidez-completitud-decidibilidad})}}\\
-- sintácticamente completo {\small{(\cref{sec:consistencia})}}\\
-- sem./sint. consistente {\small{(\cref{sec:consistencia})}}
\end{minipage}};
% Text Node
\draw (414.84,645.25) node  [color={rgb, 255:red, 0; green, 0; blue, 0 }  ,opacity=1 ] [align=left] {\begin{minipage}[lt]{534.7pt}\setlength\topsep{0pt}
\begin{center}
Esquema resumen de este Trabajo de Fin de Grado
\end{center}
\end{minipage}};
% Text Node
\draw (380.15,407.8) node  [font=\small,color={rgb, 255:red, 0; green, 0; blue, 0 }  ,opacity=1 ] [align=left] {\begin{minipage}[lt]{108.93pt}\setlength\topsep{0pt}
\begin{center}
\cref{sec:primera-aproximacion}
\end{center}

\end{minipage}};

\end{tikzpicture}}
}
%\caption{Esquema resumen (extendido)}

%\label{fig:resumen-extendido}
\end{figure}
% ====================
\end{landscape}

\endinput


% --------------------------------------------------------------------
% MAINMATTER
% --------------------------------------------------------------------
\mainmatter % activa la numeración de capítulos, resetea la numeración de las páginas y usa números arábigos

\setpartpreamble[c][0.75\linewidth]{%
	\bigskip % Deja un espacio vertical en la parte superior
  Hablaremos de la motivación que subyace en este trabajo, concretando los objetivos que se pretenden conseguir. A continuación, daremos una contextualización histórica de los temas tratados.
}
\part{Introducción}\label{part:introduccion} % Dividir un libro en partes OPCIONAL

% !TeX root = ../libro.tex
% !TeX encoding = utf8

\setchapterpreamble[c][0.75\linewidth]{%
	\sffamily
  \emph{El objetivo de la computación es el conocimiento, no los números.}
 \begin{flushright} — Richard Hamming, \cite{Hamming1986} \end{flushright}
% \\[8pt]
	\par\bigskip
}
%\vspace{28pt}

\chapter{Motivación}\label{ch:motivacion}

Los teoremas de incompletitud de Gödel son uno de los resultados más populares de las matemáticas y la filosofía, y han sido objeto de fascinación y estudio durante décadas. Estos teoremas revolucionaron los mismos fundamentos de ambas disciplinas. En este trabajo, abordaremos el primero de estos teoremas desde una perspectiva computacional. Esto es así debido a una serie de experiencias y observaciones que despertaron en mí un gran interés y que revelaron una oportunidad para abordar este tema desde una perspectiva computacional pero, sobre todo, matemáticamente rigurosa.

Este interés surge de las clases de Modelos de Computación impartidas por el tutor de este trabajo, en las que mencionó de forma breve este resultado. Tras ello sentí mucha curiosidad por profundizar en el mundo de la incompletitud, hasta llegar a un excepcional libro de Douglas Hofstadter: \emph{``Gödel, Escher, Bach: An Eternal Golden Braid''} \cite{Hofstadter1999}. Esta obra maestra interdisciplinaria combina elementos de matemáticas, música y arte, y me incitó a profundizar más en estos teoremas.

La mayor parte de los libros que tratan este tema lo hacen tal y como lo hizo Gödel: desde una perspectiva aritmética, usando la numeración de Gödel en las demostraciones, o variantes. Sin embargo, en múltiples cursos de teoría de la computación se exponen demostraciones haciendo uso de algoritmos y máquinas de Turing.

Para mi sorpresa, a pesar de que esta perspectiva se estudia extensivamente y hay una gran literatura al respecto, muchas de las referencias que encontraba eran bien poco rigurosas o bien poco consistentes en la nomenclatura y en la forma de probar los teoremas. En concreto, muchas de las fuentes no distinguen de forma precisa entre los conceptos de \emph{demostración} y \emph{verdad}.

Es aquí donde vi la oportunidad de hacer un trabajo que intentase proporcionar una perspectiva computacional a los teoremas de incompletitud de Gödel, sin perder rigor y concreción por el camino. Este enfoque permitirá una comprensión más precisa y estructurada de los resultados más relevantes al respecto, así como su relación con los modelos de computación y la teoría de la computabilidad.

\pagebreak

\section*{Objetivos}

Los objetivos de este trabajo son:

\subsection*{Objetivos teóricos}

\begin{enumerate}[label={[}T\arabic*{]}]
    \item Probar el Primer Teorema de Incompletitud de Gödel desde una perspectiva computacional.
    \item Usar para ello programas en Python en lugar de descripciones de máquinas de Turing, probando que esto es posible y riguroso.
    \item Exponer las consecuencias del teorema de incompletitud en las disciplinas matemática y filosófica.
    \item Contextualizar el teorema de incompletitud en el ámbito histórico.
    \item Encontrar la relación entre el teorema de incompletitud con la tesis de Church-Turing.
    \item Precisar conceptos lógicos relacionados con la teoría de la demostración. En concreto, distinguir entre las nociones de \emph{demostración} y \emph{verdad}, y estudiar las propiedades de los sistemas de demostración.
\end{enumerate}

\subsection*{Objetivos prácticos}

\begin{enumerate}[label={[}P\arabic*{]}]
    \item Crear contenido que pueda servir de herramienta de apoyo a la docencia.
    \item Crear una librería de programas en Python ejecutables por el usuario que permitan guiar las demostraciones que aparecen en este trabajo.
\end{enumerate}

\endinput

% !TeX root = ../libro.tex
% !TeX encoding = utf8

\setchapterpreamble[c][0.75\linewidth]{%
	\sffamily
  \emph{Si una ``religión'' se define como un sistema de ideas que contiene afirmaciones indemostrables, entonces Gödel nos enseñó que las matemáticas no sólo son una religión, sino que son la única religión que puede demostrar que lo es.}
 \begin{flushright} — John D. Barrow, \cite{Barrow2011} \end{flushright}
% \\[8pt]
	\par\bigskip
}
%\vspace{28pt}

\chapter{El Primer Teorema de Incompletitud}\label{ch:teorema-incompletitud}

Los matemáticos de principios del siglo \textsc{xx}, tal y como Leibniz, Boole y Hilbert, imaginaron que sería posible determinar la veracidad de fórmulas matemáticas de forma automática. Nosotros ahora nos preguntamos lo mismo: hemos trabajado con sistemas lógicos sencillos que son decidibles, es decir, cuyas fórmulas pueden ser comprobadas verdaderas mediante programas. En el caso de sistemas lógicos de más complejidad, como la aritmética de Peano, esto no es posible. Esto se debe a que la existencia de problemas computacionales indecidibles conduce directamente a la existencia de afirmaciones matemáticas no decidibles.

%En la \cref{sec:primera-aproximacion}, probaremos la incompletitud de \textbf{Peano}. Llamaremos a este resultado \emph{Primer Teorema de Incompletitud}. A continuación, generalizaremos este teorema para otros sistemas lógicos. Finalmente, en la \cref{sec:diferencias-godel}, veremos el enunciado original del teorema y lo compararemos con el resultado obtenido.

\section{Una primera aproximación}\label{sec:primera-aproximacion}

%\begin{teorema}[Versión no constructiva del Primer Teorema de Incompletitud]\label{teo:inco1-no-constr}
%La aritmética de Peano es incompleta: existen fórmulas verdaderas de \textbf{Peano} que no son teoremas.
%\end{teorema}
%\begin{proof}
%Asumiremos que \textbf{Peano} es completo y entraremos en contradicción. En concreto, demostraremos que si la aritmética de Peano es completa, podremos decidir \textbf{ParadaEnVacío}, lo cual contradice la \cref{prop:parada-en-vacio-no-decidible}.

%Sea $P$ un programa cualquiera. Lo usaremos como entrada al \cref{lst:parada-en-vacio-desde-peano-completo}. En la \cref{line:parada-en-vacio-desde-peano-completo-import1} importamos la función \texttt{es\_prueba\_peano}, que dada una fórmula y una demostración en \textbf{Peano} comprueba si la demostración prueba que la fórmula es un teorema. Claramente, esta función no es un oráculo. De hecho, su implementación es sencilla: basta probar si cada una de las fórmulas de la demostración es un axioma o se sigue de las fórmulas anteriores a través de las reglas de inferencia del sistema, comprobando que la última fórmula coincida con la fórmula que queremos probar.

%A continuación, importamos en la \cref{line:parada-en-vacio-desde-peano-completo-import2} la función \texttt{parada\_a\_peano} cuya existencia comentamos tras la \cref{prop:parada-a-peano}, que pasa la afirmación ``$P$ para con entrada vacía'' a una fórmula en \textbf{Peano}.

%Usamos esta función en la \cref{line:parada-en-vacio-desde-peano-completo-para-en-peano} para generar la fórmula mencionada anteriormente. En la \cref{line:parada-en-vacio-desde-peano-completo-no-para-en-peano} creamos una nueva fórmula negando la anterior, por tanto \texttt{no\_para\_en\_peano} contiene la fórmula de \textbf{Peano} correspondiente a la afirmación ``$P$ no para con entrada vacía''.

%Una de estas afirmaciones ha de ser verdadera. Como asumimos que \textbf{Peano} es completo, una de estas fórmulas tiene una prueba (es un teorema) en \textbf{Peano}. El bucle de la \cref{line:parada-en-vacio-desde-peano-completo-while} va iterando a través de todos los \emph{strings} posibles, comprobando cada uno para saber si es una prueba de \texttt{para\_en\_peano} o de \texttt{no\_para\_en\_peano}. Ya que uno de ellos debe tener una prueba, el programa terminará, decidiendo si $P$ para con entrada vacía o no. El \cref{lst:parada-en-vacio-desde-peano-completo} resuelve \textsc{ParadaEnVacío}, lo que entra en contradicción con la \cref{prop:parada-en-vacio-no-decidible}.
%\end{proof}

%\begin{lstlisting}[language=Python, caption=\lstinline{parada_en_vacio_desde_peano_completo.py},label={lst:parada-en-vacio-desde-peano-completo}]
%import utilidades

%from es_prueba_peano import es_prueba_peano # NO es un oráculo|\label{line:parada-en-vacio-desde-peano-completo-import1}|
%from parada_a_peano import parada_a_peano # NO es un oráculo|\label{line:parada-en-vacio-desde-peano-completo-import2}|

%def parada_en_vacio_desde_peano_completo(programa):|\label{line:parada-en-vacio-desde-peano-completo-main}|
%    para_en_peano = parada_a_peano(programa)|\label{line:parada-en-vacio-desde-peano-completo-para-en-peano}|
%    no_para_en_peano = 'NOT (' + para_en_peano + ')'|\label{line:parada-en-vacio-desde-peano-completo-no-para-en-peano}|
%    prueba = ''|\label{line:parada-en-vacio-desde-peano-completo-prueba}|
    
%    while True:|\label{line:parada-en-vacio-desde-peano-completo-while}|
%        if es_teorema_peano(prueba, para_en_peano) == 'sí':|\label{line:parada-en-vacio-desde-peano-completo-if1}|
%            return 'sí'
%        if es_teorema_peano(prueba, no_para_en_peano) == 'sí':|\label{line:parada-en-vacio-desde-peano-completo-if2}|
%            return 'no'
%        prueba = utilidades.siguiente_string(prueba)|\label{line:parada-en-vacio-desde-peano-completo-else}|
%\end{lstlisting}


Antes de proceder la demostración de una primera versión del Primer Teorema de Incompletitud, discutiremos de forma detallada el \cref{lst:godel}. En la \cref{line:godel-peano-import1} importamos la función \texttt{es\_teorema\_peano}, cuya existencia probamos en la \cref{prop:es-teorema-peano-semidecidible}. Recuerda que, como los axiomas de \textbf{Peano} son recursivamente enumerables (es decir, \textbf{Peano} es recursivamente axiomatizable), podemos escribir el programa \texttt{es\_teorema\_peano(fórmula)} que devuelve \palabra{sí} si \texttt{fórmula} es un teorema en \textbf{Peano}, y que no para en caso contrario (ver \cref{prop:es-teorema-peano-semidecidible}).
\vspace{8pt}
\begin{lstlisting}[language=Python, caption=\lstinline{godel_peano.py},label={lst:godel-peano}]
import utilidades

from es_teorema_peano import es_teorema_peano # NO es un oráculo|\label{line:godel-peano-import1}|
from parada_a_peano import parada_a_peano # NO es un oráculo|\label{line:godel-peano-import2}|

def godel_peano(entrada):|\label{line:godel-peano-main}|
    programa_godel = utilidades.leer('godel\_peano.py')|\label{line:godel-peano-programa}|
    para_en_peano = parada_a_peano(programa_godel)|\label{line:godel-peano-formula}|
    no_para_en_peano = 'NOT (' + para_en_peano + ')'|\label{line:godel-peano-formulaneg}|

    if es_teorema_peano(no_para_en_peano) == 'sí':|\label{line:godel-peano-es-teorema}|
        return 'para'|\label{line:godel-peano-para}|
    else:
        utilidades.ciclar()|\label{line:godel-peano-cicla}|
\end{lstlisting}

La \cref{line:godel-peano-import2} importa la función \texttt{parada\_a\_peano}, cuya existencia ya discutimos tras la \cref{prop:parada-a-peano}. Esta función traduce, dado un programa $P$, la afirmación ``$P$ para con entrada vacía'' en una fórmula de \textbf{Peano}.

En la \cref{line:godel-peano-programa} guardamos el propio programa \texttt{godel\_peano.py} en la variable \texttt{programa\_godel}. En la \cref{line:godel-peano-formula} usamos la función \texttt{parada\_a\_peano} recién comentada, y en la \cref{line:godel-peano-formulaneg} negamos la fórmula \texttt{para\_en\_peano}. De este modo:
\begin{adjustwidth}{30pt}{}
    \texttt{no\_para\_en\_peano} significa que ``\texttt{godel\_peano.py} no para con entrada vacía'' \hfill $(\star1)$
\end{adjustwidth}
Ahora es cuando las cosas se ponen interesantes. Si \texttt{no\_para\_en\_peano} es un teorema, la función \texttt{es\_teorema\_peano} devolverá \palabra{sí}, y de este modo el programa verificará la condición de la \cref{line:godel-peano-es-teorema} y devolverá \palabra{para}, haciendo que el programa pare. Si \texttt{no\_para\_en\_peano} no es un teorema, se entrará en un bucle infinito (\texttt{es\_teorema\_peano} no parará, y por lo tanto tampoco \texttt{godel\_peano.py}). Es claro que la \cref{line:godel-peano-cicla} nunca se ejecutará, pero la incluimos para enfatizar este hecho.

Ahora, nos preguntamos: ¿cuál es el comportamiento de \texttt{godel\_peano.py} con entrada vacía?

\begin{proposicion}\label{prop:godel-peano-no-para}
El \cref{lst:godel-peano} no para con entrada vacía.
\end{proposicion}
\begin{proof}
Asumimos que \texttt{godel\_peano.py} para con entrada vacía, y entraremos en contradicción. La única forma de que pare es en la \cref{line:godel-peano-para}, algo que ocurre solamente si:
\begin{adjustwidth}{30pt}{}
    \texttt{es\_teorema\_peano(no\_para\_en\_peano) == \palabra{sí}} \hfill $(\star2)$
\end{adjustwidth}
$(\star2)$ es equivalente a decir que \texttt{no\_para\_en\_peano} es un teorema. Por la solidez de \textbf{Peano}, \texttt{no\_para\_en\_peano} es verdadero.

Pero el hecho de que \texttt{no\_para\_en\_peano} sea verdadero, junto con $(\star1)$, nos dice que \texttt{godel\_peano.py} no para con entrada vacía, encontrando la contradicción esperada.
\end{proof}

Ahora estamos en condiciones de probar el resultado principal de este trabajo.

\begin{teorema}[Incompletitud de la aritmética de Peano]\label{teo:incompletitud-peano}
\index{primer teorema de incompletitud}\index{teorema de incompletitud}La aritmética de Peano es incompleta: existen fórmulas verdaderas de \normalfont{\textbf{Peano}} que no son teoremas. Concretamente, la fórmula guardada en la variable \texttt{no\_para\_en\_peano} en la \cref{line:godel-peano-formulaneg} del \cref{lst:godel} es un ejemplo de una fórmula verdadera pero no demostrable (no es un teorema) de \normalfont{\textbf{Peano}}.
\end{teorema}
\begin{proof}
Debemos mostrar que \texttt{no\_para\_en\_peano} es verdadero y no es un teorema.

Para probar que es verdadero, por la \cref{prop:godel-peano-no-para}, sabemos que \texttt{godel\_peano.py} no para con entrada vacía. Recordando $(\star1)$, inmediatamente obtenemos que \texttt{no\_para\_en\_peano} es verdadero.

Para probar que no es un teorema, de nuevo usamos el hecho de que \texttt{godel\_peano.py} no para con entrada vacía. Esto nos indica que, en la \cref{line:godel-peano-es-teorema}, la llamada a la función
\begin{adjustwidth}{30pt}{}
    \texttt{es\_teorema\_peano(no\_para\_en\_peano)}
\end{adjustwidth}
no puede devolver \palabra{sí}. Pero esto inmediatamente implica que \texttt{no\_para\_en\_peano} no es demostrable, lo que completa la prueba.
\end{proof}

%\section{El Segundo Teorema de Incompletitud}\label{sec:teo-incompletitud-2}

%\begin{teorema}[Segundo Teorema de Incompletitud]
%La consistencia de la aritmética de Peano no puede probarse en la aritmética de Peano. En concreto, la afirmación
%\begin{adjustwidth}{30pt}{}
%    ``\textbf{Peano} es consistente''
%\end{adjustwidth}
%que puede codificarse mediante una fórmula de \textbf{Peano}, no es un teorema.
%\end{teorema}

Acabamos de demostrar que \textbf{Peano} es incompleto, como ya anticipamos en la \cref{tab:sistemas-logicos}. Las consecuencias de este hecho las exploramos en detalle en el \cref{ch:conclusion}. Sin embargo, paremos un momento a analizar qué es lo que acabamos de probar.

\section{Versión semántica}\label{sec:version-semantica}

Para encontrar una fórmula verdadera y no demostrable, hemos necesitado un sistema lógico que sea, en primer lugar, recursivamente axiomatizable. Si no, no hubiesemos podido crear un programa capaz de saber cuándo una fórmula del sistema es un teorema. Por otra parte, hemos asumido que el sistema es sólido, lo cual es algo esperable. Finalmente, hemos exigido que sea capaz de expresar una cierta afirmación de las máquinas de Turing (véase \cref{prop:parada-a-peano}). Para evitar repeticiones innecesarias en esta última hipótesis, llamaremos a un sistema lógico con esta propiedad ``de tipo Turing'', algo que definimos a continuación.

\begin{definicion}[Sistema de tipo Turing]
Sea $\mathcal{S}$ un sistema formal. Diremos que $\mathcal{S}$ es \emph{de tipo Turing}\footnote{Esta definición no se encuentra en la literatura: se ha creado por conveniencia.} si puede expresar afirmaciones sobre máquinas de Turing o, equivalentemente, sobre programas;\footnote{La equivalencia la probamos en el \cref{teo:equivalencia}.} en concreto, dada una máquina de Turing $M$, puede expresar que ``$M$ para con entrada vacía''; o, equivalentemente, dado un programa $P$, puede expresar que ``$P$ para con entrada vacía''.

Análogamente podemos definirlo para el caso de sistemas lógicos.
\end{definicion}

Observemos que, sin embargo, en ningún momento hemos utilizado el sistema \textbf{Peano} en sí. Esto nos indica que podríamos realizar esta demostración con otro sistema que no sea \textbf{Peano}, siempre y cuando verifique tales condiciones. Indicamos este resultado a continuación.

\begin{teorema}[Versión semántica del Primer Teorema de Incompletitud]\label{teo:incompletitud-semantica}
Sea $\mathcal{G}$ un sistema lógico que verifique las siguientes condiciones:
\begin{enumerate}
    \item $\mathcal{G}$ es sólido.
    \item $\mathcal{G}$ es recursivamente axiomatizable.
    \item $\mathcal{G}$ es de tipo Turing.
\end{enumerate}
Entonces, $\mathcal{G}$ es incompleto, es decir, existen fórmulas verdaderas de $\mathcal{G}$ que no son teoremas.
\end{teorema}

Para reforzar este importante resultado, incluimos un resumen del proceso que hemos seguido para \textbf{Peano}, de forma general, en la demostración de este teorema.

\begin{proof}
En el \cref{lst:godel} importamos dos funciones. Por una parte, en la \cref{line:godel-import1} importamos la función \texttt{es\_teorema\_en\_G}, que comprueba si una fórmula es un teorema en $\mathcal{G}$. Esta función parará si la fórmula es un teorema, y se ejecutará indefinidamente en caso de que no lo sea. La existencia de esta función la garantiza (2).

Por otra, en la \cref{line:godel-import2} importamos la función \texttt{parada\_a\_G} que, dado un programa $P$, convierte la afirmación ``$P$ para con entrada vacía'' en una fórmula de $\mathcal{G}$, algo que es posible por (3).

En la línea \cref{line:godel-formulaneg} tenemos que:
\begin{adjustwidth}{30pt}{}
    \texttt{no\_para\_en\_G} significa que \emph{``\texttt{godel.py} no para con entrada vacía''} (en $\mathcal{G}$)  \hfill $(\star3)$
\end{adjustwidth}
Si \texttt{no\_para\_en\_G} es un teorema, la función \texttt{es\_teorema\_en\_G(no\_para\_en\_G)} (\cref{line:godel-es-teorema}) parará y devolverá \palabra{sí}, mientras que la misma función ciclará en caso de que no lo sea.

La pregunta que debemos hacernos es: ¿qué le pasa a \texttt{godel.py} cuando \texttt{entrada} es vacía?

Si asumimos que \texttt{godel.py} para con entrada vacía, la única forma de que lo haga es porque en la \cref{line:godel-es-teorema} es \texttt{es\_teorema\_en\_G(no\_para\_en\_G) == \palabra{sí}}, lo cual solo puede pasar si \texttt{no\_para\_en\_G} es un teorema. Pero, por (1), y dado que $\mathcal{G}$ es sólido, tenemos que \texttt{no\_para\_en\_G} es verdadero.

Pero si \texttt{no\_para\_en\_G} es verdadero, por $(\star3)$, esto querría decir que \texttt{godel.py} no para con entrada vacía, entrando en contradicción, por lo que debe ser:
\begin{adjustwidth}{30pt}{}
    \texttt{godel.py} no para con entrada vacía  \hfill $(\star4)$
\end{adjustwidth}
Comprobamos que \texttt{no\_para\_en\_G} es verdadero pero no es un teorema. En efecto, es verdadero por $(\star4)$. Sin embargo, no es un teorema. Para ello, usamos de nuevo $(\star4)$. El hecho de que \texttt{godel.py} no pare indica que es imposible que la condición de la \cref{line:godel-es-teorema}, \linebreak\texttt{es\_teorema\_en\_G(no\_para\_en\_G) == \palabra{sí}}, se verifique -- en otras palabras, \texttt{no\_para\_en\_G} no es un teorema.
\end{proof}

%\vspace{8pt}
\begin{lstlisting}[language=Python, caption=\lstinline{godel.py},label={lst:godel}]
import utilidades

from es_teorema_G import es_teorema_G # NO es un oráculo|\label{line:godel-import1}|
from parada_a_G import parada_a_G # NO es un oráculo|\label{line:godel-import2}|

def incompleto_semantico(entrada):|\label{line:godel-main}|
    programa = utilidades.leer('incompleto_semantico.py')|\label{line:godel-programa}|
    para_en_G = parada_a_G(programa)|\label{line:godel-formula}|
    no_para_en_G = 'NOT (' + para_en_G + ')'|\label{line:godel-formulaneg}|

    if es_teorema_G(no_para_en_G) == 'sí':|\label{line:godel-es-teorema}|
        return 'para'|\label{line:godel-para}|
    else:
        utilidades.ciclar()|\label{line:godel-cicla}|
\end{lstlisting}


Nuestro resultado nos dice que, si tenemos un sistema lógico que sea recursivamente axiomatizable y de tipo Turing, no podemos esperar que pueda ser sólido y completo al mismo tiempo. Es decir, si todos los teoremas del sistema son verdaderos, siempre podremos encontrar fórmulas del sistema que sean verdaderas pero no sean teoremas.

Podemos aplicar el \cref{teo:incompletitud-semantica} a otros sistemas lógicos, siempre que verifiquen las hipótesis. Uno de ellos es, por ejemplo, el sistema de axiomas de Zermelo-Fraenkl.

\section{Versión sintáctica}\label{sec:version-sintactica}

% WIP relajar condición de solidez?
El resultado que acabamos de probar es semántico. Sin embargo, es posible probar otro resultado, esta vez sintáctico, y más parecido al demostrado por Gödel.

Para poder probar esta variación, debemos de tener en cuenta una sutileza del problema \textsc{EsTeorema} en el caso de sistemas formales sintácticamente consistentes.

\begin{proposicion}\label{prop:es-teorema-consistente-decidible}
Sea $\mathcal{S}$ un sistema formal recursivamente axiomatizable y sintácticamente consistente. Entonces, $\textsc{EsTeorema}_\mathcal{S}$ es decidible.
\end{proposicion}
\begin{proof}
Veamos que el \cref{lst:es-teorema-s} decide $\textsc{EsTeorema}_\mathcal{S}$. En efecto, vemos en primer lugar que en la \cref{line:es-teorema-s-import} importamos la función \texttt{es\_prueba\_S}, la cual es computable por ser $\mathcal{S}$ recursivamente axiomatizable.

A continuación, en la \cref{line:es-teorema-s-negacion} almacenamos la negación de \texttt{formula} en \texttt{neg\_formula}. En la \cref{line:es-teorema-s-demostracion} declaramos la demostración.

Por ser $\mathcal{S}$ consistente, bien \texttt{formula} bien \texttt{neg\_formula} tiene una demostración. En el bucle de la \cref{line:es-teorema-s-while}, vamos comprobando si \texttt{demostracion} es una demostración de \texttt{formula} o de \texttt{neg\_formula}. En el primer caso, devolvemos \palabra{sí}, mientras que en el segundo, devolvemos \palabra{no}. En caso de que \texttt{demostracion} no sea una demostración de ninguna de ambas fórmulas, procedemos en la \cref{line:es-teorema-s-siguiente-string} a enumerar la siguiente demostración en orden \emph{shortlex} (ver \cref{prop:es-teorema-peano-semidecidible}).

Por la consistencia de $\mathcal{S}$, debe existir una demostración bien de \texttt{formula} o de \texttt{neg\_formula}, por lo que el programa debe parar al llegar a ser \texttt{demostracion} tal demostración.
\end{proof}

\begin{lstlisting}[language=Python, caption=\lstinline{es_teorema_S.py},label={lst:es-teorema-s}]
import utilidades
from es_prueba_S import es_prueba_S # NO es un oráculo|\label{line:es-teorema-s-import}|

def es_teorema_S(formula):|\label{line:es-teorema-s-main}|
    negacion_formula = 'NOT (' + formula + ')'|\label{line:es-teorema-s-negacion}|
    demostracion = ''|\label{line:es-teorema-s-demostracion}|

    while True:|\label{line:es-teorema-s-while}|
        if es_prueba_S(demostracion, formula) == 'sí':|\label{line:es-teorema-s-es-prueba}|
            return 'sí'
        if es_prueba_S(demostracion, negacion_formula) == 'sí':
            return 'no'
        demostracion = utilidades.siguiente_string(demostracion)|\label{line:es-teorema-s-siguiente-string}|
\end{lstlisting}

Ahora estamos en condiciones de probar el resultado sintáctico, que dependerá de la no decidibilidad de \textsc{AdivinaConsistente} (\cref{prob:adivina-consistente}). \cite{Aaronson2017}

\begin{teorema}[Versión sintáctica del Primer Teorema de Incompletitud]\label{teo:incompletitud-sintactica}
Sea $\mathcal{S}$ un sistema formal que verifique las siguientes condiciones:
\begin{enumerate}
    \item $\mathcal{S}$ es sintácticamente consistente.
    \item $\mathcal{S}$ es recursivamente axiomatizable.
    \item $\mathcal{S}$ es de tipo Turing.
\end{enumerate}
Entonces, $\mathcal{S}$ es sintácticamente incompleto, es decir, existen fórmulas $\phi$ de $\mathcal{S}$ tal que ni $\phi$ ni $\neg\phi$ son teoremas.
\end{teorema}
\begin{proof}
Vamos a suponer que $\mathcal{S}$ es sintácticamente completo y veremos que, en tal caso, podríamos decidir \textsc{AdivinaConsistente} (\cref{prob:adivina-consistente}), lo cual entraría en contradicción con la \cref{prop:adivina-consistente-no-decidible}.

Comencemos comentando las funciones importadas. En la \cref{line:adivina-consistente-import-1} importamos la función \texttt{oarada\_en\_S}, que traduce para un programa $P$ la afirmación ``$P$ para con entrada vacía'' en una fórmula de $\mathcal{}S$, y cuya existencia garantiza la hipótesis (3). En la \cref{line:adivina-consistente-import-2} importamos la función \texttt{es\_teorema\_S}, que decide el problema $\textsc{EsTeorema}_\mathcal{S}$, y que, por las hipótesis (1) y (2), siempre para (ver \cref{prop:es-teorema-consistente-decidible}) y devuelve \palabra{sí} si la fórmula de entrada es un teorema en $\mathcal{S}$ y \palabra{no} en caso contrario.

En la \cref{line:adivina-consistente-programa-para} usamos la función \texttt{parada\_a\_S} para almacenar en \texttt{programa\_para} la fórmula de $\mathcal{S}$ equivalente a ``\texttt{programa} para con entrada vacía''. En la \cref{line:adivina-consistente-es-teorema}, asignamos a la variable \texttt{es\_teorema} uno de dos valores posibles: \palabra{sí} o \palabra{no}, como ya comentamos anteriormente.

Llegado a la \cref{line:adivina-consistente-si}, implementamos la siguiente lógica:
\begin{itemize}
    \item Si \texttt{programa} acepta con entrada vacía (para y devuelve \palabra{sí}), la condición de la \cref{line:adivina-consistente-si} se verificará, así como la de la \cref{line:adivina-consistente-si-si}, y \texttt{adivina\_consistente.py} devolverá \palabra{sí}.
    \item Si \texttt{programa} rechaza con entrada vacía (para y devuelve \palabra{no}, la condición de la \cref{line:adivina-consistente-si} se verificará, mientras que no la de la \cref{line:adivina-consistente-si-si}, y \texttt{adivina\_consistente.py} devolverá \palabra{no}.
    \item Si \texttt{programa} cicla con entrada vacía, la condición de la \cref{line:adivina-consistente-si} no se verificará, y \texttt{adivina\_consistente.py} devolverá \palabra{no}.
\end{itemize}
Dicha lógica hace evidente que el programa decide \textsc{AdivinaConsistente}.
\end{proof}

\begin{lstlisting}[language=Python, caption=\lstinline{adivina_consistente.py},label={lst:adivina-consistente}]
import utilidades
from parada_a_S import parada_a_S # NO es un oráculo|\label{line:adivina-consistente-import-1}|
from es_teorema_S import es_teorema_S # NO es un oráculo|\label{line:adivina-consistente-import-2}|

def adivina_consistente(programa):|\label{line:adivina-consistente-main}|
    programa_para = parada_a_S(programa)|\label{line:adivina-consistente-programa-para}|

    es_teorema = es_teorema_S(programa_para)|\label{line:adivina-consistente-es-teorema}|

    if es_teorema == 'sí':|\label{line:adivina-consistente-si}|
        salida = maquina_universal(programa, '')|\label{line:adivina-consistente-universal}|
        if salida == 'sí':|\label{line:adivina-consistente-si-si}|
            return 'sí'
        else:|\label{line:adivina-consistente-si-no}|
            return 'no'
    else:|\label{line:adivina-consistente-no}|
        return 'no'
\end{lstlisting}

Hemos probado dos versiones del Primer Teorema de Incompletitud. Resumimos las diferencias entre ambas a continuación.
\vspace{8pt}
\begin{figure}[H]
$$
\begin{array}{ccc}
    {
        \begin{matrix}
            \text{\textbf{versión semántica}} \\
            \text{\small{(\cref{teo:incompletitud-semantica})}}\\
            \small{\;}
        \end{matrix}
    }
    & \;\;\; &
    {
        \begin{matrix}
            \text{\textbf{versión sintáctica}} \\
            \text{\small{(\cref{teo:incompletitud-sintactica})}}\\
            \small{\;}
        \end{matrix}
    }
    \\
    {
        \begin{matrix}
            \mathcal{G}\text{ sistema lógico}\\
            \text{rec. axiom., de tipo Turing}\\
            \text{y sólido} \\
            \small{\phi \text{ teorema} \Rightarrow \phi\text{ verdadero}}\\
            \tiny{\;}
        \end{matrix}
    }
    &  & 
    {
        \begin{matrix}
            \mathcal{S}\text{ sistema formal}\\
            \text{rec. axiom., de tipo Turing}\\
            \text{y sintácticamente consistente} \\
            \small{\nexists \phi : \phi\text{ y }\neg\phi \text{ son teoremas}}\\
            \tiny{\;}
        \end{matrix}
    } \\
    \Huge{\Downarrow} & & \Huge{\Downarrow} \\
    {
        \begin{matrix}
            \tiny{\;}\\
            \mathcal{G}\text{ no (semánticamente) completo}\\
            \small{\exists\phi:\phi\text{ es verdadero y no es teorema}}
        \end{matrix}
    }
    & &
    {
        \begin{matrix}
            \tiny{\;}\\
            \mathcal{S}\text{ no sintácticamente completo}\\
            \small{\exists\phi:\text{ni }\phi\text{ ni }\neg\phi\text{ son teoremas}}
        \end{matrix}
    }
    \\
\end{array}
$$
\caption{Comparativa entre las dos versiones del Primer Teorema de Incompletitud}
\vspace{-8pt}
\label{fig:comparativa-incompletitud}
\end{figure}

La demostración original de Gödel es puramente sintáctica, y muy parecida a la probada en el \cref{teo:incompletitud-sintactica}. La enunciamos a continuación.

\begin{teorema}[Primer Teorema de Incompletitud de Gödel-Rosser]\label{teo:incompletitud-godel}
Sea $\mathcal{S}$ un sistema formal que verifique las siguientes condiciones:
\begin{enumerate}
    \item $\mathcal{S}$ es sintácticamente consistente.
    \item $\mathcal{S}$ es recursivamente axiomatizable.
    \item $\mathcal{S}$ contiene una cierta cantidad de aritmética elemental.
\end{enumerate}

Entonces, $\mathcal{S}$ es sintácticamente incompleto, es decir, hay fórmulas $\phi$ de $\mathcal{S}$ de modo que ni $\phi$ ni $\neg\phi$ son teoremas.
\end{teorema}

Tal y como en nuestra versión sintáctica, requerimos que el sistema sea sintácticamente consistente y recursivamente axiomatizable. El teorema original de Gödel exigía como hipótesis una forma más fuerte de consistencia sintáctica llamada \emph{$\omega$-consistencia}. Esta forma más general del teorema se debe al trabajo de John B. Rosser, autor del hoy conocido como \emph{truco de Rosser}.\footnote{Es por ello por lo que lo llamamos teorema de ``Gödel-Rosser'', para distinguirlo del resultado original de Gödel.} \cite{Rosser1936}

La tercera hipótesis exige que el sistema sea capaz de probar ciertas afirmaciones de aritmética elemental. Esto se debe a que Gödel pretendía probar que la teoría de números no era axiomatizable (ver \cref{ch:historia}), por lo que usó en su demostración una forma de numerar los teoremas que hoy en día se conoce como \emph{numeración de Gödel}.


\section{Consecuencias}\label{sec:consecuencias}

El Primer Teorema de Incompletitud tiene profundas implicaciones matemáticas y filosóficas.

\subsection{Relación con el segundo problema de Hilbert}\label{subsec:segundo-problema-hilbert}

Ya introdujimos en el \cref{ch:historia} los célebres problemas de Hilbert, y en concreto, el segundo de ellos, que volveremos a referenciar a continuación:

\begin{enumerate}[label=(\arabic*)]
    \item ¿Son completas las matemáticas? Esto es, ¿puede probarse o no cada sentencia matemática?
    \item ¿Son consistentes las matemáticas? Esto es, ¿no es posible probar simultáneamente una afirmación y su negación?
    \item ¿Son decidibles las matemáticas? Esto es, ¿existe un método automático que pueda aplicarse a cualquier afirmación matemática, y que determine si es cierta? Este problema se conoce como \emph{Entscheidungsproblem}.\index{Entscheidungsproblem}
\end{enumerate}

Es necesario puntualizar que la formulación de Hilbert de este problema es relativa a la aritmética. Es claro que, para tal caso, el \cref{teo:incompletitud-godel} nos garantiza la existencia de fórmulas no demostrables, respondiendo negativamente a (1) \cite{Hellman1981}. Por otra parte, en la \cref{prop:es-teorema-consistente-decidible} probamos que \textbf{Peano}, un sistema de aritmética, es no decidible, probando que la respuesta a la pregunta (3) es negativa. La pregunta (2) queda fuera del objeto de este trabajo.

Sin embargo, es interesante ver estas preguntas sin necesidad de limitarnos a la aritmética. Hemos introducido los sistemas lógicos justamente para poder formalizar el concepto de verdad y demostración, que son el alma de las matemáticas.

Podremos encontrar respuestas análogas para sistemas formales recursivamente enumerables y de tipo Turing. En tal caso, el \cref{teo:incompletitud-sintactica} responde negativamente a la pregunta (1), mientras que una generalización de la \cref{prop:es-teorema-peano-no-decidible} (similar a la que hicimos en el \cref{teo:incompletitud-semantica} a partir del \cref{teo:incompletitud-peano}) nos permite responder negativamente a la pregunta (3).

%Tenemos respuestas a las preguntas (1) y (3) para unas hipótesis concretas pero, ¿qué pasa en el caso de sistemas no recursivamente enumerables y que no son de tipo Turing? Algunos matemáticos no están convencidos de que los resultados de Gödel resuelvan definitivamente el segundo problema de Hilbert, pues los teoremas de Gödel no se aplican en otros sistemas de lógica.

\subsection{Relación con la mente humana y con la tesis de Church-Turing}\label{subsec:mente-humana-tesis-church-turing}

Algunos autores entran en conclusiones mucho más epistemológicas, y que conciernen a la mente humana. Roger Penrose y John Lucas \cite{Penrose1989, Lucas1961} argumentan que el papel de la conciencia es no algorítmico cuando formamos juicios matemáticos, pudiendo generalizarlo al papel de la conciencia humana.

Para ello, los autores comentan que en cualquier sistema formal existen proposiciones indemostrables, esto es, para las que su algoritmo no puede proporcionar una respuesta. Si el funcionamiento de la mente humana fuese completamente algorítmico, entonces el algoritmo (o sistema formal) que utilice para formar sus juicios no podría tratar tal proposición. Sin embargo, la mente humana puede deducir que tal proposición es verdadera. Esta contradicción nos puede indicar que la mente humana, en cuanto al razonamiento matemático, no estaba utilizando un algoritmo en primer lugar.

Douglas Hofstadter \cite{Hofstadter1999, Hofstadter2007} presenta los teoremas de Gödel como un ejemplo de ``extraño bucle'', una estructura jerárquica y autoreferencial que existe en un sistema formal. Argumenta que esta es la misma estructura que origina la conciencia, el sentido del ``yo'', en la mente humana: la autoreferencia proviene de la forma en que el cerebro abstrae y categoriza los estímulos en ``símbolos'', o grupos de neuronas, que corresponden a conceptos, en lo que es de facto un sistema formal, que eventualmente forma símbolos que corresponden a la propia entidad que realiza la percepción. 

La diferencia entre las matemáticas y la conciencia reside, según Hofstdadter, en el sentido en el que se realiza el razonamiento: si es ``hacia arriba'' o ``hacia abajo''., desde los axiomas hacia las proposiciones y teoremas más complejas, o ``hacia abajo'', en el sentido contrario.

Las matemáticas funcionan hacia arriba: para derivar si una fórmula es un teorema debemos partir de los axiomas e ir razonando hasta llegar a ella. La mente humana, según el físico, funciona a hacia abajo, desde un nivel más alto --de deseos, conceptos, personalidades, ideas y pensamientos-- hasta el nivel más bajo de interacciones entre neuronas e incluso partículas fundamentales. Los seres humanos estamos hechos para percibir ``cosas grandes'' en lugar de ``cosas pequeñas'', incluso a pesar de que en el dominio de lo pequeño parece residir el motor de la realidad.

Ambos argumentos están íntimamente relacionados con la tesis de Church-Turing (véase la \cref{sec:church-turing}). Si la mente humana es una máquina de Turing, todos los procedimientos de nuestro cerebro no serían más que algoritmos, y mediante los teoremas de Gödel, nuestra forma de razonar sería limitada --¡existen afirmaciones que no podemos demostrar!--. Esto puede inducir en conclusiones más fantásticas, hablando algunos autores de que los teoremas de Gödel, aunque no permitan probarlo, son una fuerte evidencia de la existencia de dios.\footnote{No existe consenso respecto si esto es cierto, pero es sencillo argumentar los fallos de este razonamiento (véase \cite{StanfordGodel2022}).} \cite{StanfordGodel2022}

%Cerebros y máquinas

%- Autores incluyendo J. R. Lucas y Roger Penrose debaten que los teoremas tienen implicaciones sobre la inteligencia humana. El debate se centra sobre si la mente humana es equivalente a una máquina de Turing o (mediante la tesis de Church-Turing) a una máquina finita. Si la mente lo es, y los teoremas de Gödel se aplican a esta, los teoremas de incompletitud se aplicarían a ella.

%Hilary Putnam 

%Avi Widgerson

%Douglas Hofstadter 

%- 

%Consecuencias severas para el logicismo de Gottlob Frege y Bertrand Russell, que pretendían definir los números naturales en términos de lógica

%\begin{tabla}
%\begin{table}[H]
%\centering
%\begin{adjustwidth}{-12pt}{}
%\begin{tabular}{@{}ll@{}}
%\toprule
%Máquinas de Turing & Numeración Gödel \\
%(\cref{teo:incompletitud-semantica}) & (\cref{teo:incompletitud-godel}) \\ \midrule
%para sistemas lógicos (semántica) & para sistemas formales (sintaxis) \\[4pt]
%asume que el sistema lógico es sólido & asume que el sistema formal es $\omega$-consistente \\[4pt]
%usa sistemas lógicos capaces de probar & usa sistemas formales capaces de probar \\
%afirmaciones del tipo ``$P$ para con &  ciertos teoremas de aritmética elemental \\
%entrada vacía'' \\ \bottomrule
%\end{tabular}
%\end{adjustwidth}
%\end{table}
%\vspace{-8pt}
%\caption{Diferencias en los teoremas de incompletitud mostrados en este capítulo}
%\label{tab:diferencias-turing-godel}
%\end{tabla}

%\section{El Segundo Teorema de Incompletitud}\label{sec:teo-incompletitud-2}

%El Segundo Teorema de Incompletitud, en su versión sintáctica, puede obtenerse inmediatamente a partir del Primer Teorema de Incompletitud.

%\begin{lema}\label{lema:red-prob}
%Sea la siguiente fórmula un teorema:
%$$\phi_1\Rightarrow\phi_2$$

%Entonces, si $\phi_2$ no es un teorema, $\phi_1$ tampoco lo es.
%\end{lema}
%\begin{proof}
%Este resultado se obtiene de la implicación inmediata siguiente: si $\phi_1$ es un teorema, entonces $\phi_2$ ha de serlo también, pues se deduce inmediatamente de $\phi_1\Rightarrow\phi_2$ mediante \emph{modus ponens}. Concluimos la prueba por el contrarrecíproco.
%\end{proof}

%\begin{teorema}[Versión semántica del Segundo Teorema de Incompletitud]
%Sea $\mathcal{G}$ un sistema lógico que verifique las siguientes condiciones:
%\begin{itemize}
%    \item[(A2)] $\mathcal{G}$ es recursivamente axiomatizable.
%    \item[(A3)] $\mathcal{G}$ puede expresar afirmaciones sobre máquinas de Turing o, equivalentemente, sobre programas; en concreto, dada una máquina de Turing $M$, puede expresar que ``$M$ para con entrada vacía'' o, equivalentemente, dado un programa $P$, puede expresar que ``$P$ para con entrada vacía''.
%\end{itemize}
%Entonces, la fórmula correspondiente a ``$\mathcal{G}$ es sólido'' no es un teorema en $\mathcal{G}$.
%\end{teorema}
%\begin{proof}
%El \cref{teo:incompletitud-semantica} nos dice que, suponiendo (A2) y (A3), en caso de que $\mathcal{G}$ fuese sólido, entonces tendríamos una fórmula $\phi$ que (1) es verdadera pero (2) no es un teorema.

%Como hemos logrado probar (1) en $\mathcal{G}$, tenemos el siguiente teorema:
%$$``\mathcal{G}\text{ es sólido}''\Rightarrow\phi$$
%Ahora bien, aplicando (2) y el \cref{lema:red-prob}, tenemos que ``$\mathcal{G}$ es sólido'' no es un teorema.
%\end{proof}

%\emph{\textbf{Nota: 6.1} para probar la versión sintáctica (que es el verdadero segundo teorema de incompletitud) sería necesario encontrar una versión constructiva del \cref{teo:incompletitud-sintactica}, que nos diese una fórmula $\phi$ que fuese verdadera de modo que ni $\phi$ ni $\neg\phi$ sean teoremas. Pero esto es un tanto extraño porque si el teorema es sintáctico, ¿qué quiere decir que una fórmula sea verdadera?}

%\section{Consecuencias}\label{sec:teo-incompletitud-consecuencias}


\endinput



\setpartpreamble[c][0.75\linewidth]{%
	\bigskip % Deja un espacio vertical en la parte superior
  Crearemos un fundamento matemático que nos permitirá usar programas en Python como parte de demostraciones formales, y comentaremos la relación que esto tiene con una de las tesis principales de la teoría de la computación. Hablaremos sobre problemas de decisión, universalidad y el problema de parada. Finalmente, probaremos la no decidibilidad del problema de parada.
}
\cleardoublepage\part{Fundamentos}
% !TeX root = ../libro.tex
% !TeX encoding = utf8

\setchapterpreamble[c][0.75\linewidth]{%
	\sffamily
  \emph{El objetivo de la computación es el conocimiento, no los números.}
 \begin{flushright} — Richard Hamming, \cite{Hamming1986} \end{flushright}
% \\[8pt]
	\par\bigskip
}
%\vspace{28pt}

\chapter{Motivación}\label{ch:motivacion}

Los teoremas de incompletitud de Gödel son uno de los resultados más populares de las matemáticas y la filosofía, y han sido objeto de fascinación y estudio durante décadas. Estos teoremas revolucionaron los mismos fundamentos de ambas disciplinas. En este trabajo, abordaremos el primero de estos teoremas desde una perspectiva computacional. Esto es así debido a una serie de experiencias y observaciones que despertaron en mí un gran interés y que revelaron una oportunidad para abordar este tema desde una perspectiva computacional pero, sobre todo, matemáticamente rigurosa.

Este interés surge de las clases de Modelos de Computación impartidas por el tutor de este trabajo, en las que mencionó de forma breve este resultado. Tras ello sentí mucha curiosidad por profundizar en el mundo de la incompletitud, hasta llegar a un excepcional libro de Douglas Hofstadter: \emph{``Gödel, Escher, Bach: An Eternal Golden Braid''} \cite{Hofstadter1999}. Esta obra maestra interdisciplinaria combina elementos de matemáticas, música y arte, y me incitó a profundizar más en estos teoremas.

La mayor parte de los libros que tratan este tema lo hacen tal y como lo hizo Gödel: desde una perspectiva aritmética, usando la numeración de Gödel en las demostraciones, o variantes. Sin embargo, en múltiples cursos de teoría de la computación se exponen demostraciones haciendo uso de algoritmos y máquinas de Turing.

Para mi sorpresa, a pesar de que esta perspectiva se estudia extensivamente y hay una gran literatura al respecto, muchas de las referencias que encontraba eran bien poco rigurosas o bien poco consistentes en la nomenclatura y en la forma de probar los teoremas. En concreto, muchas de las fuentes no distinguen de forma precisa entre los conceptos de \emph{demostración} y \emph{verdad}.

Es aquí donde vi la oportunidad de hacer un trabajo que intentase proporcionar una perspectiva computacional a los teoremas de incompletitud de Gödel, sin perder rigor y concreción por el camino. Este enfoque permitirá una comprensión más precisa y estructurada de los resultados más relevantes al respecto, así como su relación con los modelos de computación y la teoría de la computabilidad.

\pagebreak

\section*{Objetivos}

Los objetivos de este trabajo son:

\subsection*{Objetivos teóricos}

\begin{enumerate}[label={[}T\arabic*{]}]
    \item Probar el Primer Teorema de Incompletitud de Gödel desde una perspectiva computacional.
    \item Usar para ello programas en Python en lugar de descripciones de máquinas de Turing, probando que esto es posible y riguroso.
    \item Exponer las consecuencias del teorema de incompletitud en las disciplinas matemática y filosófica.
    \item Contextualizar el teorema de incompletitud en el ámbito histórico.
    \item Encontrar la relación entre el teorema de incompletitud con la tesis de Church-Turing.
    \item Precisar conceptos lógicos relacionados con la teoría de la demostración. En concreto, distinguir entre las nociones de \emph{demostración} y \emph{verdad}, y estudiar las propiedades de los sistemas de demostración.
\end{enumerate}

\subsection*{Objetivos prácticos}

\begin{enumerate}[label={[}P\arabic*{]}]
    \item Crear contenido que pueda servir de herramienta de apoyo a la docencia.
    \item Crear una librería de programas en Python ejecutables por el usuario que permitan guiar las demostraciones que aparecen en este trabajo.
\end{enumerate}

\endinput

% !TeX root = ../libro.tex
% !TeX encoding = utf8

\setchapterpreamble[c][0.75\linewidth]{%
	\sffamily
  \emph{Si una ``religión'' se define como un sistema de ideas que contiene afirmaciones indemostrables, entonces Gödel nos enseñó que las matemáticas no sólo son una religión, sino que son la única religión que puede demostrar que lo es.}
 \begin{flushright} — John D. Barrow, \cite{Barrow2011} \end{flushright}
% \\[8pt]
	\par\bigskip
}
%\vspace{28pt}

\chapter{El Primer Teorema de Incompletitud}\label{ch:teorema-incompletitud}

Los matemáticos de principios del siglo \textsc{xx}, tal y como Leibniz, Boole y Hilbert, imaginaron que sería posible determinar la veracidad de fórmulas matemáticas de forma automática. Nosotros ahora nos preguntamos lo mismo: hemos trabajado con sistemas lógicos sencillos que son decidibles, es decir, cuyas fórmulas pueden ser comprobadas verdaderas mediante programas. En el caso de sistemas lógicos de más complejidad, como la aritmética de Peano, esto no es posible. Esto se debe a que la existencia de problemas computacionales indecidibles conduce directamente a la existencia de afirmaciones matemáticas no decidibles.

%En la \cref{sec:primera-aproximacion}, probaremos la incompletitud de \textbf{Peano}. Llamaremos a este resultado \emph{Primer Teorema de Incompletitud}. A continuación, generalizaremos este teorema para otros sistemas lógicos. Finalmente, en la \cref{sec:diferencias-godel}, veremos el enunciado original del teorema y lo compararemos con el resultado obtenido.

\section{Una primera aproximación}\label{sec:primera-aproximacion}

%\begin{teorema}[Versión no constructiva del Primer Teorema de Incompletitud]\label{teo:inco1-no-constr}
%La aritmética de Peano es incompleta: existen fórmulas verdaderas de \textbf{Peano} que no son teoremas.
%\end{teorema}
%\begin{proof}
%Asumiremos que \textbf{Peano} es completo y entraremos en contradicción. En concreto, demostraremos que si la aritmética de Peano es completa, podremos decidir \textbf{ParadaEnVacío}, lo cual contradice la \cref{prop:parada-en-vacio-no-decidible}.

%Sea $P$ un programa cualquiera. Lo usaremos como entrada al \cref{lst:parada-en-vacio-desde-peano-completo}. En la \cref{line:parada-en-vacio-desde-peano-completo-import1} importamos la función \texttt{es\_prueba\_peano}, que dada una fórmula y una demostración en \textbf{Peano} comprueba si la demostración prueba que la fórmula es un teorema. Claramente, esta función no es un oráculo. De hecho, su implementación es sencilla: basta probar si cada una de las fórmulas de la demostración es un axioma o se sigue de las fórmulas anteriores a través de las reglas de inferencia del sistema, comprobando que la última fórmula coincida con la fórmula que queremos probar.

%A continuación, importamos en la \cref{line:parada-en-vacio-desde-peano-completo-import2} la función \texttt{parada\_a\_peano} cuya existencia comentamos tras la \cref{prop:parada-a-peano}, que pasa la afirmación ``$P$ para con entrada vacía'' a una fórmula en \textbf{Peano}.

%Usamos esta función en la \cref{line:parada-en-vacio-desde-peano-completo-para-en-peano} para generar la fórmula mencionada anteriormente. En la \cref{line:parada-en-vacio-desde-peano-completo-no-para-en-peano} creamos una nueva fórmula negando la anterior, por tanto \texttt{no\_para\_en\_peano} contiene la fórmula de \textbf{Peano} correspondiente a la afirmación ``$P$ no para con entrada vacía''.

%Una de estas afirmaciones ha de ser verdadera. Como asumimos que \textbf{Peano} es completo, una de estas fórmulas tiene una prueba (es un teorema) en \textbf{Peano}. El bucle de la \cref{line:parada-en-vacio-desde-peano-completo-while} va iterando a través de todos los \emph{strings} posibles, comprobando cada uno para saber si es una prueba de \texttt{para\_en\_peano} o de \texttt{no\_para\_en\_peano}. Ya que uno de ellos debe tener una prueba, el programa terminará, decidiendo si $P$ para con entrada vacía o no. El \cref{lst:parada-en-vacio-desde-peano-completo} resuelve \textsc{ParadaEnVacío}, lo que entra en contradicción con la \cref{prop:parada-en-vacio-no-decidible}.
%\end{proof}

%\begin{lstlisting}[language=Python, caption=\lstinline{parada_en_vacio_desde_peano_completo.py},label={lst:parada-en-vacio-desde-peano-completo}]
%import utilidades

%from es_prueba_peano import es_prueba_peano # NO es un oráculo|\label{line:parada-en-vacio-desde-peano-completo-import1}|
%from parada_a_peano import parada_a_peano # NO es un oráculo|\label{line:parada-en-vacio-desde-peano-completo-import2}|

%def parada_en_vacio_desde_peano_completo(programa):|\label{line:parada-en-vacio-desde-peano-completo-main}|
%    para_en_peano = parada_a_peano(programa)|\label{line:parada-en-vacio-desde-peano-completo-para-en-peano}|
%    no_para_en_peano = 'NOT (' + para_en_peano + ')'|\label{line:parada-en-vacio-desde-peano-completo-no-para-en-peano}|
%    prueba = ''|\label{line:parada-en-vacio-desde-peano-completo-prueba}|
    
%    while True:|\label{line:parada-en-vacio-desde-peano-completo-while}|
%        if es_teorema_peano(prueba, para_en_peano) == 'sí':|\label{line:parada-en-vacio-desde-peano-completo-if1}|
%            return 'sí'
%        if es_teorema_peano(prueba, no_para_en_peano) == 'sí':|\label{line:parada-en-vacio-desde-peano-completo-if2}|
%            return 'no'
%        prueba = utilidades.siguiente_string(prueba)|\label{line:parada-en-vacio-desde-peano-completo-else}|
%\end{lstlisting}


Antes de proceder la demostración de una primera versión del Primer Teorema de Incompletitud, discutiremos de forma detallada el \cref{lst:godel}. En la \cref{line:godel-peano-import1} importamos la función \texttt{es\_teorema\_peano}, cuya existencia probamos en la \cref{prop:es-teorema-peano-semidecidible}. Recuerda que, como los axiomas de \textbf{Peano} son recursivamente enumerables (es decir, \textbf{Peano} es recursivamente axiomatizable), podemos escribir el programa \texttt{es\_teorema\_peano(fórmula)} que devuelve \palabra{sí} si \texttt{fórmula} es un teorema en \textbf{Peano}, y que no para en caso contrario (ver \cref{prop:es-teorema-peano-semidecidible}).
\vspace{8pt}
\begin{lstlisting}[language=Python, caption=\lstinline{godel_peano.py},label={lst:godel-peano}]
import utilidades

from es_teorema_peano import es_teorema_peano # NO es un oráculo|\label{line:godel-peano-import1}|
from parada_a_peano import parada_a_peano # NO es un oráculo|\label{line:godel-peano-import2}|

def godel_peano(entrada):|\label{line:godel-peano-main}|
    programa_godel = utilidades.leer('godel\_peano.py')|\label{line:godel-peano-programa}|
    para_en_peano = parada_a_peano(programa_godel)|\label{line:godel-peano-formula}|
    no_para_en_peano = 'NOT (' + para_en_peano + ')'|\label{line:godel-peano-formulaneg}|

    if es_teorema_peano(no_para_en_peano) == 'sí':|\label{line:godel-peano-es-teorema}|
        return 'para'|\label{line:godel-peano-para}|
    else:
        utilidades.ciclar()|\label{line:godel-peano-cicla}|
\end{lstlisting}

La \cref{line:godel-peano-import2} importa la función \texttt{parada\_a\_peano}, cuya existencia ya discutimos tras la \cref{prop:parada-a-peano}. Esta función traduce, dado un programa $P$, la afirmación ``$P$ para con entrada vacía'' en una fórmula de \textbf{Peano}.

En la \cref{line:godel-peano-programa} guardamos el propio programa \texttt{godel\_peano.py} en la variable \texttt{programa\_godel}. En la \cref{line:godel-peano-formula} usamos la función \texttt{parada\_a\_peano} recién comentada, y en la \cref{line:godel-peano-formulaneg} negamos la fórmula \texttt{para\_en\_peano}. De este modo:
\begin{adjustwidth}{30pt}{}
    \texttt{no\_para\_en\_peano} significa que ``\texttt{godel\_peano.py} no para con entrada vacía'' \hfill $(\star1)$
\end{adjustwidth}
Ahora es cuando las cosas se ponen interesantes. Si \texttt{no\_para\_en\_peano} es un teorema, la función \texttt{es\_teorema\_peano} devolverá \palabra{sí}, y de este modo el programa verificará la condición de la \cref{line:godel-peano-es-teorema} y devolverá \palabra{para}, haciendo que el programa pare. Si \texttt{no\_para\_en\_peano} no es un teorema, se entrará en un bucle infinito (\texttt{es\_teorema\_peano} no parará, y por lo tanto tampoco \texttt{godel\_peano.py}). Es claro que la \cref{line:godel-peano-cicla} nunca se ejecutará, pero la incluimos para enfatizar este hecho.

Ahora, nos preguntamos: ¿cuál es el comportamiento de \texttt{godel\_peano.py} con entrada vacía?

\begin{proposicion}\label{prop:godel-peano-no-para}
El \cref{lst:godel-peano} no para con entrada vacía.
\end{proposicion}
\begin{proof}
Asumimos que \texttt{godel\_peano.py} para con entrada vacía, y entraremos en contradicción. La única forma de que pare es en la \cref{line:godel-peano-para}, algo que ocurre solamente si:
\begin{adjustwidth}{30pt}{}
    \texttt{es\_teorema\_peano(no\_para\_en\_peano) == \palabra{sí}} \hfill $(\star2)$
\end{adjustwidth}
$(\star2)$ es equivalente a decir que \texttt{no\_para\_en\_peano} es un teorema. Por la solidez de \textbf{Peano}, \texttt{no\_para\_en\_peano} es verdadero.

Pero el hecho de que \texttt{no\_para\_en\_peano} sea verdadero, junto con $(\star1)$, nos dice que \texttt{godel\_peano.py} no para con entrada vacía, encontrando la contradicción esperada.
\end{proof}

Ahora estamos en condiciones de probar el resultado principal de este trabajo.

\begin{teorema}[Incompletitud de la aritmética de Peano]\label{teo:incompletitud-peano}
\index{primer teorema de incompletitud}\index{teorema de incompletitud}La aritmética de Peano es incompleta: existen fórmulas verdaderas de \normalfont{\textbf{Peano}} que no son teoremas. Concretamente, la fórmula guardada en la variable \texttt{no\_para\_en\_peano} en la \cref{line:godel-peano-formulaneg} del \cref{lst:godel} es un ejemplo de una fórmula verdadera pero no demostrable (no es un teorema) de \normalfont{\textbf{Peano}}.
\end{teorema}
\begin{proof}
Debemos mostrar que \texttt{no\_para\_en\_peano} es verdadero y no es un teorema.

Para probar que es verdadero, por la \cref{prop:godel-peano-no-para}, sabemos que \texttt{godel\_peano.py} no para con entrada vacía. Recordando $(\star1)$, inmediatamente obtenemos que \texttt{no\_para\_en\_peano} es verdadero.

Para probar que no es un teorema, de nuevo usamos el hecho de que \texttt{godel\_peano.py} no para con entrada vacía. Esto nos indica que, en la \cref{line:godel-peano-es-teorema}, la llamada a la función
\begin{adjustwidth}{30pt}{}
    \texttt{es\_teorema\_peano(no\_para\_en\_peano)}
\end{adjustwidth}
no puede devolver \palabra{sí}. Pero esto inmediatamente implica que \texttt{no\_para\_en\_peano} no es demostrable, lo que completa la prueba.
\end{proof}

%\section{El Segundo Teorema de Incompletitud}\label{sec:teo-incompletitud-2}

%\begin{teorema}[Segundo Teorema de Incompletitud]
%La consistencia de la aritmética de Peano no puede probarse en la aritmética de Peano. En concreto, la afirmación
%\begin{adjustwidth}{30pt}{}
%    ``\textbf{Peano} es consistente''
%\end{adjustwidth}
%que puede codificarse mediante una fórmula de \textbf{Peano}, no es un teorema.
%\end{teorema}

Acabamos de demostrar que \textbf{Peano} es incompleto, como ya anticipamos en la \cref{tab:sistemas-logicos}. Las consecuencias de este hecho las exploramos en detalle en el \cref{ch:conclusion}. Sin embargo, paremos un momento a analizar qué es lo que acabamos de probar.

\section{Versión semántica}\label{sec:version-semantica}

Para encontrar una fórmula verdadera y no demostrable, hemos necesitado un sistema lógico que sea, en primer lugar, recursivamente axiomatizable. Si no, no hubiesemos podido crear un programa capaz de saber cuándo una fórmula del sistema es un teorema. Por otra parte, hemos asumido que el sistema es sólido, lo cual es algo esperable. Finalmente, hemos exigido que sea capaz de expresar una cierta afirmación de las máquinas de Turing (véase \cref{prop:parada-a-peano}). Para evitar repeticiones innecesarias en esta última hipótesis, llamaremos a un sistema lógico con esta propiedad ``de tipo Turing'', algo que definimos a continuación.

\begin{definicion}[Sistema de tipo Turing]
Sea $\mathcal{S}$ un sistema formal. Diremos que $\mathcal{S}$ es \emph{de tipo Turing}\footnote{Esta definición no se encuentra en la literatura: se ha creado por conveniencia.} si puede expresar afirmaciones sobre máquinas de Turing o, equivalentemente, sobre programas;\footnote{La equivalencia la probamos en el \cref{teo:equivalencia}.} en concreto, dada una máquina de Turing $M$, puede expresar que ``$M$ para con entrada vacía''; o, equivalentemente, dado un programa $P$, puede expresar que ``$P$ para con entrada vacía''.

Análogamente podemos definirlo para el caso de sistemas lógicos.
\end{definicion}

Observemos que, sin embargo, en ningún momento hemos utilizado el sistema \textbf{Peano} en sí. Esto nos indica que podríamos realizar esta demostración con otro sistema que no sea \textbf{Peano}, siempre y cuando verifique tales condiciones. Indicamos este resultado a continuación.

\begin{teorema}[Versión semántica del Primer Teorema de Incompletitud]\label{teo:incompletitud-semantica}
Sea $\mathcal{G}$ un sistema lógico que verifique las siguientes condiciones:
\begin{enumerate}
    \item $\mathcal{G}$ es sólido.
    \item $\mathcal{G}$ es recursivamente axiomatizable.
    \item $\mathcal{G}$ es de tipo Turing.
\end{enumerate}
Entonces, $\mathcal{G}$ es incompleto, es decir, existen fórmulas verdaderas de $\mathcal{G}$ que no son teoremas.
\end{teorema}

Para reforzar este importante resultado, incluimos un resumen del proceso que hemos seguido para \textbf{Peano}, de forma general, en la demostración de este teorema.

\begin{proof}
En el \cref{lst:godel} importamos dos funciones. Por una parte, en la \cref{line:godel-import1} importamos la función \texttt{es\_teorema\_en\_G}, que comprueba si una fórmula es un teorema en $\mathcal{G}$. Esta función parará si la fórmula es un teorema, y se ejecutará indefinidamente en caso de que no lo sea. La existencia de esta función la garantiza (2).

Por otra, en la \cref{line:godel-import2} importamos la función \texttt{parada\_a\_G} que, dado un programa $P$, convierte la afirmación ``$P$ para con entrada vacía'' en una fórmula de $\mathcal{G}$, algo que es posible por (3).

En la línea \cref{line:godel-formulaneg} tenemos que:
\begin{adjustwidth}{30pt}{}
    \texttt{no\_para\_en\_G} significa que \emph{``\texttt{godel.py} no para con entrada vacía''} (en $\mathcal{G}$)  \hfill $(\star3)$
\end{adjustwidth}
Si \texttt{no\_para\_en\_G} es un teorema, la función \texttt{es\_teorema\_en\_G(no\_para\_en\_G)} (\cref{line:godel-es-teorema}) parará y devolverá \palabra{sí}, mientras que la misma función ciclará en caso de que no lo sea.

La pregunta que debemos hacernos es: ¿qué le pasa a \texttt{godel.py} cuando \texttt{entrada} es vacía?

Si asumimos que \texttt{godel.py} para con entrada vacía, la única forma de que lo haga es porque en la \cref{line:godel-es-teorema} es \texttt{es\_teorema\_en\_G(no\_para\_en\_G) == \palabra{sí}}, lo cual solo puede pasar si \texttt{no\_para\_en\_G} es un teorema. Pero, por (1), y dado que $\mathcal{G}$ es sólido, tenemos que \texttt{no\_para\_en\_G} es verdadero.

Pero si \texttt{no\_para\_en\_G} es verdadero, por $(\star3)$, esto querría decir que \texttt{godel.py} no para con entrada vacía, entrando en contradicción, por lo que debe ser:
\begin{adjustwidth}{30pt}{}
    \texttt{godel.py} no para con entrada vacía  \hfill $(\star4)$
\end{adjustwidth}
Comprobamos que \texttt{no\_para\_en\_G} es verdadero pero no es un teorema. En efecto, es verdadero por $(\star4)$. Sin embargo, no es un teorema. Para ello, usamos de nuevo $(\star4)$. El hecho de que \texttt{godel.py} no pare indica que es imposible que la condición de la \cref{line:godel-es-teorema}, \linebreak\texttt{es\_teorema\_en\_G(no\_para\_en\_G) == \palabra{sí}}, se verifique -- en otras palabras, \texttt{no\_para\_en\_G} no es un teorema.
\end{proof}

%\vspace{8pt}
\begin{lstlisting}[language=Python, caption=\lstinline{godel.py},label={lst:godel}]
import utilidades

from es_teorema_G import es_teorema_G # NO es un oráculo|\label{line:godel-import1}|
from parada_a_G import parada_a_G # NO es un oráculo|\label{line:godel-import2}|

def incompleto_semantico(entrada):|\label{line:godel-main}|
    programa = utilidades.leer('incompleto_semantico.py')|\label{line:godel-programa}|
    para_en_G = parada_a_G(programa)|\label{line:godel-formula}|
    no_para_en_G = 'NOT (' + para_en_G + ')'|\label{line:godel-formulaneg}|

    if es_teorema_G(no_para_en_G) == 'sí':|\label{line:godel-es-teorema}|
        return 'para'|\label{line:godel-para}|
    else:
        utilidades.ciclar()|\label{line:godel-cicla}|
\end{lstlisting}


Nuestro resultado nos dice que, si tenemos un sistema lógico que sea recursivamente axiomatizable y de tipo Turing, no podemos esperar que pueda ser sólido y completo al mismo tiempo. Es decir, si todos los teoremas del sistema son verdaderos, siempre podremos encontrar fórmulas del sistema que sean verdaderas pero no sean teoremas.

Podemos aplicar el \cref{teo:incompletitud-semantica} a otros sistemas lógicos, siempre que verifiquen las hipótesis. Uno de ellos es, por ejemplo, el sistema de axiomas de Zermelo-Fraenkl.

\section{Versión sintáctica}\label{sec:version-sintactica}

% WIP relajar condición de solidez?
El resultado que acabamos de probar es semántico. Sin embargo, es posible probar otro resultado, esta vez sintáctico, y más parecido al demostrado por Gödel.

Para poder probar esta variación, debemos de tener en cuenta una sutileza del problema \textsc{EsTeorema} en el caso de sistemas formales sintácticamente consistentes.

\begin{proposicion}\label{prop:es-teorema-consistente-decidible}
Sea $\mathcal{S}$ un sistema formal recursivamente axiomatizable y sintácticamente consistente. Entonces, $\textsc{EsTeorema}_\mathcal{S}$ es decidible.
\end{proposicion}
\begin{proof}
Veamos que el \cref{lst:es-teorema-s} decide $\textsc{EsTeorema}_\mathcal{S}$. En efecto, vemos en primer lugar que en la \cref{line:es-teorema-s-import} importamos la función \texttt{es\_prueba\_S}, la cual es computable por ser $\mathcal{S}$ recursivamente axiomatizable.

A continuación, en la \cref{line:es-teorema-s-negacion} almacenamos la negación de \texttt{formula} en \texttt{neg\_formula}. En la \cref{line:es-teorema-s-demostracion} declaramos la demostración.

Por ser $\mathcal{S}$ consistente, bien \texttt{formula} bien \texttt{neg\_formula} tiene una demostración. En el bucle de la \cref{line:es-teorema-s-while}, vamos comprobando si \texttt{demostracion} es una demostración de \texttt{formula} o de \texttt{neg\_formula}. En el primer caso, devolvemos \palabra{sí}, mientras que en el segundo, devolvemos \palabra{no}. En caso de que \texttt{demostracion} no sea una demostración de ninguna de ambas fórmulas, procedemos en la \cref{line:es-teorema-s-siguiente-string} a enumerar la siguiente demostración en orden \emph{shortlex} (ver \cref{prop:es-teorema-peano-semidecidible}).

Por la consistencia de $\mathcal{S}$, debe existir una demostración bien de \texttt{formula} o de \texttt{neg\_formula}, por lo que el programa debe parar al llegar a ser \texttt{demostracion} tal demostración.
\end{proof}

\begin{lstlisting}[language=Python, caption=\lstinline{es_teorema_S.py},label={lst:es-teorema-s}]
import utilidades
from es_prueba_S import es_prueba_S # NO es un oráculo|\label{line:es-teorema-s-import}|

def es_teorema_S(formula):|\label{line:es-teorema-s-main}|
    negacion_formula = 'NOT (' + formula + ')'|\label{line:es-teorema-s-negacion}|
    demostracion = ''|\label{line:es-teorema-s-demostracion}|

    while True:|\label{line:es-teorema-s-while}|
        if es_prueba_S(demostracion, formula) == 'sí':|\label{line:es-teorema-s-es-prueba}|
            return 'sí'
        if es_prueba_S(demostracion, negacion_formula) == 'sí':
            return 'no'
        demostracion = utilidades.siguiente_string(demostracion)|\label{line:es-teorema-s-siguiente-string}|
\end{lstlisting}

Ahora estamos en condiciones de probar el resultado sintáctico, que dependerá de la no decidibilidad de \textsc{AdivinaConsistente} (\cref{prob:adivina-consistente}). \cite{Aaronson2017}

\begin{teorema}[Versión sintáctica del Primer Teorema de Incompletitud]\label{teo:incompletitud-sintactica}
Sea $\mathcal{S}$ un sistema formal que verifique las siguientes condiciones:
\begin{enumerate}
    \item $\mathcal{S}$ es sintácticamente consistente.
    \item $\mathcal{S}$ es recursivamente axiomatizable.
    \item $\mathcal{S}$ es de tipo Turing.
\end{enumerate}
Entonces, $\mathcal{S}$ es sintácticamente incompleto, es decir, existen fórmulas $\phi$ de $\mathcal{S}$ tal que ni $\phi$ ni $\neg\phi$ son teoremas.
\end{teorema}
\begin{proof}
Vamos a suponer que $\mathcal{S}$ es sintácticamente completo y veremos que, en tal caso, podríamos decidir \textsc{AdivinaConsistente} (\cref{prob:adivina-consistente}), lo cual entraría en contradicción con la \cref{prop:adivina-consistente-no-decidible}.

Comencemos comentando las funciones importadas. En la \cref{line:adivina-consistente-import-1} importamos la función \texttt{oarada\_en\_S}, que traduce para un programa $P$ la afirmación ``$P$ para con entrada vacía'' en una fórmula de $\mathcal{}S$, y cuya existencia garantiza la hipótesis (3). En la \cref{line:adivina-consistente-import-2} importamos la función \texttt{es\_teorema\_S}, que decide el problema $\textsc{EsTeorema}_\mathcal{S}$, y que, por las hipótesis (1) y (2), siempre para (ver \cref{prop:es-teorema-consistente-decidible}) y devuelve \palabra{sí} si la fórmula de entrada es un teorema en $\mathcal{S}$ y \palabra{no} en caso contrario.

En la \cref{line:adivina-consistente-programa-para} usamos la función \texttt{parada\_a\_S} para almacenar en \texttt{programa\_para} la fórmula de $\mathcal{S}$ equivalente a ``\texttt{programa} para con entrada vacía''. En la \cref{line:adivina-consistente-es-teorema}, asignamos a la variable \texttt{es\_teorema} uno de dos valores posibles: \palabra{sí} o \palabra{no}, como ya comentamos anteriormente.

Llegado a la \cref{line:adivina-consistente-si}, implementamos la siguiente lógica:
\begin{itemize}
    \item Si \texttt{programa} acepta con entrada vacía (para y devuelve \palabra{sí}), la condición de la \cref{line:adivina-consistente-si} se verificará, así como la de la \cref{line:adivina-consistente-si-si}, y \texttt{adivina\_consistente.py} devolverá \palabra{sí}.
    \item Si \texttt{programa} rechaza con entrada vacía (para y devuelve \palabra{no}, la condición de la \cref{line:adivina-consistente-si} se verificará, mientras que no la de la \cref{line:adivina-consistente-si-si}, y \texttt{adivina\_consistente.py} devolverá \palabra{no}.
    \item Si \texttt{programa} cicla con entrada vacía, la condición de la \cref{line:adivina-consistente-si} no se verificará, y \texttt{adivina\_consistente.py} devolverá \palabra{no}.
\end{itemize}
Dicha lógica hace evidente que el programa decide \textsc{AdivinaConsistente}.
\end{proof}

\begin{lstlisting}[language=Python, caption=\lstinline{adivina_consistente.py},label={lst:adivina-consistente}]
import utilidades
from parada_a_S import parada_a_S # NO es un oráculo|\label{line:adivina-consistente-import-1}|
from es_teorema_S import es_teorema_S # NO es un oráculo|\label{line:adivina-consistente-import-2}|

def adivina_consistente(programa):|\label{line:adivina-consistente-main}|
    programa_para = parada_a_S(programa)|\label{line:adivina-consistente-programa-para}|

    es_teorema = es_teorema_S(programa_para)|\label{line:adivina-consistente-es-teorema}|

    if es_teorema == 'sí':|\label{line:adivina-consistente-si}|
        salida = maquina_universal(programa, '')|\label{line:adivina-consistente-universal}|
        if salida == 'sí':|\label{line:adivina-consistente-si-si}|
            return 'sí'
        else:|\label{line:adivina-consistente-si-no}|
            return 'no'
    else:|\label{line:adivina-consistente-no}|
        return 'no'
\end{lstlisting}

Hemos probado dos versiones del Primer Teorema de Incompletitud. Resumimos las diferencias entre ambas a continuación.
\vspace{8pt}
\begin{figure}[H]
$$
\begin{array}{ccc}
    {
        \begin{matrix}
            \text{\textbf{versión semántica}} \\
            \text{\small{(\cref{teo:incompletitud-semantica})}}\\
            \small{\;}
        \end{matrix}
    }
    & \;\;\; &
    {
        \begin{matrix}
            \text{\textbf{versión sintáctica}} \\
            \text{\small{(\cref{teo:incompletitud-sintactica})}}\\
            \small{\;}
        \end{matrix}
    }
    \\
    {
        \begin{matrix}
            \mathcal{G}\text{ sistema lógico}\\
            \text{rec. axiom., de tipo Turing}\\
            \text{y sólido} \\
            \small{\phi \text{ teorema} \Rightarrow \phi\text{ verdadero}}\\
            \tiny{\;}
        \end{matrix}
    }
    &  & 
    {
        \begin{matrix}
            \mathcal{S}\text{ sistema formal}\\
            \text{rec. axiom., de tipo Turing}\\
            \text{y sintácticamente consistente} \\
            \small{\nexists \phi : \phi\text{ y }\neg\phi \text{ son teoremas}}\\
            \tiny{\;}
        \end{matrix}
    } \\
    \Huge{\Downarrow} & & \Huge{\Downarrow} \\
    {
        \begin{matrix}
            \tiny{\;}\\
            \mathcal{G}\text{ no (semánticamente) completo}\\
            \small{\exists\phi:\phi\text{ es verdadero y no es teorema}}
        \end{matrix}
    }
    & &
    {
        \begin{matrix}
            \tiny{\;}\\
            \mathcal{S}\text{ no sintácticamente completo}\\
            \small{\exists\phi:\text{ni }\phi\text{ ni }\neg\phi\text{ son teoremas}}
        \end{matrix}
    }
    \\
\end{array}
$$
\caption{Comparativa entre las dos versiones del Primer Teorema de Incompletitud}
\vspace{-8pt}
\label{fig:comparativa-incompletitud}
\end{figure}

La demostración original de Gödel es puramente sintáctica, y muy parecida a la probada en el \cref{teo:incompletitud-sintactica}. La enunciamos a continuación.

\begin{teorema}[Primer Teorema de Incompletitud de Gödel-Rosser]\label{teo:incompletitud-godel}
Sea $\mathcal{S}$ un sistema formal que verifique las siguientes condiciones:
\begin{enumerate}
    \item $\mathcal{S}$ es sintácticamente consistente.
    \item $\mathcal{S}$ es recursivamente axiomatizable.
    \item $\mathcal{S}$ contiene una cierta cantidad de aritmética elemental.
\end{enumerate}

Entonces, $\mathcal{S}$ es sintácticamente incompleto, es decir, hay fórmulas $\phi$ de $\mathcal{S}$ de modo que ni $\phi$ ni $\neg\phi$ son teoremas.
\end{teorema}

Tal y como en nuestra versión sintáctica, requerimos que el sistema sea sintácticamente consistente y recursivamente axiomatizable. El teorema original de Gödel exigía como hipótesis una forma más fuerte de consistencia sintáctica llamada \emph{$\omega$-consistencia}. Esta forma más general del teorema se debe al trabajo de John B. Rosser, autor del hoy conocido como \emph{truco de Rosser}.\footnote{Es por ello por lo que lo llamamos teorema de ``Gödel-Rosser'', para distinguirlo del resultado original de Gödel.} \cite{Rosser1936}

La tercera hipótesis exige que el sistema sea capaz de probar ciertas afirmaciones de aritmética elemental. Esto se debe a que Gödel pretendía probar que la teoría de números no era axiomatizable (ver \cref{ch:historia}), por lo que usó en su demostración una forma de numerar los teoremas que hoy en día se conoce como \emph{numeración de Gödel}.


\section{Consecuencias}\label{sec:consecuencias}

El Primer Teorema de Incompletitud tiene profundas implicaciones matemáticas y filosóficas.

\subsection{Relación con el segundo problema de Hilbert}\label{subsec:segundo-problema-hilbert}

Ya introdujimos en el \cref{ch:historia} los célebres problemas de Hilbert, y en concreto, el segundo de ellos, que volveremos a referenciar a continuación:

\begin{enumerate}[label=(\arabic*)]
    \item ¿Son completas las matemáticas? Esto es, ¿puede probarse o no cada sentencia matemática?
    \item ¿Son consistentes las matemáticas? Esto es, ¿no es posible probar simultáneamente una afirmación y su negación?
    \item ¿Son decidibles las matemáticas? Esto es, ¿existe un método automático que pueda aplicarse a cualquier afirmación matemática, y que determine si es cierta? Este problema se conoce como \emph{Entscheidungsproblem}.\index{Entscheidungsproblem}
\end{enumerate}

Es necesario puntualizar que la formulación de Hilbert de este problema es relativa a la aritmética. Es claro que, para tal caso, el \cref{teo:incompletitud-godel} nos garantiza la existencia de fórmulas no demostrables, respondiendo negativamente a (1) \cite{Hellman1981}. Por otra parte, en la \cref{prop:es-teorema-consistente-decidible} probamos que \textbf{Peano}, un sistema de aritmética, es no decidible, probando que la respuesta a la pregunta (3) es negativa. La pregunta (2) queda fuera del objeto de este trabajo.

Sin embargo, es interesante ver estas preguntas sin necesidad de limitarnos a la aritmética. Hemos introducido los sistemas lógicos justamente para poder formalizar el concepto de verdad y demostración, que son el alma de las matemáticas.

Podremos encontrar respuestas análogas para sistemas formales recursivamente enumerables y de tipo Turing. En tal caso, el \cref{teo:incompletitud-sintactica} responde negativamente a la pregunta (1), mientras que una generalización de la \cref{prop:es-teorema-peano-no-decidible} (similar a la que hicimos en el \cref{teo:incompletitud-semantica} a partir del \cref{teo:incompletitud-peano}) nos permite responder negativamente a la pregunta (3).

%Tenemos respuestas a las preguntas (1) y (3) para unas hipótesis concretas pero, ¿qué pasa en el caso de sistemas no recursivamente enumerables y que no son de tipo Turing? Algunos matemáticos no están convencidos de que los resultados de Gödel resuelvan definitivamente el segundo problema de Hilbert, pues los teoremas de Gödel no se aplican en otros sistemas de lógica.

\subsection{Relación con la mente humana y con la tesis de Church-Turing}\label{subsec:mente-humana-tesis-church-turing}

Algunos autores entran en conclusiones mucho más epistemológicas, y que conciernen a la mente humana. Roger Penrose y John Lucas \cite{Penrose1989, Lucas1961} argumentan que el papel de la conciencia es no algorítmico cuando formamos juicios matemáticos, pudiendo generalizarlo al papel de la conciencia humana.

Para ello, los autores comentan que en cualquier sistema formal existen proposiciones indemostrables, esto es, para las que su algoritmo no puede proporcionar una respuesta. Si el funcionamiento de la mente humana fuese completamente algorítmico, entonces el algoritmo (o sistema formal) que utilice para formar sus juicios no podría tratar tal proposición. Sin embargo, la mente humana puede deducir que tal proposición es verdadera. Esta contradicción nos puede indicar que la mente humana, en cuanto al razonamiento matemático, no estaba utilizando un algoritmo en primer lugar.

Douglas Hofstadter \cite{Hofstadter1999, Hofstadter2007} presenta los teoremas de Gödel como un ejemplo de ``extraño bucle'', una estructura jerárquica y autoreferencial que existe en un sistema formal. Argumenta que esta es la misma estructura que origina la conciencia, el sentido del ``yo'', en la mente humana: la autoreferencia proviene de la forma en que el cerebro abstrae y categoriza los estímulos en ``símbolos'', o grupos de neuronas, que corresponden a conceptos, en lo que es de facto un sistema formal, que eventualmente forma símbolos que corresponden a la propia entidad que realiza la percepción. 

La diferencia entre las matemáticas y la conciencia reside, según Hofstdadter, en el sentido en el que se realiza el razonamiento: si es ``hacia arriba'' o ``hacia abajo''., desde los axiomas hacia las proposiciones y teoremas más complejas, o ``hacia abajo'', en el sentido contrario.

Las matemáticas funcionan hacia arriba: para derivar si una fórmula es un teorema debemos partir de los axiomas e ir razonando hasta llegar a ella. La mente humana, según el físico, funciona a hacia abajo, desde un nivel más alto --de deseos, conceptos, personalidades, ideas y pensamientos-- hasta el nivel más bajo de interacciones entre neuronas e incluso partículas fundamentales. Los seres humanos estamos hechos para percibir ``cosas grandes'' en lugar de ``cosas pequeñas'', incluso a pesar de que en el dominio de lo pequeño parece residir el motor de la realidad.

Ambos argumentos están íntimamente relacionados con la tesis de Church-Turing (véase la \cref{sec:church-turing}). Si la mente humana es una máquina de Turing, todos los procedimientos de nuestro cerebro no serían más que algoritmos, y mediante los teoremas de Gödel, nuestra forma de razonar sería limitada --¡existen afirmaciones que no podemos demostrar!--. Esto puede inducir en conclusiones más fantásticas, hablando algunos autores de que los teoremas de Gödel, aunque no permitan probarlo, son una fuerte evidencia de la existencia de dios.\footnote{No existe consenso respecto si esto es cierto, pero es sencillo argumentar los fallos de este razonamiento (véase \cite{StanfordGodel2022}).} \cite{StanfordGodel2022}

%Cerebros y máquinas

%- Autores incluyendo J. R. Lucas y Roger Penrose debaten que los teoremas tienen implicaciones sobre la inteligencia humana. El debate se centra sobre si la mente humana es equivalente a una máquina de Turing o (mediante la tesis de Church-Turing) a una máquina finita. Si la mente lo es, y los teoremas de Gödel se aplican a esta, los teoremas de incompletitud se aplicarían a ella.

%Hilary Putnam 

%Avi Widgerson

%Douglas Hofstadter 

%- 

%Consecuencias severas para el logicismo de Gottlob Frege y Bertrand Russell, que pretendían definir los números naturales en términos de lógica

%\begin{tabla}
%\begin{table}[H]
%\centering
%\begin{adjustwidth}{-12pt}{}
%\begin{tabular}{@{}ll@{}}
%\toprule
%Máquinas de Turing & Numeración Gödel \\
%(\cref{teo:incompletitud-semantica}) & (\cref{teo:incompletitud-godel}) \\ \midrule
%para sistemas lógicos (semántica) & para sistemas formales (sintaxis) \\[4pt]
%asume que el sistema lógico es sólido & asume que el sistema formal es $\omega$-consistente \\[4pt]
%usa sistemas lógicos capaces de probar & usa sistemas formales capaces de probar \\
%afirmaciones del tipo ``$P$ para con &  ciertos teoremas de aritmética elemental \\
%entrada vacía'' \\ \bottomrule
%\end{tabular}
%\end{adjustwidth}
%\end{table}
%\vspace{-8pt}
%\caption{Diferencias en los teoremas de incompletitud mostrados en este capítulo}
%\label{tab:diferencias-turing-godel}
%\end{tabla}

%\section{El Segundo Teorema de Incompletitud}\label{sec:teo-incompletitud-2}

%El Segundo Teorema de Incompletitud, en su versión sintáctica, puede obtenerse inmediatamente a partir del Primer Teorema de Incompletitud.

%\begin{lema}\label{lema:red-prob}
%Sea la siguiente fórmula un teorema:
%$$\phi_1\Rightarrow\phi_2$$

%Entonces, si $\phi_2$ no es un teorema, $\phi_1$ tampoco lo es.
%\end{lema}
%\begin{proof}
%Este resultado se obtiene de la implicación inmediata siguiente: si $\phi_1$ es un teorema, entonces $\phi_2$ ha de serlo también, pues se deduce inmediatamente de $\phi_1\Rightarrow\phi_2$ mediante \emph{modus ponens}. Concluimos la prueba por el contrarrecíproco.
%\end{proof}

%\begin{teorema}[Versión semántica del Segundo Teorema de Incompletitud]
%Sea $\mathcal{G}$ un sistema lógico que verifique las siguientes condiciones:
%\begin{itemize}
%    \item[(A2)] $\mathcal{G}$ es recursivamente axiomatizable.
%    \item[(A3)] $\mathcal{G}$ puede expresar afirmaciones sobre máquinas de Turing o, equivalentemente, sobre programas; en concreto, dada una máquina de Turing $M$, puede expresar que ``$M$ para con entrada vacía'' o, equivalentemente, dado un programa $P$, puede expresar que ``$P$ para con entrada vacía''.
%\end{itemize}
%Entonces, la fórmula correspondiente a ``$\mathcal{G}$ es sólido'' no es un teorema en $\mathcal{G}$.
%\end{teorema}
%\begin{proof}
%El \cref{teo:incompletitud-semantica} nos dice que, suponiendo (A2) y (A3), en caso de que $\mathcal{G}$ fuese sólido, entonces tendríamos una fórmula $\phi$ que (1) es verdadera pero (2) no es un teorema.

%Como hemos logrado probar (1) en $\mathcal{G}$, tenemos el siguiente teorema:
%$$``\mathcal{G}\text{ es sólido}''\Rightarrow\phi$$
%Ahora bien, aplicando (2) y el \cref{lema:red-prob}, tenemos que ``$\mathcal{G}$ es sólido'' no es un teorema.
%\end{proof}

%\emph{\textbf{Nota: 6.1} para probar la versión sintáctica (que es el verdadero segundo teorema de incompletitud) sería necesario encontrar una versión constructiva del \cref{teo:incompletitud-sintactica}, que nos diese una fórmula $\phi$ que fuese verdadera de modo que ni $\phi$ ni $\neg\phi$ sean teoremas. Pero esto es un tanto extraño porque si el teorema es sintáctico, ¿qué quiere decir que una fórmula sea verdadera?}

%\section{Consecuencias}\label{sec:teo-incompletitud-consecuencias}


\endinput


\setpartpreamble[c][0.75\linewidth]{%
	\bigskip % Deja un espacio vertical en la parte superior
  A partir de la decidibilidad del problema de parada, probaremos que la aritmética de Peano es incompleta, bajo unas ciertas condiciones. Veremos que la elección de la aritmética de Peano no restringe sacar conclusiones de otros sistemas lógicos. A continuación, compraremos la demostración de este trabajo con la original de Gödel. Finalmente, veremos las limitaciones de nuestro teorema, y aprenderemos un truco que nos permitirá relajar las condiciones del primer teorema.
}
\cleardoublepage\part{Incompletitud}
% !TeX root = ../libro.tex
% !TeX encoding = utf8

\setchapterpreamble[c][0.75\linewidth]{%
	\sffamily
  \emph{El objetivo de la computación es el conocimiento, no los números.}
 \begin{flushright} — Richard Hamming, \cite{Hamming1986} \end{flushright}
% \\[8pt]
	\par\bigskip
}
%\vspace{28pt}

\chapter{Motivación}\label{ch:motivacion}

Los teoremas de incompletitud de Gödel son uno de los resultados más populares de las matemáticas y la filosofía, y han sido objeto de fascinación y estudio durante décadas. Estos teoremas revolucionaron los mismos fundamentos de ambas disciplinas. En este trabajo, abordaremos el primero de estos teoremas desde una perspectiva computacional. Esto es así debido a una serie de experiencias y observaciones que despertaron en mí un gran interés y que revelaron una oportunidad para abordar este tema desde una perspectiva computacional pero, sobre todo, matemáticamente rigurosa.

Este interés surge de las clases de Modelos de Computación impartidas por el tutor de este trabajo, en las que mencionó de forma breve este resultado. Tras ello sentí mucha curiosidad por profundizar en el mundo de la incompletitud, hasta llegar a un excepcional libro de Douglas Hofstadter: \emph{``Gödel, Escher, Bach: An Eternal Golden Braid''} \cite{Hofstadter1999}. Esta obra maestra interdisciplinaria combina elementos de matemáticas, música y arte, y me incitó a profundizar más en estos teoremas.

La mayor parte de los libros que tratan este tema lo hacen tal y como lo hizo Gödel: desde una perspectiva aritmética, usando la numeración de Gödel en las demostraciones, o variantes. Sin embargo, en múltiples cursos de teoría de la computación se exponen demostraciones haciendo uso de algoritmos y máquinas de Turing.

Para mi sorpresa, a pesar de que esta perspectiva se estudia extensivamente y hay una gran literatura al respecto, muchas de las referencias que encontraba eran bien poco rigurosas o bien poco consistentes en la nomenclatura y en la forma de probar los teoremas. En concreto, muchas de las fuentes no distinguen de forma precisa entre los conceptos de \emph{demostración} y \emph{verdad}.

Es aquí donde vi la oportunidad de hacer un trabajo que intentase proporcionar una perspectiva computacional a los teoremas de incompletitud de Gödel, sin perder rigor y concreción por el camino. Este enfoque permitirá una comprensión más precisa y estructurada de los resultados más relevantes al respecto, así como su relación con los modelos de computación y la teoría de la computabilidad.

\pagebreak

\section*{Objetivos}

Los objetivos de este trabajo son:

\subsection*{Objetivos teóricos}

\begin{enumerate}[label={[}T\arabic*{]}]
    \item Probar el Primer Teorema de Incompletitud de Gödel desde una perspectiva computacional.
    \item Usar para ello programas en Python en lugar de descripciones de máquinas de Turing, probando que esto es posible y riguroso.
    \item Exponer las consecuencias del teorema de incompletitud en las disciplinas matemática y filosófica.
    \item Contextualizar el teorema de incompletitud en el ámbito histórico.
    \item Encontrar la relación entre el teorema de incompletitud con la tesis de Church-Turing.
    \item Precisar conceptos lógicos relacionados con la teoría de la demostración. En concreto, distinguir entre las nociones de \emph{demostración} y \emph{verdad}, y estudiar las propiedades de los sistemas de demostración.
\end{enumerate}

\subsection*{Objetivos prácticos}

\begin{enumerate}[label={[}P\arabic*{]}]
    \item Crear contenido que pueda servir de herramienta de apoyo a la docencia.
    \item Crear una librería de programas en Python ejecutables por el usuario que permitan guiar las demostraciones que aparecen en este trabajo.
\end{enumerate}

\endinput

% !TeX root = ../libro.tex
% !TeX encoding = utf8

\setchapterpreamble[c][0.75\linewidth]{%
	\sffamily
  \emph{Si una ``religión'' se define como un sistema de ideas que contiene afirmaciones indemostrables, entonces Gödel nos enseñó que las matemáticas no sólo son una religión, sino que son la única religión que puede demostrar que lo es.}
 \begin{flushright} — John D. Barrow, \cite{Barrow2011} \end{flushright}
% \\[8pt]
	\par\bigskip
}
%\vspace{28pt}

\chapter{El Primer Teorema de Incompletitud}\label{ch:teorema-incompletitud}

Los matemáticos de principios del siglo \textsc{xx}, tal y como Leibniz, Boole y Hilbert, imaginaron que sería posible determinar la veracidad de fórmulas matemáticas de forma automática. Nosotros ahora nos preguntamos lo mismo: hemos trabajado con sistemas lógicos sencillos que son decidibles, es decir, cuyas fórmulas pueden ser comprobadas verdaderas mediante programas. En el caso de sistemas lógicos de más complejidad, como la aritmética de Peano, esto no es posible. Esto se debe a que la existencia de problemas computacionales indecidibles conduce directamente a la existencia de afirmaciones matemáticas no decidibles.

%En la \cref{sec:primera-aproximacion}, probaremos la incompletitud de \textbf{Peano}. Llamaremos a este resultado \emph{Primer Teorema de Incompletitud}. A continuación, generalizaremos este teorema para otros sistemas lógicos. Finalmente, en la \cref{sec:diferencias-godel}, veremos el enunciado original del teorema y lo compararemos con el resultado obtenido.

\section{Una primera aproximación}\label{sec:primera-aproximacion}

%\begin{teorema}[Versión no constructiva del Primer Teorema de Incompletitud]\label{teo:inco1-no-constr}
%La aritmética de Peano es incompleta: existen fórmulas verdaderas de \textbf{Peano} que no son teoremas.
%\end{teorema}
%\begin{proof}
%Asumiremos que \textbf{Peano} es completo y entraremos en contradicción. En concreto, demostraremos que si la aritmética de Peano es completa, podremos decidir \textbf{ParadaEnVacío}, lo cual contradice la \cref{prop:parada-en-vacio-no-decidible}.

%Sea $P$ un programa cualquiera. Lo usaremos como entrada al \cref{lst:parada-en-vacio-desde-peano-completo}. En la \cref{line:parada-en-vacio-desde-peano-completo-import1} importamos la función \texttt{es\_prueba\_peano}, que dada una fórmula y una demostración en \textbf{Peano} comprueba si la demostración prueba que la fórmula es un teorema. Claramente, esta función no es un oráculo. De hecho, su implementación es sencilla: basta probar si cada una de las fórmulas de la demostración es un axioma o se sigue de las fórmulas anteriores a través de las reglas de inferencia del sistema, comprobando que la última fórmula coincida con la fórmula que queremos probar.

%A continuación, importamos en la \cref{line:parada-en-vacio-desde-peano-completo-import2} la función \texttt{parada\_a\_peano} cuya existencia comentamos tras la \cref{prop:parada-a-peano}, que pasa la afirmación ``$P$ para con entrada vacía'' a una fórmula en \textbf{Peano}.

%Usamos esta función en la \cref{line:parada-en-vacio-desde-peano-completo-para-en-peano} para generar la fórmula mencionada anteriormente. En la \cref{line:parada-en-vacio-desde-peano-completo-no-para-en-peano} creamos una nueva fórmula negando la anterior, por tanto \texttt{no\_para\_en\_peano} contiene la fórmula de \textbf{Peano} correspondiente a la afirmación ``$P$ no para con entrada vacía''.

%Una de estas afirmaciones ha de ser verdadera. Como asumimos que \textbf{Peano} es completo, una de estas fórmulas tiene una prueba (es un teorema) en \textbf{Peano}. El bucle de la \cref{line:parada-en-vacio-desde-peano-completo-while} va iterando a través de todos los \emph{strings} posibles, comprobando cada uno para saber si es una prueba de \texttt{para\_en\_peano} o de \texttt{no\_para\_en\_peano}. Ya que uno de ellos debe tener una prueba, el programa terminará, decidiendo si $P$ para con entrada vacía o no. El \cref{lst:parada-en-vacio-desde-peano-completo} resuelve \textsc{ParadaEnVacío}, lo que entra en contradicción con la \cref{prop:parada-en-vacio-no-decidible}.
%\end{proof}

%\begin{lstlisting}[language=Python, caption=\lstinline{parada_en_vacio_desde_peano_completo.py},label={lst:parada-en-vacio-desde-peano-completo}]
%import utilidades

%from es_prueba_peano import es_prueba_peano # NO es un oráculo|\label{line:parada-en-vacio-desde-peano-completo-import1}|
%from parada_a_peano import parada_a_peano # NO es un oráculo|\label{line:parada-en-vacio-desde-peano-completo-import2}|

%def parada_en_vacio_desde_peano_completo(programa):|\label{line:parada-en-vacio-desde-peano-completo-main}|
%    para_en_peano = parada_a_peano(programa)|\label{line:parada-en-vacio-desde-peano-completo-para-en-peano}|
%    no_para_en_peano = 'NOT (' + para_en_peano + ')'|\label{line:parada-en-vacio-desde-peano-completo-no-para-en-peano}|
%    prueba = ''|\label{line:parada-en-vacio-desde-peano-completo-prueba}|
    
%    while True:|\label{line:parada-en-vacio-desde-peano-completo-while}|
%        if es_teorema_peano(prueba, para_en_peano) == 'sí':|\label{line:parada-en-vacio-desde-peano-completo-if1}|
%            return 'sí'
%        if es_teorema_peano(prueba, no_para_en_peano) == 'sí':|\label{line:parada-en-vacio-desde-peano-completo-if2}|
%            return 'no'
%        prueba = utilidades.siguiente_string(prueba)|\label{line:parada-en-vacio-desde-peano-completo-else}|
%\end{lstlisting}


Antes de proceder la demostración de una primera versión del Primer Teorema de Incompletitud, discutiremos de forma detallada el \cref{lst:godel}. En la \cref{line:godel-peano-import1} importamos la función \texttt{es\_teorema\_peano}, cuya existencia probamos en la \cref{prop:es-teorema-peano-semidecidible}. Recuerda que, como los axiomas de \textbf{Peano} son recursivamente enumerables (es decir, \textbf{Peano} es recursivamente axiomatizable), podemos escribir el programa \texttt{es\_teorema\_peano(fórmula)} que devuelve \palabra{sí} si \texttt{fórmula} es un teorema en \textbf{Peano}, y que no para en caso contrario (ver \cref{prop:es-teorema-peano-semidecidible}).
\vspace{8pt}
\begin{lstlisting}[language=Python, caption=\lstinline{godel_peano.py},label={lst:godel-peano}]
import utilidades

from es_teorema_peano import es_teorema_peano # NO es un oráculo|\label{line:godel-peano-import1}|
from parada_a_peano import parada_a_peano # NO es un oráculo|\label{line:godel-peano-import2}|

def godel_peano(entrada):|\label{line:godel-peano-main}|
    programa_godel = utilidades.leer('godel\_peano.py')|\label{line:godel-peano-programa}|
    para_en_peano = parada_a_peano(programa_godel)|\label{line:godel-peano-formula}|
    no_para_en_peano = 'NOT (' + para_en_peano + ')'|\label{line:godel-peano-formulaneg}|

    if es_teorema_peano(no_para_en_peano) == 'sí':|\label{line:godel-peano-es-teorema}|
        return 'para'|\label{line:godel-peano-para}|
    else:
        utilidades.ciclar()|\label{line:godel-peano-cicla}|
\end{lstlisting}

La \cref{line:godel-peano-import2} importa la función \texttt{parada\_a\_peano}, cuya existencia ya discutimos tras la \cref{prop:parada-a-peano}. Esta función traduce, dado un programa $P$, la afirmación ``$P$ para con entrada vacía'' en una fórmula de \textbf{Peano}.

En la \cref{line:godel-peano-programa} guardamos el propio programa \texttt{godel\_peano.py} en la variable \texttt{programa\_godel}. En la \cref{line:godel-peano-formula} usamos la función \texttt{parada\_a\_peano} recién comentada, y en la \cref{line:godel-peano-formulaneg} negamos la fórmula \texttt{para\_en\_peano}. De este modo:
\begin{adjustwidth}{30pt}{}
    \texttt{no\_para\_en\_peano} significa que ``\texttt{godel\_peano.py} no para con entrada vacía'' \hfill $(\star1)$
\end{adjustwidth}
Ahora es cuando las cosas se ponen interesantes. Si \texttt{no\_para\_en\_peano} es un teorema, la función \texttt{es\_teorema\_peano} devolverá \palabra{sí}, y de este modo el programa verificará la condición de la \cref{line:godel-peano-es-teorema} y devolverá \palabra{para}, haciendo que el programa pare. Si \texttt{no\_para\_en\_peano} no es un teorema, se entrará en un bucle infinito (\texttt{es\_teorema\_peano} no parará, y por lo tanto tampoco \texttt{godel\_peano.py}). Es claro que la \cref{line:godel-peano-cicla} nunca se ejecutará, pero la incluimos para enfatizar este hecho.

Ahora, nos preguntamos: ¿cuál es el comportamiento de \texttt{godel\_peano.py} con entrada vacía?

\begin{proposicion}\label{prop:godel-peano-no-para}
El \cref{lst:godel-peano} no para con entrada vacía.
\end{proposicion}
\begin{proof}
Asumimos que \texttt{godel\_peano.py} para con entrada vacía, y entraremos en contradicción. La única forma de que pare es en la \cref{line:godel-peano-para}, algo que ocurre solamente si:
\begin{adjustwidth}{30pt}{}
    \texttt{es\_teorema\_peano(no\_para\_en\_peano) == \palabra{sí}} \hfill $(\star2)$
\end{adjustwidth}
$(\star2)$ es equivalente a decir que \texttt{no\_para\_en\_peano} es un teorema. Por la solidez de \textbf{Peano}, \texttt{no\_para\_en\_peano} es verdadero.

Pero el hecho de que \texttt{no\_para\_en\_peano} sea verdadero, junto con $(\star1)$, nos dice que \texttt{godel\_peano.py} no para con entrada vacía, encontrando la contradicción esperada.
\end{proof}

Ahora estamos en condiciones de probar el resultado principal de este trabajo.

\begin{teorema}[Incompletitud de la aritmética de Peano]\label{teo:incompletitud-peano}
\index{primer teorema de incompletitud}\index{teorema de incompletitud}La aritmética de Peano es incompleta: existen fórmulas verdaderas de \normalfont{\textbf{Peano}} que no son teoremas. Concretamente, la fórmula guardada en la variable \texttt{no\_para\_en\_peano} en la \cref{line:godel-peano-formulaneg} del \cref{lst:godel} es un ejemplo de una fórmula verdadera pero no demostrable (no es un teorema) de \normalfont{\textbf{Peano}}.
\end{teorema}
\begin{proof}
Debemos mostrar que \texttt{no\_para\_en\_peano} es verdadero y no es un teorema.

Para probar que es verdadero, por la \cref{prop:godel-peano-no-para}, sabemos que \texttt{godel\_peano.py} no para con entrada vacía. Recordando $(\star1)$, inmediatamente obtenemos que \texttt{no\_para\_en\_peano} es verdadero.

Para probar que no es un teorema, de nuevo usamos el hecho de que \texttt{godel\_peano.py} no para con entrada vacía. Esto nos indica que, en la \cref{line:godel-peano-es-teorema}, la llamada a la función
\begin{adjustwidth}{30pt}{}
    \texttt{es\_teorema\_peano(no\_para\_en\_peano)}
\end{adjustwidth}
no puede devolver \palabra{sí}. Pero esto inmediatamente implica que \texttt{no\_para\_en\_peano} no es demostrable, lo que completa la prueba.
\end{proof}

%\section{El Segundo Teorema de Incompletitud}\label{sec:teo-incompletitud-2}

%\begin{teorema}[Segundo Teorema de Incompletitud]
%La consistencia de la aritmética de Peano no puede probarse en la aritmética de Peano. En concreto, la afirmación
%\begin{adjustwidth}{30pt}{}
%    ``\textbf{Peano} es consistente''
%\end{adjustwidth}
%que puede codificarse mediante una fórmula de \textbf{Peano}, no es un teorema.
%\end{teorema}

Acabamos de demostrar que \textbf{Peano} es incompleto, como ya anticipamos en la \cref{tab:sistemas-logicos}. Las consecuencias de este hecho las exploramos en detalle en el \cref{ch:conclusion}. Sin embargo, paremos un momento a analizar qué es lo que acabamos de probar.

\section{Versión semántica}\label{sec:version-semantica}

Para encontrar una fórmula verdadera y no demostrable, hemos necesitado un sistema lógico que sea, en primer lugar, recursivamente axiomatizable. Si no, no hubiesemos podido crear un programa capaz de saber cuándo una fórmula del sistema es un teorema. Por otra parte, hemos asumido que el sistema es sólido, lo cual es algo esperable. Finalmente, hemos exigido que sea capaz de expresar una cierta afirmación de las máquinas de Turing (véase \cref{prop:parada-a-peano}). Para evitar repeticiones innecesarias en esta última hipótesis, llamaremos a un sistema lógico con esta propiedad ``de tipo Turing'', algo que definimos a continuación.

\begin{definicion}[Sistema de tipo Turing]
Sea $\mathcal{S}$ un sistema formal. Diremos que $\mathcal{S}$ es \emph{de tipo Turing}\footnote{Esta definición no se encuentra en la literatura: se ha creado por conveniencia.} si puede expresar afirmaciones sobre máquinas de Turing o, equivalentemente, sobre programas;\footnote{La equivalencia la probamos en el \cref{teo:equivalencia}.} en concreto, dada una máquina de Turing $M$, puede expresar que ``$M$ para con entrada vacía''; o, equivalentemente, dado un programa $P$, puede expresar que ``$P$ para con entrada vacía''.

Análogamente podemos definirlo para el caso de sistemas lógicos.
\end{definicion}

Observemos que, sin embargo, en ningún momento hemos utilizado el sistema \textbf{Peano} en sí. Esto nos indica que podríamos realizar esta demostración con otro sistema que no sea \textbf{Peano}, siempre y cuando verifique tales condiciones. Indicamos este resultado a continuación.

\begin{teorema}[Versión semántica del Primer Teorema de Incompletitud]\label{teo:incompletitud-semantica}
Sea $\mathcal{G}$ un sistema lógico que verifique las siguientes condiciones:
\begin{enumerate}
    \item $\mathcal{G}$ es sólido.
    \item $\mathcal{G}$ es recursivamente axiomatizable.
    \item $\mathcal{G}$ es de tipo Turing.
\end{enumerate}
Entonces, $\mathcal{G}$ es incompleto, es decir, existen fórmulas verdaderas de $\mathcal{G}$ que no son teoremas.
\end{teorema}

Para reforzar este importante resultado, incluimos un resumen del proceso que hemos seguido para \textbf{Peano}, de forma general, en la demostración de este teorema.

\begin{proof}
En el \cref{lst:godel} importamos dos funciones. Por una parte, en la \cref{line:godel-import1} importamos la función \texttt{es\_teorema\_en\_G}, que comprueba si una fórmula es un teorema en $\mathcal{G}$. Esta función parará si la fórmula es un teorema, y se ejecutará indefinidamente en caso de que no lo sea. La existencia de esta función la garantiza (2).

Por otra, en la \cref{line:godel-import2} importamos la función \texttt{parada\_a\_G} que, dado un programa $P$, convierte la afirmación ``$P$ para con entrada vacía'' en una fórmula de $\mathcal{G}$, algo que es posible por (3).

En la línea \cref{line:godel-formulaneg} tenemos que:
\begin{adjustwidth}{30pt}{}
    \texttt{no\_para\_en\_G} significa que \emph{``\texttt{godel.py} no para con entrada vacía''} (en $\mathcal{G}$)  \hfill $(\star3)$
\end{adjustwidth}
Si \texttt{no\_para\_en\_G} es un teorema, la función \texttt{es\_teorema\_en\_G(no\_para\_en\_G)} (\cref{line:godel-es-teorema}) parará y devolverá \palabra{sí}, mientras que la misma función ciclará en caso de que no lo sea.

La pregunta que debemos hacernos es: ¿qué le pasa a \texttt{godel.py} cuando \texttt{entrada} es vacía?

Si asumimos que \texttt{godel.py} para con entrada vacía, la única forma de que lo haga es porque en la \cref{line:godel-es-teorema} es \texttt{es\_teorema\_en\_G(no\_para\_en\_G) == \palabra{sí}}, lo cual solo puede pasar si \texttt{no\_para\_en\_G} es un teorema. Pero, por (1), y dado que $\mathcal{G}$ es sólido, tenemos que \texttt{no\_para\_en\_G} es verdadero.

Pero si \texttt{no\_para\_en\_G} es verdadero, por $(\star3)$, esto querría decir que \texttt{godel.py} no para con entrada vacía, entrando en contradicción, por lo que debe ser:
\begin{adjustwidth}{30pt}{}
    \texttt{godel.py} no para con entrada vacía  \hfill $(\star4)$
\end{adjustwidth}
Comprobamos que \texttt{no\_para\_en\_G} es verdadero pero no es un teorema. En efecto, es verdadero por $(\star4)$. Sin embargo, no es un teorema. Para ello, usamos de nuevo $(\star4)$. El hecho de que \texttt{godel.py} no pare indica que es imposible que la condición de la \cref{line:godel-es-teorema}, \linebreak\texttt{es\_teorema\_en\_G(no\_para\_en\_G) == \palabra{sí}}, se verifique -- en otras palabras, \texttt{no\_para\_en\_G} no es un teorema.
\end{proof}

%\vspace{8pt}
\begin{lstlisting}[language=Python, caption=\lstinline{godel.py},label={lst:godel}]
import utilidades

from es_teorema_G import es_teorema_G # NO es un oráculo|\label{line:godel-import1}|
from parada_a_G import parada_a_G # NO es un oráculo|\label{line:godel-import2}|

def incompleto_semantico(entrada):|\label{line:godel-main}|
    programa = utilidades.leer('incompleto_semantico.py')|\label{line:godel-programa}|
    para_en_G = parada_a_G(programa)|\label{line:godel-formula}|
    no_para_en_G = 'NOT (' + para_en_G + ')'|\label{line:godel-formulaneg}|

    if es_teorema_G(no_para_en_G) == 'sí':|\label{line:godel-es-teorema}|
        return 'para'|\label{line:godel-para}|
    else:
        utilidades.ciclar()|\label{line:godel-cicla}|
\end{lstlisting}


Nuestro resultado nos dice que, si tenemos un sistema lógico que sea recursivamente axiomatizable y de tipo Turing, no podemos esperar que pueda ser sólido y completo al mismo tiempo. Es decir, si todos los teoremas del sistema son verdaderos, siempre podremos encontrar fórmulas del sistema que sean verdaderas pero no sean teoremas.

Podemos aplicar el \cref{teo:incompletitud-semantica} a otros sistemas lógicos, siempre que verifiquen las hipótesis. Uno de ellos es, por ejemplo, el sistema de axiomas de Zermelo-Fraenkl.

\section{Versión sintáctica}\label{sec:version-sintactica}

% WIP relajar condición de solidez?
El resultado que acabamos de probar es semántico. Sin embargo, es posible probar otro resultado, esta vez sintáctico, y más parecido al demostrado por Gödel.

Para poder probar esta variación, debemos de tener en cuenta una sutileza del problema \textsc{EsTeorema} en el caso de sistemas formales sintácticamente consistentes.

\begin{proposicion}\label{prop:es-teorema-consistente-decidible}
Sea $\mathcal{S}$ un sistema formal recursivamente axiomatizable y sintácticamente consistente. Entonces, $\textsc{EsTeorema}_\mathcal{S}$ es decidible.
\end{proposicion}
\begin{proof}
Veamos que el \cref{lst:es-teorema-s} decide $\textsc{EsTeorema}_\mathcal{S}$. En efecto, vemos en primer lugar que en la \cref{line:es-teorema-s-import} importamos la función \texttt{es\_prueba\_S}, la cual es computable por ser $\mathcal{S}$ recursivamente axiomatizable.

A continuación, en la \cref{line:es-teorema-s-negacion} almacenamos la negación de \texttt{formula} en \texttt{neg\_formula}. En la \cref{line:es-teorema-s-demostracion} declaramos la demostración.

Por ser $\mathcal{S}$ consistente, bien \texttt{formula} bien \texttt{neg\_formula} tiene una demostración. En el bucle de la \cref{line:es-teorema-s-while}, vamos comprobando si \texttt{demostracion} es una demostración de \texttt{formula} o de \texttt{neg\_formula}. En el primer caso, devolvemos \palabra{sí}, mientras que en el segundo, devolvemos \palabra{no}. En caso de que \texttt{demostracion} no sea una demostración de ninguna de ambas fórmulas, procedemos en la \cref{line:es-teorema-s-siguiente-string} a enumerar la siguiente demostración en orden \emph{shortlex} (ver \cref{prop:es-teorema-peano-semidecidible}).

Por la consistencia de $\mathcal{S}$, debe existir una demostración bien de \texttt{formula} o de \texttt{neg\_formula}, por lo que el programa debe parar al llegar a ser \texttt{demostracion} tal demostración.
\end{proof}

\begin{lstlisting}[language=Python, caption=\lstinline{es_teorema_S.py},label={lst:es-teorema-s}]
import utilidades
from es_prueba_S import es_prueba_S # NO es un oráculo|\label{line:es-teorema-s-import}|

def es_teorema_S(formula):|\label{line:es-teorema-s-main}|
    negacion_formula = 'NOT (' + formula + ')'|\label{line:es-teorema-s-negacion}|
    demostracion = ''|\label{line:es-teorema-s-demostracion}|

    while True:|\label{line:es-teorema-s-while}|
        if es_prueba_S(demostracion, formula) == 'sí':|\label{line:es-teorema-s-es-prueba}|
            return 'sí'
        if es_prueba_S(demostracion, negacion_formula) == 'sí':
            return 'no'
        demostracion = utilidades.siguiente_string(demostracion)|\label{line:es-teorema-s-siguiente-string}|
\end{lstlisting}

Ahora estamos en condiciones de probar el resultado sintáctico, que dependerá de la no decidibilidad de \textsc{AdivinaConsistente} (\cref{prob:adivina-consistente}). \cite{Aaronson2017}

\begin{teorema}[Versión sintáctica del Primer Teorema de Incompletitud]\label{teo:incompletitud-sintactica}
Sea $\mathcal{S}$ un sistema formal que verifique las siguientes condiciones:
\begin{enumerate}
    \item $\mathcal{S}$ es sintácticamente consistente.
    \item $\mathcal{S}$ es recursivamente axiomatizable.
    \item $\mathcal{S}$ es de tipo Turing.
\end{enumerate}
Entonces, $\mathcal{S}$ es sintácticamente incompleto, es decir, existen fórmulas $\phi$ de $\mathcal{S}$ tal que ni $\phi$ ni $\neg\phi$ son teoremas.
\end{teorema}
\begin{proof}
Vamos a suponer que $\mathcal{S}$ es sintácticamente completo y veremos que, en tal caso, podríamos decidir \textsc{AdivinaConsistente} (\cref{prob:adivina-consistente}), lo cual entraría en contradicción con la \cref{prop:adivina-consistente-no-decidible}.

Comencemos comentando las funciones importadas. En la \cref{line:adivina-consistente-import-1} importamos la función \texttt{oarada\_en\_S}, que traduce para un programa $P$ la afirmación ``$P$ para con entrada vacía'' en una fórmula de $\mathcal{}S$, y cuya existencia garantiza la hipótesis (3). En la \cref{line:adivina-consistente-import-2} importamos la función \texttt{es\_teorema\_S}, que decide el problema $\textsc{EsTeorema}_\mathcal{S}$, y que, por las hipótesis (1) y (2), siempre para (ver \cref{prop:es-teorema-consistente-decidible}) y devuelve \palabra{sí} si la fórmula de entrada es un teorema en $\mathcal{S}$ y \palabra{no} en caso contrario.

En la \cref{line:adivina-consistente-programa-para} usamos la función \texttt{parada\_a\_S} para almacenar en \texttt{programa\_para} la fórmula de $\mathcal{S}$ equivalente a ``\texttt{programa} para con entrada vacía''. En la \cref{line:adivina-consistente-es-teorema}, asignamos a la variable \texttt{es\_teorema} uno de dos valores posibles: \palabra{sí} o \palabra{no}, como ya comentamos anteriormente.

Llegado a la \cref{line:adivina-consistente-si}, implementamos la siguiente lógica:
\begin{itemize}
    \item Si \texttt{programa} acepta con entrada vacía (para y devuelve \palabra{sí}), la condición de la \cref{line:adivina-consistente-si} se verificará, así como la de la \cref{line:adivina-consistente-si-si}, y \texttt{adivina\_consistente.py} devolverá \palabra{sí}.
    \item Si \texttt{programa} rechaza con entrada vacía (para y devuelve \palabra{no}, la condición de la \cref{line:adivina-consistente-si} se verificará, mientras que no la de la \cref{line:adivina-consistente-si-si}, y \texttt{adivina\_consistente.py} devolverá \palabra{no}.
    \item Si \texttt{programa} cicla con entrada vacía, la condición de la \cref{line:adivina-consistente-si} no se verificará, y \texttt{adivina\_consistente.py} devolverá \palabra{no}.
\end{itemize}
Dicha lógica hace evidente que el programa decide \textsc{AdivinaConsistente}.
\end{proof}

\begin{lstlisting}[language=Python, caption=\lstinline{adivina_consistente.py},label={lst:adivina-consistente}]
import utilidades
from parada_a_S import parada_a_S # NO es un oráculo|\label{line:adivina-consistente-import-1}|
from es_teorema_S import es_teorema_S # NO es un oráculo|\label{line:adivina-consistente-import-2}|

def adivina_consistente(programa):|\label{line:adivina-consistente-main}|
    programa_para = parada_a_S(programa)|\label{line:adivina-consistente-programa-para}|

    es_teorema = es_teorema_S(programa_para)|\label{line:adivina-consistente-es-teorema}|

    if es_teorema == 'sí':|\label{line:adivina-consistente-si}|
        salida = maquina_universal(programa, '')|\label{line:adivina-consistente-universal}|
        if salida == 'sí':|\label{line:adivina-consistente-si-si}|
            return 'sí'
        else:|\label{line:adivina-consistente-si-no}|
            return 'no'
    else:|\label{line:adivina-consistente-no}|
        return 'no'
\end{lstlisting}

Hemos probado dos versiones del Primer Teorema de Incompletitud. Resumimos las diferencias entre ambas a continuación.
\vspace{8pt}
\begin{figure}[H]
$$
\begin{array}{ccc}
    {
        \begin{matrix}
            \text{\textbf{versión semántica}} \\
            \text{\small{(\cref{teo:incompletitud-semantica})}}\\
            \small{\;}
        \end{matrix}
    }
    & \;\;\; &
    {
        \begin{matrix}
            \text{\textbf{versión sintáctica}} \\
            \text{\small{(\cref{teo:incompletitud-sintactica})}}\\
            \small{\;}
        \end{matrix}
    }
    \\
    {
        \begin{matrix}
            \mathcal{G}\text{ sistema lógico}\\
            \text{rec. axiom., de tipo Turing}\\
            \text{y sólido} \\
            \small{\phi \text{ teorema} \Rightarrow \phi\text{ verdadero}}\\
            \tiny{\;}
        \end{matrix}
    }
    &  & 
    {
        \begin{matrix}
            \mathcal{S}\text{ sistema formal}\\
            \text{rec. axiom., de tipo Turing}\\
            \text{y sintácticamente consistente} \\
            \small{\nexists \phi : \phi\text{ y }\neg\phi \text{ son teoremas}}\\
            \tiny{\;}
        \end{matrix}
    } \\
    \Huge{\Downarrow} & & \Huge{\Downarrow} \\
    {
        \begin{matrix}
            \tiny{\;}\\
            \mathcal{G}\text{ no (semánticamente) completo}\\
            \small{\exists\phi:\phi\text{ es verdadero y no es teorema}}
        \end{matrix}
    }
    & &
    {
        \begin{matrix}
            \tiny{\;}\\
            \mathcal{S}\text{ no sintácticamente completo}\\
            \small{\exists\phi:\text{ni }\phi\text{ ni }\neg\phi\text{ son teoremas}}
        \end{matrix}
    }
    \\
\end{array}
$$
\caption{Comparativa entre las dos versiones del Primer Teorema de Incompletitud}
\vspace{-8pt}
\label{fig:comparativa-incompletitud}
\end{figure}

La demostración original de Gödel es puramente sintáctica, y muy parecida a la probada en el \cref{teo:incompletitud-sintactica}. La enunciamos a continuación.

\begin{teorema}[Primer Teorema de Incompletitud de Gödel-Rosser]\label{teo:incompletitud-godel}
Sea $\mathcal{S}$ un sistema formal que verifique las siguientes condiciones:
\begin{enumerate}
    \item $\mathcal{S}$ es sintácticamente consistente.
    \item $\mathcal{S}$ es recursivamente axiomatizable.
    \item $\mathcal{S}$ contiene una cierta cantidad de aritmética elemental.
\end{enumerate}

Entonces, $\mathcal{S}$ es sintácticamente incompleto, es decir, hay fórmulas $\phi$ de $\mathcal{S}$ de modo que ni $\phi$ ni $\neg\phi$ son teoremas.
\end{teorema}

Tal y como en nuestra versión sintáctica, requerimos que el sistema sea sintácticamente consistente y recursivamente axiomatizable. El teorema original de Gödel exigía como hipótesis una forma más fuerte de consistencia sintáctica llamada \emph{$\omega$-consistencia}. Esta forma más general del teorema se debe al trabajo de John B. Rosser, autor del hoy conocido como \emph{truco de Rosser}.\footnote{Es por ello por lo que lo llamamos teorema de ``Gödel-Rosser'', para distinguirlo del resultado original de Gödel.} \cite{Rosser1936}

La tercera hipótesis exige que el sistema sea capaz de probar ciertas afirmaciones de aritmética elemental. Esto se debe a que Gödel pretendía probar que la teoría de números no era axiomatizable (ver \cref{ch:historia}), por lo que usó en su demostración una forma de numerar los teoremas que hoy en día se conoce como \emph{numeración de Gödel}.


\section{Consecuencias}\label{sec:consecuencias}

El Primer Teorema de Incompletitud tiene profundas implicaciones matemáticas y filosóficas.

\subsection{Relación con el segundo problema de Hilbert}\label{subsec:segundo-problema-hilbert}

Ya introdujimos en el \cref{ch:historia} los célebres problemas de Hilbert, y en concreto, el segundo de ellos, que volveremos a referenciar a continuación:

\begin{enumerate}[label=(\arabic*)]
    \item ¿Son completas las matemáticas? Esto es, ¿puede probarse o no cada sentencia matemática?
    \item ¿Son consistentes las matemáticas? Esto es, ¿no es posible probar simultáneamente una afirmación y su negación?
    \item ¿Son decidibles las matemáticas? Esto es, ¿existe un método automático que pueda aplicarse a cualquier afirmación matemática, y que determine si es cierta? Este problema se conoce como \emph{Entscheidungsproblem}.\index{Entscheidungsproblem}
\end{enumerate}

Es necesario puntualizar que la formulación de Hilbert de este problema es relativa a la aritmética. Es claro que, para tal caso, el \cref{teo:incompletitud-godel} nos garantiza la existencia de fórmulas no demostrables, respondiendo negativamente a (1) \cite{Hellman1981}. Por otra parte, en la \cref{prop:es-teorema-consistente-decidible} probamos que \textbf{Peano}, un sistema de aritmética, es no decidible, probando que la respuesta a la pregunta (3) es negativa. La pregunta (2) queda fuera del objeto de este trabajo.

Sin embargo, es interesante ver estas preguntas sin necesidad de limitarnos a la aritmética. Hemos introducido los sistemas lógicos justamente para poder formalizar el concepto de verdad y demostración, que son el alma de las matemáticas.

Podremos encontrar respuestas análogas para sistemas formales recursivamente enumerables y de tipo Turing. En tal caso, el \cref{teo:incompletitud-sintactica} responde negativamente a la pregunta (1), mientras que una generalización de la \cref{prop:es-teorema-peano-no-decidible} (similar a la que hicimos en el \cref{teo:incompletitud-semantica} a partir del \cref{teo:incompletitud-peano}) nos permite responder negativamente a la pregunta (3).

%Tenemos respuestas a las preguntas (1) y (3) para unas hipótesis concretas pero, ¿qué pasa en el caso de sistemas no recursivamente enumerables y que no son de tipo Turing? Algunos matemáticos no están convencidos de que los resultados de Gödel resuelvan definitivamente el segundo problema de Hilbert, pues los teoremas de Gödel no se aplican en otros sistemas de lógica.

\subsection{Relación con la mente humana y con la tesis de Church-Turing}\label{subsec:mente-humana-tesis-church-turing}

Algunos autores entran en conclusiones mucho más epistemológicas, y que conciernen a la mente humana. Roger Penrose y John Lucas \cite{Penrose1989, Lucas1961} argumentan que el papel de la conciencia es no algorítmico cuando formamos juicios matemáticos, pudiendo generalizarlo al papel de la conciencia humana.

Para ello, los autores comentan que en cualquier sistema formal existen proposiciones indemostrables, esto es, para las que su algoritmo no puede proporcionar una respuesta. Si el funcionamiento de la mente humana fuese completamente algorítmico, entonces el algoritmo (o sistema formal) que utilice para formar sus juicios no podría tratar tal proposición. Sin embargo, la mente humana puede deducir que tal proposición es verdadera. Esta contradicción nos puede indicar que la mente humana, en cuanto al razonamiento matemático, no estaba utilizando un algoritmo en primer lugar.

Douglas Hofstadter \cite{Hofstadter1999, Hofstadter2007} presenta los teoremas de Gödel como un ejemplo de ``extraño bucle'', una estructura jerárquica y autoreferencial que existe en un sistema formal. Argumenta que esta es la misma estructura que origina la conciencia, el sentido del ``yo'', en la mente humana: la autoreferencia proviene de la forma en que el cerebro abstrae y categoriza los estímulos en ``símbolos'', o grupos de neuronas, que corresponden a conceptos, en lo que es de facto un sistema formal, que eventualmente forma símbolos que corresponden a la propia entidad que realiza la percepción. 

La diferencia entre las matemáticas y la conciencia reside, según Hofstdadter, en el sentido en el que se realiza el razonamiento: si es ``hacia arriba'' o ``hacia abajo''., desde los axiomas hacia las proposiciones y teoremas más complejas, o ``hacia abajo'', en el sentido contrario.

Las matemáticas funcionan hacia arriba: para derivar si una fórmula es un teorema debemos partir de los axiomas e ir razonando hasta llegar a ella. La mente humana, según el físico, funciona a hacia abajo, desde un nivel más alto --de deseos, conceptos, personalidades, ideas y pensamientos-- hasta el nivel más bajo de interacciones entre neuronas e incluso partículas fundamentales. Los seres humanos estamos hechos para percibir ``cosas grandes'' en lugar de ``cosas pequeñas'', incluso a pesar de que en el dominio de lo pequeño parece residir el motor de la realidad.

Ambos argumentos están íntimamente relacionados con la tesis de Church-Turing (véase la \cref{sec:church-turing}). Si la mente humana es una máquina de Turing, todos los procedimientos de nuestro cerebro no serían más que algoritmos, y mediante los teoremas de Gödel, nuestra forma de razonar sería limitada --¡existen afirmaciones que no podemos demostrar!--. Esto puede inducir en conclusiones más fantásticas, hablando algunos autores de que los teoremas de Gödel, aunque no permitan probarlo, son una fuerte evidencia de la existencia de dios.\footnote{No existe consenso respecto si esto es cierto, pero es sencillo argumentar los fallos de este razonamiento (véase \cite{StanfordGodel2022}).} \cite{StanfordGodel2022}

%Cerebros y máquinas

%- Autores incluyendo J. R. Lucas y Roger Penrose debaten que los teoremas tienen implicaciones sobre la inteligencia humana. El debate se centra sobre si la mente humana es equivalente a una máquina de Turing o (mediante la tesis de Church-Turing) a una máquina finita. Si la mente lo es, y los teoremas de Gödel se aplican a esta, los teoremas de incompletitud se aplicarían a ella.

%Hilary Putnam 

%Avi Widgerson

%Douglas Hofstadter 

%- 

%Consecuencias severas para el logicismo de Gottlob Frege y Bertrand Russell, que pretendían definir los números naturales en términos de lógica

%\begin{tabla}
%\begin{table}[H]
%\centering
%\begin{adjustwidth}{-12pt}{}
%\begin{tabular}{@{}ll@{}}
%\toprule
%Máquinas de Turing & Numeración Gödel \\
%(\cref{teo:incompletitud-semantica}) & (\cref{teo:incompletitud-godel}) \\ \midrule
%para sistemas lógicos (semántica) & para sistemas formales (sintaxis) \\[4pt]
%asume que el sistema lógico es sólido & asume que el sistema formal es $\omega$-consistente \\[4pt]
%usa sistemas lógicos capaces de probar & usa sistemas formales capaces de probar \\
%afirmaciones del tipo ``$P$ para con &  ciertos teoremas de aritmética elemental \\
%entrada vacía'' \\ \bottomrule
%\end{tabular}
%\end{adjustwidth}
%\end{table}
%\vspace{-8pt}
%\caption{Diferencias en los teoremas de incompletitud mostrados en este capítulo}
%\label{tab:diferencias-turing-godel}
%\end{tabla}

%\section{El Segundo Teorema de Incompletitud}\label{sec:teo-incompletitud-2}

%El Segundo Teorema de Incompletitud, en su versión sintáctica, puede obtenerse inmediatamente a partir del Primer Teorema de Incompletitud.

%\begin{lema}\label{lema:red-prob}
%Sea la siguiente fórmula un teorema:
%$$\phi_1\Rightarrow\phi_2$$

%Entonces, si $\phi_2$ no es un teorema, $\phi_1$ tampoco lo es.
%\end{lema}
%\begin{proof}
%Este resultado se obtiene de la implicación inmediata siguiente: si $\phi_1$ es un teorema, entonces $\phi_2$ ha de serlo también, pues se deduce inmediatamente de $\phi_1\Rightarrow\phi_2$ mediante \emph{modus ponens}. Concluimos la prueba por el contrarrecíproco.
%\end{proof}

%\begin{teorema}[Versión semántica del Segundo Teorema de Incompletitud]
%Sea $\mathcal{G}$ un sistema lógico que verifique las siguientes condiciones:
%\begin{itemize}
%    \item[(A2)] $\mathcal{G}$ es recursivamente axiomatizable.
%    \item[(A3)] $\mathcal{G}$ puede expresar afirmaciones sobre máquinas de Turing o, equivalentemente, sobre programas; en concreto, dada una máquina de Turing $M$, puede expresar que ``$M$ para con entrada vacía'' o, equivalentemente, dado un programa $P$, puede expresar que ``$P$ para con entrada vacía''.
%\end{itemize}
%Entonces, la fórmula correspondiente a ``$\mathcal{G}$ es sólido'' no es un teorema en $\mathcal{G}$.
%\end{teorema}
%\begin{proof}
%El \cref{teo:incompletitud-semantica} nos dice que, suponiendo (A2) y (A3), en caso de que $\mathcal{G}$ fuese sólido, entonces tendríamos una fórmula $\phi$ que (1) es verdadera pero (2) no es un teorema.

%Como hemos logrado probar (1) en $\mathcal{G}$, tenemos el siguiente teorema:
%$$``\mathcal{G}\text{ es sólido}''\Rightarrow\phi$$
%Ahora bien, aplicando (2) y el \cref{lema:red-prob}, tenemos que ``$\mathcal{G}$ es sólido'' no es un teorema.
%\end{proof}

%\emph{\textbf{Nota: 6.1} para probar la versión sintáctica (que es el verdadero segundo teorema de incompletitud) sería necesario encontrar una versión constructiva del \cref{teo:incompletitud-sintactica}, que nos diese una fórmula $\phi$ que fuese verdadera de modo que ni $\phi$ ni $\neg\phi$ sean teoremas. Pero esto es un tanto extraño porque si el teorema es sintáctico, ¿qué quiere decir que una fórmula sea verdadera?}

%\section{Consecuencias}\label{sec:teo-incompletitud-consecuencias}


\endinput

%\input{capitulos/parte03/capitulo03}
%\input{capitulos/parte03/capitulo04}

\setpartpreamble[c][0.75\linewidth]{%
	\bigskip % Deja un espacio vertical en la parte superior
  Concluimos el trabajo revisando los objetivos iniciales y asentando las posibles vías de investigación y desarrollo en un futuro
}
\cleardoublepage\part{Conclusión}\label{part:conclusion}
% !TeX root = ../libro.tex
% !TeX encoding = utf8

\setchapterpreamble[c][0.75\linewidth]{%
	\sffamily
  \emph{El objetivo de la computación es el conocimiento, no los números.}
 \begin{flushright} — Richard Hamming, \cite{Hamming1986} \end{flushright}
% \\[8pt]
	\par\bigskip
}
%\vspace{28pt}

\chapter{Motivación}\label{ch:motivacion}

Los teoremas de incompletitud de Gödel son uno de los resultados más populares de las matemáticas y la filosofía, y han sido objeto de fascinación y estudio durante décadas. Estos teoremas revolucionaron los mismos fundamentos de ambas disciplinas. En este trabajo, abordaremos el primero de estos teoremas desde una perspectiva computacional. Esto es así debido a una serie de experiencias y observaciones que despertaron en mí un gran interés y que revelaron una oportunidad para abordar este tema desde una perspectiva computacional pero, sobre todo, matemáticamente rigurosa.

Este interés surge de las clases de Modelos de Computación impartidas por el tutor de este trabajo, en las que mencionó de forma breve este resultado. Tras ello sentí mucha curiosidad por profundizar en el mundo de la incompletitud, hasta llegar a un excepcional libro de Douglas Hofstadter: \emph{``Gödel, Escher, Bach: An Eternal Golden Braid''} \cite{Hofstadter1999}. Esta obra maestra interdisciplinaria combina elementos de matemáticas, música y arte, y me incitó a profundizar más en estos teoremas.

La mayor parte de los libros que tratan este tema lo hacen tal y como lo hizo Gödel: desde una perspectiva aritmética, usando la numeración de Gödel en las demostraciones, o variantes. Sin embargo, en múltiples cursos de teoría de la computación se exponen demostraciones haciendo uso de algoritmos y máquinas de Turing.

Para mi sorpresa, a pesar de que esta perspectiva se estudia extensivamente y hay una gran literatura al respecto, muchas de las referencias que encontraba eran bien poco rigurosas o bien poco consistentes en la nomenclatura y en la forma de probar los teoremas. En concreto, muchas de las fuentes no distinguen de forma precisa entre los conceptos de \emph{demostración} y \emph{verdad}.

Es aquí donde vi la oportunidad de hacer un trabajo que intentase proporcionar una perspectiva computacional a los teoremas de incompletitud de Gödel, sin perder rigor y concreción por el camino. Este enfoque permitirá una comprensión más precisa y estructurada de los resultados más relevantes al respecto, así como su relación con los modelos de computación y la teoría de la computabilidad.

\pagebreak

\section*{Objetivos}

Los objetivos de este trabajo son:

\subsection*{Objetivos teóricos}

\begin{enumerate}[label={[}T\arabic*{]}]
    \item Probar el Primer Teorema de Incompletitud de Gödel desde una perspectiva computacional.
    \item Usar para ello programas en Python en lugar de descripciones de máquinas de Turing, probando que esto es posible y riguroso.
    \item Exponer las consecuencias del teorema de incompletitud en las disciplinas matemática y filosófica.
    \item Contextualizar el teorema de incompletitud en el ámbito histórico.
    \item Encontrar la relación entre el teorema de incompletitud con la tesis de Church-Turing.
    \item Precisar conceptos lógicos relacionados con la teoría de la demostración. En concreto, distinguir entre las nociones de \emph{demostración} y \emph{verdad}, y estudiar las propiedades de los sistemas de demostración.
\end{enumerate}

\subsection*{Objetivos prácticos}

\begin{enumerate}[label={[}P\arabic*{]}]
    \item Crear contenido que pueda servir de herramienta de apoyo a la docencia.
    \item Crear una librería de programas en Python ejecutables por el usuario que permitan guiar las demostraciones que aparecen en este trabajo.
\end{enumerate}

\endinput

% !TeX root = ../libro.tex
% !TeX encoding = utf8

\setchapterpreamble[c][0.75\linewidth]{%
	\sffamily
  \emph{Si una ``religión'' se define como un sistema de ideas que contiene afirmaciones indemostrables, entonces Gödel nos enseñó que las matemáticas no sólo son una religión, sino que son la única religión que puede demostrar que lo es.}
 \begin{flushright} — John D. Barrow, \cite{Barrow2011} \end{flushright}
% \\[8pt]
	\par\bigskip
}
%\vspace{28pt}

\chapter{El Primer Teorema de Incompletitud}\label{ch:teorema-incompletitud}

Los matemáticos de principios del siglo \textsc{xx}, tal y como Leibniz, Boole y Hilbert, imaginaron que sería posible determinar la veracidad de fórmulas matemáticas de forma automática. Nosotros ahora nos preguntamos lo mismo: hemos trabajado con sistemas lógicos sencillos que son decidibles, es decir, cuyas fórmulas pueden ser comprobadas verdaderas mediante programas. En el caso de sistemas lógicos de más complejidad, como la aritmética de Peano, esto no es posible. Esto se debe a que la existencia de problemas computacionales indecidibles conduce directamente a la existencia de afirmaciones matemáticas no decidibles.

%En la \cref{sec:primera-aproximacion}, probaremos la incompletitud de \textbf{Peano}. Llamaremos a este resultado \emph{Primer Teorema de Incompletitud}. A continuación, generalizaremos este teorema para otros sistemas lógicos. Finalmente, en la \cref{sec:diferencias-godel}, veremos el enunciado original del teorema y lo compararemos con el resultado obtenido.

\section{Una primera aproximación}\label{sec:primera-aproximacion}

%\begin{teorema}[Versión no constructiva del Primer Teorema de Incompletitud]\label{teo:inco1-no-constr}
%La aritmética de Peano es incompleta: existen fórmulas verdaderas de \textbf{Peano} que no son teoremas.
%\end{teorema}
%\begin{proof}
%Asumiremos que \textbf{Peano} es completo y entraremos en contradicción. En concreto, demostraremos que si la aritmética de Peano es completa, podremos decidir \textbf{ParadaEnVacío}, lo cual contradice la \cref{prop:parada-en-vacio-no-decidible}.

%Sea $P$ un programa cualquiera. Lo usaremos como entrada al \cref{lst:parada-en-vacio-desde-peano-completo}. En la \cref{line:parada-en-vacio-desde-peano-completo-import1} importamos la función \texttt{es\_prueba\_peano}, que dada una fórmula y una demostración en \textbf{Peano} comprueba si la demostración prueba que la fórmula es un teorema. Claramente, esta función no es un oráculo. De hecho, su implementación es sencilla: basta probar si cada una de las fórmulas de la demostración es un axioma o se sigue de las fórmulas anteriores a través de las reglas de inferencia del sistema, comprobando que la última fórmula coincida con la fórmula que queremos probar.

%A continuación, importamos en la \cref{line:parada-en-vacio-desde-peano-completo-import2} la función \texttt{parada\_a\_peano} cuya existencia comentamos tras la \cref{prop:parada-a-peano}, que pasa la afirmación ``$P$ para con entrada vacía'' a una fórmula en \textbf{Peano}.

%Usamos esta función en la \cref{line:parada-en-vacio-desde-peano-completo-para-en-peano} para generar la fórmula mencionada anteriormente. En la \cref{line:parada-en-vacio-desde-peano-completo-no-para-en-peano} creamos una nueva fórmula negando la anterior, por tanto \texttt{no\_para\_en\_peano} contiene la fórmula de \textbf{Peano} correspondiente a la afirmación ``$P$ no para con entrada vacía''.

%Una de estas afirmaciones ha de ser verdadera. Como asumimos que \textbf{Peano} es completo, una de estas fórmulas tiene una prueba (es un teorema) en \textbf{Peano}. El bucle de la \cref{line:parada-en-vacio-desde-peano-completo-while} va iterando a través de todos los \emph{strings} posibles, comprobando cada uno para saber si es una prueba de \texttt{para\_en\_peano} o de \texttt{no\_para\_en\_peano}. Ya que uno de ellos debe tener una prueba, el programa terminará, decidiendo si $P$ para con entrada vacía o no. El \cref{lst:parada-en-vacio-desde-peano-completo} resuelve \textsc{ParadaEnVacío}, lo que entra en contradicción con la \cref{prop:parada-en-vacio-no-decidible}.
%\end{proof}

%\begin{lstlisting}[language=Python, caption=\lstinline{parada_en_vacio_desde_peano_completo.py},label={lst:parada-en-vacio-desde-peano-completo}]
%import utilidades

%from es_prueba_peano import es_prueba_peano # NO es un oráculo|\label{line:parada-en-vacio-desde-peano-completo-import1}|
%from parada_a_peano import parada_a_peano # NO es un oráculo|\label{line:parada-en-vacio-desde-peano-completo-import2}|

%def parada_en_vacio_desde_peano_completo(programa):|\label{line:parada-en-vacio-desde-peano-completo-main}|
%    para_en_peano = parada_a_peano(programa)|\label{line:parada-en-vacio-desde-peano-completo-para-en-peano}|
%    no_para_en_peano = 'NOT (' + para_en_peano + ')'|\label{line:parada-en-vacio-desde-peano-completo-no-para-en-peano}|
%    prueba = ''|\label{line:parada-en-vacio-desde-peano-completo-prueba}|
    
%    while True:|\label{line:parada-en-vacio-desde-peano-completo-while}|
%        if es_teorema_peano(prueba, para_en_peano) == 'sí':|\label{line:parada-en-vacio-desde-peano-completo-if1}|
%            return 'sí'
%        if es_teorema_peano(prueba, no_para_en_peano) == 'sí':|\label{line:parada-en-vacio-desde-peano-completo-if2}|
%            return 'no'
%        prueba = utilidades.siguiente_string(prueba)|\label{line:parada-en-vacio-desde-peano-completo-else}|
%\end{lstlisting}


Antes de proceder la demostración de una primera versión del Primer Teorema de Incompletitud, discutiremos de forma detallada el \cref{lst:godel}. En la \cref{line:godel-peano-import1} importamos la función \texttt{es\_teorema\_peano}, cuya existencia probamos en la \cref{prop:es-teorema-peano-semidecidible}. Recuerda que, como los axiomas de \textbf{Peano} son recursivamente enumerables (es decir, \textbf{Peano} es recursivamente axiomatizable), podemos escribir el programa \texttt{es\_teorema\_peano(fórmula)} que devuelve \palabra{sí} si \texttt{fórmula} es un teorema en \textbf{Peano}, y que no para en caso contrario (ver \cref{prop:es-teorema-peano-semidecidible}).
\vspace{8pt}
\begin{lstlisting}[language=Python, caption=\lstinline{godel_peano.py},label={lst:godel-peano}]
import utilidades

from es_teorema_peano import es_teorema_peano # NO es un oráculo|\label{line:godel-peano-import1}|
from parada_a_peano import parada_a_peano # NO es un oráculo|\label{line:godel-peano-import2}|

def godel_peano(entrada):|\label{line:godel-peano-main}|
    programa_godel = utilidades.leer('godel\_peano.py')|\label{line:godel-peano-programa}|
    para_en_peano = parada_a_peano(programa_godel)|\label{line:godel-peano-formula}|
    no_para_en_peano = 'NOT (' + para_en_peano + ')'|\label{line:godel-peano-formulaneg}|

    if es_teorema_peano(no_para_en_peano) == 'sí':|\label{line:godel-peano-es-teorema}|
        return 'para'|\label{line:godel-peano-para}|
    else:
        utilidades.ciclar()|\label{line:godel-peano-cicla}|
\end{lstlisting}

La \cref{line:godel-peano-import2} importa la función \texttt{parada\_a\_peano}, cuya existencia ya discutimos tras la \cref{prop:parada-a-peano}. Esta función traduce, dado un programa $P$, la afirmación ``$P$ para con entrada vacía'' en una fórmula de \textbf{Peano}.

En la \cref{line:godel-peano-programa} guardamos el propio programa \texttt{godel\_peano.py} en la variable \texttt{programa\_godel}. En la \cref{line:godel-peano-formula} usamos la función \texttt{parada\_a\_peano} recién comentada, y en la \cref{line:godel-peano-formulaneg} negamos la fórmula \texttt{para\_en\_peano}. De este modo:
\begin{adjustwidth}{30pt}{}
    \texttt{no\_para\_en\_peano} significa que ``\texttt{godel\_peano.py} no para con entrada vacía'' \hfill $(\star1)$
\end{adjustwidth}
Ahora es cuando las cosas se ponen interesantes. Si \texttt{no\_para\_en\_peano} es un teorema, la función \texttt{es\_teorema\_peano} devolverá \palabra{sí}, y de este modo el programa verificará la condición de la \cref{line:godel-peano-es-teorema} y devolverá \palabra{para}, haciendo que el programa pare. Si \texttt{no\_para\_en\_peano} no es un teorema, se entrará en un bucle infinito (\texttt{es\_teorema\_peano} no parará, y por lo tanto tampoco \texttt{godel\_peano.py}). Es claro que la \cref{line:godel-peano-cicla} nunca se ejecutará, pero la incluimos para enfatizar este hecho.

Ahora, nos preguntamos: ¿cuál es el comportamiento de \texttt{godel\_peano.py} con entrada vacía?

\begin{proposicion}\label{prop:godel-peano-no-para}
El \cref{lst:godel-peano} no para con entrada vacía.
\end{proposicion}
\begin{proof}
Asumimos que \texttt{godel\_peano.py} para con entrada vacía, y entraremos en contradicción. La única forma de que pare es en la \cref{line:godel-peano-para}, algo que ocurre solamente si:
\begin{adjustwidth}{30pt}{}
    \texttt{es\_teorema\_peano(no\_para\_en\_peano) == \palabra{sí}} \hfill $(\star2)$
\end{adjustwidth}
$(\star2)$ es equivalente a decir que \texttt{no\_para\_en\_peano} es un teorema. Por la solidez de \textbf{Peano}, \texttt{no\_para\_en\_peano} es verdadero.

Pero el hecho de que \texttt{no\_para\_en\_peano} sea verdadero, junto con $(\star1)$, nos dice que \texttt{godel\_peano.py} no para con entrada vacía, encontrando la contradicción esperada.
\end{proof}

Ahora estamos en condiciones de probar el resultado principal de este trabajo.

\begin{teorema}[Incompletitud de la aritmética de Peano]\label{teo:incompletitud-peano}
\index{primer teorema de incompletitud}\index{teorema de incompletitud}La aritmética de Peano es incompleta: existen fórmulas verdaderas de \normalfont{\textbf{Peano}} que no son teoremas. Concretamente, la fórmula guardada en la variable \texttt{no\_para\_en\_peano} en la \cref{line:godel-peano-formulaneg} del \cref{lst:godel} es un ejemplo de una fórmula verdadera pero no demostrable (no es un teorema) de \normalfont{\textbf{Peano}}.
\end{teorema}
\begin{proof}
Debemos mostrar que \texttt{no\_para\_en\_peano} es verdadero y no es un teorema.

Para probar que es verdadero, por la \cref{prop:godel-peano-no-para}, sabemos que \texttt{godel\_peano.py} no para con entrada vacía. Recordando $(\star1)$, inmediatamente obtenemos que \texttt{no\_para\_en\_peano} es verdadero.

Para probar que no es un teorema, de nuevo usamos el hecho de que \texttt{godel\_peano.py} no para con entrada vacía. Esto nos indica que, en la \cref{line:godel-peano-es-teorema}, la llamada a la función
\begin{adjustwidth}{30pt}{}
    \texttt{es\_teorema\_peano(no\_para\_en\_peano)}
\end{adjustwidth}
no puede devolver \palabra{sí}. Pero esto inmediatamente implica que \texttt{no\_para\_en\_peano} no es demostrable, lo que completa la prueba.
\end{proof}

%\section{El Segundo Teorema de Incompletitud}\label{sec:teo-incompletitud-2}

%\begin{teorema}[Segundo Teorema de Incompletitud]
%La consistencia de la aritmética de Peano no puede probarse en la aritmética de Peano. En concreto, la afirmación
%\begin{adjustwidth}{30pt}{}
%    ``\textbf{Peano} es consistente''
%\end{adjustwidth}
%que puede codificarse mediante una fórmula de \textbf{Peano}, no es un teorema.
%\end{teorema}

Acabamos de demostrar que \textbf{Peano} es incompleto, como ya anticipamos en la \cref{tab:sistemas-logicos}. Las consecuencias de este hecho las exploramos en detalle en el \cref{ch:conclusion}. Sin embargo, paremos un momento a analizar qué es lo que acabamos de probar.

\section{Versión semántica}\label{sec:version-semantica}

Para encontrar una fórmula verdadera y no demostrable, hemos necesitado un sistema lógico que sea, en primer lugar, recursivamente axiomatizable. Si no, no hubiesemos podido crear un programa capaz de saber cuándo una fórmula del sistema es un teorema. Por otra parte, hemos asumido que el sistema es sólido, lo cual es algo esperable. Finalmente, hemos exigido que sea capaz de expresar una cierta afirmación de las máquinas de Turing (véase \cref{prop:parada-a-peano}). Para evitar repeticiones innecesarias en esta última hipótesis, llamaremos a un sistema lógico con esta propiedad ``de tipo Turing'', algo que definimos a continuación.

\begin{definicion}[Sistema de tipo Turing]
Sea $\mathcal{S}$ un sistema formal. Diremos que $\mathcal{S}$ es \emph{de tipo Turing}\footnote{Esta definición no se encuentra en la literatura: se ha creado por conveniencia.} si puede expresar afirmaciones sobre máquinas de Turing o, equivalentemente, sobre programas;\footnote{La equivalencia la probamos en el \cref{teo:equivalencia}.} en concreto, dada una máquina de Turing $M$, puede expresar que ``$M$ para con entrada vacía''; o, equivalentemente, dado un programa $P$, puede expresar que ``$P$ para con entrada vacía''.

Análogamente podemos definirlo para el caso de sistemas lógicos.
\end{definicion}

Observemos que, sin embargo, en ningún momento hemos utilizado el sistema \textbf{Peano} en sí. Esto nos indica que podríamos realizar esta demostración con otro sistema que no sea \textbf{Peano}, siempre y cuando verifique tales condiciones. Indicamos este resultado a continuación.

\begin{teorema}[Versión semántica del Primer Teorema de Incompletitud]\label{teo:incompletitud-semantica}
Sea $\mathcal{G}$ un sistema lógico que verifique las siguientes condiciones:
\begin{enumerate}
    \item $\mathcal{G}$ es sólido.
    \item $\mathcal{G}$ es recursivamente axiomatizable.
    \item $\mathcal{G}$ es de tipo Turing.
\end{enumerate}
Entonces, $\mathcal{G}$ es incompleto, es decir, existen fórmulas verdaderas de $\mathcal{G}$ que no son teoremas.
\end{teorema}

Para reforzar este importante resultado, incluimos un resumen del proceso que hemos seguido para \textbf{Peano}, de forma general, en la demostración de este teorema.

\begin{proof}
En el \cref{lst:godel} importamos dos funciones. Por una parte, en la \cref{line:godel-import1} importamos la función \texttt{es\_teorema\_en\_G}, que comprueba si una fórmula es un teorema en $\mathcal{G}$. Esta función parará si la fórmula es un teorema, y se ejecutará indefinidamente en caso de que no lo sea. La existencia de esta función la garantiza (2).

Por otra, en la \cref{line:godel-import2} importamos la función \texttt{parada\_a\_G} que, dado un programa $P$, convierte la afirmación ``$P$ para con entrada vacía'' en una fórmula de $\mathcal{G}$, algo que es posible por (3).

En la línea \cref{line:godel-formulaneg} tenemos que:
\begin{adjustwidth}{30pt}{}
    \texttt{no\_para\_en\_G} significa que \emph{``\texttt{godel.py} no para con entrada vacía''} (en $\mathcal{G}$)  \hfill $(\star3)$
\end{adjustwidth}
Si \texttt{no\_para\_en\_G} es un teorema, la función \texttt{es\_teorema\_en\_G(no\_para\_en\_G)} (\cref{line:godel-es-teorema}) parará y devolverá \palabra{sí}, mientras que la misma función ciclará en caso de que no lo sea.

La pregunta que debemos hacernos es: ¿qué le pasa a \texttt{godel.py} cuando \texttt{entrada} es vacía?

Si asumimos que \texttt{godel.py} para con entrada vacía, la única forma de que lo haga es porque en la \cref{line:godel-es-teorema} es \texttt{es\_teorema\_en\_G(no\_para\_en\_G) == \palabra{sí}}, lo cual solo puede pasar si \texttt{no\_para\_en\_G} es un teorema. Pero, por (1), y dado que $\mathcal{G}$ es sólido, tenemos que \texttt{no\_para\_en\_G} es verdadero.

Pero si \texttt{no\_para\_en\_G} es verdadero, por $(\star3)$, esto querría decir que \texttt{godel.py} no para con entrada vacía, entrando en contradicción, por lo que debe ser:
\begin{adjustwidth}{30pt}{}
    \texttt{godel.py} no para con entrada vacía  \hfill $(\star4)$
\end{adjustwidth}
Comprobamos que \texttt{no\_para\_en\_G} es verdadero pero no es un teorema. En efecto, es verdadero por $(\star4)$. Sin embargo, no es un teorema. Para ello, usamos de nuevo $(\star4)$. El hecho de que \texttt{godel.py} no pare indica que es imposible que la condición de la \cref{line:godel-es-teorema}, \linebreak\texttt{es\_teorema\_en\_G(no\_para\_en\_G) == \palabra{sí}}, se verifique -- en otras palabras, \texttt{no\_para\_en\_G} no es un teorema.
\end{proof}

%\vspace{8pt}
\begin{lstlisting}[language=Python, caption=\lstinline{godel.py},label={lst:godel}]
import utilidades

from es_teorema_G import es_teorema_G # NO es un oráculo|\label{line:godel-import1}|
from parada_a_G import parada_a_G # NO es un oráculo|\label{line:godel-import2}|

def incompleto_semantico(entrada):|\label{line:godel-main}|
    programa = utilidades.leer('incompleto_semantico.py')|\label{line:godel-programa}|
    para_en_G = parada_a_G(programa)|\label{line:godel-formula}|
    no_para_en_G = 'NOT (' + para_en_G + ')'|\label{line:godel-formulaneg}|

    if es_teorema_G(no_para_en_G) == 'sí':|\label{line:godel-es-teorema}|
        return 'para'|\label{line:godel-para}|
    else:
        utilidades.ciclar()|\label{line:godel-cicla}|
\end{lstlisting}


Nuestro resultado nos dice que, si tenemos un sistema lógico que sea recursivamente axiomatizable y de tipo Turing, no podemos esperar que pueda ser sólido y completo al mismo tiempo. Es decir, si todos los teoremas del sistema son verdaderos, siempre podremos encontrar fórmulas del sistema que sean verdaderas pero no sean teoremas.

Podemos aplicar el \cref{teo:incompletitud-semantica} a otros sistemas lógicos, siempre que verifiquen las hipótesis. Uno de ellos es, por ejemplo, el sistema de axiomas de Zermelo-Fraenkl.

\section{Versión sintáctica}\label{sec:version-sintactica}

% WIP relajar condición de solidez?
El resultado que acabamos de probar es semántico. Sin embargo, es posible probar otro resultado, esta vez sintáctico, y más parecido al demostrado por Gödel.

Para poder probar esta variación, debemos de tener en cuenta una sutileza del problema \textsc{EsTeorema} en el caso de sistemas formales sintácticamente consistentes.

\begin{proposicion}\label{prop:es-teorema-consistente-decidible}
Sea $\mathcal{S}$ un sistema formal recursivamente axiomatizable y sintácticamente consistente. Entonces, $\textsc{EsTeorema}_\mathcal{S}$ es decidible.
\end{proposicion}
\begin{proof}
Veamos que el \cref{lst:es-teorema-s} decide $\textsc{EsTeorema}_\mathcal{S}$. En efecto, vemos en primer lugar que en la \cref{line:es-teorema-s-import} importamos la función \texttt{es\_prueba\_S}, la cual es computable por ser $\mathcal{S}$ recursivamente axiomatizable.

A continuación, en la \cref{line:es-teorema-s-negacion} almacenamos la negación de \texttt{formula} en \texttt{neg\_formula}. En la \cref{line:es-teorema-s-demostracion} declaramos la demostración.

Por ser $\mathcal{S}$ consistente, bien \texttt{formula} bien \texttt{neg\_formula} tiene una demostración. En el bucle de la \cref{line:es-teorema-s-while}, vamos comprobando si \texttt{demostracion} es una demostración de \texttt{formula} o de \texttt{neg\_formula}. En el primer caso, devolvemos \palabra{sí}, mientras que en el segundo, devolvemos \palabra{no}. En caso de que \texttt{demostracion} no sea una demostración de ninguna de ambas fórmulas, procedemos en la \cref{line:es-teorema-s-siguiente-string} a enumerar la siguiente demostración en orden \emph{shortlex} (ver \cref{prop:es-teorema-peano-semidecidible}).

Por la consistencia de $\mathcal{S}$, debe existir una demostración bien de \texttt{formula} o de \texttt{neg\_formula}, por lo que el programa debe parar al llegar a ser \texttt{demostracion} tal demostración.
\end{proof}

\begin{lstlisting}[language=Python, caption=\lstinline{es_teorema_S.py},label={lst:es-teorema-s}]
import utilidades
from es_prueba_S import es_prueba_S # NO es un oráculo|\label{line:es-teorema-s-import}|

def es_teorema_S(formula):|\label{line:es-teorema-s-main}|
    negacion_formula = 'NOT (' + formula + ')'|\label{line:es-teorema-s-negacion}|
    demostracion = ''|\label{line:es-teorema-s-demostracion}|

    while True:|\label{line:es-teorema-s-while}|
        if es_prueba_S(demostracion, formula) == 'sí':|\label{line:es-teorema-s-es-prueba}|
            return 'sí'
        if es_prueba_S(demostracion, negacion_formula) == 'sí':
            return 'no'
        demostracion = utilidades.siguiente_string(demostracion)|\label{line:es-teorema-s-siguiente-string}|
\end{lstlisting}

Ahora estamos en condiciones de probar el resultado sintáctico, que dependerá de la no decidibilidad de \textsc{AdivinaConsistente} (\cref{prob:adivina-consistente}). \cite{Aaronson2017}

\begin{teorema}[Versión sintáctica del Primer Teorema de Incompletitud]\label{teo:incompletitud-sintactica}
Sea $\mathcal{S}$ un sistema formal que verifique las siguientes condiciones:
\begin{enumerate}
    \item $\mathcal{S}$ es sintácticamente consistente.
    \item $\mathcal{S}$ es recursivamente axiomatizable.
    \item $\mathcal{S}$ es de tipo Turing.
\end{enumerate}
Entonces, $\mathcal{S}$ es sintácticamente incompleto, es decir, existen fórmulas $\phi$ de $\mathcal{S}$ tal que ni $\phi$ ni $\neg\phi$ son teoremas.
\end{teorema}
\begin{proof}
Vamos a suponer que $\mathcal{S}$ es sintácticamente completo y veremos que, en tal caso, podríamos decidir \textsc{AdivinaConsistente} (\cref{prob:adivina-consistente}), lo cual entraría en contradicción con la \cref{prop:adivina-consistente-no-decidible}.

Comencemos comentando las funciones importadas. En la \cref{line:adivina-consistente-import-1} importamos la función \texttt{oarada\_en\_S}, que traduce para un programa $P$ la afirmación ``$P$ para con entrada vacía'' en una fórmula de $\mathcal{}S$, y cuya existencia garantiza la hipótesis (3). En la \cref{line:adivina-consistente-import-2} importamos la función \texttt{es\_teorema\_S}, que decide el problema $\textsc{EsTeorema}_\mathcal{S}$, y que, por las hipótesis (1) y (2), siempre para (ver \cref{prop:es-teorema-consistente-decidible}) y devuelve \palabra{sí} si la fórmula de entrada es un teorema en $\mathcal{S}$ y \palabra{no} en caso contrario.

En la \cref{line:adivina-consistente-programa-para} usamos la función \texttt{parada\_a\_S} para almacenar en \texttt{programa\_para} la fórmula de $\mathcal{S}$ equivalente a ``\texttt{programa} para con entrada vacía''. En la \cref{line:adivina-consistente-es-teorema}, asignamos a la variable \texttt{es\_teorema} uno de dos valores posibles: \palabra{sí} o \palabra{no}, como ya comentamos anteriormente.

Llegado a la \cref{line:adivina-consistente-si}, implementamos la siguiente lógica:
\begin{itemize}
    \item Si \texttt{programa} acepta con entrada vacía (para y devuelve \palabra{sí}), la condición de la \cref{line:adivina-consistente-si} se verificará, así como la de la \cref{line:adivina-consistente-si-si}, y \texttt{adivina\_consistente.py} devolverá \palabra{sí}.
    \item Si \texttt{programa} rechaza con entrada vacía (para y devuelve \palabra{no}, la condición de la \cref{line:adivina-consistente-si} se verificará, mientras que no la de la \cref{line:adivina-consistente-si-si}, y \texttt{adivina\_consistente.py} devolverá \palabra{no}.
    \item Si \texttt{programa} cicla con entrada vacía, la condición de la \cref{line:adivina-consistente-si} no se verificará, y \texttt{adivina\_consistente.py} devolverá \palabra{no}.
\end{itemize}
Dicha lógica hace evidente que el programa decide \textsc{AdivinaConsistente}.
\end{proof}

\begin{lstlisting}[language=Python, caption=\lstinline{adivina_consistente.py},label={lst:adivina-consistente}]
import utilidades
from parada_a_S import parada_a_S # NO es un oráculo|\label{line:adivina-consistente-import-1}|
from es_teorema_S import es_teorema_S # NO es un oráculo|\label{line:adivina-consistente-import-2}|

def adivina_consistente(programa):|\label{line:adivina-consistente-main}|
    programa_para = parada_a_S(programa)|\label{line:adivina-consistente-programa-para}|

    es_teorema = es_teorema_S(programa_para)|\label{line:adivina-consistente-es-teorema}|

    if es_teorema == 'sí':|\label{line:adivina-consistente-si}|
        salida = maquina_universal(programa, '')|\label{line:adivina-consistente-universal}|
        if salida == 'sí':|\label{line:adivina-consistente-si-si}|
            return 'sí'
        else:|\label{line:adivina-consistente-si-no}|
            return 'no'
    else:|\label{line:adivina-consistente-no}|
        return 'no'
\end{lstlisting}

Hemos probado dos versiones del Primer Teorema de Incompletitud. Resumimos las diferencias entre ambas a continuación.
\vspace{8pt}
\begin{figure}[H]
$$
\begin{array}{ccc}
    {
        \begin{matrix}
            \text{\textbf{versión semántica}} \\
            \text{\small{(\cref{teo:incompletitud-semantica})}}\\
            \small{\;}
        \end{matrix}
    }
    & \;\;\; &
    {
        \begin{matrix}
            \text{\textbf{versión sintáctica}} \\
            \text{\small{(\cref{teo:incompletitud-sintactica})}}\\
            \small{\;}
        \end{matrix}
    }
    \\
    {
        \begin{matrix}
            \mathcal{G}\text{ sistema lógico}\\
            \text{rec. axiom., de tipo Turing}\\
            \text{y sólido} \\
            \small{\phi \text{ teorema} \Rightarrow \phi\text{ verdadero}}\\
            \tiny{\;}
        \end{matrix}
    }
    &  & 
    {
        \begin{matrix}
            \mathcal{S}\text{ sistema formal}\\
            \text{rec. axiom., de tipo Turing}\\
            \text{y sintácticamente consistente} \\
            \small{\nexists \phi : \phi\text{ y }\neg\phi \text{ son teoremas}}\\
            \tiny{\;}
        \end{matrix}
    } \\
    \Huge{\Downarrow} & & \Huge{\Downarrow} \\
    {
        \begin{matrix}
            \tiny{\;}\\
            \mathcal{G}\text{ no (semánticamente) completo}\\
            \small{\exists\phi:\phi\text{ es verdadero y no es teorema}}
        \end{matrix}
    }
    & &
    {
        \begin{matrix}
            \tiny{\;}\\
            \mathcal{S}\text{ no sintácticamente completo}\\
            \small{\exists\phi:\text{ni }\phi\text{ ni }\neg\phi\text{ son teoremas}}
        \end{matrix}
    }
    \\
\end{array}
$$
\caption{Comparativa entre las dos versiones del Primer Teorema de Incompletitud}
\vspace{-8pt}
\label{fig:comparativa-incompletitud}
\end{figure}

La demostración original de Gödel es puramente sintáctica, y muy parecida a la probada en el \cref{teo:incompletitud-sintactica}. La enunciamos a continuación.

\begin{teorema}[Primer Teorema de Incompletitud de Gödel-Rosser]\label{teo:incompletitud-godel}
Sea $\mathcal{S}$ un sistema formal que verifique las siguientes condiciones:
\begin{enumerate}
    \item $\mathcal{S}$ es sintácticamente consistente.
    \item $\mathcal{S}$ es recursivamente axiomatizable.
    \item $\mathcal{S}$ contiene una cierta cantidad de aritmética elemental.
\end{enumerate}

Entonces, $\mathcal{S}$ es sintácticamente incompleto, es decir, hay fórmulas $\phi$ de $\mathcal{S}$ de modo que ni $\phi$ ni $\neg\phi$ son teoremas.
\end{teorema}

Tal y como en nuestra versión sintáctica, requerimos que el sistema sea sintácticamente consistente y recursivamente axiomatizable. El teorema original de Gödel exigía como hipótesis una forma más fuerte de consistencia sintáctica llamada \emph{$\omega$-consistencia}. Esta forma más general del teorema se debe al trabajo de John B. Rosser, autor del hoy conocido como \emph{truco de Rosser}.\footnote{Es por ello por lo que lo llamamos teorema de ``Gödel-Rosser'', para distinguirlo del resultado original de Gödel.} \cite{Rosser1936}

La tercera hipótesis exige que el sistema sea capaz de probar ciertas afirmaciones de aritmética elemental. Esto se debe a que Gödel pretendía probar que la teoría de números no era axiomatizable (ver \cref{ch:historia}), por lo que usó en su demostración una forma de numerar los teoremas que hoy en día se conoce como \emph{numeración de Gödel}.


\section{Consecuencias}\label{sec:consecuencias}

El Primer Teorema de Incompletitud tiene profundas implicaciones matemáticas y filosóficas.

\subsection{Relación con el segundo problema de Hilbert}\label{subsec:segundo-problema-hilbert}

Ya introdujimos en el \cref{ch:historia} los célebres problemas de Hilbert, y en concreto, el segundo de ellos, que volveremos a referenciar a continuación:

\begin{enumerate}[label=(\arabic*)]
    \item ¿Son completas las matemáticas? Esto es, ¿puede probarse o no cada sentencia matemática?
    \item ¿Son consistentes las matemáticas? Esto es, ¿no es posible probar simultáneamente una afirmación y su negación?
    \item ¿Son decidibles las matemáticas? Esto es, ¿existe un método automático que pueda aplicarse a cualquier afirmación matemática, y que determine si es cierta? Este problema se conoce como \emph{Entscheidungsproblem}.\index{Entscheidungsproblem}
\end{enumerate}

Es necesario puntualizar que la formulación de Hilbert de este problema es relativa a la aritmética. Es claro que, para tal caso, el \cref{teo:incompletitud-godel} nos garantiza la existencia de fórmulas no demostrables, respondiendo negativamente a (1) \cite{Hellman1981}. Por otra parte, en la \cref{prop:es-teorema-consistente-decidible} probamos que \textbf{Peano}, un sistema de aritmética, es no decidible, probando que la respuesta a la pregunta (3) es negativa. La pregunta (2) queda fuera del objeto de este trabajo.

Sin embargo, es interesante ver estas preguntas sin necesidad de limitarnos a la aritmética. Hemos introducido los sistemas lógicos justamente para poder formalizar el concepto de verdad y demostración, que son el alma de las matemáticas.

Podremos encontrar respuestas análogas para sistemas formales recursivamente enumerables y de tipo Turing. En tal caso, el \cref{teo:incompletitud-sintactica} responde negativamente a la pregunta (1), mientras que una generalización de la \cref{prop:es-teorema-peano-no-decidible} (similar a la que hicimos en el \cref{teo:incompletitud-semantica} a partir del \cref{teo:incompletitud-peano}) nos permite responder negativamente a la pregunta (3).

%Tenemos respuestas a las preguntas (1) y (3) para unas hipótesis concretas pero, ¿qué pasa en el caso de sistemas no recursivamente enumerables y que no son de tipo Turing? Algunos matemáticos no están convencidos de que los resultados de Gödel resuelvan definitivamente el segundo problema de Hilbert, pues los teoremas de Gödel no se aplican en otros sistemas de lógica.

\subsection{Relación con la mente humana y con la tesis de Church-Turing}\label{subsec:mente-humana-tesis-church-turing}

Algunos autores entran en conclusiones mucho más epistemológicas, y que conciernen a la mente humana. Roger Penrose y John Lucas \cite{Penrose1989, Lucas1961} argumentan que el papel de la conciencia es no algorítmico cuando formamos juicios matemáticos, pudiendo generalizarlo al papel de la conciencia humana.

Para ello, los autores comentan que en cualquier sistema formal existen proposiciones indemostrables, esto es, para las que su algoritmo no puede proporcionar una respuesta. Si el funcionamiento de la mente humana fuese completamente algorítmico, entonces el algoritmo (o sistema formal) que utilice para formar sus juicios no podría tratar tal proposición. Sin embargo, la mente humana puede deducir que tal proposición es verdadera. Esta contradicción nos puede indicar que la mente humana, en cuanto al razonamiento matemático, no estaba utilizando un algoritmo en primer lugar.

Douglas Hofstadter \cite{Hofstadter1999, Hofstadter2007} presenta los teoremas de Gödel como un ejemplo de ``extraño bucle'', una estructura jerárquica y autoreferencial que existe en un sistema formal. Argumenta que esta es la misma estructura que origina la conciencia, el sentido del ``yo'', en la mente humana: la autoreferencia proviene de la forma en que el cerebro abstrae y categoriza los estímulos en ``símbolos'', o grupos de neuronas, que corresponden a conceptos, en lo que es de facto un sistema formal, que eventualmente forma símbolos que corresponden a la propia entidad que realiza la percepción. 

La diferencia entre las matemáticas y la conciencia reside, según Hofstdadter, en el sentido en el que se realiza el razonamiento: si es ``hacia arriba'' o ``hacia abajo''., desde los axiomas hacia las proposiciones y teoremas más complejas, o ``hacia abajo'', en el sentido contrario.

Las matemáticas funcionan hacia arriba: para derivar si una fórmula es un teorema debemos partir de los axiomas e ir razonando hasta llegar a ella. La mente humana, según el físico, funciona a hacia abajo, desde un nivel más alto --de deseos, conceptos, personalidades, ideas y pensamientos-- hasta el nivel más bajo de interacciones entre neuronas e incluso partículas fundamentales. Los seres humanos estamos hechos para percibir ``cosas grandes'' en lugar de ``cosas pequeñas'', incluso a pesar de que en el dominio de lo pequeño parece residir el motor de la realidad.

Ambos argumentos están íntimamente relacionados con la tesis de Church-Turing (véase la \cref{sec:church-turing}). Si la mente humana es una máquina de Turing, todos los procedimientos de nuestro cerebro no serían más que algoritmos, y mediante los teoremas de Gödel, nuestra forma de razonar sería limitada --¡existen afirmaciones que no podemos demostrar!--. Esto puede inducir en conclusiones más fantásticas, hablando algunos autores de que los teoremas de Gödel, aunque no permitan probarlo, son una fuerte evidencia de la existencia de dios.\footnote{No existe consenso respecto si esto es cierto, pero es sencillo argumentar los fallos de este razonamiento (véase \cite{StanfordGodel2022}).} \cite{StanfordGodel2022}

%Cerebros y máquinas

%- Autores incluyendo J. R. Lucas y Roger Penrose debaten que los teoremas tienen implicaciones sobre la inteligencia humana. El debate se centra sobre si la mente humana es equivalente a una máquina de Turing o (mediante la tesis de Church-Turing) a una máquina finita. Si la mente lo es, y los teoremas de Gödel se aplican a esta, los teoremas de incompletitud se aplicarían a ella.

%Hilary Putnam 

%Avi Widgerson

%Douglas Hofstadter 

%- 

%Consecuencias severas para el logicismo de Gottlob Frege y Bertrand Russell, que pretendían definir los números naturales en términos de lógica

%\begin{tabla}
%\begin{table}[H]
%\centering
%\begin{adjustwidth}{-12pt}{}
%\begin{tabular}{@{}ll@{}}
%\toprule
%Máquinas de Turing & Numeración Gödel \\
%(\cref{teo:incompletitud-semantica}) & (\cref{teo:incompletitud-godel}) \\ \midrule
%para sistemas lógicos (semántica) & para sistemas formales (sintaxis) \\[4pt]
%asume que el sistema lógico es sólido & asume que el sistema formal es $\omega$-consistente \\[4pt]
%usa sistemas lógicos capaces de probar & usa sistemas formales capaces de probar \\
%afirmaciones del tipo ``$P$ para con &  ciertos teoremas de aritmética elemental \\
%entrada vacía'' \\ \bottomrule
%\end{tabular}
%\end{adjustwidth}
%\end{table}
%\vspace{-8pt}
%\caption{Diferencias en los teoremas de incompletitud mostrados en este capítulo}
%\label{tab:diferencias-turing-godel}
%\end{tabla}

%\section{El Segundo Teorema de Incompletitud}\label{sec:teo-incompletitud-2}

%El Segundo Teorema de Incompletitud, en su versión sintáctica, puede obtenerse inmediatamente a partir del Primer Teorema de Incompletitud.

%\begin{lema}\label{lema:red-prob}
%Sea la siguiente fórmula un teorema:
%$$\phi_1\Rightarrow\phi_2$$

%Entonces, si $\phi_2$ no es un teorema, $\phi_1$ tampoco lo es.
%\end{lema}
%\begin{proof}
%Este resultado se obtiene de la implicación inmediata siguiente: si $\phi_1$ es un teorema, entonces $\phi_2$ ha de serlo también, pues se deduce inmediatamente de $\phi_1\Rightarrow\phi_2$ mediante \emph{modus ponens}. Concluimos la prueba por el contrarrecíproco.
%\end{proof}

%\begin{teorema}[Versión semántica del Segundo Teorema de Incompletitud]
%Sea $\mathcal{G}$ un sistema lógico que verifique las siguientes condiciones:
%\begin{itemize}
%    \item[(A2)] $\mathcal{G}$ es recursivamente axiomatizable.
%    \item[(A3)] $\mathcal{G}$ puede expresar afirmaciones sobre máquinas de Turing o, equivalentemente, sobre programas; en concreto, dada una máquina de Turing $M$, puede expresar que ``$M$ para con entrada vacía'' o, equivalentemente, dado un programa $P$, puede expresar que ``$P$ para con entrada vacía''.
%\end{itemize}
%Entonces, la fórmula correspondiente a ``$\mathcal{G}$ es sólido'' no es un teorema en $\mathcal{G}$.
%\end{teorema}
%\begin{proof}
%El \cref{teo:incompletitud-semantica} nos dice que, suponiendo (A2) y (A3), en caso de que $\mathcal{G}$ fuese sólido, entonces tendríamos una fórmula $\phi$ que (1) es verdadera pero (2) no es un teorema.

%Como hemos logrado probar (1) en $\mathcal{G}$, tenemos el siguiente teorema:
%$$``\mathcal{G}\text{ es sólido}''\Rightarrow\phi$$
%Ahora bien, aplicando (2) y el \cref{lema:red-prob}, tenemos que ``$\mathcal{G}$ es sólido'' no es un teorema.
%\end{proof}

%\emph{\textbf{Nota: 6.1} para probar la versión sintáctica (que es el verdadero segundo teorema de incompletitud) sería necesario encontrar una versión constructiva del \cref{teo:incompletitud-sintactica}, que nos diese una fórmula $\phi$ que fuese verdadera de modo que ni $\phi$ ni $\neg\phi$ sean teoremas. Pero esto es un tanto extraño porque si el teorema es sintáctico, ¿qué quiere decir que una fórmula sea verdadera?}

%\section{Consecuencias}\label{sec:teo-incompletitud-consecuencias}


\endinput


% --------------------------------------------------------------------
% APPENDIX: Opcional
% --------------------------------------------------------------------

\appendix % Reinicia la numeración de los capítulos y usa letras para numerarlos
%\pdfbookmark[-1]{Apéndices}{appendix} % Alternativamente podemos agrupar los apéndices con un nuevo \part{Apéndices}
\cleardoublepage\part*{Apéndices}
\addcontentsline{toc}{part}{Apéndices}

% !TeX root = ../libro.tex
% !TeX encoding = utf8

\chapter{Sobre el código de este trabajo}\label{ap:codigo-trabajo}

Junto a esta memoria, el trabajo incluye un repositorio de código accesible desde:
\begin{adjustwidth}{30pt}{}
    \href{https://github.com/mianfg-DGIIM/TFG}{https://github.com/mianfg-DGIIM/TFG}
\end{adjustwidth}
El repositorio tiene la siguiente estructura de archivos y directorios:
\vspace{8pt}
\dirtree{%
.1 /TFG.
.2 memoria\DTcomment{Memoria del trabajo}.
.3 memoria.pdf\DTcomment{Archivo PDF de esta memoria}.
.3 tex\DTcomment{Contiene el código en \LaTeX{ }{ }de la memoria}.
.4 memoria.tex.
.4 ....
.2 codigo\DTcomment{Todo el código del trabajo}.
.3 maquinas\_turing\DTcomment{Contiene descripciones de máquinas de Turing}.
.4 mas\_a\_que\_b.mt.
.4 ....
.3 adivina\_consistente.py.
.3 ....
}
\vspace{8pt}
En la carpeta \texttt{codigo} se encuentran:
\begin{itemize}
    \item Por una parte, las descripciones de máquinas de Turing comentadas en la \cref{prop:python-a-monocinta}, como archivos \texttt{.mt} dentro de la carpeta \texttt{maquinas\_turing}.
    \item Los programas desarrollados en este trabajo, junto con algunos otros.
\end{itemize}

\section{Cómo ejecutar los programas}

Los programas de la carpeta \texttt{codigo} pueden ejecutarse importando las funciones necesarias en el intérprete de Python. En la librería \texttt{utilidades.py} aparecen múltiples funciones que se usan en varios programas, y que son útiles para poder ejecutar otros. A modo de ejemplo, si queremos ejecutar la función \texttt{simula\_turing\_mult} de \texttt{simula\_turing.py}, haremos:
\begin{lstlisting}[numbers=none,frame=none]
TFG/codigo$ python
>>> from utilidades import leer
>>> from simula_turing import simula_turing_mult
>>> codificacion = leer('./maquinas_turing/mas_a_que_b.mt')
>>> entrada = 'abaaabbb'
>>> simula_turing_mult(codificacion, entrada)
'q_0 : X X X X X X X X [_] (rechaza)'
\end{lstlisting}

\section{Índice de programas}

%\footnotesize
\begin{tabularx}{\textwidth}{L{0.3} L{0.5} L{0.2}}
\midrule
Programa & Descripción & ¿Ejecutable? \\
\midrule
\texttt{adivina\_consistente.py} \linebreak \small{(\cref{lst:adivina-consistente})} & Usado en la prueba del \cref{teo:incompletitud-sintactica}: de ser el sistema lógico completo, este programa decidiría \hyperref[prob:adivina-consistente]{\textsc{AdivinaConsistente}} \vspace{4pt} & no [NE1] \\

\texttt{c\_d.py} \linebreak \small{(\cref{lst:c-d})} & Usado para probar que si un programa es decidible, su complementario también (\cref{prop:c-decidible}) \vspace{4pt} & no [NE2] \\

\texttt{c\_diagonal.py}$^\lozenge$ \linebreak \small{(\cref{lst:c-diagonal})} & Si \hyperref[prob:diagonal]{\textsc{Diagonal}} fuese decidible, este programa decidiría el problema \hyperref[prob:c-diagonal]{\textsc{C-Diagonal}} (\cref{prop:universal-no-decidible}) \vspace{4pt} & sí \\

\texttt{diagonal.py}$^\lozenge$ \linebreak \small{(\cref{lst:diagonal})} & Si \hyperref[prob:universal]{\textsc{Universal}} fuese decidible, este programa decidiría \hyperref[prob:diagonal]{\textsc{Diagonal}} (\cref{prop:universal-no-decidible}) \vspace{4pt} & sí \\

\texttt{diagonal\_a\_universal.py} \linebreak \small{(\cref{lst:diagonal-a-universal})} & Implementa la reducción de \hyperref[prob:diagonal]{\textsc{Diagonal}} a \hyperref[prob:universal]{\textsc{Universal}} \vspace{4pt} & sí \\

\texttt{es\_prueba\_peano.py} & Este programa es capaz de comprobar si una secuencia de fórmulas es prueba de otra fórmula en \hyperref[sl:peano]{\textbf{Peano}} \vspace{4pt} & no impl. [NI1] \\

\texttt{es\_teorema\_peano.py}$^\lozenge$ \linebreak \small{(\cref{lst:es-teorema-peano})} & Este programa hace \hyperref[prob:es-teorema-peano]{\textsc{EsTeoremaPeano}} semidecidible (\cref{prop:es-teorema-peano-semidecidible}) \vspace{4pt} & no [NE3] \\

\texttt{es\_teorema\_S.py} \linebreak \small{(\cref{lst:es-teorema-s})} & Prueba que \hyperref[prob:es-teorema]{\textsc{EsTeorema}} es decidible para sistemas lógicos sintácticamente consistentes (\cref{prop:es-teorema-consistente-decidible}) \vspace{4pt} & no [NE4] \\

\texttt{funcion\_con\_objeto}\linebreak\texttt{\_codificado.py} \linebreak \small{(\cref{lst:funcion-con-objeto-codificado})} & Este programa muestra que los programas SISO no imponen una restricción sobre los programas que podemos decidir: podemos aceptar como entrada y devolver como salida cualquier programa debidamente codificado (en este caso, mediante la librería nativa \texttt{picke}) \vspace{4pt} & sí \\

\texttt{godel.py} \linebreak \small{(\cref{lst:godel})} & Programa central en la demostración de la versión semántica del Primer Teorema de Incompletitud en el caso general (\cref{teo:incompletitud-semantica}) \vspace{4pt} & no [NE5] \\

\texttt{godel\_peano.py}$^\lozenge$ \linebreak \small{(\cref{lst:godel-peano})} & Programa central en la demostración de la versión semántica del Primer Teorema de Incompletitud para \hyperref[sl:peano]{\textbf{Peano}} (\cref{teo:incompletitud-peano}) \vspace{2pt} & no [NE6] \\

\midrule
\end{tabularx}
\vspace*{-0.2cm}
\begin{tabla}
\caption{Índice de programas usados en este trabajo}
\label{tab:indice-programas}
\end{tabla}
\vspace*{-1cm}

\newpage

\vspace*{0.9cm}

\begin{tabularx}{\textwidth}{L{0.3} L{0.5} L{0.2}}
\midrule
Programa & Descripción & ¿Ejecutable? \\
\midrule
\texttt{ignora\_entrada.py}$^\lozenge$ \linebreak \small{(\cref{lst:ignora-entrada})} \vspace{4pt} & Permite ignorar la entrada y ejecutar programas y entradas almacenados en disco & sí \\

\texttt{maquina\_universal.py} \linebreak \small{(\cref{lst:maquina-universal})} & Implementa una máquina universal en Python, es decir, una función que acepta como entradas un programa en Python y una entrada, y ejecuta el programa en Python con dicha entrada \vspace{4pt} & sí \\

\texttt{maquina\_universal}\linebreak\texttt{\_parada.py}$^\lozenge$ \linebreak \small{(\cref{lst:maquina-universal-parada})} \vspace{4pt} & Programa usado para demostrar que \hyperref[prob:parada]{\textsc{Parada}} es no decidible (\cref{prop:parada-no-decidible}) & sí \\

\texttt{mas\_a\_que\_b.py} \linebreak \small{(\cref{lst:mas-a-que-b})} & Este programa decide el problema \hyperref[prob:mas-a-que-b]{\textsc{MásAQueB}} (\cref{prop:masaqueb-decidible}) \vspace{4pt} & sí \\

\texttt{mas\_a\_que\_b\_v2.py} \linebreak \small{(\cref{lst:mas-a-que-b-v2})} & Este programa decide el problema \hyperref[prob:mas-a-que-b]{\textsc{MásAQueB}}, aceptando cualquier palabra como entrada \vspace{4pt} & sí \\

\texttt{modifica\_adivina}\linebreak\texttt{\_consistente.py}  \linebreak \small{(\cref{lst:modifica-adivina-consistente})} & Usado para probar la no decidibilidad de \hyperref[prob:adivina-consistente]{\textsc{AdivinaConsistente}} (\cref{prop:adivina-consistente-no-decidible}) \vspace{4pt} & no [NE7] \\

\texttt{parada\_a\_parada\_en}\linebreak\texttt{\_vacio.py} \linebreak \small{(\cref{lst:parada-a-parada-en-vacio})} \vspace{4pt} & Implementa la reducción de \hyperref[prob:parada]{\textsc{Parada}} a \hyperref[prob:parada-en-vacio]{\textsc{ParadaEnVacío}} & no [NE8] \\

\texttt{parada\_a\_peano.py} & Traduce una afirmación del tipo \emph{``$P$ para con entrada vacía''} a una fórmula de \hyperref[sl:peano]{\textbf{Peano}} \vspace{4pt} & no impl. [NI2] \\

\texttt{parada\_en\_vacio\_a\_es}\linebreak\texttt{\_verdadero\_peano.py}$^\lozenge$ \linebreak \small{(\cref{lst:parada-en-vacio-a-es-verdadero-peano})} \vspace{4pt} & Implementa la reducción de \hyperref[prob:parada-en-vacio]{\textsc{ParadaEnVacío}} a \hyperref[prob:es-verdadero-peano]{\textsc{EsVerdaderoPeano}} & no [NE9] \\

\texttt{si.py} \linebreak \small{(\cref{lst:si})} \vspace{4pt} & Este programa decide el problema \hyperref[prob:si]{\textsc{Sí}} & sí \\

\texttt{simula\_turing.py} \linebreak \small{(\cref{lst:simula-turing})} & Este programa simula una máquina de Turing. Puede usarse con cualquiera de las máquinas de la carpeta \texttt{maquinas\_turing} \vspace{4pt} & sí \\

\texttt{simula\_turing}\linebreak\texttt{\_decision.py} & Este programa es una variación de \texttt{simula\_turing.py} para hacer que el programa sea de decisión (siempre se devuelve \palabra{sí} o \palabra{no}). Este programa no se usa en el trabajo, pero se deja para experimentación \vspace{2pt} & sí \\
\midrule
\end{tabularx}
\vspace*{-0.2cm}
\begin{tabla}
\caption*{Tabla A.1. (cont.): Índice de programas usados en este trabajo}
\label{tab:indice-programas}
\end{tabla}


\newpage

\vspace*{0.9cm}

\begin{tabularx}{\textwidth}{L{0.3} L{0.5} L{0.2}}
\midrule
Programa & Descripción & ¿Ejecutable? \\
\midrule
\texttt{suma\_binaria.py}$^\lozenge$ & Este programa decide los problemas \hyperref[prob:es-teorema]{\textsc{EsTeorema}} y \hyperref[prob:es-verdadero]{\textsc{EsVerdadero}} para todos los sistemas lógicos derivados del sistema formal \hyperref[sf:suma-binaria]{\textsc{SumaBinaria}} (incluyendo la decidibilidad del problema \hyperref[prob:es-teorema]{\textsc{EsTeorema}} de este sistema formal) \vspace{4pt} & sí \\

\texttt{turing.py} & Este programa incluye una clase \texttt{Turing} que encapsula la simulación de \texttt{simula\_turing.py} \vspace{4pt} & sí \\

\texttt{universal.py} & Este programa hace el problema \hyperref[prob:universal]{\textsc{Universal}} semidecidible. Esto no es probado en el trabajo, pero se deja aquí para permitir la experimentación \vspace{4pt} & sí \\

\texttt{universal\_a\_parada.py}$^\lozenge$ \linebreak \small{(\cref{lst:universal-a-parada})} & Implementa la reducción de \hyperref[prob:universal]{\textsc{Universal}} a \hyperref[prob:parada]{\textsc{Parada}} \vspace{4pt} & no [NE10] \\

\texttt{utilidades.py} & Este programa implementa funciones que se usan en el resto de programas (``utilidades'') \vspace{2pt} & sí \\

\midrule
\end{tabularx}
\vspace*{-0.2cm}
\begin{tabla}
\caption*{Tabla A.1. (cont.): Índice de programas usados en este trabajo}
\label{tab:indice-programas}
\end{tabla}

Los programas indicados con $^\lozenge$ están basados en programas de \cite{MacCormick2018}; en concreto:

\begin{itemize}
    \item Los programas \texttt{c\_diagonal.py} y \texttt{diagonal.py} son inmediatos de implementar, y son prácticamente idénticos a los programas de la figura 3.7 de \cite{MacCormick2018}.
    \item El programa \texttt{es\_teorema\_peano.py} se basa en el programa de la figura 16.5 de \cite{MacCormick2018}.
    \item El programa \texttt{godel\_peano.py} se basa en el programa de la figura 16.8 de \cite{MacCormick2018}. Usamos la nomenclatura en \emph{strings} de la artimética de Peano que definimos en la \cref{tab:peano-string}.
    \item El programa \texttt{ignora\_entrada.py} se basa en el programa de la figura 6.9 de \cite{MacCormick2018}.
    \item El programa \texttt{maquina\_universal\_parada.py} se basa en el primer programa de la figura 7.10 de \cite{MacCormick2018}. Devolvemos el valor \palabra{sí} para hacer que el programa sea de decisión.
    \item El programa \texttt{parada\_en\_vacio\_a\_es\_verdadero\_peano.py} se basa en el programa de la figura 16.6 de \cite{MacCormick2018}.
    \item El programa \texttt{suma\_binaria.py} usa las expresiones regulares que figuran en el programa \texttt{binAd.py} de la librería de código de \cite{MacCormick2018}. Además, se incluye el código correspondiente a las reglas del resto de sistemas lógicos del capítulo, y se reestructura el código para poder evaluar las fórmulas en cualquiera de estos sistemas.
    \item El programa \texttt{universal\_a\_parada.py} se basa en el segundo programa de la figura 7.10 de \cite{MacCormick2018}. Reestructuramos el código para hacerlo más legible.
\end{itemize}

\subsection*{Justificación de programas no ejecutables}

Algunos de los programas que aparecen en la carpeta \texttt{codigo} no son ejecutables. Los motivos de esto se explican a continuación.

\begin{enumerate}[label={[}NE\arabic*{]},wide = 0pt,widest={100}, leftmargin =*]
    \item Las funciones \texttt{parada\_a\_S} y \texttt{es\_prueba\_S} no están implementadas, pues corresponden a un sistema lógico cualquiera (sería necesario crear una implementación concreta para cada sistema lógico), pero su existencia es probada.
    \item La función \texttt{d} no existe (este programa es ilustrativo). Basta probar a sustituir \texttt{d} por cualquier otro programa de decisión.
    \item La función \texttt{es\_prueba\_peano} no está implementada (véase [NI1]).
    \item La función \texttt{es\_prueba\_S} no está implementada (véase [NE1]).
    \item Las funciones \texttt{es\_teorema\_G} y \texttt{parada\_a\_G} no están implementadas, pues corresponden a un sistema formal cualquiera (sería necesario crear una implementación concreta para cada sistema lógico), pero su existencia se supone.
    \item La función \texttt{parada\_a\_peano} no está implementada (véase [NI2]).
    \item No puede escribirse dado que \textsc{AdivinaConsistente} es no decidible (y no está implementado).
    \item Es una reducción en la que \texttt{parada\_en\_vacio} es un oráculo.
    \item Es una reducción en la que \texttt{es\_verdadero\_peano} es un oráculo, y donde \texttt{parada\_a\_peano} no está implementada (véase [NI2]), pero NO es un oráculo (su existencia se prueba).
    \item Es una reducción en la que \texttt{parada} es un oráculo (de hecho no está implementada porque no es computable).
\end{enumerate}

\subsection*{Justificación de programas no implementados}

Dos de los programas a los que hacemos referencia en este trabajo no han sido implementados. A continuación se concreta el por qué.

\begin{enumerate}[label={[}NI\arabic*{]},wide = 0pt,widest={NI1,NI}, leftmargin =*]
    \item[{[}NI1, NI2{]}] Implementar estas funciones sería un trabajo muy arduo, que se deja como posible en un futuro (véase el \cref{ch:trabajo-futuro}).
\end{enumerate}

\endinput

% !TeX root = ../libro.tex
% !TeX encoding = utf8

\chapter{Demostraciones adicionales}\label{ap:demostraciones-adicionales}

% ===================+
\begin{proposicion}\label{prop:ap-mtm}
Toda máquina de Turing multicinta puede ser simulada por una máquina de Turing de una sola cinta.
\end{proposicion}
\begin{proof}
Sea $M=(Q, A, B, \delta, q_0, \Vtextvisiblespace\:, F)$ una máquina de Turing con $k$ cintas. Definiremos la máquina $\widetilde{M}=(\widetilde{Q}, A, \widetilde{B}, \widetilde{\delta}, q_0, [\Vtextvisiblespace^k], F)$. Para ello, definimos en primer lugar $\widecheck{B}$ como:
$$
    \widecheck{B} = \{\widecheck{b} : b \in B\}
$$
Para definir $\widetilde{B}$, observemos que un símbolo de un alfabeto no tiene por qué ser únicamente un carácter. Para representar un único símbolo, lo rodearemos de corchetes ($[\:]$). Por tanto, el siguiente alfabeto tiene sentido:
$$
    \widetilde{B} = \{[u_1\:u_2\:...\:u_k]:u_i\in B\cup\widecheck{B}\}\cup\{\#\}
$$
En $\widetilde{B}$ insertamos, además de los símbolos que hemos comentado, un símbolo delimitador $\#$.

El conjunto de estados $\widetilde{Q}$ lo hacemos más amplio que $Q$. Para ello, numeraremos todos los estados de $Q$:
$$
    Q=\{q_0,q_1,...,q_m\}
$$
Y definimos el nuevo conjunto de estados como:
$$
    \widetilde{Q}=Q\cup\left(\bigcup_{i=0}^m\bigcup_{j=1}^k\left\{\widetilde{q}_{i,j},\:\widetilde{q}_{i,j}^{(R)},\:\widetilde{q}_{i,j}^{(I)},\:\widetilde{q}_{i,j}^{(Is)},\:\widetilde{q}_{i,j}^{(D)},\:\widetilde{q}_{i,j}^{(Ds)}\right\}\right)
$$
Definimos como configuración inicial de la máquina con entrada $u=u_1\:u_2\:...\:u_n\in A^*$ a:\footnote{Siempre es posible cambiar la posición de la configuración inicial de una máquina de Turing, ejecutando las transiciones necesarias.}
$$
    q_0\;:\;\fbox{\#}\:[\widecheck{u_1}\:\underbrace{\widecheck{\Vtextvisiblespace} \:... \:\widecheck{\Vtextvisiblespace}}_{k-1)}]\:[u_2\:\underbrace{{\Vtextvisiblespace} \:... \:{\Vtextvisiblespace}}_{k-1)}]\:...\:[u_n\:\underbrace{{\Vtextvisiblespace} \:... \:{\Vtextvisiblespace}}_{k-1)}]\:\#
$$
Describiremos $\widetilde{\delta}$ poco a poco. En primer lugar, queremos partir del mismo estado que la máquina multicinta (esto nos permitirá que el lenguaje aceptado sea el mismo, como veremos más adelante). En el momento en el que la máquina se encuentre en este estado, nos movemos a la derecha y comenzaremos a buscar la posición en la cinta del primer cabezal. Usaremos el estado $\widetilde{q}_{i,j}$ para buscar el cabezal $j$-ésimo, estando la máquina multicinta en el estado $q_i$. Añadimos para hacer esto las siguientes transiciones:
\begin{equation}
    \widetilde{\delta}(q_i, \#) = (\widetilde{q}_{i,1}, \#, D)\;\;\text{ para cada }\;\;i=0,1,...,m
\end{equation}
Para buscar el cabezal $j$-ésimo, estando la máquina multicinta en el estado $q_i$, nos iremos moviendo a la derecha mientras que no hayamos encontrado un carácter de $\widecheck{B}$ en la posición $j$-ésima. Para ello, añadimos para cada $i=0,1,...,m$ y $j=1,2,...,k$:
\begin{equation}
    \widetilde{\delta}(\widetilde{q}_{i,j}, [ w_1\:w_2\:...\:w_l\:...\:w_k]) = (\widetilde{q}_{i,j}, [ w_1\:w_2\:...\:w_l\:...\:w_k], D)\;\;\text{ con }\;\;w_j\in B, w_l \in B\cup\widecheck{B} \text{ para } l \neq j
\end{equation}
Nótese que, de forma explícita, hacemos que el símbolo $w_j$ no pueda estar en $\widecheck{B}$ (de aquí en adelante, si no explicitamos, será $w_l\in B\cup\widecheck{B}$ para cada $l=0,1,...,k$). Esto es para continuar moviéndonos a la derecha hasta encontrar un carácter con cabezal en la posición $j$-ésima de cada símbolo.

Ahora estamos en condiciones de incluir las transiciones de $\delta$ en $\widetilde{\delta}$. Sea una transición de $\delta$:
$$
    \delta(q_i, b_1, b_2, ..., b_k) = (q_j, c_1, M_1, c_2, M_2, ..., c_k, M_k)
$$
Para cada una de estas transiciones, deberemos insertar reglas en $\widetilde{\delta}$. Estas transiciones dependerán de cómo se mueva el cabezal, ya que si lo hace a la derecha o a la izquierda hasta toparse con el delimitador $\#$, deberemos extender la cinta de $\widetilde{M}$ con un símbolo blanco de $\widetilde{M}$ (que es $[\Vtextvisiblespace^k]$), y desplazar el delimitador $\#$ una posición. Dada la transición anterior, y para cada $l=1, 2, ..., k$:
\begin{itemize}
    \item Si $M_l=D$, insertamos varias transiciones. En primer lugar, la transición correspondiente a sustituir el símbolo en el que se encuentra el cabezal, $b_l$, por el correspondiente, $c_l$:
    \begin{equation}
        \widetilde{\delta}(\widetilde{q}_{i,l}, [w_1\:w_2\:...\:w_{l-1}\:\widecheck{b_l}\:w_{l+1}\:...\:w_k]) = (\widetilde{q}_{i,l}^{(D)}, [w_1\:w_2\:...\:w_{l-1}\:c_l\:w_{l+1}\:...\:w_k], D)
    \end{equation}
    A continuación simulamos que el cabezal $l$-ésimo se mueve a la derecha. En caso de que no hayamos alcanzado el separador $\#$, la transición es directa, simplemente indicamos que el cabezal está en la posición determinada, y volvemos al primer delimitador $\#$ mediante el estado $^{(R)}$ (que veremos más adelante):
    \begin{equation}
        \widetilde{\delta}(\widetilde{q}_{i,l}^{(D)}, [w_1\:w_2\:...\:w_{l-1}\:w_l\:w_{l+1}\:...\:w_k]) = (\widetilde{q}_{i,l}^{(R)}, [w_1\:w_2\:...\:w_{l-1}\:\widecheck{w_l}\:w_{l+1}\:...\:w_k], I)
    \end{equation}
    En caso de que nos topemos con el delimitador $\#$, iniciaremos una ``subrutina'' indicada con los estados $^{(Ds)}$, en el que añadiremos un símbolo blanco y moveremos el delimitador a la derecha, siguiendo al estado $^{(R)}$:
    \begin{align}
        \widetilde{\delta}(\widetilde{q}_{i,l}^{(D)}, \#) &= (\widetilde{q}_{i,l}^{(Ds)}, [\underbrace{{\Vtextvisiblespace} \:... \:{\Vtextvisiblespace}}_{l-1)}\:\widecheck{\Vtextvisiblespace}\:\underbrace{{\Vtextvisiblespace} \:... \:{\Vtextvisiblespace}}_{k-l)}], D)\\
        \widetilde{\delta}(\widetilde{q}_{i,l}^{(Ds)}, [\Vtextvisiblespace^k]) &= (\widetilde{q}_{i,l}^{(R)}, \#, I)\\
    \end{align}
    \item Si $M_l=I$, añadimos de forma análoga las transiciones:
    \begin{align}
        \widetilde{\delta}(\widetilde{q}_{i,l}, [w_1\:w_2\:...\:w_{l-1}\:\widecheck{b_l}\:w_{l+1}\:...\:w_k]) &= (\widetilde{q}_{i,l}^{(I)}, [w_1\:w_2\:...\:w_{l-1}\:c_l\:w_{l+1}\:...\:w_k], I)\\
        \widetilde{\delta}(\widetilde{q}_{i,l}^{(I)}, [w_1\:w_2\:...\:w_{l-1}\:w_l\:w_{l+1}\:...\:w_k]) &= (\widetilde{q}_{i,l}^{(R)}, [w_1\:w_2\:...\:w_{l-1}\:\widecheck{w_l}\:w_{l+1}\:...\:w_k], I)\\
        \widetilde{\delta}(\widetilde{q}_{i,l}^{(I)}, \#) &= (\widetilde{q}_{i,l}^{(Is)}, [\underbrace{{\Vtextvisiblespace} \:... \:{\Vtextvisiblespace}}_{l-1)}\:\widecheck{\Vtextvisiblespace}\:\underbrace{{\Vtextvisiblespace} \:... \:{\Vtextvisiblespace}}_{k-l)}], I)\\
        \widetilde{\delta}(\widetilde{q}_{i,l}^{(Is)}, [\Vtextvisiblespace^k]) &= (\widetilde{q}_{i,l}^{(R)}, \#, D)\\
    \end{align}
    \item Si $M_l=S$, la transición es más inmediata, dado que no debemos mover el cabezal, simplemente reemplazar el símbolo y cambiar a $^{(R)}$:
    \begin{align}
        \widetilde{\delta}(\widetilde{q}_{i,l}, [w_1\:w_2\:...\:w_{l-1}\:\widecheck{b_l}\:w_{l+1}\:...\:w_k]) &= (\widetilde{q}_{i,l}^{(R)}, [w_1\:w_2\:...\:w_{l-1}\:\widecheck{c_l}\:w_{l+1}\:...\:w_k], I)
    \end{align}
\end{itemize}
Ahora, nos queda una forma de cambiar la cinta que estamos analizando. Comenzaremos añadiendo para ello un estado de retorno, que posiciona el cabezal en el primer delimitador $\#$. Para cada $i=0,1,...,m$ y $j=1,2,...,k$ añadimos las transiciones:
\begin{equation}
    \widetilde{\delta}(\widetilde{q}_{i,j}^{(R)}, [w_1\:w_2\:...\:w_l\:...\:w_k]) = (\widetilde{q}_{i,j}^{(R)}, [w_1\:w_2\:...\:w_l\:...\:w_k], I)
\end{equation}
Para leer la cinta siguiente, simplemente añadimos estas transiciones para cada $i=0,1,...,m$ y $j=1,2,...,k-1$:
\begin{equation}
    \widetilde{\delta}(\widetilde{q}_{i,j}^{(R)}, \#) = (\widetilde{q}_{i,j+1}^{(R)}, \#, D)
\end{equation}
Si ya hemos leído la última cinta, cambiamos el estado a uno de los de $Q$ de acuerdo a la función de transición. De nuevo revisitamos $\delta$. Para cada transición:
$$
    \delta(q_i, b_1, b_2, ..., b_k) = (q_j, c_1, M_1, c_2, M_2, ..., c_k, M_k)
$$
añadimos la transición de cambio de estado:
\begin{equation}
    \widetilde{\delta}(\widetilde{q}_{i,k}^{(R)}, \#) = (q_j, \#, S)
\end{equation}
Ya hemos definido todas las transiciones necesarias. Ahora, comprobaremos que la simulación es correcta, es decir, que $L(M)=L(\widetilde{M})$. En efecto, hemos construido una máquina $\widetilde{M}$ que ejecuta el mismo proceso de cálculo que $M$, y que además pasa por los mismos estados $Q$ de $M$ tras modificar la cinta. En otras palabras, la máquina $M$  Es precisamente por esto por lo que, en caso de la máquina parar, se encontrará en un estado $q\in Q$ (la hemos construido para que ocurra esto). De este modo, el lenguaje aceptado por ambas máquinas coincide, dado que usamos el mismo conjunto de estados finales $F$.
\end{proof}
% ===================+

\pagebreak

% ===================+
\begin{proposicion}\label{prop:f-s}
Para cada problema $F:X\longrightarrow Y$, existe un problema equivalente $\widehat{F}:S\longrightarrow S$.
\end{proposicion}
\begin{proof}
Si queremos resolver un problema de ordenador, deberemos codificar las entradas y las salidas mediante un formato que pueda entender. La idea es poder traducir un problema $F$ a otro problema $\widehat{F}$ de \emph{strings}, de modo que procese la misma función $F$ para entradas correctamente codificadas, y se devuelva \palabra{no} en caso de que la codificación de entrada sea incorrecta.

Para formalizar este concepto convenientemente, definiremos los conjuntos $X^\oslash = X \cup \{\oslash\}$ y $Y^\oslash = Y \cup \{\oslash\}$. A $\oslash$ lo llamaremos el \emph{elemento error}.

Expandiremos el problema $F$ para aceptar el elemento error como entrada, de modo que su solución sea el elemento error. Lo llamaremos $F^\oslash : X^\oslash \longrightarrow Y^\oslash$, definido como:
$$
    F^\oslash(x) = \begin{cases} 
      F(x), & \text{si} \;\;\; x \in X \\
      \oslash, & \text{si} \;\;\; x = \oslash 
   \end{cases}
$$
Sea $\mathcal{C}_X : S \longrightarrow X^\oslash$ la función que, para cada codificación en $S$, asigna la entrada de $X$ correspondiente. Si la codificación es inválida, asigna $\oslash$.

Sea $\mathcal{C}_Y^* : Y^\oslash \longrightarrow S$ la función que, para cada solución de $Y$, asigna la codificación en $S$ correspondiente. Para la solución error $\oslash$, asigna la codificación \palabra{no}. Del mismo modo que para las entradas, también podemos definir $\mathcal{C}_Y : S \longrightarrow Y^\oslash$, que para cada codificación de una solución en $S$, obtiene su solución en $Y$. Para codificaciones inválidas, asigna $\oslash$.

Ahora estamos en condiciones de definir $\widehat{F} = \mathcal{C}_Y^* \circ F^\oslash \circ \mathcal{C}_X$.
\begin{figure}[H]
\centering

\tikzset{every picture/.style={line width=0.75pt}} %set default line width to 0.75pt        

\begin{tikzpicture}[x=0.75pt,y=0.75pt,yscale=-1,xscale=1]
%uncomment if require: \path (0,629); %set diagram left start at 0, and has height of 629

%Straight Lines [id:da560650477859397] 
\draw    (210.17,60.86) -- (278,60.86) ;
\draw [shift={(280,60.86)}, rotate = 180] [color={rgb, 255:red, 0; green, 0; blue, 0 }  ][line width=0.75]    (10.93,-4.9) .. controls (6.95,-2.3) and (3.31,-0.67) .. (0,0) .. controls (3.31,0.67) and (6.95,2.3) .. (10.93,4.9)   ;
%Curve Lines [id:da8170296687430567] 
\draw    (210,73.17) .. controls (295.24,123.25) and (414.47,123.33) .. (498.73,73.92) ;
\draw [shift={(500,73.17)}, rotate = 149.25] [color={rgb, 255:red, 0; green, 0; blue, 0 }  ][line width=0.75]    (10.93,-3.29) .. controls (6.95,-1.4) and (3.31,-0.3) .. (0,0) .. controls (3.31,0.3) and (6.95,1.4) .. (10.93,3.29)   ;
%Straight Lines [id:da24312545548628428] 
\draw    (320.17,60.86) -- (388,60.86) ;
\draw [shift={(390,60.86)}, rotate = 180] [color={rgb, 255:red, 0; green, 0; blue, 0 }  ][line width=0.75]    (10.93,-4.9) .. controls (6.95,-2.3) and (3.31,-0.67) .. (0,0) .. controls (3.31,0.67) and (6.95,2.3) .. (10.93,4.9)   ;
%Straight Lines [id:da8363047188359607] 
\draw    (430.17,60.86) -- (498,60.86) ;
\draw [shift={(500,60.86)}, rotate = 180] [color={rgb, 255:red, 0; green, 0; blue, 0 }  ][line width=0.75]    (10.93,-4.9) .. controls (6.95,-2.3) and (3.31,-0.67) .. (0,0) .. controls (3.31,0.67) and (6.95,2.3) .. (10.93,4.9)   ;

% Text Node
\draw (194.75,61.35) node  [  ] [align=left] {\begin{minipage}[lt]{20.51pt}\setlength\topsep{0pt}
\begin{center}
$S$
\end{center}

\end{minipage}};
% Text Node
\draw (300,61.35) node  [  ] [align=left] {\begin{minipage}[lt]{27.2pt}\setlength\topsep{0pt}
\begin{center}
$X^\oslash$
\end{center}

\end{minipage}};
% Text Node
\draw (409.83,61.35) node  [  ] [align=left] {\begin{minipage}[lt]{26.97pt}\setlength\topsep{0pt}
\begin{center}
$Y^\oslash$
\end{center}

\end{minipage}};
% Text Node
\draw (514.42,61.35) node  [  ] [align=left] {\begin{minipage}[lt]{20.51pt}\setlength\topsep{0pt}
\begin{center}
$S$
\end{center}

\end{minipage}};
% Text Node
\draw (244.83,45.51) node  [  ] [align=left] {\begin{minipage}[lt]{47.83pt}\setlength\topsep{0pt}
\begin{center}
$\mathcal{C}_X$
\end{center}

\end{minipage}};
% Text Node
\draw (354.83,45.51) node  [  ] [align=left] {\begin{minipage}[lt]{47.83pt}\setlength\topsep{0pt}
\begin{center}
$F^\oslash$
\end{center}

\end{minipage}};
% Text Node
\draw (464.83,45.51) node  [  ] [align=left] {\begin{minipage}[lt]{47.83pt}\setlength\topsep{0pt}
\begin{center}
$\mathcal{C}_Y^*$
\end{center}

\end{minipage}};
% Text Node
\draw (354.83,95.51) node  [  ] [align=left] {\begin{minipage}[lt]{47.83pt}\setlength\topsep{0pt}
\begin{center}
$\widehat{F}$
\end{center}

\end{minipage}};


\end{tikzpicture}

\end{figure}
El problema $\widehat{F}$ resuelve el mismo problema que $F$, y es capaz de tratar con entradas erróneas. Basta ver que, si $I\in S$ es una codificación correcta de una entrada $x\in X$ (es decir, $\mathcal{C}_X(I)=x$), entonces es $\widehat{F}(I) = \mathcal{C}_Y^*(F^\oslash(\mathcal{C}_X(I))) = \mathcal{C}_Y^*(F^\oslash(x)) = \mathcal{C}_Y^*(F(x)) = O$, siendo $O$ la codificación de $F(x)$, $O=\mathcal{C}_Y^*(F(x))$. Dada la codificación de la solución $O$, es fácil obtener la solución $F(x) = \mathcal{C}_Y(O)$.

En el caso de que $I \in S$ sea una codificación de entrada incorrecta, será $\widehat{F}(I)=\mathcal{C}_Y^*(F^\oslash(\mathcal{C}_X(I)))=\mathcal{C}_Y^*(F^\oslash(\oslash))=\mathcal{C}_Y^*(\oslash)=\palabra{no}$.
\end{proof}
% ===================+

\pagebreak

\begin{proposicion}\label{prop:ap-cons}
El sistema lógico \normalfont{\textbf{SumaBinariaLógico}} es consistente.
\end{proposicion}
\begin{proof}
Para probar que el sistema es consistente, veremos que todas las reglas de inferencia mapean fórmulas verdaderas a fórmulas verdaderas. En primer lugar, vemos que el axioma \texttt{1=1} es verdadero, pues es cierto que $1=1$.

Ahora, comprobemos cada regla:

$$
    \begin{matrix}
    \text{(R1)} & A\texttt{=}B & \mapsto & \texttt{1+}A\texttt{=}\texttt{1+}B 
    \end{matrix}
$$

En este caso, si $A$ y $B$ son números binarios y asumimos que la fórmula $A\texttt{=}B$ es verdadera, entonces tenemos que $A=B$. Por tanto, la fórmula $\texttt{1+}A\texttt{=1+}B$ es verdadera, ya que es cierto que $1+A=1+B$.

$$
    \begin{matrix}
    \text{(R2)} & A\underline{N}\texttt{+}\underline{N}B & \mapsto & AN\texttt{0}B
    \end{matrix}
$$

El hecho de que las $N$ tengan coincidencias maximales nos lleva a comprobar si dado $N$ un número binario, asumiendo que $N\texttt{+}N$ es verdadera tenemos que $N\texttt{0}$ es verdadera, o lo que es lo mismo, debemos comprobar que $N\texttt{+}N=N\texttt{0}$.

Un número binario $N$ se representa mediante $m$ dígitos, ya sean $0$ o $1$, siendo \linebreak $N=n_{m-1}\:n_{m-2}\:...\:n_1\:n_0$, y se verifica que:
$$N = n_{m-1}\cdot2^{m-1} + n_{m-2}\cdot2^{m-2}+...+n_1\cdot2+n_0$$

Basta ver que:
$$N+N=2N=n_{m-1}\cdot2^{m} + n_{m-2}\cdot2^{m-1}+...+n_1\cdot2^2+n_0\cdot 2$$
y que el número $N+N$ corresponde exactamente a $n_{m-1}\:n_{m-2}\:...\:n_1\:n_0\:0=N\:0$.
$$
    \begin{matrix}
    \text{(R3)} & A\texttt{\underline{1}+}\underline{N\texttt{0}}B & \mapsto & AN\texttt{1}B
    \end{matrix}
$$
De forma análoga a la regla anterior, comprobamos que:
$$1+N\:0=N\:0+1=n_{m-1}\cdot2^{m} + n_{m-2}\cdot2^{m-1}+...+n_1\cdot2^2+n_0\cdot 2+1=N\:1$$
\end{proof}

\begin{proposicion}\label{prop:ap-comp}
El sistema lógico \normalfont{\textbf{SumaBinariaLógico}} es completo.
\end{proposicion}
\begin{proof}
Veremos que toda fórmula verdadera es un teorema. Vemos en primer lugar que toda fórmula verdadera $\phi$ es del tipo
$$\phi=N_1\texttt{+}N_2\texttt{+}...\texttt{+}N_k\texttt{=}M_1\texttt{+}M_2\texttt{+}...\texttt{+}M_l$$
donde cada $N_i, M_j\;\forall i\in\{1, 2, ..., k\}, j\in\{1, 2, ..., l\}$ es un número binario (que verifica la expresión regular $\texttt{1(0|1)}^\texttt{*}$, y de modo que la siguiente fórmula aritmética es cierta:
$$N_1+N_2+...+N_k=M_1+M_2+...+M_l$$
Nuestro objetivo es probar que la fórmula $\phi$ es un teorema. Para ello, vamos a probar que, dado un número binario $n$ que sea compilación maximal de la expresión regular \texttt{1(0|1)*}, éste puede derivarse a partir de la suma de $n$ unos, es decir, debemos probar que:
$$A\underline{n}B\xmapsto{*}A\underbrace{\texttt{1}\texttt{+}\texttt{1}\texttt{+}...\texttt{+}\texttt{1}}_{n)}B$$
Dado que, por ser esto cierto, tenemos que:
$$\psi=\underbrace{\texttt{1}\texttt{+}\texttt{1}\texttt{+}...\texttt{+}\texttt{1}}_{N_1)}\texttt{+}\underbrace{\texttt{1}\texttt{+}\texttt{1}\texttt{+}...\texttt{+}\texttt{1}}_{N_2)}\texttt{+}...\texttt{+}\underbrace{\texttt{1}\texttt{+}\texttt{1}\texttt{+}...\texttt{+}\texttt{1}}_{N_k)}\texttt{=}\underbrace{\texttt{1}\texttt{+}\texttt{1}\texttt{+}...\texttt{+}\texttt{1}}_{M_1)}\texttt{+}\underbrace{\texttt{1}\texttt{+}\texttt{1}\texttt{+}...\texttt{+}\texttt{1}}_{M_2)}\texttt{+}...\texttt{+}\underbrace{\texttt{1}\texttt{+}\texttt{1}\texttt{+}...\texttt{+}\texttt{1}}_{M_l)}\xmapsto{*}\phi$$
Y como tenemos $N_1+N_2+...+N_k$ unos a la izquierda de \texttt{=} y $M_1+M_2+...+M_l$ a la derecha, y es
$$N_1+N_2+...+N_k=M_1+M_2+...+M_l=n$$
tenemos que
$$\psi=\underbrace{\texttt{1}\texttt{+}\texttt{1}\texttt{+}...\texttt{+}\texttt{1}}_{n)}\texttt{=}\underbrace{\texttt{1}\texttt{+}\texttt{1}\texttt{+}...\texttt{+}\texttt{1}}_{n)}$$
pero es claro que podemos obtener $\psi$ aplicando la regla (R1) $n$ veces a partir del axioma $1=1$, luego es
$$\texttt{1=1}\xmapsto{*}\psi=\underbrace{\texttt{1}\texttt{+}\texttt{1}\texttt{+}...\texttt{+}\texttt{1}}_{n)}\texttt{=}\underbrace{\texttt{1}\texttt{+}\texttt{1}\texttt{+}...\texttt{+}\texttt{1}}_{n)}\xmapsto{*}\phi$$
y como $\texttt{1=1}\xmapsto{*}\phi$ y es \texttt{1=1} un axioma, $\phi$ es un teorema.

Probemos que, en efecto, es
$$A\underline{n}B\xmapsto{*}A\underbrace{\texttt{1}\texttt{+}\texttt{1}\texttt{+}...\texttt{+}\texttt{1}}_{n)}B$$

Para ello, primero observemos que, si $m$ es un número binario positivo, entonces puede que sea par o impar.

Si $m$ es impar, entonces puede ser escrito con $k$ dígitos: $m=b_{k-1}\:b_{k-2}\:...\:b_1\:\texttt{1}$. Sea $\varphi$ una fórmula de \textbf{SumaBinariaLógico}, y supongamos que $\varphi$ contiene un \emph{substring} $m$ como compilación maximal de $\texttt{1(0|1)}^\texttt{*}$, es decir, $\phi=A\:\underline{m}\:B$. Entonces, $\varphi$ puede ser derivado de un cierto $\varphi'$, $\varphi'\mapsto\varphi$ mediante la regla (R3), siendo$$\varphi'=A\:\underline{b_{k-1}\:b_{k-2}\:...\:b_1\:\texttt{0}}\:B$$

En caso de que $m$ sea par, es claro que puede expresarse como la suma $p+p$, cierto $p$ un cierto binario positivo. Entonces, si $\varphi$ contiene un \emph{substring} $m$ como compilación maximal de $\texttt{1(0|1)}^\texttt{*}$, $\varphi=A\:\underline{m}\:B$, es $\varphi'\mapsto\varphi$ con $$\varphi'=A\:\underline{p}\texttt{+}\underline{p}\:B$$

Sabiendo esto, cualquier ocurrencia de un entero positivo $m$ en una fórmula de \textbf{SumaBinariaLógico} puede ser reducida a la suma de $m$ unos, es decir, dado $\varphi=A\:\underline{m}\:B$, se tiene que:
$$A\:\underbrace{\texttt{1}\texttt{+}\texttt{1}\texttt{+}...\texttt{+}\texttt{1}}_{m)}\:B\xmapsto{*}\varphi$$
Esto es debido a que todo número natural $m$ puede ser reducido a la suma de $m$ unos aplicando recursivamente las operaciones anteriormente descritas (dividir por $2$ en caso de ser par y restar $1$ en caso de ser impar).
\end{proof}

\begin{proposicion}\label{prop:ap-restr-arr}
El sistema lógico \normalfont{\textbf{SumaBinariaLógicoArreglado}} es completo.
\end{proposicion}
\begin{proof}
Para demostrar esto, usaremos el resultado obtenido en la proposición anterior. Toda fórmula verdadera $\phi$ es del tipo
$$\phi=N_1\texttt{+}N_2\texttt{+}...\texttt{+}N_k\texttt{=}M_1\texttt{+}M_2\texttt{+}...\texttt{+}M_l$$
donde cada $N_i, M_j\;\forall i\in\{1, 2, ..., k\}, j\in\{1, 2, ..., l\}$ es un número binario (que verifica la expresión regular $\texttt{1(0|1)}^\texttt{*}$, y de modo que la siguiente fórmula aritmética es cierta:
$$N_1+N_2+...+N_k\geq M_1+M_2+...+M_l$$
Vamos a llamar a ambas partes con $n$ y $m$:
$$n=N_1+N_2+...+N_k, \;\;\;\;m=M_1+M_2+...+M_l$$
Como $n\geq m$, existe un $p\geq 0$ de modo que $n=m+p$. Ya probamos en la proposición anterior que:
$$\psi=\underbrace{\texttt{1}\texttt{+}\texttt{1}\texttt{+}...\texttt{+}\texttt{1}}_{p)}\texttt{+}\underbrace{\texttt{1}\texttt{+}\texttt{1}\texttt{+}...\texttt{+}\texttt{1}}_{m)}\texttt{=}\underbrace{\texttt{1}\texttt{+}\texttt{1}\texttt{+}...\texttt{+}\texttt{1}}_{m)}\xmapsto{*}\phi$$
Ahora basta ver que:
$$\texttt{1=1}\underset{(1)}{\overset{*}{\longmapsto}}\underbrace{\texttt{1}\texttt{+}\texttt{1}\texttt{+}...\texttt{+}\texttt{1}}_{m)}\texttt{=}\underbrace{\texttt{1}\texttt{+}\texttt{1}\texttt{+}...\texttt{+}\texttt{1}}_{m)}\underset{(2)}{\overset{*}{\longmapsto}}\psi\xmapsto{*}\phi$$
Donde en $(1)$ aplicamos $m$ veces la regla (R1) y en $(2)$ aplicamos $p$ veces la regla (R1b).
\end{proof}

\endinput
%------------------------------------------------------------------------------------
% FIN DEL APÉNDICE. 
%------------------------------------------------------------------------------------

% !TeX root = ../libro.tex
% !TeX encoding = utf8

\chapter{Descripción de la aritmética de Peano}\label{ap:peano}

A continuación se especifica con detalle el sistema lógico \textbf{Peano}, correspondiente a la aritmética de Peano de \emph{Principia Mathematica} \cite{Whitehead1927} y descrito por Gödel en su publicación en la que demostró el Primer Teorema de Incompletitud \cite{Godel1931}. Los contenidos del \cref{sl:peano} se corresponden a una traducción de tal artículo en \cite{Godel1992}, con la ayuda de \cite{Hirzel2000}.

En el artículo original de Gödel, los símbolos usados son diferentes. En la \cref{tab:peano-simbolos} se especifican tales diferencias.

\begin{sistemalogico}[Aritmética de Peano (\textbf{Peano})]\label{sl:peano}
Describimos las fórmulas del \emph{lenguaje} de \textbf{Peano} como sigue:

\begin{itemize}
    \item[I.] Las \emph{constantes}: $\neg$, $\vee$, $\forall$, $0$, $S$, $($ y $)$ (véase la \cref{tab:peano-simbolos}).
    \item[II.] Las \emph{variables de tipo uno} (individuales, son números naturales incluyendo el $0$): $x_1, y_1, z_1, ...$
    
    Las \emph{variables de tipo dos} (clases de individuos, subconjuntos de $\mathbb{N}$): $x_2, y_2, z_2, ...$

    Las \emph{variables de tipo tres} (clases de clases de individuos, subconjuntos de subconjuntos de $\mathbb{N}$): $x_3, y_3, z_3, ...$

    Y sigue con un tipo para cada número natural.
\end{itemize}

Una nota: las variables para las funciones de dos o más términos son superfluas como signos básicos, ya que las relaciones pueden ser definidas como clases de pares ordenados y los pares ordenados pueden ser definidos como clases de clases, por ejemplo, el par ordenado $a,b$ por $((a), (a,b))$ donde $(x,y)$ significa la clase cuyos elementos son $x$ e $y$, y $x$ la clase cuyo único elemento es $x$.

Por un \emph{signo de primer tipo} entendemos una combinación de signos de la forma:$$a, S(a), S(S(a)), S(S(S(a))), ...$$
donde $a$ es bien $0$ o una variable de primer tipo. En el primer caso llamamos a tal símbolo un \emph{símbolo numérico}. Para $n>1$ entendemos un \emph{símbolo de tipo $n$-ésimo} lo mismo que una \emph{variable de tipo $n$-ésimo}. Las combinaciones de los símbolos de la forma $a(b)$, donde $b$ es un símbolo de tipo $n$ y $a$ es un símbolo de tipo $n+1$, se llaman \emph{fórmulas elementales}.

Las \emph{fórmulas} [del sistema lógico] las definimos como la menor clase que contiene todas las fórmulas elementales y, además, con cualquier $a$ y $b$ las siguientes: $\neg (a)$, $(a) \vee (b)$, $\forall x : (a)$ (donde $x$ es cualquier variable). Llamamos a $(a)\vee(b)$ la \emph{disyunción} de $a$ y $b$, a $\neg(a)$ la \emph{negación} y a $\forall x : (a)$ la \emph{generalización} de $a$. Una fórmula en la que no hay ninguna variable libre\footnote{Las variables libres son aquellas que no son cuantificadas por un $\forall$ o un $\exists$.} se llama \emph{fórmula proposicional}. [...]

Denotamos como $\text{Subst } a\binom{v}{b}$ (donde $a$ es una fórmula, $v$ es una variable y $b$ es un símbolo del mismo tipo que $v$) a la fórmula que se deriva a partir de $a$ reemplazando todas las ocurrencias de $v$, donde sean libres, por $b$.

Decimos que una fórmula $a$ es una \emph{elevación de tipo} de otra fórmula $b$ si podemos obtener $b$ incrementando el tipo de todas las variables de $a$ por el mismo número.

Las fórmulas siguientes (I-V) son los \emph{axiomas} [de nuestro sistema lógico]. Los indicamos con las abreviaturas $\wedge$, $\Rightarrow$, $\iff$, $\exists$ y $=$\footnote{Observa que: $p\wedge q$ se define como $\not(\not p \vee \not q)$, $p\Rightarrow q$ como $\neg p \vee q$, $p \iff q$ como $(\neg p \vee q) \wedge (\neg q \vee p)$, $\exists x : p$ como $\neg\forall x: \neg p$, y $=$ como $\forall x_2 : (x_2(x_1) \Rightarrow x_2(y_1))$.} y sujetos a las convenciones usuales para eliminar paréntesis:\footnote{Eliminamos los paréntesis cuando no dé lugar a confusión.}
\begin{itemize}
    \item[I.] \begin{align*}
        \text{(P1)} &\;\;\;\; \neg(S(x_1)=0) \\
        \text{(P2)} &\;\;\;\; S(x_1)=S(y_1) \Rightarrow x_1=y_1 \\
        \text{(P3)} &\;\;\;\; \left(x_2(0)\wedge\forall x_1 : x_2(x_1) \Rightarrow x_2(S(x_1))\right) \Rightarrow \forall x_1 : x_2(x_1)
    \end{align*}
    \item[II.] Toda fórmula derivada de insertar cualquier fórmula para $p$, $q$ y $r$ en los esquemas siguientes:
    \begin{align}
        \text{(P4)} &\;\;\;\; p \vee p \Rightarrow p \\
        \text{(P5)} &\;\;\;\; p \Rightarrow p \vee q \\
        \text{(P6)} &\;\;\;\; p \vee q \Rightarrow q \vee p \\
        \text{(P7)} &\;\;\;\; (p \Rightarrow q) \Rightarrow (r \vee p \Rightarrow r \vee q)
    \end{align}
    \item[III.] Cada fórmula obtenida de los dos esquemas:
    \begin{align}
        \text{(P8)} &\;\;\;\; (\forall v : a) \Rightarrow \text{Subst } a\binom{v}{c} \\
        \text{(P9)} &\;\;\;\; (\forall v : b \vee a) \Rightarrow (b \vee \forall v : a)
    \end{align}

    haciendo las siguientes sustituciones para $a$, $v$, $b$ y $c$ (y realizando en (P8) la operación denotada por Subst): para cada cualquier fórmula $a$, cualquier variable $v$, para una fórmula cualquiera $b$ en la que $v$ no aparece de forma libre, y para $c$ un signo del mismo tipo que $v$, de modo que $c$ no contenga una variable que esté enlazada a $a$ en un lugar donde $v$ sea libre.\footnote{$c$ es bien $0$, bien un signo de la forma $S(S(S...(u)..)$ donde $u$ es $0$ o una variable de tipo $1$.}

    \item[IV.] Cada fórmula obtenida del esquema:
    $$
    \text{(P10)} \;\;\;\; \exists u : \forall v : (u(v) \iff a)
    $$
    insertando en el lugar de $v$ y $u$ cualquier variable de tipo $n$ y $n+1$, respectivamente, y en $a$ una fórmula sin ocurrencias libres de $u$. Este axioma es el de \emph{reducibilidad}.\footnote{Un axioma de la teoría de conjuntos.}

    \item[V.] Cualquier fórmula obtenida de la siguiente mediante elevación de tipo (y la fórmula en sí misma):
    $$
    \text{(P11)} \;\;\;\; \left(\forall x_1 : (x_2(x_1)\iff y_2(x_1))\right)\Rightarrow x_2=y_2
    $$

    Este axioma nos indica que una clase está completamente determinada por sus elementos.
\end{itemize}

Una fórmula $c$ se llama \emph{consecuencia inmediata} de $a$ y $b$ si $a$ es la fórmula $(\neg (b))\vee (c)$, y una \emph{consecuencia inmediata} de [únicamente] $a$ si $c$ es la fórmula $\forall v : a$, donde $v$ es cualquier variable. La clase de [\emph{teoremas}] es definida como la menor clase de fórmulas que contiene a los axiomas y que es cerrada bajo la operación ``consecuencia inmediata''.

Una fórmula de \textbf{Peano} es \emph{verdadera} si es cierta interpretándola como una afirmación lógica sobre los números naturales.
\end{sistemalogico}

\vfill
% ====================
\begin{tabla}
\begin{table}[H]
\centering
\begin{tabular}{@{}ccl@{}}
\toprule
Símbolo & Notación original (Gödel) & Descripción \\ \midrule
$\neg$ & $\sim$ & Negación \\
$\vee$ & $\vee$ & Disyunción \\
$\forall$ & $\Pi$ & Cuantificador universal \\
$S$ & $f$ & Sucesor \\
$($ & $($ & Paréntesis de apertura \\
$)$ & $)$ & Paréntesis de cierre \\
$\exists$ & $(Ex)$ & Cuantificador existencial \\
$:$ & $\textbf{.}$ & Tal que \\
$\Rightarrow$ & $\supset$ & Implicación \\
$\iff$ & $\equiv$ & Coimplicación \\
$=$ & $=$ & Igualdad \\\bottomrule
\end{tabular}
\end{table}
\vspace{-8pt}
\caption{Correspondencia entre los símbolos usados en este trabajo y los símbolos usados originalmente por Gödel \cite{Godel1931}}
\label{tab:peano-simbolos}
\end{tabla}
% ====================
\vfill
\endinput
%------------------------------------------------------------------------------------
% FIN DEL APÉNDICE. 
%------------------------------------------------------------------------------------

% Añadir tantos apéndices como sea necesario 

% --------------------------------------------------------------------
% GLOSARIO: Opcional
% --------------------------------------------------------------------

%\input{glosario} 

% -------------------------------------------------------------------
% BACKMATTER
% -------------------------------------------------------------------
\cleardoublepage
\phantomsection

\backmatter % Desactiva la numeración de los capítulos
\cleardoublepage
\phantomsection

%\addcontentsline{toc}{chapter}{Índice alfabético}
%\cleardoublepage\part{Referencias e Índices}
%\pdfbookmark[-1]{Referencias e Índices}{BM-Referencias}

%\printindex

% BIBLIOGRAFÍA
%-------------------------------------------------------------------

%\setbibpreamble{References are listed in alphabetical order, by author. The references with more than one author are ordered according to the first author.\par\bigskip}
%\bibliographystyle{apalike} 
%\nocite{*}
%\
\begin{small} % Normalmente la bibliografía se imprime en un tamaño de letra más pequeño.

\cleardoublepage
\phantomsection
%\addcontentsline{toc}{part}{Bibliografía}
%\printbibliography[heading=subbibliography,title={Bibliografía básica},segment=1]
\part*{Referencias e Índices}
\addcontentsline{toc}{part}{Referencias e Índices}

\phantomsection
\vspace*{1.5cm}
{\sffamily \LARGE \textbf{Bibliografía}}\\[8pt]
\thispagestyle{plain}
\addcontentsline{toc}{chapter}{Bibliografía}
\markboth{Bibliografía}{Bibliografía}

{\sffamily \Large {Bibliografía básica}}\\

{\normalsize El contenido de este trabajo sigue el orden expuesto en \cite{MacCormick2018}, profundizando en los conceptos explicados en tal obra.}\\

%\pdfbookmark[-1]{Bibliografía}{Bibliografia}

\printbibliography[heading=subbibliography,keyword={basica},heading=none]
\vspace{18pt}
\printbibliography[heading=subbibliography,notkeyword={basica},heading=none]

\end{small}

\printindex
\markboth{Índice alfabético}{Índice alfabético}


%\printbibliography

%\bibliography{library.bib}

\phantomsection 
\listoffigures
\markboth{Lista de figuras}{Lista de figuras}


\phantomsection 
\listoftablas

\phantomsection 
\listofproblemas

\markboth{Lista de figuras}{Lista de figuras}

\listoftheorems[ignoreall,show={definicion},title={Índice de definiciones}]
%\addcontentsline{toc}{chapter}{Índice de definiciones}

\listoftheorems[ignoreall,show={sistemaformal,sistemalogico},title={Índice de sistemas formales y lógicos}]
%\addcontentsline{toc}{chapter}{Índice de sistemas formales y lógicos}


%\listoftheorems

% ÍNDICE TERMINOLÓGICO  (Opcional) 
%------------------------------------------------------------------- 

% Para incluir el índice terminológico es necesario descomentar los siguientes comandos. Incluir un índice terminológico es opcional

% \cleardoublepage 
% \begin{footnotesize} % Normalmente el índice se imprime en un tamaño de letra más pequeño.
% \printindex 
% \end{footnotesize}

\end{document}
